\subsection*{Textbook Formal Inventory}
This subsection records the exact formal material in \file{HAETAE_HopGames.ec}.  It is generated from the proof-surface EasyCrypt file, so the line numbers and statements are meant to be read next to the checked source.

\noindent\textbf{Chapter role.} HAETAE-specific hybrid games, simulator transitions, and hop bounds.

\noindent\textbf{Inventory size.} 88 context declarations, 148 definitions/game objects, 271 lemmas or relational theorems, and 12 explicit Hoare/Phoare judgments were found.

\subsubsection*{Theory Context, Imports, and Quantified Modules}
\paragraph{Context 1 (line 1)}
\noindent\textbf{EasyCrypt name.}
\begin{lstlisting}[language=EasyCrypt,basicstyle=\ttfamily\scriptsize]
imported theories
\end{lstlisting}
\noindent\textbf{Checked EasyCrypt statement.}
\begin{lstlisting}[language=EasyCrypt,basicstyle=\ttfamily\scriptsize]
require import AllCore Distr List Real FSet FMap StdOrder FelTactic Mu_mem.
\end{lstlisting}
\noindent\textbf{Formal mathematical reading.}
Mathematical context. This line imports or specializes earlier theories, making their carrier sets, functions, modules, and theorems available to the current file.

\noindent\textbf{Role in the proof.}
Use in this chapter. It fixes the proof context, imported mathematical language, or quantified module parameters for the following statements.

\paragraph{Context 2 (line 2)}
\noindent\textbf{EasyCrypt name.}
\begin{lstlisting}[language=EasyCrypt,basicstyle=\ttfamily\scriptsize]
imported theories
\end{lstlisting}
\noindent\textbf{Checked EasyCrypt statement.}
\begin{lstlisting}[language=EasyCrypt,basicstyle=\ttfamily\scriptsize]
require import Sig_ROM HAETAE_Scheme HAETAE_Params HAETAE_Algebra.
\end{lstlisting}
\noindent\textbf{Formal mathematical reading.}
Mathematical context. This line imports or specializes earlier theories, making their carrier sets, functions, modules, and theorems available to the current file.

\noindent\textbf{Role in the proof.}
Use in this chapter. It fixes the proof context, imported mathematical language, or quantified module parameters for the following statements.

\paragraph{Context 3 (line 3)}
\noindent\textbf{EasyCrypt name.}
\begin{lstlisting}[language=EasyCrypt,basicstyle=\ttfamily\scriptsize]
imported theories
\end{lstlisting}
\noindent\textbf{Checked EasyCrypt statement.}
\begin{lstlisting}[language=EasyCrypt,basicstyle=\ttfamily\scriptsize]
require import HAETAE_Distributions.
\end{lstlisting}
\noindent\textbf{Formal mathematical reading.}
Mathematical context. This line imports or specializes earlier theories, making their carrier sets, functions, modules, and theorems available to the current file.

\noindent\textbf{Role in the proof.}
Use in this chapter. It fixes the proof context, imported mathematical language, or quantified module parameters for the following statements.

\paragraph{Context 4 (line 4)}
\noindent\textbf{EasyCrypt name.}
\begin{lstlisting}[language=EasyCrypt,basicstyle=\ttfamily\scriptsize]
imported theories
\end{lstlisting}
\noindent\textbf{Checked EasyCrypt statement.}
\begin{lstlisting}[language=EasyCrypt,basicstyle=\ttfamily\scriptsize]
require import HAETAE_Assumptions.
\end{lstlisting}
\noindent\textbf{Formal mathematical reading.}
Mathematical context. This line imports or specializes earlier theories, making their carrier sets, functions, modules, and theorems available to the current file.

\noindent\textbf{Role in the proof.}
Use in this chapter. It fixes the proof context, imported mathematical language, or quantified module parameters for the following statements.

\paragraph{Context 5 (line 5)}
\noindent\textbf{EasyCrypt name.}
\begin{lstlisting}[language=EasyCrypt,basicstyle=\ttfamily\scriptsize]
imported theories
\end{lstlisting}
\noindent\textbf{Checked EasyCrypt statement.}
\begin{lstlisting}[language=EasyCrypt,basicstyle=\ttfamily\scriptsize]
require import HAETAE_Transcript.
\end{lstlisting}
\noindent\textbf{Formal mathematical reading.}
Mathematical context. This line imports or specializes earlier theories, making their carrier sets, functions, modules, and theorems available to the current file.

\noindent\textbf{Role in the proof.}
Use in this chapter. It fixes the proof context, imported mathematical language, or quantified module parameters for the following statements.

\paragraph{Context 6 (line 6)}
\noindent\textbf{EasyCrypt name.}
\begin{lstlisting}[language=EasyCrypt,basicstyle=\ttfamily\scriptsize]
imported theories
\end{lstlisting}
\noindent\textbf{Checked EasyCrypt statement.}
\begin{lstlisting}[language=EasyCrypt,basicstyle=\ttfamily\scriptsize]
require import HAETAE_Reductions.
\end{lstlisting}
\noindent\textbf{Formal mathematical reading.}
Mathematical context. This line imports or specializes earlier theories, making their carrier sets, functions, modules, and theorems available to the current file.

\noindent\textbf{Role in the proof.}
Use in this chapter. It fixes the proof context, imported mathematical language, or quantified module parameters for the following statements.

\paragraph{Context 7 (line 7)}
\noindent\textbf{EasyCrypt name.}
\begin{lstlisting}[language=EasyCrypt,basicstyle=\ttfamily\scriptsize]
imported theories
\end{lstlisting}
\noindent\textbf{Checked EasyCrypt statement.}
\begin{lstlisting}[language=EasyCrypt,basicstyle=\ttfamily\scriptsize]
require import HAETAE_Rejection HAETAE_ROM HAETAE_ROM_Programming.
\end{lstlisting}
\noindent\textbf{Formal mathematical reading.}
Mathematical context. This line imports or specializes earlier theories, making their carrier sets, functions, modules, and theorems available to the current file.

\noindent\textbf{Role in the proof.}
Use in this chapter. It fixes the proof context, imported mathematical language, or quantified module parameters for the following statements.

\paragraph{Context 8 (line 9)}
\noindent\textbf{EasyCrypt name.}
\begin{lstlisting}[language=EasyCrypt,basicstyle=\ttfamily\scriptsize]
HAETAE_HopGames
\end{lstlisting}
\noindent\textbf{Checked EasyCrypt statement.}
\begin{lstlisting}[language=EasyCrypt,basicstyle=\ttfamily\scriptsize]
theory HAETAE_HopGames.
\end{lstlisting}
\noindent\textbf{Formal mathematical reading.}
Mathematical context. This begins a namespace for the formal development in this file.

\noindent\textbf{Role in the proof.}
Use in this chapter. It fixes the proof context, imported mathematical language, or quantified module parameters for the following statements.

\paragraph{Context 9 (line 11)}
\noindent\textbf{EasyCrypt name.}
\begin{lstlisting}[language=EasyCrypt,basicstyle=\ttfamily\scriptsize]
opened namespace
\end{lstlisting}
\noindent\textbf{Checked EasyCrypt statement.}
\begin{lstlisting}[language=EasyCrypt,basicstyle=\ttfamily\scriptsize]
import HAETAE_Scheme.
\end{lstlisting}
\noindent\textbf{Formal mathematical reading.}
Mathematical context. This line imports or specializes earlier theories, making their carrier sets, functions, modules, and theorems available to the current file.

\noindent\textbf{Role in the proof.}
Use in this chapter. It fixes the proof context, imported mathematical language, or quantified module parameters for the following statements.

\paragraph{Context 10 (line 12)}
\noindent\textbf{EasyCrypt name.}
\begin{lstlisting}[language=EasyCrypt,basicstyle=\ttfamily\scriptsize]
opened namespace
\end{lstlisting}
\noindent\textbf{Checked EasyCrypt statement.}
\begin{lstlisting}[language=EasyCrypt,basicstyle=\ttfamily\scriptsize]
import HAETAE_Params.
\end{lstlisting}
\noindent\textbf{Formal mathematical reading.}
Mathematical context. This line imports or specializes earlier theories, making their carrier sets, functions, modules, and theorems available to the current file.

\noindent\textbf{Role in the proof.}
Use in this chapter. It fixes the proof context, imported mathematical language, or quantified module parameters for the following statements.

\paragraph{Context 11 (line 13)}
\noindent\textbf{EasyCrypt name.}
\begin{lstlisting}[language=EasyCrypt,basicstyle=\ttfamily\scriptsize]
opened namespace
\end{lstlisting}
\noindent\textbf{Checked EasyCrypt statement.}
\begin{lstlisting}[language=EasyCrypt,basicstyle=\ttfamily\scriptsize]
import HAETAE_Algebra.
\end{lstlisting}
\noindent\textbf{Formal mathematical reading.}
Mathematical context. This line imports or specializes earlier theories, making their carrier sets, functions, modules, and theorems available to the current file.

\noindent\textbf{Role in the proof.}
Use in this chapter. It fixes the proof context, imported mathematical language, or quantified module parameters for the following statements.

\paragraph{Context 12 (line 14)}
\noindent\textbf{EasyCrypt name.}
\begin{lstlisting}[language=EasyCrypt,basicstyle=\ttfamily\scriptsize]
opened namespace
\end{lstlisting}
\noindent\textbf{Checked EasyCrypt statement.}
\begin{lstlisting}[language=EasyCrypt,basicstyle=\ttfamily\scriptsize]
import HAETAE_Distributions.
\end{lstlisting}
\noindent\textbf{Formal mathematical reading.}
Mathematical context. This line imports or specializes earlier theories, making their carrier sets, functions, modules, and theorems available to the current file.

\noindent\textbf{Role in the proof.}
Use in this chapter. It fixes the proof context, imported mathematical language, or quantified module parameters for the following statements.

\paragraph{Context 13 (line 15)}
\noindent\textbf{EasyCrypt name.}
\begin{lstlisting}[language=EasyCrypt,basicstyle=\ttfamily\scriptsize]
opened namespace
\end{lstlisting}
\noindent\textbf{Checked EasyCrypt statement.}
\begin{lstlisting}[language=EasyCrypt,basicstyle=\ttfamily\scriptsize]
import HAETAE_Assumptions.
\end{lstlisting}
\noindent\textbf{Formal mathematical reading.}
Mathematical context. This line imports or specializes earlier theories, making their carrier sets, functions, modules, and theorems available to the current file.

\noindent\textbf{Role in the proof.}
Use in this chapter. It fixes the proof context, imported mathematical language, or quantified module parameters for the following statements.

\paragraph{Context 14 (line 16)}
\noindent\textbf{EasyCrypt name.}
\begin{lstlisting}[language=EasyCrypt,basicstyle=\ttfamily\scriptsize]
opened namespace
\end{lstlisting}
\noindent\textbf{Checked EasyCrypt statement.}
\begin{lstlisting}[language=EasyCrypt,basicstyle=\ttfamily\scriptsize]
import HAETAE_Transcript.
\end{lstlisting}
\noindent\textbf{Formal mathematical reading.}
Mathematical context. This line imports or specializes earlier theories, making their carrier sets, functions, modules, and theorems available to the current file.

\noindent\textbf{Role in the proof.}
Use in this chapter. It fixes the proof context, imported mathematical language, or quantified module parameters for the following statements.

\paragraph{Context 15 (line 17)}
\noindent\textbf{EasyCrypt name.}
\begin{lstlisting}[language=EasyCrypt,basicstyle=\ttfamily\scriptsize]
opened namespace
\end{lstlisting}
\noindent\textbf{Checked EasyCrypt statement.}
\begin{lstlisting}[language=EasyCrypt,basicstyle=\ttfamily\scriptsize]
import HAETAE_Reductions.
\end{lstlisting}
\noindent\textbf{Formal mathematical reading.}
Mathematical context. This line imports or specializes earlier theories, making their carrier sets, functions, modules, and theorems available to the current file.

\noindent\textbf{Role in the proof.}
Use in this chapter. It fixes the proof context, imported mathematical language, or quantified module parameters for the following statements.

\paragraph{Context 16 (line 18)}
\noindent\textbf{EasyCrypt name.}
\begin{lstlisting}[language=EasyCrypt,basicstyle=\ttfamily\scriptsize]
opened namespace
\end{lstlisting}
\noindent\textbf{Checked EasyCrypt statement.}
\begin{lstlisting}[language=EasyCrypt,basicstyle=\ttfamily\scriptsize]
import HAETAE_Rejection.
\end{lstlisting}
\noindent\textbf{Formal mathematical reading.}
Mathematical context. This line imports or specializes earlier theories, making their carrier sets, functions, modules, and theorems available to the current file.

\noindent\textbf{Role in the proof.}
Use in this chapter. It fixes the proof context, imported mathematical language, or quantified module parameters for the following statements.

\paragraph{Context 17 (line 19)}
\noindent\textbf{EasyCrypt name.}
\begin{lstlisting}[language=EasyCrypt,basicstyle=\ttfamily\scriptsize]
opened namespace
\end{lstlisting}
\noindent\textbf{Checked EasyCrypt statement.}
\begin{lstlisting}[language=EasyCrypt,basicstyle=\ttfamily\scriptsize]
import HAETAE_ROM.
\end{lstlisting}
\noindent\textbf{Formal mathematical reading.}
Mathematical context. This line imports or specializes earlier theories, making their carrier sets, functions, modules, and theorems available to the current file.

\noindent\textbf{Role in the proof.}
Use in this chapter. It fixes the proof context, imported mathematical language, or quantified module parameters for the following statements.

\paragraph{Context 18 (line 20)}
\noindent\textbf{EasyCrypt name.}
\begin{lstlisting}[language=EasyCrypt,basicstyle=\ttfamily\scriptsize]
opened namespace
\end{lstlisting}
\noindent\textbf{Checked EasyCrypt statement.}
\begin{lstlisting}[language=EasyCrypt,basicstyle=\ttfamily\scriptsize]
import HAETAE_ROM_Programming.
\end{lstlisting}
\noindent\textbf{Formal mathematical reading.}
Mathematical context. This line imports or specializes earlier theories, making their carrier sets, functions, modules, and theorems available to the current file.

\noindent\textbf{Role in the proof.}
Use in this chapter. It fixes the proof context, imported mathematical language, or quantified module parameters for the following statements.

\paragraph{Context 19 (line 21)}
\noindent\textbf{EasyCrypt name.}
\begin{lstlisting}[language=EasyCrypt,basicstyle=\ttfamily\scriptsize]
opened namespace
\end{lstlisting}
\noindent\textbf{Checked EasyCrypt statement.}
\begin{lstlisting}[language=EasyCrypt,basicstyle=\ttfamily\scriptsize]
import RealOrder.
\end{lstlisting}
\noindent\textbf{Formal mathematical reading.}
Mathematical context. This line imports or specializes earlier theories, making their carrier sets, functions, modules, and theorems available to the current file.

\noindent\textbf{Role in the proof.}
Use in this chapter. It fixes the proof context, imported mathematical language, or quantified module parameters for the following statements.

\paragraph{Context 20 (line 1766)}
\noindent\textbf{EasyCrypt name.}
\begin{lstlisting}[language=EasyCrypt,basicstyle=\ttfamily\scriptsize]
ROMInternalTranscriptBudgetedPaperSimNMAQueryBudget
\end{lstlisting}
\noindent\textbf{Checked EasyCrypt statement.}
\begin{lstlisting}[language=EasyCrypt,basicstyle=\ttfamily\scriptsize]
section ROMInternalTranscriptBudgetedPaperSimNMAQueryBudget.
\end{lstlisting}
\noindent\textbf{Formal mathematical reading.}
Mathematical context. This opens a scoped block of hypotheses, module parameters, and lemmas.

\noindent\textbf{Role in the proof.}
Use in this chapter. It fixes the proof context, imported mathematical language, or quantified module parameters for the following statements.

\paragraph{Context 21 (line 1768)}
\noindent\textbf{EasyCrypt name.}
\begin{lstlisting}[language=EasyCrypt,basicstyle=\ttfamily\scriptsize]
H
\end{lstlisting}
\noindent\textbf{Checked EasyCrypt statement.}
\begin{lstlisting}[language=EasyCrypt,basicstyle=\ttfamily\scriptsize]
declare module H <: SIG.Oracle {-ROMInternalTranscriptBudgetedPaperSimAsNMA}.
\end{lstlisting}
\noindent\textbf{Formal mathematical reading.}
Mathematical context. This is universal quantification over an adversary, oracle, scheme, sampler, or reduction module satisfying the stated interface and memory restrictions.

\noindent\textbf{Role in the proof.}
Use in this chapter. It fixes the proof context, imported mathematical language, or quantified module parameters for the following statements.

\paragraph{Context 22 (line 1769)}
\noindent\textbf{EasyCrypt name.}
\begin{lstlisting}[language=EasyCrypt,basicstyle=\ttfamily\scriptsize]
A
\end{lstlisting}
\noindent\textbf{Checked EasyCrypt statement.}
\begin{lstlisting}[language=EasyCrypt,basicstyle=\ttfamily\scriptsize]
declare module A <: SIG.Adversary {-H,
                                   -ROMInternalTranscriptBudgetedPaperSimAsNMA}.
\end{lstlisting}
\noindent\textbf{Formal mathematical reading.}
Mathematical context. This is universal quantification over an adversary, oracle, scheme, sampler, or reduction module satisfying the stated interface and memory restrictions.

\noindent\textbf{Role in the proof.}
Use in this chapter. It fixes the proof context, imported mathematical language, or quantified module parameters for the following statements.

\paragraph{Context 23 (line 1771)}
\noindent\textbf{EasyCrypt name.}
\begin{lstlisting}[language=EasyCrypt,basicstyle=\ttfamily\scriptsize]
Samp
\end{lstlisting}
\noindent\textbf{Checked EasyCrypt statement.}
\begin{lstlisting}[language=EasyCrypt,basicstyle=\ttfamily\scriptsize]
declare module Samp <: PaperSimSigningSampler
  {-A, -ROMInternalTranscriptBudgetedPaperSimAsNMA}.
\end{lstlisting}
\noindent\textbf{Formal mathematical reading.}
Mathematical context. This is universal quantification over an adversary, oracle, scheme, sampler, or reduction module satisfying the stated interface and memory restrictions.

\noindent\textbf{Role in the proof.}
Use in this chapter. It fixes the proof context, imported mathematical language, or quantified module parameters for the following statements.

\paragraph{Context 24 (line 2979)}
\noindent\textbf{EasyCrypt name.}
\begin{lstlisting}[language=EasyCrypt,basicstyle=\ttfamily\scriptsize]
InternalTranscriptSignSound
\end{lstlisting}
\noindent\textbf{Checked EasyCrypt statement.}
\begin{lstlisting}[language=EasyCrypt,basicstyle=\ttfamily\scriptsize]
section InternalTranscriptSignSound.
\end{lstlisting}
\noindent\textbf{Formal mathematical reading.}
Mathematical context. This opens a scoped block of hypotheses, module parameters, and lemmas.

\noindent\textbf{Role in the proof.}
Use in this chapter. It fixes the proof context, imported mathematical language, or quantified module parameters for the following statements.

\paragraph{Context 25 (line 2981)}
\noindent\textbf{EasyCrypt name.}
\begin{lstlisting}[language=EasyCrypt,basicstyle=\ttfamily\scriptsize]
H
\end{lstlisting}
\noindent\textbf{Checked EasyCrypt statement.}
\begin{lstlisting}[language=EasyCrypt,basicstyle=\ttfamily\scriptsize]
declare module H <: SIG.Oracle {-EUF_CMA_InternalTranscriptSign,
                                -EUF_CMA_ROMInternalTranscriptSign}.
\end{lstlisting}
\noindent\textbf{Formal mathematical reading.}
Mathematical context. This is universal quantification over an adversary, oracle, scheme, sampler, or reduction module satisfying the stated interface and memory restrictions.

\noindent\textbf{Role in the proof.}
Use in this chapter. It fixes the proof context, imported mathematical language, or quantified module parameters for the following statements.

\paragraph{Context 26 (line 2983)}
\noindent\textbf{EasyCrypt name.}
\begin{lstlisting}[language=EasyCrypt,basicstyle=\ttfamily\scriptsize]
A
\end{lstlisting}
\noindent\textbf{Checked EasyCrypt statement.}
\begin{lstlisting}[language=EasyCrypt,basicstyle=\ttfamily\scriptsize]
declare module A <: SIG.Adversary {-H, -EUF_CMA_InternalTranscriptSign,
                                   -EUF_CMA_ROMInternalTranscriptSign}.
\end{lstlisting}
\noindent\textbf{Formal mathematical reading.}
Mathematical context. This is universal quantification over an adversary, oracle, scheme, sampler, or reduction module satisfying the stated interface and memory restrictions.

\noindent\textbf{Role in the proof.}
Use in this chapter. It fixes the proof context, imported mathematical language, or quantified module parameters for the following statements.

\paragraph{Context 27 (line 3576)}
\noindent\textbf{EasyCrypt name.}
\begin{lstlisting}[language=EasyCrypt,basicstyle=\ttfamily\scriptsize]
CountedROMInternalTranscriptDiscipline
\end{lstlisting}
\noindent\textbf{Checked EasyCrypt statement.}
\begin{lstlisting}[language=EasyCrypt,basicstyle=\ttfamily\scriptsize]
section CountedROMInternalTranscriptDiscipline.
\end{lstlisting}
\noindent\textbf{Formal mathematical reading.}
Mathematical context. This opens a scoped block of hypotheses, module parameters, and lemmas.

\noindent\textbf{Role in the proof.}
Use in this chapter. It fixes the proof context, imported mathematical language, or quantified module parameters for the following statements.

\paragraph{Context 28 (line 3578)}
\noindent\textbf{EasyCrypt name.}
\begin{lstlisting}[language=EasyCrypt,basicstyle=\ttfamily\scriptsize]
H
\end{lstlisting}
\noindent\textbf{Checked EasyCrypt statement.}
\begin{lstlisting}[language=EasyCrypt,basicstyle=\ttfamily\scriptsize]
declare module H <: Oracle {-CountedLazyROM,
                             -EUF_CMA_ROMInternalTranscriptCountedSign}.
\end{lstlisting}
\noindent\textbf{Formal mathematical reading.}
Mathematical context. This is universal quantification over an adversary, oracle, scheme, sampler, or reduction module satisfying the stated interface and memory restrictions.

\noindent\textbf{Role in the proof.}
Use in this chapter. It fixes the proof context, imported mathematical language, or quantified module parameters for the following statements.

\paragraph{Context 29 (line 3580)}
\noindent\textbf{EasyCrypt name.}
\begin{lstlisting}[language=EasyCrypt,basicstyle=\ttfamily\scriptsize]
A
\end{lstlisting}
\noindent\textbf{Checked EasyCrypt statement.}
\begin{lstlisting}[language=EasyCrypt,basicstyle=\ttfamily\scriptsize]
declare module A <: SIG.Adversary {-H, -CountedLazyROM,
                                   -EUF_CMA_ROMInternalTranscriptCountedSign}.
\end{lstlisting}
\noindent\textbf{Formal mathematical reading.}
Mathematical context. This is universal quantification over an adversary, oracle, scheme, sampler, or reduction module satisfying the stated interface and memory restrictions.

\noindent\textbf{Role in the proof.}
Use in this chapter. It fixes the proof context, imported mathematical language, or quantified module parameters for the following statements.

\paragraph{Context 30 (line 3771)}
\noindent\textbf{EasyCrypt name.}
\begin{lstlisting}[language=EasyCrypt,basicstyle=\ttfamily\scriptsize]
ProgrammedROMTranscriptDiscipline
\end{lstlisting}
\noindent\textbf{Checked EasyCrypt statement.}
\begin{lstlisting}[language=EasyCrypt,basicstyle=\ttfamily\scriptsize]
section ProgrammedROMTranscriptDiscipline.
\end{lstlisting}
\noindent\textbf{Formal mathematical reading.}
Mathematical context. This opens a scoped block of hypotheses, module parameters, and lemmas.

\noindent\textbf{Role in the proof.}
Use in this chapter. It fixes the proof context, imported mathematical language, or quantified module parameters for the following statements.

\paragraph{Context 31 (line 3773)}
\noindent\textbf{EasyCrypt name.}
\begin{lstlisting}[language=EasyCrypt,basicstyle=\ttfamily\scriptsize]
H
\end{lstlisting}
\noindent\textbf{Checked EasyCrypt statement.}
\begin{lstlisting}[language=EasyCrypt,basicstyle=\ttfamily\scriptsize]
declare module H <: Oracle {-CountedLazyROM,
                             -EUF_CMA_ROMProgrammedTranscriptCountedSign}.
\end{lstlisting}
\noindent\textbf{Formal mathematical reading.}
Mathematical context. This is universal quantification over an adversary, oracle, scheme, sampler, or reduction module satisfying the stated interface and memory restrictions.

\noindent\textbf{Role in the proof.}
Use in this chapter. It fixes the proof context, imported mathematical language, or quantified module parameters for the following statements.

\paragraph{Context 32 (line 3775)}
\noindent\textbf{EasyCrypt name.}
\begin{lstlisting}[language=EasyCrypt,basicstyle=\ttfamily\scriptsize]
A
\end{lstlisting}
\noindent\textbf{Checked EasyCrypt statement.}
\begin{lstlisting}[language=EasyCrypt,basicstyle=\ttfamily\scriptsize]
declare module A <: SIG.Adversary {-H, -CountedLazyROM,
                                   -EUF_CMA_ROMProgrammedTranscriptCountedSign}.
\end{lstlisting}
\noindent\textbf{Formal mathematical reading.}
Mathematical context. This is universal quantification over an adversary, oracle, scheme, sampler, or reduction module satisfying the stated interface and memory restrictions.

\noindent\textbf{Role in the proof.}
Use in this chapter. It fixes the proof context, imported mathematical language, or quantified module parameters for the following statements.

\paragraph{Context 33 (line 3906)}
\noindent\textbf{EasyCrypt name.}
\begin{lstlisting}[language=EasyCrypt,basicstyle=\ttfamily\scriptsize]
BudgetedProgrammedROMTranscriptDiscipline
\end{lstlisting}
\noindent\textbf{Checked EasyCrypt statement.}
\begin{lstlisting}[language=EasyCrypt,basicstyle=\ttfamily\scriptsize]
section BudgetedProgrammedROMTranscriptDiscipline.
\end{lstlisting}
\noindent\textbf{Formal mathematical reading.}
Mathematical context. This opens a scoped block of hypotheses, module parameters, and lemmas.

\noindent\textbf{Role in the proof.}
Use in this chapter. It fixes the proof context, imported mathematical language, or quantified module parameters for the following statements.

\paragraph{Context 34 (line 3908)}
\noindent\textbf{EasyCrypt name.}
\begin{lstlisting}[language=EasyCrypt,basicstyle=\ttfamily\scriptsize]
H
\end{lstlisting}
\noindent\textbf{Checked EasyCrypt statement.}
\begin{lstlisting}[language=EasyCrypt,basicstyle=\ttfamily\scriptsize]
declare module H <: Oracle {-CountedLazyROM,
                             -DirectSigningCoinSampler,
                             -EUF_CMA_ROMProgrammedTranscriptBudgetedCountedSign}.
\end{lstlisting}
\noindent\textbf{Formal mathematical reading.}
Mathematical context. This is universal quantification over an adversary, oracle, scheme, sampler, or reduction module satisfying the stated interface and memory restrictions.

\noindent\textbf{Role in the proof.}
Use in this chapter. It fixes the proof context, imported mathematical language, or quantified module parameters for the following statements.

\paragraph{Context 35 (line 3911)}
\noindent\textbf{EasyCrypt name.}
\begin{lstlisting}[language=EasyCrypt,basicstyle=\ttfamily\scriptsize]
A
\end{lstlisting}
\noindent\textbf{Checked EasyCrypt statement.}
\begin{lstlisting}[language=EasyCrypt,basicstyle=\ttfamily\scriptsize]
declare module A <: SIG.Adversary {-H, -CountedLazyROM,
                                   -DirectSigningCoinSampler,
                                   -EUF_CMA_ROMProgrammedTranscriptBudgetedCountedSign}.
\end{lstlisting}
\noindent\textbf{Formal mathematical reading.}
Mathematical context. This is universal quantification over an adversary, oracle, scheme, sampler, or reduction module satisfying the stated interface and memory restrictions.

\noindent\textbf{Role in the proof.}
Use in this chapter. It fixes the proof context, imported mathematical language, or quantified module parameters for the following statements.

\paragraph{Context 36 (line 4520)}
\noindent\textbf{EasyCrypt name.}
\begin{lstlisting}[language=EasyCrypt,basicstyle=\ttfamily\scriptsize]
BudgetedSampledROMTranscriptDiscipline
\end{lstlisting}
\noindent\textbf{Checked EasyCrypt statement.}
\begin{lstlisting}[language=EasyCrypt,basicstyle=\ttfamily\scriptsize]
section BudgetedSampledROMTranscriptDiscipline.
\end{lstlisting}
\noindent\textbf{Formal mathematical reading.}
Mathematical context. This opens a scoped block of hypotheses, module parameters, and lemmas.

\noindent\textbf{Role in the proof.}
Use in this chapter. It fixes the proof context, imported mathematical language, or quantified module parameters for the following statements.

\paragraph{Context 37 (line 4522)}
\noindent\textbf{EasyCrypt name.}
\begin{lstlisting}[language=EasyCrypt,basicstyle=\ttfamily\scriptsize]
H
\end{lstlisting}
\noindent\textbf{Checked EasyCrypt statement.}
\begin{lstlisting}[language=EasyCrypt,basicstyle=\ttfamily\scriptsize]
declare module H <: Oracle {-CountedLazyROM,
                             -DirectSigningCoinSampler,
                             -EUF_CMA_ROMProgrammedTranscriptBudgetedSampledSign}.
\end{lstlisting}
\noindent\textbf{Formal mathematical reading.}
Mathematical context. This is universal quantification over an adversary, oracle, scheme, sampler, or reduction module satisfying the stated interface and memory restrictions.

\noindent\textbf{Role in the proof.}
Use in this chapter. It fixes the proof context, imported mathematical language, or quantified module parameters for the following statements.

\paragraph{Context 38 (line 4525)}
\noindent\textbf{EasyCrypt name.}
\begin{lstlisting}[language=EasyCrypt,basicstyle=\ttfamily\scriptsize]
A
\end{lstlisting}
\noindent\textbf{Checked EasyCrypt statement.}
\begin{lstlisting}[language=EasyCrypt,basicstyle=\ttfamily\scriptsize]
declare module A <: SIG.Adversary {-H, -CountedLazyROM,
                                   -DirectSigningCoinSampler,
                                   -EUF_CMA_ROMProgrammedTranscriptBudgetedSampledSign}.
\end{lstlisting}
\noindent\textbf{Formal mathematical reading.}
Mathematical context. This is universal quantification over an adversary, oracle, scheme, sampler, or reduction module satisfying the stated interface and memory restrictions.

\noindent\textbf{Role in the proof.}
Use in this chapter. It fixes the proof context, imported mathematical language, or quantified module parameters for the following statements.

\paragraph{Context 39 (line 5230)}
\noindent\textbf{EasyCrypt name.}
\begin{lstlisting}[language=EasyCrypt,basicstyle=\ttfamily\scriptsize]
BudgetedProgrammedToSampledDirectBridge
\end{lstlisting}
\noindent\textbf{Checked EasyCrypt statement.}
\begin{lstlisting}[language=EasyCrypt,basicstyle=\ttfamily\scriptsize]
section BudgetedProgrammedToSampledDirectBridge.
\end{lstlisting}
\noindent\textbf{Formal mathematical reading.}
Mathematical context. This opens a scoped block of hypotheses, module parameters, and lemmas.

\noindent\textbf{Role in the proof.}
Use in this chapter. It fixes the proof context, imported mathematical language, or quantified module parameters for the following statements.

\paragraph{Context 40 (line 5232)}
\noindent\textbf{EasyCrypt name.}
\begin{lstlisting}[language=EasyCrypt,basicstyle=\ttfamily\scriptsize]
H
\end{lstlisting}
\noindent\textbf{Checked EasyCrypt statement.}
\begin{lstlisting}[language=EasyCrypt,basicstyle=\ttfamily\scriptsize]
declare module H <: Oracle {-CountedLazyROM,
                             -DirectSigningCoinSampler,
                             -EUF_CMA_ROMProgrammedTranscriptBudgetedCountedSign,
                             -EUF_CMA_ROMProgrammedTranscriptBudgetedSampledSign}.
\end{lstlisting}
\noindent\textbf{Formal mathematical reading.}
Mathematical context. This is universal quantification over an adversary, oracle, scheme, sampler, or reduction module satisfying the stated interface and memory restrictions.

\noindent\textbf{Role in the proof.}
Use in this chapter. It fixes the proof context, imported mathematical language, or quantified module parameters for the following statements.

\paragraph{Context 41 (line 5236)}
\noindent\textbf{EasyCrypt name.}
\begin{lstlisting}[language=EasyCrypt,basicstyle=\ttfamily\scriptsize]
A
\end{lstlisting}
\noindent\textbf{Checked EasyCrypt statement.}
\begin{lstlisting}[language=EasyCrypt,basicstyle=\ttfamily\scriptsize]
declare module A <: SIG.Adversary {-H, -CountedLazyROM,
                                   -DirectSigningCoinSampler,
                                   -EUF_CMA_ROMProgrammedTranscriptBudgetedCountedSign,
                                   -EUF_CMA_ROMProgrammedTranscriptBudgetedSampledSign}.
\end{lstlisting}
\noindent\textbf{Formal mathematical reading.}
Mathematical context. This is universal quantification over an adversary, oracle, scheme, sampler, or reduction module satisfying the stated interface and memory restrictions.

\noindent\textbf{Role in the proof.}
Use in this chapter. It fixes the proof context, imported mathematical language, or quantified module parameters for the following statements.

\paragraph{Context 42 (line 5601)}
\noindent\textbf{EasyCrypt name.}
\begin{lstlisting}[language=EasyCrypt,basicstyle=\ttfamily\scriptsize]
HAETAEROMInternalTranscriptExact
\end{lstlisting}
\noindent\textbf{Checked EasyCrypt statement.}
\begin{lstlisting}[language=EasyCrypt,basicstyle=\ttfamily\scriptsize]
section HAETAEROMInternalTranscriptExact.
\end{lstlisting}
\noindent\textbf{Formal mathematical reading.}
Mathematical context. This opens a scoped block of hypotheses, module parameters, and lemmas.

\noindent\textbf{Role in the proof.}
Use in this chapter. It fixes the proof context, imported mathematical language, or quantified module parameters for the following statements.

\paragraph{Context 43 (line 5603)}
\noindent\textbf{EasyCrypt name.}
\begin{lstlisting}[language=EasyCrypt,basicstyle=\ttfamily\scriptsize]
H
\end{lstlisting}
\noindent\textbf{Checked EasyCrypt statement.}
\begin{lstlisting}[language=EasyCrypt,basicstyle=\ttfamily\scriptsize]
declare module H <: SIG.Oracle {-SIG.EUF_CMA,
                                -EUF_CMA_ROMInternalTranscriptSign}.
\end{lstlisting}
\noindent\textbf{Formal mathematical reading.}
Mathematical context. This is universal quantification over an adversary, oracle, scheme, sampler, or reduction module satisfying the stated interface and memory restrictions.

\noindent\textbf{Role in the proof.}
Use in this chapter. It fixes the proof context, imported mathematical language, or quantified module parameters for the following statements.

\paragraph{Context 44 (line 5605)}
\noindent\textbf{EasyCrypt name.}
\begin{lstlisting}[language=EasyCrypt,basicstyle=\ttfamily\scriptsize]
A
\end{lstlisting}
\noindent\textbf{Checked EasyCrypt statement.}
\begin{lstlisting}[language=EasyCrypt,basicstyle=\ttfamily\scriptsize]
declare module A <: SIG.Adversary {-H, -SIG.EUF_CMA,
                                   -EUF_CMA_ROMInternalTranscriptSign}.
\end{lstlisting}
\noindent\textbf{Formal mathematical reading.}
Mathematical context. This is universal quantification over an adversary, oracle, scheme, sampler, or reduction module satisfying the stated interface and memory restrictions.

\noindent\textbf{Role in the proof.}
Use in this chapter. It fixes the proof context, imported mathematical language, or quantified module parameters for the following statements.

\paragraph{Context 45 (line 5706)}
\noindent\textbf{EasyCrypt name.}
\begin{lstlisting}[language=EasyCrypt,basicstyle=\ttfamily\scriptsize]
ROMInternalTranscriptStructuralNMA
\end{lstlisting}
\noindent\textbf{Checked EasyCrypt statement.}
\begin{lstlisting}[language=EasyCrypt,basicstyle=\ttfamily\scriptsize]
section ROMInternalTranscriptStructuralNMA.
\end{lstlisting}
\noindent\textbf{Formal mathematical reading.}
Mathematical context. This opens a scoped block of hypotheses, module parameters, and lemmas.

\noindent\textbf{Role in the proof.}
Use in this chapter. It fixes the proof context, imported mathematical language, or quantified module parameters for the following statements.

\paragraph{Context 46 (line 5708)}
\noindent\textbf{EasyCrypt name.}
\begin{lstlisting}[language=EasyCrypt,basicstyle=\ttfamily\scriptsize]
H
\end{lstlisting}
\noindent\textbf{Checked EasyCrypt statement.}
\begin{lstlisting}[language=EasyCrypt,basicstyle=\ttfamily\scriptsize]
declare module H <: SIG.Oracle {-EUF_CMA_ROMInternalTranscriptSign,
                                -ROMInternalTranscriptAsNMA}.
\end{lstlisting}
\noindent\textbf{Formal mathematical reading.}
Mathematical context. This is universal quantification over an adversary, oracle, scheme, sampler, or reduction module satisfying the stated interface and memory restrictions.

\noindent\textbf{Role in the proof.}
Use in this chapter. It fixes the proof context, imported mathematical language, or quantified module parameters for the following statements.

\paragraph{Context 47 (line 5710)}
\noindent\textbf{EasyCrypt name.}
\begin{lstlisting}[language=EasyCrypt,basicstyle=\ttfamily\scriptsize]
A
\end{lstlisting}
\noindent\textbf{Checked EasyCrypt statement.}
\begin{lstlisting}[language=EasyCrypt,basicstyle=\ttfamily\scriptsize]
declare module A <: SIG.Adversary {-H, -EUF_CMA_ROMInternalTranscriptSign,
                                   -ROMInternalTranscriptAsNMA}.
\end{lstlisting}
\noindent\textbf{Formal mathematical reading.}
Mathematical context. This is universal quantification over an adversary, oracle, scheme, sampler, or reduction module satisfying the stated interface and memory restrictions.

\noindent\textbf{Role in the proof.}
Use in this chapter. It fixes the proof context, imported mathematical language, or quantified module parameters for the following statements.

\paragraph{Context 48 (line 5874)}
\noindent\textbf{EasyCrypt name.}
\begin{lstlisting}[language=EasyCrypt,basicstyle=\ttfamily\scriptsize]
ROMInternalTranscriptPublicSimNMA
\end{lstlisting}
\noindent\textbf{Checked EasyCrypt statement.}
\begin{lstlisting}[language=EasyCrypt,basicstyle=\ttfamily\scriptsize]
section ROMInternalTranscriptPublicSimNMA.
\end{lstlisting}
\noindent\textbf{Formal mathematical reading.}
Mathematical context. This opens a scoped block of hypotheses, module parameters, and lemmas.

\noindent\textbf{Role in the proof.}
Use in this chapter. It fixes the proof context, imported mathematical language, or quantified module parameters for the following statements.

\paragraph{Context 49 (line 5876)}
\noindent\textbf{EasyCrypt name.}
\begin{lstlisting}[language=EasyCrypt,basicstyle=\ttfamily\scriptsize]
H
\end{lstlisting}
\noindent\textbf{Checked EasyCrypt statement.}
\begin{lstlisting}[language=EasyCrypt,basicstyle=\ttfamily\scriptsize]
declare module H <: SIG.Oracle {-ROMInternalTranscriptAsNMA,
                                -ROMInternalTranscriptPublicSimAsNMA}.
\end{lstlisting}
\noindent\textbf{Formal mathematical reading.}
Mathematical context. This is universal quantification over an adversary, oracle, scheme, sampler, or reduction module satisfying the stated interface and memory restrictions.

\noindent\textbf{Role in the proof.}
Use in this chapter. It fixes the proof context, imported mathematical language, or quantified module parameters for the following statements.

\paragraph{Context 50 (line 5878)}
\noindent\textbf{EasyCrypt name.}
\begin{lstlisting}[language=EasyCrypt,basicstyle=\ttfamily\scriptsize]
A
\end{lstlisting}
\noindent\textbf{Checked EasyCrypt statement.}
\begin{lstlisting}[language=EasyCrypt,basicstyle=\ttfamily\scriptsize]
declare module A <: SIG.Adversary {-H, -ROMInternalTranscriptAsNMA,
                                   -ROMInternalTranscriptPublicSimAsNMA}.
\end{lstlisting}
\noindent\textbf{Formal mathematical reading.}
Mathematical context. This is universal quantification over an adversary, oracle, scheme, sampler, or reduction module satisfying the stated interface and memory restrictions.

\noindent\textbf{Role in the proof.}
Use in this chapter. It fixes the proof context, imported mathematical language, or quantified module parameters for the following statements.

\paragraph{Context 51 (line 6041)}
\noindent\textbf{EasyCrypt name.}
\begin{lstlisting}[language=EasyCrypt,basicstyle=\ttfamily\scriptsize]
ROMInternalTranscriptROMPaperSimNMA
\end{lstlisting}
\noindent\textbf{Checked EasyCrypt statement.}
\begin{lstlisting}[language=EasyCrypt,basicstyle=\ttfamily\scriptsize]
section ROMInternalTranscriptROMPaperSimNMA.
\end{lstlisting}
\noindent\textbf{Formal mathematical reading.}
Mathematical context. This opens a scoped block of hypotheses, module parameters, and lemmas.

\noindent\textbf{Role in the proof.}
Use in this chapter. It fixes the proof context, imported mathematical language, or quantified module parameters for the following statements.

\paragraph{Context 52 (line 6043)}
\noindent\textbf{EasyCrypt name.}
\begin{lstlisting}[language=EasyCrypt,basicstyle=\ttfamily\scriptsize]
H
\end{lstlisting}
\noindent\textbf{Checked EasyCrypt statement.}
\begin{lstlisting}[language=EasyCrypt,basicstyle=\ttfamily\scriptsize]
declare module H <: SIG.Oracle {-ROMInternalTranscriptPublicSimAsNMA,
                                -ROMInternalTranscriptPaperSimAsNMA,
                                -ROMPaperSimSigningSampler}.
\end{lstlisting}
\noindent\textbf{Formal mathematical reading.}
Mathematical context. This is universal quantification over an adversary, oracle, scheme, sampler, or reduction module satisfying the stated interface and memory restrictions.

\noindent\textbf{Role in the proof.}
Use in this chapter. It fixes the proof context, imported mathematical language, or quantified module parameters for the following statements.

\paragraph{Context 53 (line 6046)}
\noindent\textbf{EasyCrypt name.}
\begin{lstlisting}[language=EasyCrypt,basicstyle=\ttfamily\scriptsize]
A
\end{lstlisting}
\noindent\textbf{Checked EasyCrypt statement.}
\begin{lstlisting}[language=EasyCrypt,basicstyle=\ttfamily\scriptsize]
declare module A <: SIG.Adversary {-H, -ROMInternalTranscriptPublicSimAsNMA,
                                   -ROMInternalTranscriptPaperSimAsNMA,
                                   -ROMPaperSimSigningSampler}.
\end{lstlisting}
\noindent\textbf{Formal mathematical reading.}
Mathematical context. This is universal quantification over an adversary, oracle, scheme, sampler, or reduction module satisfying the stated interface and memory restrictions.

\noindent\textbf{Role in the proof.}
Use in this chapter. It fixes the proof context, imported mathematical language, or quantified module parameters for the following statements.

\paragraph{Context 54 (line 6212)}
\noindent\textbf{EasyCrypt name.}
\begin{lstlisting}[language=EasyCrypt,basicstyle=\ttfamily\scriptsize]
ROMInternalTranscriptRealSigningPaperSimNMA
\end{lstlisting}
\noindent\textbf{Checked EasyCrypt statement.}
\begin{lstlisting}[language=EasyCrypt,basicstyle=\ttfamily\scriptsize]
section ROMInternalTranscriptRealSigningPaperSimNMA.
\end{lstlisting}
\noindent\textbf{Formal mathematical reading.}
Mathematical context. This opens a scoped block of hypotheses, module parameters, and lemmas.

\noindent\textbf{Role in the proof.}
Use in this chapter. It fixes the proof context, imported mathematical language, or quantified module parameters for the following statements.

\paragraph{Context 55 (line 6214)}
\noindent\textbf{EasyCrypt name.}
\begin{lstlisting}[language=EasyCrypt,basicstyle=\ttfamily\scriptsize]
H
\end{lstlisting}
\noindent\textbf{Checked EasyCrypt statement.}
\begin{lstlisting}[language=EasyCrypt,basicstyle=\ttfamily\scriptsize]
declare module H <: SIG.Oracle {-ROMInternalTranscriptAsNMA,
                                -ROMInternalTranscriptPaperSimAsNMA,
                                -RealSigningPaperSimSampler}.
\end{lstlisting}
\noindent\textbf{Formal mathematical reading.}
Mathematical context. This is universal quantification over an adversary, oracle, scheme, sampler, or reduction module satisfying the stated interface and memory restrictions.

\noindent\textbf{Role in the proof.}
Use in this chapter. It fixes the proof context, imported mathematical language, or quantified module parameters for the following statements.

\paragraph{Context 56 (line 6217)}
\noindent\textbf{EasyCrypt name.}
\begin{lstlisting}[language=EasyCrypt,basicstyle=\ttfamily\scriptsize]
A
\end{lstlisting}
\noindent\textbf{Checked EasyCrypt statement.}
\begin{lstlisting}[language=EasyCrypt,basicstyle=\ttfamily\scriptsize]
declare module A <: SIG.Adversary {-H, -ROMInternalTranscriptAsNMA,
                                   -ROMInternalTranscriptPaperSimAsNMA,
                                   -RealSigningPaperSimSampler}.
\end{lstlisting}
\noindent\textbf{Formal mathematical reading.}
Mathematical context. This is universal quantification over an adversary, oracle, scheme, sampler, or reduction module satisfying the stated interface and memory restrictions.

\noindent\textbf{Role in the proof.}
Use in this chapter. It fixes the proof context, imported mathematical language, or quantified module parameters for the following statements.

\paragraph{Context 57 (line 6418)}
\noindent\textbf{EasyCrypt name.}
\begin{lstlisting}[language=EasyCrypt,basicstyle=\ttfamily\scriptsize]
ROMInternalTranscriptROSigningAttemptPaperSimNMA
\end{lstlisting}
\noindent\textbf{Checked EasyCrypt statement.}
\begin{lstlisting}[language=EasyCrypt,basicstyle=\ttfamily\scriptsize]
section ROMInternalTranscriptROSigningAttemptPaperSimNMA.
\end{lstlisting}
\noindent\textbf{Formal mathematical reading.}
Mathematical context. This opens a scoped block of hypotheses, module parameters, and lemmas.

\noindent\textbf{Role in the proof.}
Use in this chapter. It fixes the proof context, imported mathematical language, or quantified module parameters for the following statements.

\paragraph{Context 58 (line 6420)}
\noindent\textbf{EasyCrypt name.}
\begin{lstlisting}[language=EasyCrypt,basicstyle=\ttfamily\scriptsize]
H
\end{lstlisting}
\noindent\textbf{Checked EasyCrypt statement.}
\begin{lstlisting}[language=EasyCrypt,basicstyle=\ttfamily\scriptsize]
declare module H <: SIG.Oracle {-ROMInternalTranscriptPaperSimAsNMA,
                                -RealSigningPaperSimSampler,
                                -ROSigningAttemptPaperSimSampler}.
\end{lstlisting}
\noindent\textbf{Formal mathematical reading.}
Mathematical context. This is universal quantification over an adversary, oracle, scheme, sampler, or reduction module satisfying the stated interface and memory restrictions.

\noindent\textbf{Role in the proof.}
Use in this chapter. It fixes the proof context, imported mathematical language, or quantified module parameters for the following statements.

\paragraph{Context 59 (line 6423)}
\noindent\textbf{EasyCrypt name.}
\begin{lstlisting}[language=EasyCrypt,basicstyle=\ttfamily\scriptsize]
A
\end{lstlisting}
\noindent\textbf{Checked EasyCrypt statement.}
\begin{lstlisting}[language=EasyCrypt,basicstyle=\ttfamily\scriptsize]
declare module A <: SIG.Adversary {-H, -ROMInternalTranscriptPaperSimAsNMA,
                                   -RealSigningPaperSimSampler,
                                   -ROSigningAttemptPaperSimSampler}.
\end{lstlisting}
\noindent\textbf{Formal mathematical reading.}
Mathematical context. This is universal quantification over an adversary, oracle, scheme, sampler, or reduction module satisfying the stated interface and memory restrictions.

\noindent\textbf{Role in the proof.}
Use in this chapter. It fixes the proof context, imported mathematical language, or quantified module parameters for the following statements.

\paragraph{Context 60 (line 7002)}
\noindent\textbf{EasyCrypt name.}
\begin{lstlisting}[language=EasyCrypt,basicstyle=\ttfamily\scriptsize]
ROMInternalTranscriptROSigningAttemptExactHyperballPaperSimNMA
\end{lstlisting}
\noindent\textbf{Checked EasyCrypt statement.}
\begin{lstlisting}[language=EasyCrypt,basicstyle=\ttfamily\scriptsize]
section ROMInternalTranscriptROSigningAttemptExactHyperballPaperSimNMA.
\end{lstlisting}
\noindent\textbf{Formal mathematical reading.}
Mathematical context. This opens a scoped block of hypotheses, module parameters, and lemmas.

\noindent\textbf{Role in the proof.}
Use in this chapter. It fixes the proof context, imported mathematical language, or quantified module parameters for the following statements.

\paragraph{Context 61 (line 7004)}
\noindent\textbf{EasyCrypt name.}
\begin{lstlisting}[language=EasyCrypt,basicstyle=\ttfamily\scriptsize]
H
\end{lstlisting}
\noindent\textbf{Checked EasyCrypt statement.}
\begin{lstlisting}[language=EasyCrypt,basicstyle=\ttfamily\scriptsize]
declare module H <: SIG.Oracle {-ROMInternalTranscriptPaperSimAsNMA,
                                -ROMInternalTranscriptBudgetedPaperSimAsNMA,
                                -ROSigningAttemptPaperSimSampler,
                                -ROExactHyperballPaperSimSampler}.
\end{lstlisting}
\noindent\textbf{Formal mathematical reading.}
Mathematical context. This is universal quantification over an adversary, oracle, scheme, sampler, or reduction module satisfying the stated interface and memory restrictions.

\noindent\textbf{Role in the proof.}
Use in this chapter. It fixes the proof context, imported mathematical language, or quantified module parameters for the following statements.

\paragraph{Context 62 (line 7008)}
\noindent\textbf{EasyCrypt name.}
\begin{lstlisting}[language=EasyCrypt,basicstyle=\ttfamily\scriptsize]
A
\end{lstlisting}
\noindent\textbf{Checked EasyCrypt statement.}
\begin{lstlisting}[language=EasyCrypt,basicstyle=\ttfamily\scriptsize]
declare module A <: SIG.Adversary {-H, -ROMInternalTranscriptPaperSimAsNMA,
                                   -ROMInternalTranscriptBudgetedPaperSimAsNMA,
                                   -ROSigningAttemptPaperSimSampler,
                                   -ROExactHyperballPaperSimSampler}.
\end{lstlisting}
\noindent\textbf{Formal mathematical reading.}
Mathematical context. This is universal quantification over an adversary, oracle, scheme, sampler, or reduction module satisfying the stated interface and memory restrictions.

\noindent\textbf{Role in the proof.}
Use in this chapter. It fixes the proof context, imported mathematical language, or quantified module parameters for the following statements.

\paragraph{Context 63 (line 7318)}
\noindent\textbf{EasyCrypt name.}
\begin{lstlisting}[language=EasyCrypt,basicstyle=\ttfamily\scriptsize]
ConcreteROMInternalTranscriptROSigningAttemptExactHyperballPaperSimNMA
\end{lstlisting}
\noindent\textbf{Checked EasyCrypt statement.}
\begin{lstlisting}[language=EasyCrypt,basicstyle=\ttfamily\scriptsize]
section ConcreteROMInternalTranscriptROSigningAttemptExactHyperballPaperSimNMA.
\end{lstlisting}
\noindent\textbf{Formal mathematical reading.}
Mathematical context. This opens a scoped block of hypotheses, module parameters, and lemmas.

\noindent\textbf{Role in the proof.}
Use in this chapter. It fixes the proof context, imported mathematical language, or quantified module parameters for the following statements.

\paragraph{Context 64 (line 7320)}
\noindent\textbf{EasyCrypt name.}
\begin{lstlisting}[language=EasyCrypt,basicstyle=\ttfamily\scriptsize]
A
\end{lstlisting}
\noindent\textbf{Checked EasyCrypt statement.}
\begin{lstlisting}[language=EasyCrypt,basicstyle=\ttfamily\scriptsize]
declare module A <: SIG.Adversary {-HAETAE_RO.FRO,
                                   -ROMInternalTranscriptPaperSimAsNMA,
                                   -ROMInternalTranscriptBudgetedPaperSimAsNMA,
                                   -ROSigningAttemptPaperSimSampler,
                                   -ROExactHyperballPaperSimSampler}.
\end{lstlisting}
\noindent\textbf{Formal mathematical reading.}
Mathematical context. This is universal quantification over an adversary, oracle, scheme, sampler, or reduction module satisfying the stated interface and memory restrictions.

\noindent\textbf{Role in the proof.}
Use in this chapter. It fixes the proof context, imported mathematical language, or quantified module parameters for the following statements.

\paragraph{Context 65 (line 7819)}
\noindent\textbf{EasyCrypt name.}
\begin{lstlisting}[language=EasyCrypt,basicstyle=\ttfamily\scriptsize]
ROMInternalTranscriptHAETAERejectionPaperSimNMA
\end{lstlisting}
\noindent\textbf{Checked EasyCrypt statement.}
\begin{lstlisting}[language=EasyCrypt,basicstyle=\ttfamily\scriptsize]
section ROMInternalTranscriptHAETAERejectionPaperSimNMA.
\end{lstlisting}
\noindent\textbf{Formal mathematical reading.}
Mathematical context. This opens a scoped block of hypotheses, module parameters, and lemmas.

\noindent\textbf{Role in the proof.}
Use in this chapter. It fixes the proof context, imported mathematical language, or quantified module parameters for the following statements.

\paragraph{Context 66 (line 7821)}
\noindent\textbf{EasyCrypt name.}
\begin{lstlisting}[language=EasyCrypt,basicstyle=\ttfamily\scriptsize]
H
\end{lstlisting}
\noindent\textbf{Checked EasyCrypt statement.}
\begin{lstlisting}[language=EasyCrypt,basicstyle=\ttfamily\scriptsize]
declare module H <: SIG.Oracle {-ROMInternalTranscriptPaperSimAsNMA,
                                 -RealSigningPaperSimSampler,
                                -HAETAERejectionPaperSimSampler}.
\end{lstlisting}
\noindent\textbf{Formal mathematical reading.}
Mathematical context. This is universal quantification over an adversary, oracle, scheme, sampler, or reduction module satisfying the stated interface and memory restrictions.

\noindent\textbf{Role in the proof.}
Use in this chapter. It fixes the proof context, imported mathematical language, or quantified module parameters for the following statements.

\paragraph{Context 67 (line 7824)}
\noindent\textbf{EasyCrypt name.}
\begin{lstlisting}[language=EasyCrypt,basicstyle=\ttfamily\scriptsize]
A
\end{lstlisting}
\noindent\textbf{Checked EasyCrypt statement.}
\begin{lstlisting}[language=EasyCrypt,basicstyle=\ttfamily\scriptsize]
declare module A <: SIG.Adversary {-H, -ROMInternalTranscriptPaperSimAsNMA,
                                   -RealSigningPaperSimSampler,
                                   -HAETAERejectionPaperSimSampler}.
\end{lstlisting}
\noindent\textbf{Formal mathematical reading.}
Mathematical context. This is universal quantification over an adversary, oracle, scheme, sampler, or reduction module satisfying the stated interface and memory restrictions.

\noindent\textbf{Role in the proof.}
Use in this chapter. It fixes the proof context, imported mathematical language, or quantified module parameters for the following statements.

\paragraph{Context 68 (line 8030)}
\noindent\textbf{EasyCrypt name.}
\begin{lstlisting}[language=EasyCrypt,basicstyle=\ttfamily\scriptsize]
ROMInternalTranscriptExactHyperballPaperSimNMA
\end{lstlisting}
\noindent\textbf{Checked EasyCrypt statement.}
\begin{lstlisting}[language=EasyCrypt,basicstyle=\ttfamily\scriptsize]
section ROMInternalTranscriptExactHyperballPaperSimNMA.
\end{lstlisting}
\noindent\textbf{Formal mathematical reading.}
Mathematical context. This opens a scoped block of hypotheses, module parameters, and lemmas.

\noindent\textbf{Role in the proof.}
Use in this chapter. It fixes the proof context, imported mathematical language, or quantified module parameters for the following statements.

\paragraph{Context 69 (line 8032)}
\noindent\textbf{EasyCrypt name.}
\begin{lstlisting}[language=EasyCrypt,basicstyle=\ttfamily\scriptsize]
H
\end{lstlisting}
\noindent\textbf{Checked EasyCrypt statement.}
\begin{lstlisting}[language=EasyCrypt,basicstyle=\ttfamily\scriptsize]
declare module H <: SIG.Oracle {-ROMInternalTranscriptPaperSimAsNMA,
                                 -ExactHyperballHAETAERejectionPaperSimSampler,
                                 -ExactHyperballPaperSimSampler}.
\end{lstlisting}
\noindent\textbf{Formal mathematical reading.}
Mathematical context. This is universal quantification over an adversary, oracle, scheme, sampler, or reduction module satisfying the stated interface and memory restrictions.

\noindent\textbf{Role in the proof.}
Use in this chapter. It fixes the proof context, imported mathematical language, or quantified module parameters for the following statements.

\paragraph{Context 70 (line 8035)}
\noindent\textbf{EasyCrypt name.}
\begin{lstlisting}[language=EasyCrypt,basicstyle=\ttfamily\scriptsize]
A
\end{lstlisting}
\noindent\textbf{Checked EasyCrypt statement.}
\begin{lstlisting}[language=EasyCrypt,basicstyle=\ttfamily\scriptsize]
declare module A <: SIG.Adversary {-H, -ROMInternalTranscriptPaperSimAsNMA,
                                   -ExactHyperballHAETAERejectionPaperSimSampler,
                                   -ExactHyperballPaperSimSampler}.
\end{lstlisting}
\noindent\textbf{Formal mathematical reading.}
Mathematical context. This is universal quantification over an adversary, oracle, scheme, sampler, or reduction module satisfying the stated interface and memory restrictions.

\noindent\textbf{Role in the proof.}
Use in this chapter. It fixes the proof context, imported mathematical language, or quantified module parameters for the following statements.

\paragraph{Context 71 (line 8232)}
\noindent\textbf{EasyCrypt name.}
\begin{lstlisting}[language=EasyCrypt,basicstyle=\ttfamily\scriptsize]
ROMInternalTranscriptPaperSimHop
\end{lstlisting}
\noindent\textbf{Checked EasyCrypt statement.}
\begin{lstlisting}[language=EasyCrypt,basicstyle=\ttfamily\scriptsize]
section ROMInternalTranscriptPaperSimHop.
\end{lstlisting}
\noindent\textbf{Formal mathematical reading.}
Mathematical context. This opens a scoped block of hypotheses, module parameters, and lemmas.

\noindent\textbf{Role in the proof.}
Use in this chapter. It fixes the proof context, imported mathematical language, or quantified module parameters for the following statements.

\paragraph{Context 72 (line 8234)}
\noindent\textbf{EasyCrypt name.}
\begin{lstlisting}[language=EasyCrypt,basicstyle=\ttfamily\scriptsize]
H
\end{lstlisting}
\noindent\textbf{Checked EasyCrypt statement.}
\begin{lstlisting}[language=EasyCrypt,basicstyle=\ttfamily\scriptsize]
declare module H <: SIG.Oracle {-ROMInternalTranscriptPublicSimAsNMA,
                                -ROMInternalTranscriptPaperSimAsNMA}.
\end{lstlisting}
\noindent\textbf{Formal mathematical reading.}
Mathematical context. This is universal quantification over an adversary, oracle, scheme, sampler, or reduction module satisfying the stated interface and memory restrictions.

\noindent\textbf{Role in the proof.}
Use in this chapter. It fixes the proof context, imported mathematical language, or quantified module parameters for the following statements.

\paragraph{Context 73 (line 8236)}
\noindent\textbf{EasyCrypt name.}
\begin{lstlisting}[language=EasyCrypt,basicstyle=\ttfamily\scriptsize]
A
\end{lstlisting}
\noindent\textbf{Checked EasyCrypt statement.}
\begin{lstlisting}[language=EasyCrypt,basicstyle=\ttfamily\scriptsize]
declare module A <: SIG.Adversary {-H, -ROMInternalTranscriptPublicSimAsNMA,
                                   -ROMInternalTranscriptPaperSimAsNMA}.
\end{lstlisting}
\noindent\textbf{Formal mathematical reading.}
Mathematical context. This is universal quantification over an adversary, oracle, scheme, sampler, or reduction module satisfying the stated interface and memory restrictions.

\noindent\textbf{Role in the proof.}
Use in this chapter. It fixes the proof context, imported mathematical language, or quantified module parameters for the following statements.

\paragraph{Context 74 (line 8238)}
\noindent\textbf{EasyCrypt name.}
\begin{lstlisting}[language=EasyCrypt,basicstyle=\ttfamily\scriptsize]
Samp
\end{lstlisting}
\noindent\textbf{Checked EasyCrypt statement.}
\begin{lstlisting}[language=EasyCrypt,basicstyle=\ttfamily\scriptsize]
declare module Samp <: PaperSimSigningSampler {-H, -A,
                                               -ROMInternalTranscriptPaperSimAsNMA}.
\end{lstlisting}
\noindent\textbf{Formal mathematical reading.}
Mathematical context. This is universal quantification over an adversary, oracle, scheme, sampler, or reduction module satisfying the stated interface and memory restrictions.

\noindent\textbf{Role in the proof.}
Use in this chapter. It fixes the proof context, imported mathematical language, or quantified module parameters for the following statements.

\paragraph{Context 75 (line 8256)}
\noindent\textbf{EasyCrypt name.}
\begin{lstlisting}[language=EasyCrypt,basicstyle=\ttfamily\scriptsize]
TranscriptLoggingErasure
\end{lstlisting}
\noindent\textbf{Checked EasyCrypt statement.}
\begin{lstlisting}[language=EasyCrypt,basicstyle=\ttfamily\scriptsize]
section TranscriptLoggingErasure.
\end{lstlisting}
\noindent\textbf{Formal mathematical reading.}
Mathematical context. This opens a scoped block of hypotheses, module parameters, and lemmas.

\noindent\textbf{Role in the proof.}
Use in this chapter. It fixes the proof context, imported mathematical language, or quantified module parameters for the following statements.

\paragraph{Context 76 (line 8258)}
\noindent\textbf{EasyCrypt name.}
\begin{lstlisting}[language=EasyCrypt,basicstyle=\ttfamily\scriptsize]
H
\end{lstlisting}
\noindent\textbf{Checked EasyCrypt statement.}
\begin{lstlisting}[language=EasyCrypt,basicstyle=\ttfamily\scriptsize]
declare module H <: SIG.Oracle {-EUF_CMA_SimulatedSign,
                                 -EUF_CMA_TranscriptSimulatedSign}.
\end{lstlisting}
\noindent\textbf{Formal mathematical reading.}
Mathematical context. This is universal quantification over an adversary, oracle, scheme, sampler, or reduction module satisfying the stated interface and memory restrictions.

\noindent\textbf{Role in the proof.}
Use in this chapter. It fixes the proof context, imported mathematical language, or quantified module parameters for the following statements.

\paragraph{Context 77 (line 8260)}
\noindent\textbf{EasyCrypt name.}
\begin{lstlisting}[language=EasyCrypt,basicstyle=\ttfamily\scriptsize]
S
\end{lstlisting}
\noindent\textbf{Checked EasyCrypt statement.}
\begin{lstlisting}[language=EasyCrypt,basicstyle=\ttfamily\scriptsize]
declare module S <: SIG.Scheme {-H, -EUF_CMA_SimulatedSign,
                                 -EUF_CMA_TranscriptSimulatedSign}.
\end{lstlisting}
\noindent\textbf{Formal mathematical reading.}
Mathematical context. This is universal quantification over an adversary, oracle, scheme, sampler, or reduction module satisfying the stated interface and memory restrictions.

\noindent\textbf{Role in the proof.}
Use in this chapter. It fixes the proof context, imported mathematical language, or quantified module parameters for the following statements.

\paragraph{Context 78 (line 8262)}
\noindent\textbf{EasyCrypt name.}
\begin{lstlisting}[language=EasyCrypt,basicstyle=\ttfamily\scriptsize]
A
\end{lstlisting}
\noindent\textbf{Checked EasyCrypt statement.}
\begin{lstlisting}[language=EasyCrypt,basicstyle=\ttfamily\scriptsize]
declare module A <: SIG.Adversary {-H, -S, -EUF_CMA_SimulatedSign,
                                    -EUF_CMA_TranscriptSimulatedSign}.
\end{lstlisting}
\noindent\textbf{Formal mathematical reading.}
Mathematical context. This is universal quantification over an adversary, oracle, scheme, sampler, or reduction module satisfying the stated interface and memory restrictions.

\noindent\textbf{Role in the proof.}
Use in this chapter. It fixes the proof context, imported mathematical language, or quantified module parameters for the following statements.

\paragraph{Context 79 (line 8264)}
\noindent\textbf{EasyCrypt name.}
\begin{lstlisting}[language=EasyCrypt,basicstyle=\ttfamily\scriptsize]
Sim
\end{lstlisting}
\noindent\textbf{Checked EasyCrypt statement.}
\begin{lstlisting}[language=EasyCrypt,basicstyle=\ttfamily\scriptsize]
declare module Sim <: SigningSimulator {-H, -S, -A,
                                         -EUF_CMA_SimulatedSign,
                                         -EUF_CMA_TranscriptSimulatedSign}.
\end{lstlisting}
\noindent\textbf{Formal mathematical reading.}
Mathematical context. This is universal quantification over an adversary, oracle, scheme, sampler, or reduction module satisfying the stated interface and memory restrictions.

\noindent\textbf{Role in the proof.}
Use in this chapter. It fixes the proof context, imported mathematical language, or quantified module parameters for the following statements.

\paragraph{Context 80 (line 8340)}
\noindent\textbf{EasyCrypt name.}
\begin{lstlisting}[language=EasyCrypt,basicstyle=\ttfamily\scriptsize]
RealSigningSimulatorExact
\end{lstlisting}
\noindent\textbf{Checked EasyCrypt statement.}
\begin{lstlisting}[language=EasyCrypt,basicstyle=\ttfamily\scriptsize]
section RealSigningSimulatorExact.
\end{lstlisting}
\noindent\textbf{Formal mathematical reading.}
Mathematical context. This opens a scoped block of hypotheses, module parameters, and lemmas.

\noindent\textbf{Role in the proof.}
Use in this chapter. It fixes the proof context, imported mathematical language, or quantified module parameters for the following statements.

\paragraph{Context 81 (line 8342)}
\noindent\textbf{EasyCrypt name.}
\begin{lstlisting}[language=EasyCrypt,basicstyle=\ttfamily\scriptsize]
H
\end{lstlisting}
\noindent\textbf{Checked EasyCrypt statement.}
\begin{lstlisting}[language=EasyCrypt,basicstyle=\ttfamily\scriptsize]
declare module H <: SIG.Oracle {-SIG.EUF_CMA, -EUF_CMA_SimulatedSign,
                                 -RealSigningSimulator}.
\end{lstlisting}
\noindent\textbf{Formal mathematical reading.}
Mathematical context. This is universal quantification over an adversary, oracle, scheme, sampler, or reduction module satisfying the stated interface and memory restrictions.

\noindent\textbf{Role in the proof.}
Use in this chapter. It fixes the proof context, imported mathematical language, or quantified module parameters for the following statements.

\paragraph{Context 82 (line 8344)}
\noindent\textbf{EasyCrypt name.}
\begin{lstlisting}[language=EasyCrypt,basicstyle=\ttfamily\scriptsize]
S
\end{lstlisting}
\noindent\textbf{Checked EasyCrypt statement.}
\begin{lstlisting}[language=EasyCrypt,basicstyle=\ttfamily\scriptsize]
declare module S <: SIG.Scheme {-H, -SIG.EUF_CMA, -EUF_CMA_SimulatedSign,
                                 -RealSigningSimulator}.
\end{lstlisting}
\noindent\textbf{Formal mathematical reading.}
Mathematical context. This is universal quantification over an adversary, oracle, scheme, sampler, or reduction module satisfying the stated interface and memory restrictions.

\noindent\textbf{Role in the proof.}
Use in this chapter. It fixes the proof context, imported mathematical language, or quantified module parameters for the following statements.

\paragraph{Context 83 (line 8346)}
\noindent\textbf{EasyCrypt name.}
\begin{lstlisting}[language=EasyCrypt,basicstyle=\ttfamily\scriptsize]
A
\end{lstlisting}
\noindent\textbf{Checked EasyCrypt statement.}
\begin{lstlisting}[language=EasyCrypt,basicstyle=\ttfamily\scriptsize]
declare module A <: SIG.Adversary {-H, -S, -SIG.EUF_CMA,
                                    -EUF_CMA_SimulatedSign,
                                    -RealSigningSimulator}.
\end{lstlisting}
\noindent\textbf{Formal mathematical reading.}
Mathematical context. This is universal quantification over an adversary, oracle, scheme, sampler, or reduction module satisfying the stated interface and memory restrictions.

\noindent\textbf{Role in the proof.}
Use in this chapter. It fixes the proof context, imported mathematical language, or quantified module parameters for the following statements.

\paragraph{Context 84 (line 8500)}
\noindent\textbf{EasyCrypt name.}
\begin{lstlisting}[language=EasyCrypt,basicstyle=\ttfamily\scriptsize]
BimodalSpecialSoundnessLift
\end{lstlisting}
\noindent\textbf{Checked EasyCrypt statement.}
\begin{lstlisting}[language=EasyCrypt,basicstyle=\ttfamily\scriptsize]
section BimodalSpecialSoundnessLift.
\end{lstlisting}
\noindent\textbf{Formal mathematical reading.}
Mathematical context. This opens a scoped block of hypotheses, module parameters, and lemmas.

\noindent\textbf{Role in the proof.}
Use in this chapter. It fixes the proof context, imported mathematical language, or quantified module parameters for the following statements.

\paragraph{Context 85 (line 8502)}
\noindent\textbf{EasyCrypt name.}
\begin{lstlisting}[language=EasyCrypt,basicstyle=\ttfamily\scriptsize]
H
\end{lstlisting}
\noindent\textbf{Checked EasyCrypt statement.}
\begin{lstlisting}[language=EasyCrypt,basicstyle=\ttfamily\scriptsize]
declare module H <: SIG.Oracle.
\end{lstlisting}
\noindent\textbf{Formal mathematical reading.}
Mathematical context. This is universal quantification over an adversary, oracle, scheme, sampler, or reduction module satisfying the stated interface and memory restrictions.

\noindent\textbf{Role in the proof.}
Use in this chapter. It fixes the proof context, imported mathematical language, or quantified module parameters for the following statements.

\paragraph{Context 86 (line 8503)}
\noindent\textbf{EasyCrypt name.}
\begin{lstlisting}[language=EasyCrypt,basicstyle=\ttfamily\scriptsize]
S
\end{lstlisting}
\noindent\textbf{Checked EasyCrypt statement.}
\begin{lstlisting}[language=EasyCrypt,basicstyle=\ttfamily\scriptsize]
declare module S <: SIG.Scheme {-H}.
\end{lstlisting}
\noindent\textbf{Formal mathematical reading.}
Mathematical context. This is universal quantification over an adversary, oracle, scheme, sampler, or reduction module satisfying the stated interface and memory restrictions.

\noindent\textbf{Role in the proof.}
Use in this chapter. It fixes the proof context, imported mathematical language, or quantified module parameters for the following statements.

\paragraph{Context 87 (line 8504)}
\noindent\textbf{EasyCrypt name.}
\begin{lstlisting}[language=EasyCrypt,basicstyle=\ttfamily\scriptsize]
F
\end{lstlisting}
\noindent\textbf{Checked EasyCrypt statement.}
\begin{lstlisting}[language=EasyCrypt,basicstyle=\ttfamily\scriptsize]
declare module F <: ForkingAdversary {-H, -S}.
\end{lstlisting}
\noindent\textbf{Formal mathematical reading.}
Mathematical context. This is universal quantification over an adversary, oracle, scheme, sampler, or reduction module satisfying the stated interface and memory restrictions.

\noindent\textbf{Role in the proof.}
Use in this chapter. It fixes the proof context, imported mathematical language, or quantified module parameters for the following statements.

\paragraph{Context 88 (line 8505)}
\noindent\textbf{EasyCrypt name.}
\begin{lstlisting}[language=EasyCrypt,basicstyle=\ttfamily\scriptsize]
X
\end{lstlisting}
\noindent\textbf{Checked EasyCrypt statement.}
\begin{lstlisting}[language=EasyCrypt,basicstyle=\ttfamily\scriptsize]
declare module X <: BimodalTranscriptExtractor {-H, -S, -F}.
\end{lstlisting}
\noindent\textbf{Formal mathematical reading.}
Mathematical context. This is universal quantification over an adversary, oracle, scheme, sampler, or reduction module satisfying the stated interface and memory restrictions.

\noindent\textbf{Role in the proof.}
Use in this chapter. It fixes the proof context, imported mathematical language, or quantified module parameters for the following statements.

\subsubsection*{Definitions, Carrier Sets, Operations, Games, and Procedures}
\begin{formaldefinition}[EasyCrypt line 23]
\noindent\textbf{EasyCrypt name.}
\begin{lstlisting}[language=EasyCrypt,basicstyle=\ttfamily\scriptsize]
SigningSimulator
\end{lstlisting}
\noindent\textbf{Checked EasyCrypt statement.}
\begin{lstlisting}[language=EasyCrypt,basicstyle=\ttfamily\scriptsize]
module type SigningSimulator(H : SIG.POracle) = {
\end{lstlisting}
\noindent\textbf{Formal mathematical reading.}
Mathematical definition. This is an interface for an oracle, adversary, scheme, sampler, solver, or reduction.  A later theorem quantified over this interface is universally quantified over all modules implementing these procedures.

\noindent\textbf{Role in the proof.}
Use in this chapter. It supplies an executable component of the formal probabilistic model.

\end{formaldefinition}

\begin{formaldefinition}[EasyCrypt line 24]
\noindent\textbf{EasyCrypt name.}
\begin{lstlisting}[language=EasyCrypt,basicstyle=\ttfamily\scriptsize]
init
\end{lstlisting}
\noindent\textbf{Checked EasyCrypt statement.}
\begin{lstlisting}[language=EasyCrypt,basicstyle=\ttfamily\scriptsize]
proc init(pk : pkey, sk : skey) : unit
  proc sign(m : message, ctx : context) : signature
}.
\end{lstlisting}
\noindent\textbf{Formal mathematical reading.}
Mathematical definition. This procedure is a command in the probabilistic program.  Its return value, state updates, and oracle calls are the operational content of the game.

\noindent\textbf{Role in the proof.}
Use in this chapter. It supplies an executable component of the formal probabilistic model.

\end{formaldefinition}

\begin{formaldefinition}[EasyCrypt line 28]
\noindent\textbf{EasyCrypt name.}
\begin{lstlisting}[language=EasyCrypt,basicstyle=\ttfamily\scriptsize]
RealSigningSimulator
\end{lstlisting}
\noindent\textbf{Checked EasyCrypt statement.}
\begin{lstlisting}[language=EasyCrypt,basicstyle=\ttfamily\scriptsize]
module RealSigningSimulator(S : SIG.Scheme) (H : SIG.POracle) = {
\end{lstlisting}
\noindent\textbf{Formal mathematical reading.}
Mathematical definition. This is a probabilistic program/game or a module adapter.  Its procedures induce the probability space used by the game-based proof.

\noindent\textbf{Role in the proof.}
Use in this chapter. It supplies an executable component of the formal probabilistic model.

\end{formaldefinition}

\begin{formaldefinition}[EasyCrypt line 31]
\noindent\textbf{EasyCrypt name.}
\begin{lstlisting}[language=EasyCrypt,basicstyle=\ttfamily\scriptsize]
init
\end{lstlisting}
\noindent\textbf{Checked EasyCrypt statement.}
\begin{lstlisting}[language=EasyCrypt,basicstyle=\ttfamily\scriptsize]
proc init(pk : pkey, sk0 : skey) : unit = {
\end{lstlisting}
\noindent\textbf{Formal mathematical reading.}
Mathematical definition. This procedure is a command in the probabilistic program.  Its return value, state updates, and oracle calls are the operational content of the game.

\noindent\textbf{Role in the proof.}
Use in this chapter. It supplies an executable component of the formal probabilistic model.

\end{formaldefinition}

\begin{formaldefinition}[EasyCrypt line 35]
\noindent\textbf{EasyCrypt name.}
\begin{lstlisting}[language=EasyCrypt,basicstyle=\ttfamily\scriptsize]
sign
\end{lstlisting}
\noindent\textbf{Checked EasyCrypt statement.}
\begin{lstlisting}[language=EasyCrypt,basicstyle=\ttfamily\scriptsize]
proc sign(m : message, ctx : context) : signature = {
\end{lstlisting}
\noindent\textbf{Formal mathematical reading.}
Mathematical definition. This procedure is a command in the probabilistic program.  Its return value, state updates, and oracle calls are the operational content of the game.

\noindent\textbf{Role in the proof.}
Use in this chapter. It supplies an executable component of the formal probabilistic model.

\end{formaldefinition}

\begin{formaldefinition}[EasyCrypt line 43]
\noindent\textbf{EasyCrypt name.}
\begin{lstlisting}[language=EasyCrypt,basicstyle=\ttfamily\scriptsize]
EUF_CMA_SimulatedSign
\end{lstlisting}
\noindent\textbf{Checked EasyCrypt statement.}
\begin{lstlisting}[language=EasyCrypt,basicstyle=\ttfamily\scriptsize]
module EUF_CMA_SimulatedSign(
  H : SIG.Oracle,
  S : SIG.Scheme,
  A : SIG.Adversary,
  Sim : SigningSimulator
) = {
\end{lstlisting}
\noindent\textbf{Formal mathematical reading.}
Mathematical definition. This is a probabilistic program/game or a module adapter.  Its procedures induce the probability space used by the game-based proof.

\noindent\textbf{Role in the proof.}
Use in this chapter. It defines the game or game interface whose success event is bounded by later theorems.

\end{formaldefinition}

\begin{formaldefinition}[EasyCrypt line 51]
\noindent\textbf{EasyCrypt name.}
\begin{lstlisting}[language=EasyCrypt,basicstyle=\ttfamily\scriptsize]
O
\end{lstlisting}
\noindent\textbf{Checked EasyCrypt statement.}
\begin{lstlisting}[language=EasyCrypt,basicstyle=\ttfamily\scriptsize]
module O = {
\end{lstlisting}
\noindent\textbf{Formal mathematical reading.}
Mathematical definition. This is a probabilistic program/game or a module adapter.  Its procedures induce the probability space used by the game-based proof.

\noindent\textbf{Role in the proof.}
Use in this chapter. It supplies an executable component of the formal probabilistic model.

\end{formaldefinition}

\begin{formaldefinition}[EasyCrypt line 52]
\noindent\textbf{EasyCrypt name.}
\begin{lstlisting}[language=EasyCrypt,basicstyle=\ttfamily\scriptsize]
sign
\end{lstlisting}
\noindent\textbf{Checked EasyCrypt statement.}
\begin{lstlisting}[language=EasyCrypt,basicstyle=\ttfamily\scriptsize]
proc sign(m : message, ctx : context) : signature = {
\end{lstlisting}
\noindent\textbf{Formal mathematical reading.}
Mathematical definition. This procedure is a command in the probabilistic program.  Its return value, state updates, and oracle calls are the operational content of the game.

\noindent\textbf{Role in the proof.}
Use in this chapter. It supplies an executable component of the formal probabilistic model.

\end{formaldefinition}

\begin{formaldefinition}[EasyCrypt line 61]
\noindent\textbf{EasyCrypt name.}
\begin{lstlisting}[language=EasyCrypt,basicstyle=\ttfamily\scriptsize]
A
\end{lstlisting}
\noindent\textbf{Checked EasyCrypt statement.}
\begin{lstlisting}[language=EasyCrypt,basicstyle=\ttfamily\scriptsize]
module A = A(H, O)

  proc main() : bool = {
\end{lstlisting}
\noindent\textbf{Formal mathematical reading.}
Mathematical definition. This is a probabilistic program/game or a module adapter.  Its procedures induce the probability space used by the game-based proof.

\noindent\textbf{Role in the proof.}
Use in this chapter. It supplies an executable component of the formal probabilistic model.

\end{formaldefinition}

\begin{formaldefinition}[EasyCrypt line 81]
\noindent\textbf{EasyCrypt name.}
\begin{lstlisting}[language=EasyCrypt,basicstyle=\ttfamily\scriptsize]
transcript_log_consistent
\end{lstlisting}
\noindent\textbf{Checked EasyCrypt statement.}
\begin{lstlisting}[language=EasyCrypt,basicstyle=\ttfamily\scriptsize]
op transcript_log_consistent
   (md : mode) (pk : pkey) (qs : SIG.query list)
   (trs : transcript list) (rs : signing_transcript_record list) : bool =
  qs = signing_record_queries rs /\
  trs = signing_record_transcripts rs /\
  signing_record_log_sound md pk rs.
\end{lstlisting}
\noindent\textbf{Formal mathematical reading.}
Mathematical definition. This is a total EasyCrypt operation.  The exact displayed statement gives the domain, codomain, and defining equation when present.

\noindent\textbf{Role in the proof.}
Use in this chapter. It defines transcript/extraction data for the Fiat-Shamir forking and special-soundness argument.

\end{formaldefinition}

\begin{formaldefinition}[EasyCrypt line 134]
\noindent\textbf{EasyCrypt name.}
\begin{lstlisting}[language=EasyCrypt,basicstyle=\ttfamily\scriptsize]
transcript_log_ready
\end{lstlisting}
\noindent\textbf{Checked EasyCrypt statement.}
\begin{lstlisting}[language=EasyCrypt,basicstyle=\ttfamily\scriptsize]
op transcript_log_ready (md : mode) (trs : transcript list) : bool =
  transcript_log_valid md trs /\
  transcript_log_min_entropy_clear md trs.
\end{lstlisting}
\noindent\textbf{Formal mathematical reading.}
Mathematical definition. This is a Boolean predicate.  The exact displayed EasyCrypt statement gives its arguments; mathematically it is an event, invariant, support condition, or well-formedness predicate over those arguments.

\noindent\textbf{Role in the proof.}
Use in this chapter. It defines transcript/extraction data for the Fiat-Shamir forking and special-soundness argument.

\end{formaldefinition}

\begin{formaldefinition}[EasyCrypt line 138]
\noindent\textbf{EasyCrypt name.}
\begin{lstlisting}[language=EasyCrypt,basicstyle=\ttfamily\scriptsize]
budgeted_programmed_query_count_discipline
\end{lstlisting}
\noindent\textbf{Checked EasyCrypt statement.}
\begin{lstlisting}[language=EasyCrypt,basicstyle=\ttfamily\scriptsize]
op budgeted_programmed_query_count_discipline
   (hash_count signing_count : int) : bool =
  0 <= hash_count /\
  hash_count <= hash_query_budget_count /\
  0 <= signing_count /\
  signing_count <= signature_query_budget_count.
\end{lstlisting}
\noindent\textbf{Formal mathematical reading.}
Mathematical definition. This is a Boolean predicate.  The exact displayed EasyCrypt statement gives its arguments; mathematically it is an event, invariant, support condition, or well-formedness predicate over those arguments.

\noindent\textbf{Role in the proof.}
Use in this chapter. It is part of the exact formal vocabulary used by subsequent lemmas and games.

\end{formaldefinition}

\begin{formaldefinition}[EasyCrypt line 145]
\noindent\textbf{EasyCrypt name.}
\begin{lstlisting}[language=EasyCrypt,basicstyle=\ttfamily\scriptsize]
budgeted_paper_sim_signing_count_discipline
\end{lstlisting}
\noindent\textbf{Checked EasyCrypt statement.}
\begin{lstlisting}[language=EasyCrypt,basicstyle=\ttfamily\scriptsize]
op budgeted_paper_sim_signing_count_discipline
   (signing_count : int) : bool =
  0 <= signing_count /\
  signing_count <= signature_query_budget_count.
\end{lstlisting}
\noindent\textbf{Formal mathematical reading.}
Mathematical definition. This is a Boolean predicate.  The exact displayed EasyCrypt statement gives its arguments; mathematically it is an event, invariant, support condition, or well-formedness predicate over those arguments.

\noindent\textbf{Role in the proof.}
Use in this chapter. It is part of the exact formal vocabulary used by subsequent lemmas and games.

\end{formaldefinition}

\begin{formaldefinition}[EasyCrypt line 150]
\noindent\textbf{EasyCrypt name.}
\begin{lstlisting}[language=EasyCrypt,basicstyle=\ttfamily\scriptsize]
budgeted_paper_sim_sampler_freshness_discipline
\end{lstlisting}
\noindent\textbf{Checked EasyCrypt statement.}
\begin{lstlisting}[language=EasyCrypt,basicstyle=\ttfamily\scriptsize]
op budgeted_paper_sim_sampler_freshness_discipline
   (hash_count signing_count : int)
   (hash_qs sampler_qs : ro_query list) : bool =
  0 <= hash_count /\
  hash_count <= hash_query_budget_count /\
  size hash_qs <= hash_count /\
  0 <= signing_count /\
  signing_count <= signature_query_budget_count /\
  size sampler_qs <= signing_count.
\end{lstlisting}
\noindent\textbf{Formal mathematical reading.}
Mathematical definition. This is a Boolean predicate.  The exact displayed EasyCrypt statement gives its arguments; mathematically it is an event, invariant, support condition, or well-formedness predicate over those arguments.

\noindent\textbf{Role in the proof.}
Use in this chapter. It contributes a sampler or sampler-derived object to the distributional model.

\end{formaldefinition}

\begin{formaldefinition}[EasyCrypt line 160]
\noindent\textbf{EasyCrypt name.}
\begin{lstlisting}[language=EasyCrypt,basicstyle=\ttfamily\scriptsize]
sampler_rom_covered
\end{lstlisting}
\noindent\textbf{Checked EasyCrypt statement.}
\begin{lstlisting}[language=EasyCrypt,basicstyle=\ttfamily\scriptsize]
op sampler_rom_covered
   (rom : (ro_query, ro_output * PROM.flag) fmap)
   (hash_qs sampler_qs : ro_query list) : bool =
  forall seed_coins,
    sampler_expand_query seed_coins \in rom =>
    sampler_expand_query seed_coins \in hash_qs \/
    sampler_expand_query seed_coins \in sampler_qs.
\end{lstlisting}
\noindent\textbf{Formal mathematical reading.}
Mathematical definition. This is a total EasyCrypt operation.  The exact displayed statement gives the domain, codomain, and defining equation when present.

\noindent\textbf{Role in the proof.}
Use in this chapter. It contributes a sampler or sampler-derived object to the distributional model.

\end{formaldefinition}

\begin{formaldefinition}[EasyCrypt line 205]
\noindent\textbf{EasyCrypt name.}
\begin{lstlisting}[language=EasyCrypt,basicstyle=\ttfamily\scriptsize]
budgeted_programmed_reprogram_discipline
\end{lstlisting}
\noindent\textbf{Checked EasyCrypt statement.}
\begin{lstlisting}[language=EasyCrypt,basicstyle=\ttfamily\scriptsize]
op budgeted_programmed_reprogram_discipline
   (pk : pkey) (sk : skey)
   (programmed : programming_site list) (bad_reprogram : bool) : bool =
  pk = public_key_of_secret haetae_mode sk /\
  programming_site_log_conflict_free_for_honest haetae_mode sk programmed /\
  ! bad_reprogram.
\end{lstlisting}
\noindent\textbf{Formal mathematical reading.}
Mathematical definition. This is a Boolean predicate.  The exact displayed EasyCrypt statement gives its arguments; mathematically it is an event, invariant, support condition, or well-formedness predicate over those arguments.

\noindent\textbf{Role in the proof.}
Use in this chapter. It is part of the exact formal vocabulary used by subsequent lemmas and games.

\end{formaldefinition}

\begin{formaldefinition}[EasyCrypt line 212]
\noindent\textbf{EasyCrypt name.}
\begin{lstlisting}[language=EasyCrypt,basicstyle=\ttfamily\scriptsize]
budgeted_programmed_challenge_query_discipline
\end{lstlisting}
\noindent\textbf{Checked EasyCrypt statement.}
\begin{lstlisting}[language=EasyCrypt,basicstyle=\ttfamily\scriptsize]
op budgeted_programmed_challenge_query_discipline
   (hash_count : int) (queries : ro_query list) : bool =
  0 <= hash_count /\
  hash_count <= hash_query_budget_count /\
  challenge_query_log_count queries <= hash_count + 1.
\end{lstlisting}
\noindent\textbf{Formal mathematical reading.}
Mathematical definition. This is a Boolean predicate.  The exact displayed EasyCrypt statement gives its arguments; mathematically it is an event, invariant, support condition, or well-formedness predicate over those arguments.

\noindent\textbf{Role in the proof.}
Use in this chapter. It is part of the exact formal vocabulary used by subsequent lemmas and games.

\end{formaldefinition}

\begin{formaldefinition}[EasyCrypt line 332]
\noindent\textbf{EasyCrypt name.}
\begin{lstlisting}[language=EasyCrypt,basicstyle=\ttfamily\scriptsize]
internal_transcript_state_sound
\end{lstlisting}
\noindent\textbf{Checked EasyCrypt statement.}
\begin{lstlisting}[language=EasyCrypt,basicstyle=\ttfamily\scriptsize]
op internal_transcript_state_sound
   (md : mode) (pk : pkey) (sk : skey) (qs : SIG.query list)
   (trs : transcript list) (rs : signing_transcript_record list) : bool =
  pk = public_key_of_secret md sk /\
  transcript_log_consistent md pk qs trs rs.
\end{lstlisting}
\noindent\textbf{Formal mathematical reading.}
Mathematical definition. This is a total EasyCrypt operation.  The exact displayed statement gives the domain, codomain, and defining equation when present.

\noindent\textbf{Role in the proof.}
Use in this chapter. It defines transcript/extraction data for the Fiat-Shamir forking and special-soundness argument.

\end{formaldefinition}

\begin{formaldefinition}[EasyCrypt line 442]
\noindent\textbf{EasyCrypt name.}
\begin{lstlisting}[language=EasyCrypt,basicstyle=\ttfamily\scriptsize]
EUF_CMA_TranscriptSimulatedSign
\end{lstlisting}
\noindent\textbf{Checked EasyCrypt statement.}
\begin{lstlisting}[language=EasyCrypt,basicstyle=\ttfamily\scriptsize]
module EUF_CMA_TranscriptSimulatedSign(
  H : SIG.Oracle,
  S : SIG.Scheme,
  A : SIG.Adversary,
  Sim : SigningSimulator
) = {
\end{lstlisting}
\noindent\textbf{Formal mathematical reading.}
Mathematical definition. This is a probabilistic program/game or a module adapter.  Its procedures induce the probability space used by the game-based proof.

\noindent\textbf{Role in the proof.}
Use in this chapter. It defines the game or game interface whose success event is bounded by later theorems.

\end{formaldefinition}

\begin{formaldefinition}[EasyCrypt line 453]
\noindent\textbf{EasyCrypt name.}
\begin{lstlisting}[language=EasyCrypt,basicstyle=\ttfamily\scriptsize]
O
\end{lstlisting}
\noindent\textbf{Checked EasyCrypt statement.}
\begin{lstlisting}[language=EasyCrypt,basicstyle=\ttfamily\scriptsize]
module O = {
\end{lstlisting}
\noindent\textbf{Formal mathematical reading.}
Mathematical definition. This is a probabilistic program/game or a module adapter.  Its procedures induce the probability space used by the game-based proof.

\noindent\textbf{Role in the proof.}
Use in this chapter. It supplies an executable component of the formal probabilistic model.

\end{formaldefinition}

\begin{formaldefinition}[EasyCrypt line 454]
\noindent\textbf{EasyCrypt name.}
\begin{lstlisting}[language=EasyCrypt,basicstyle=\ttfamily\scriptsize]
sign
\end{lstlisting}
\noindent\textbf{Checked EasyCrypt statement.}
\begin{lstlisting}[language=EasyCrypt,basicstyle=\ttfamily\scriptsize]
proc sign(m : message, ctx : context) : signature = {
\end{lstlisting}
\noindent\textbf{Formal mathematical reading.}
Mathematical definition. This procedure is a command in the probabilistic program.  Its return value, state updates, and oracle calls are the operational content of the game.

\noindent\textbf{Role in the proof.}
Use in this chapter. It supplies an executable component of the formal probabilistic model.

\end{formaldefinition}

\begin{formaldefinition}[EasyCrypt line 467]
\noindent\textbf{EasyCrypt name.}
\begin{lstlisting}[language=EasyCrypt,basicstyle=\ttfamily\scriptsize]
A
\end{lstlisting}
\noindent\textbf{Checked EasyCrypt statement.}
\begin{lstlisting}[language=EasyCrypt,basicstyle=\ttfamily\scriptsize]
module A = A(H, O)

  proc main() : bool = {
\end{lstlisting}
\noindent\textbf{Formal mathematical reading.}
Mathematical definition. This is a probabilistic program/game or a module adapter.  Its procedures induce the probability space used by the game-based proof.

\noindent\textbf{Role in the proof.}
Use in this chapter. It supplies an executable component of the formal probabilistic model.

\end{formaldefinition}

\begin{formaldefinition}[EasyCrypt line 490]
\noindent\textbf{EasyCrypt name.}
\begin{lstlisting}[language=EasyCrypt,basicstyle=\ttfamily\scriptsize]
EUF_CMA_InternalTranscriptSign
\end{lstlisting}
\noindent\textbf{Checked EasyCrypt statement.}
\begin{lstlisting}[language=EasyCrypt,basicstyle=\ttfamily\scriptsize]
module EUF_CMA_InternalTranscriptSign(
  H : SIG.Oracle,
  A : SIG.Adversary
) = {
\end{lstlisting}
\noindent\textbf{Formal mathematical reading.}
Mathematical definition. This is a probabilistic program/game or a module adapter.  Its procedures induce the probability space used by the game-based proof.

\noindent\textbf{Role in the proof.}
Use in this chapter. It defines the game or game interface whose success event is bounded by later theorems.

\end{formaldefinition}

\begin{formaldefinition}[EasyCrypt line 500]
\noindent\textbf{EasyCrypt name.}
\begin{lstlisting}[language=EasyCrypt,basicstyle=\ttfamily\scriptsize]
O
\end{lstlisting}
\noindent\textbf{Checked EasyCrypt statement.}
\begin{lstlisting}[language=EasyCrypt,basicstyle=\ttfamily\scriptsize]
module O = {
\end{lstlisting}
\noindent\textbf{Formal mathematical reading.}
Mathematical definition. This is a probabilistic program/game or a module adapter.  Its procedures induce the probability space used by the game-based proof.

\noindent\textbf{Role in the proof.}
Use in this chapter. It supplies an executable component of the formal probabilistic model.

\end{formaldefinition}

\begin{formaldefinition}[EasyCrypt line 501]
\noindent\textbf{EasyCrypt name.}
\begin{lstlisting}[language=EasyCrypt,basicstyle=\ttfamily\scriptsize]
sign
\end{lstlisting}
\noindent\textbf{Checked EasyCrypt statement.}
\begin{lstlisting}[language=EasyCrypt,basicstyle=\ttfamily\scriptsize]
proc sign(m : message, ctx : context) : signature = {
\end{lstlisting}
\noindent\textbf{Formal mathematical reading.}
Mathematical definition. This procedure is a command in the probabilistic program.  Its return value, state updates, and oracle calls are the operational content of the game.

\noindent\textbf{Role in the proof.}
Use in this chapter. It supplies an executable component of the formal probabilistic model.

\end{formaldefinition}

\begin{formaldefinition}[EasyCrypt line 517]
\noindent\textbf{EasyCrypt name.}
\begin{lstlisting}[language=EasyCrypt,basicstyle=\ttfamily\scriptsize]
A
\end{lstlisting}
\noindent\textbf{Checked EasyCrypt statement.}
\begin{lstlisting}[language=EasyCrypt,basicstyle=\ttfamily\scriptsize]
module A = A(H, O)

  proc main() : bool = {
\end{lstlisting}
\noindent\textbf{Formal mathematical reading.}
Mathematical definition. This is a probabilistic program/game or a module adapter.  Its procedures induce the probability space used by the game-based proof.

\noindent\textbf{Role in the proof.}
Use in this chapter. It supplies an executable component of the formal probabilistic model.

\end{formaldefinition}

\begin{formaldefinition}[EasyCrypt line 544]
\noindent\textbf{EasyCrypt name.}
\begin{lstlisting}[language=EasyCrypt,basicstyle=\ttfamily\scriptsize]
EUF_CMA_ROMInternalTranscriptSign
\end{lstlisting}
\noindent\textbf{Checked EasyCrypt statement.}
\begin{lstlisting}[language=EasyCrypt,basicstyle=\ttfamily\scriptsize]
module EUF_CMA_ROMInternalTranscriptSign(
  H : SIG.Oracle,
  A : SIG.Adversary
) = {
\end{lstlisting}
\noindent\textbf{Formal mathematical reading.}
Mathematical definition. This is a probabilistic program/game or a module adapter.  Its procedures induce the probability space used by the game-based proof.

\noindent\textbf{Role in the proof.}
Use in this chapter. It defines the game or game interface whose success event is bounded by later theorems.

\end{formaldefinition}

\begin{formaldefinition}[EasyCrypt line 554]
\noindent\textbf{EasyCrypt name.}
\begin{lstlisting}[language=EasyCrypt,basicstyle=\ttfamily\scriptsize]
O
\end{lstlisting}
\noindent\textbf{Checked EasyCrypt statement.}
\begin{lstlisting}[language=EasyCrypt,basicstyle=\ttfamily\scriptsize]
module O = {
\end{lstlisting}
\noindent\textbf{Formal mathematical reading.}
Mathematical definition. This is a probabilistic program/game or a module adapter.  Its procedures induce the probability space used by the game-based proof.

\noindent\textbf{Role in the proof.}
Use in this chapter. It supplies an executable component of the formal probabilistic model.

\end{formaldefinition}

\begin{formaldefinition}[EasyCrypt line 555]
\noindent\textbf{EasyCrypt name.}
\begin{lstlisting}[language=EasyCrypt,basicstyle=\ttfamily\scriptsize]
sign
\end{lstlisting}
\noindent\textbf{Checked EasyCrypt statement.}
\begin{lstlisting}[language=EasyCrypt,basicstyle=\ttfamily\scriptsize]
proc sign(m : message, ctx : context) : signature = {
\end{lstlisting}
\noindent\textbf{Formal mathematical reading.}
Mathematical definition. This procedure is a command in the probabilistic program.  Its return value, state updates, and oracle calls are the operational content of the game.

\noindent\textbf{Role in the proof.}
Use in this chapter. It supplies an executable component of the formal probabilistic model.

\end{formaldefinition}

\begin{formaldefinition}[EasyCrypt line 584]
\noindent\textbf{EasyCrypt name.}
\begin{lstlisting}[language=EasyCrypt,basicstyle=\ttfamily\scriptsize]
A
\end{lstlisting}
\noindent\textbf{Checked EasyCrypt statement.}
\begin{lstlisting}[language=EasyCrypt,basicstyle=\ttfamily\scriptsize]
module A = A(H, O)

  proc main() : bool = {
\end{lstlisting}
\noindent\textbf{Formal mathematical reading.}
Mathematical definition. This is a probabilistic program/game or a module adapter.  Its procedures induce the probability space used by the game-based proof.

\noindent\textbf{Role in the proof.}
Use in this chapter. It supplies an executable component of the formal probabilistic model.

\end{formaldefinition}

\begin{formaldefinition}[EasyCrypt line 613]
\noindent\textbf{EasyCrypt name.}
\begin{lstlisting}[language=EasyCrypt,basicstyle=\ttfamily\scriptsize]
ROMInternalTranscriptAsNMA
\end{lstlisting}
\noindent\textbf{Checked EasyCrypt statement.}
\begin{lstlisting}[language=EasyCrypt,basicstyle=\ttfamily\scriptsize]
module ROMInternalTranscriptAsNMA(A : SIG.Adversary) (H : SIG.POracle) = {
\end{lstlisting}
\noindent\textbf{Formal mathematical reading.}
Mathematical definition. This is a probabilistic program/game or a module adapter.  Its procedures induce the probability space used by the game-based proof.

\noindent\textbf{Role in the proof.}
Use in this chapter. It defines the game or game interface whose success event is bounded by later theorems.

\end{formaldefinition}

\begin{formaldefinition}[EasyCrypt line 620]
\noindent\textbf{EasyCrypt name.}
\begin{lstlisting}[language=EasyCrypt,basicstyle=\ttfamily\scriptsize]
O
\end{lstlisting}
\noindent\textbf{Checked EasyCrypt statement.}
\begin{lstlisting}[language=EasyCrypt,basicstyle=\ttfamily\scriptsize]
module O = {
\end{lstlisting}
\noindent\textbf{Formal mathematical reading.}
Mathematical definition. This is a probabilistic program/game or a module adapter.  Its procedures induce the probability space used by the game-based proof.

\noindent\textbf{Role in the proof.}
Use in this chapter. It supplies an executable component of the formal probabilistic model.

\end{formaldefinition}

\begin{formaldefinition}[EasyCrypt line 621]
\noindent\textbf{EasyCrypt name.}
\begin{lstlisting}[language=EasyCrypt,basicstyle=\ttfamily\scriptsize]
sign
\end{lstlisting}
\noindent\textbf{Checked EasyCrypt statement.}
\begin{lstlisting}[language=EasyCrypt,basicstyle=\ttfamily\scriptsize]
proc sign(m : message, ctx : context) : signature = {
\end{lstlisting}
\noindent\textbf{Formal mathematical reading.}
Mathematical definition. This procedure is a command in the probabilistic program.  Its return value, state updates, and oracle calls are the operational content of the game.

\noindent\textbf{Role in the proof.}
Use in this chapter. It supplies an executable component of the formal probabilistic model.

\end{formaldefinition}

\begin{formaldefinition}[EasyCrypt line 650]
\noindent\textbf{EasyCrypt name.}
\begin{lstlisting}[language=EasyCrypt,basicstyle=\ttfamily\scriptsize]
A
\end{lstlisting}
\noindent\textbf{Checked EasyCrypt statement.}
\begin{lstlisting}[language=EasyCrypt,basicstyle=\ttfamily\scriptsize]
module A = A(H, O)

  proc forge(pk : pkey) : message * context * signature = {
\end{lstlisting}
\noindent\textbf{Formal mathematical reading.}
Mathematical definition. This is a probabilistic program/game or a module adapter.  Its procedures induce the probability space used by the game-based proof.

\noindent\textbf{Role in the proof.}
Use in this chapter. It supplies an executable component of the formal probabilistic model.

\end{formaldefinition}

\begin{formaldefinition}[EasyCrypt line 665]
\noindent\textbf{EasyCrypt name.}
\begin{lstlisting}[language=EasyCrypt,basicstyle=\ttfamily\scriptsize]
public_sim_source
\end{lstlisting}
\noindent\textbf{Checked EasyCrypt statement.}
\begin{lstlisting}[language=EasyCrypt,basicstyle=\ttfamily\scriptsize]
op public_sim_source
   (pk : pkey) (m : message) (ctx : context)
   (coins : random_coins) : byte list =
  coins ++ pk.`1 ++ ctx ++ m.
\end{lstlisting}
\noindent\textbf{Formal mathematical reading.}
Mathematical definition. This is a total EasyCrypt operation.  The exact displayed statement gives the domain, codomain, and defining equation when present.

\noindent\textbf{Role in the proof.}
Use in this chapter. It is part of the exact formal vocabulary used by subsequent lemmas and games.

\end{formaldefinition}

\begin{formaldefinition}[EasyCrypt line 670]
\noindent\textbf{EasyCrypt name.}
\begin{lstlisting}[language=EasyCrypt,basicstyle=\ttfamily\scriptsize]
public_sim_response_aux_vector
\end{lstlisting}
\noindent\textbf{Checked EasyCrypt statement.}
\begin{lstlisting}[language=EasyCrypt,basicstyle=\ttfamily\scriptsize]
op public_sim_response_aux_vector :
  mode -> pkey -> message -> context -> random_coins -> polyveck =
  fun md pk m ctx coins =>
    deterministic_unit_polyveck md 29
      (public_sim_source pk m ctx coins).
\end{lstlisting}
\noindent\textbf{Formal mathematical reading.}
Mathematical definition. This is a total EasyCrypt operation.  The exact displayed statement gives the domain, codomain, and defining equation when present.

\noindent\textbf{Role in the proof.}
Use in this chapter. It is part of the exact formal vocabulary used by subsequent lemmas and games.

\end{formaldefinition}

\begin{formaldefinition}[EasyCrypt line 676]
\noindent\textbf{EasyCrypt name.}
\begin{lstlisting}[language=EasyCrypt,basicstyle=\ttfamily\scriptsize]
public_sim_response_vector
\end{lstlisting}
\noindent\textbf{Checked EasyCrypt statement.}
\begin{lstlisting}[language=EasyCrypt,basicstyle=\ttfamily\scriptsize]
op public_sim_response_vector :
  mode -> pkey -> message -> context -> random_coins -> polyvecl =
  fun md pk m ctx coins =>
    deterministic_unit_polyvecl md 17
      (public_sim_source pk m ctx coins).
\end{lstlisting}
\noindent\textbf{Formal mathematical reading.}
Mathematical definition. This is a total EasyCrypt operation.  The exact displayed statement gives the domain, codomain, and defining equation when present.

\noindent\textbf{Role in the proof.}
Use in this chapter. It is part of the exact formal vocabulary used by subsequent lemmas and games.

\end{formaldefinition}

\begin{formaldefinition}[EasyCrypt line 682]
\noindent\textbf{EasyCrypt name.}
\begin{lstlisting}[language=EasyCrypt,basicstyle=\ttfamily\scriptsize]
public_sim_hint_vector
\end{lstlisting}
\noindent\textbf{Checked EasyCrypt statement.}
\begin{lstlisting}[language=EasyCrypt,basicstyle=\ttfamily\scriptsize]
op public_sim_hint_vector :
  mode -> pkey -> message -> context -> random_coins -> hint_t =
  fun md pk m ctx coins =>
    deterministic_unit_polyveck md 41
      (public_sim_source pk m ctx coins).
\end{lstlisting}
\noindent\textbf{Formal mathematical reading.}
Mathematical definition. This is a total EasyCrypt operation.  The exact displayed statement gives the domain, codomain, and defining equation when present.

\noindent\textbf{Role in the proof.}
Use in this chapter. It is part of the exact formal vocabulary used by subsequent lemmas and games.

\end{formaldefinition}

\begin{formaldefinition}[EasyCrypt line 688]
\noindent\textbf{EasyCrypt name.}
\begin{lstlisting}[language=EasyCrypt,basicstyle=\ttfamily\scriptsize]
public_sim_commitment_lowbits
\end{lstlisting}
\noindent\textbf{Checked EasyCrypt statement.}
\begin{lstlisting}[language=EasyCrypt,basicstyle=\ttfamily\scriptsize]
op public_sim_commitment_lowbits :
  mode -> pkey -> message -> context -> random_coins -> poly =
  fun md pk m ctx coins =>
    let raw = deterministic_polyveck md
      (public_sim_source pk m ctx coins) in
    let row0 = nth poly_zero raw 0 in
    mkseq
      (fun i =>
        if i = 0 then signing_entropy_token_of_coins coins
        else poly_coeff row0 i)
      n.
\end{lstlisting}
\noindent\textbf{Formal mathematical reading.}
Mathematical definition. This is a total EasyCrypt operation.  The exact displayed statement gives the domain, codomain, and defining equation when present.

\noindent\textbf{Role in the proof.}
Use in this chapter. It is part of the exact formal vocabulary used by subsequent lemmas and games.

\end{formaldefinition}

\begin{formaldefinition}[EasyCrypt line 700]
\noindent\textbf{EasyCrypt name.}
\begin{lstlisting}[language=EasyCrypt,basicstyle=\ttfamily\scriptsize]
public_sim_commitment_challenge
\end{lstlisting}
\noindent\textbf{Checked EasyCrypt statement.}
\begin{lstlisting}[language=EasyCrypt,basicstyle=\ttfamily\scriptsize]
op public_sim_commitment_challenge :
  mode -> pkey -> message -> context -> random_coins -> challenge =
  fun md pk m ctx coins =>
    challenge_hash md
      (public_sim_response_aux_vector md pk m ctx coins)
      (public_sim_commitment_lowbits md pk m ctx coins)
      (message_hash pk ctx m).
\end{lstlisting}
\noindent\textbf{Formal mathematical reading.}
Mathematical definition. This is a total EasyCrypt operation.  The exact displayed statement gives the domain, codomain, and defining equation when present.

\noindent\textbf{Role in the proof.}
Use in this chapter. It is part of the exact formal vocabulary used by subsequent lemmas and games.

\end{formaldefinition}

\begin{formaldefinition}[EasyCrypt line 708]
\noindent\textbf{EasyCrypt name.}
\begin{lstlisting}[language=EasyCrypt,basicstyle=\ttfamily\scriptsize]
public_sim_commitment_highbits
\end{lstlisting}
\noindent\textbf{Checked EasyCrypt statement.}
\begin{lstlisting}[language=EasyCrypt,basicstyle=\ttfamily\scriptsize]
op public_sim_commitment_highbits :
  mode -> pkey -> message -> context -> random_coins -> polyveck =
  fun md pk m ctx coins =>
    let ch = public_sim_commitment_challenge md pk m ctx coins in
    reconstructed_highbits md
      (public_sim_response_aux_vector md pk m ctx coins)
      (public_key_challenge_term md pk ch).
\end{lstlisting}
\noindent\textbf{Formal mathematical reading.}
Mathematical definition. This is a total EasyCrypt operation.  The exact displayed statement gives the domain, codomain, and defining equation when present.

\noindent\textbf{Role in the proof.}
Use in this chapter. It is part of the exact formal vocabulary used by subsequent lemmas and games.

\end{formaldefinition}

\begin{formaldefinition}[EasyCrypt line 716]
\noindent\textbf{EasyCrypt name.}
\begin{lstlisting}[language=EasyCrypt,basicstyle=\ttfamily\scriptsize]
public_sim_signature_token
\end{lstlisting}
\noindent\textbf{Checked EasyCrypt statement.}
\begin{lstlisting}[language=EasyCrypt,basicstyle=\ttfamily\scriptsize]
op public_sim_signature_token :
  mode -> pkey -> message -> context -> random_coins -> sig_token =
  fun md pk m ctx coins =>
    ( public_sim_response_vector md pk m ctx coins,
      public_sim_response_aux_vector md pk m ctx coins,
      public_sim_hint_vector md pk m ctx coins).
\end{lstlisting}
\noindent\textbf{Formal mathematical reading.}
Mathematical definition. This is a total EasyCrypt operation.  The exact displayed statement gives the domain, codomain, and defining equation when present.

\noindent\textbf{Role in the proof.}
Use in this chapter. It names a well-formedness, support, or validity predicate needed by later preservation lemmas.

\end{formaldefinition}

\begin{formaldefinition}[EasyCrypt line 723]
\noindent\textbf{EasyCrypt name.}
\begin{lstlisting}[language=EasyCrypt,basicstyle=\ttfamily\scriptsize]
public_sim_signature
\end{lstlisting}
\noindent\textbf{Checked EasyCrypt statement.}
\begin{lstlisting}[language=EasyCrypt,basicstyle=\ttfamily\scriptsize]
op public_sim_signature :
  mode -> pkey -> message -> context -> random_coins -> signature =
  fun md pk m ctx coins =>
    let highbits = public_sim_commitment_highbits md pk m ctx coins in
    let lowbits = public_sim_commitment_lowbits md pk m ctx coins in
    let mu = message_hash pk ctx m in
    let ch = public_sim_commitment_challenge md pk m ctx coins in
    (highbits, lowbits, mu, ch,
     public_sim_signature_token md pk m ctx coins, 0, 0).
\end{lstlisting}
\noindent\textbf{Formal mathematical reading.}
Mathematical definition. This is a total EasyCrypt operation.  The exact displayed statement gives the domain, codomain, and defining equation when present.

\noindent\textbf{Role in the proof.}
Use in this chapter. It is part of the exact formal vocabulary used by subsequent lemmas and games.

\end{formaldefinition}

\begin{formaldefinition}[EasyCrypt line 804]
\noindent\textbf{EasyCrypt name.}
\begin{lstlisting}[language=EasyCrypt,basicstyle=\ttfamily\scriptsize]
ROMInternalTranscriptPublicSimAsNMA
\end{lstlisting}
\noindent\textbf{Checked EasyCrypt statement.}
\begin{lstlisting}[language=EasyCrypt,basicstyle=\ttfamily\scriptsize]
module ROMInternalTranscriptPublicSimAsNMA(A : SIG.Adversary)
                                         (H : SIG.POracle) = {
\end{lstlisting}
\noindent\textbf{Formal mathematical reading.}
Mathematical definition. This is a probabilistic program/game or a module adapter.  Its procedures induce the probability space used by the game-based proof.

\noindent\textbf{Role in the proof.}
Use in this chapter. It defines the game or game interface whose success event is bounded by later theorems.

\end{formaldefinition}

\begin{formaldefinition}[EasyCrypt line 811]
\noindent\textbf{EasyCrypt name.}
\begin{lstlisting}[language=EasyCrypt,basicstyle=\ttfamily\scriptsize]
O
\end{lstlisting}
\noindent\textbf{Checked EasyCrypt statement.}
\begin{lstlisting}[language=EasyCrypt,basicstyle=\ttfamily\scriptsize]
module O = {
\end{lstlisting}
\noindent\textbf{Formal mathematical reading.}
Mathematical definition. This is a probabilistic program/game or a module adapter.  Its procedures induce the probability space used by the game-based proof.

\noindent\textbf{Role in the proof.}
Use in this chapter. It supplies an executable component of the formal probabilistic model.

\end{formaldefinition}

\begin{formaldefinition}[EasyCrypt line 812]
\noindent\textbf{EasyCrypt name.}
\begin{lstlisting}[language=EasyCrypt,basicstyle=\ttfamily\scriptsize]
sign
\end{lstlisting}
\noindent\textbf{Checked EasyCrypt statement.}
\begin{lstlisting}[language=EasyCrypt,basicstyle=\ttfamily\scriptsize]
proc sign(m : message, ctx : context) : signature = {
\end{lstlisting}
\noindent\textbf{Formal mathematical reading.}
Mathematical definition. This procedure is a command in the probabilistic program.  Its return value, state updates, and oracle calls are the operational content of the game.

\noindent\textbf{Role in the proof.}
Use in this chapter. It supplies an executable component of the formal probabilistic model.

\end{formaldefinition}

\begin{formaldefinition}[EasyCrypt line 841]
\noindent\textbf{EasyCrypt name.}
\begin{lstlisting}[language=EasyCrypt,basicstyle=\ttfamily\scriptsize]
A
\end{lstlisting}
\noindent\textbf{Checked EasyCrypt statement.}
\begin{lstlisting}[language=EasyCrypt,basicstyle=\ttfamily\scriptsize]
module A = A(H, O)

  proc forge(pk : pkey) : message * context * signature = {
\end{lstlisting}
\noindent\textbf{Formal mathematical reading.}
Mathematical definition. This is a probabilistic program/game or a module adapter.  Its procedures induce the probability space used by the game-based proof.

\noindent\textbf{Role in the proof.}
Use in this chapter. It supplies an executable component of the formal probabilistic model.

\end{formaldefinition}

\begin{formaldefinition}[EasyCrypt line 855]
\noindent\textbf{EasyCrypt name.}
\begin{lstlisting}[language=EasyCrypt,basicstyle=\ttfamily\scriptsize]
paper_sim_signature_sample
\end{lstlisting}
\noindent\textbf{Checked EasyCrypt statement.}
\begin{lstlisting}[language=EasyCrypt,basicstyle=\ttfamily\scriptsize]
type paper_sim_signature_sample = sig_token * poly * int.
\end{lstlisting}
\noindent\textbf{Formal mathematical reading.}
Mathematical definition. This is a definitional type abbreviation.  The exact displayed EasyCrypt statement gives the right-hand side; later proof steps may unfold it without adding an assumption.

\noindent\textbf{Role in the proof.}
Use in this chapter. It contributes a sampler or sampler-derived object to the distributional model.

\end{formaldefinition}

\begin{formaldefinition}[EasyCrypt line 857]
\noindent\textbf{EasyCrypt name.}
\begin{lstlisting}[language=EasyCrypt,basicstyle=\ttfamily\scriptsize]
paper_sim_sample_token
\end{lstlisting}
\noindent\textbf{Checked EasyCrypt statement.}
\begin{lstlisting}[language=EasyCrypt,basicstyle=\ttfamily\scriptsize]
op paper_sim_sample_token (s : paper_sim_signature_sample) : sig_token =
  s.`1.
\end{lstlisting}
\noindent\textbf{Formal mathematical reading.}
Mathematical definition. This is a total EasyCrypt operation.  The exact displayed statement gives the domain, codomain, and defining equation when present.

\noindent\textbf{Role in the proof.}
Use in this chapter. It names a well-formedness, support, or validity predicate needed by later preservation lemmas.

\end{formaldefinition}

\begin{formaldefinition}[EasyCrypt line 859]
\noindent\textbf{EasyCrypt name.}
\begin{lstlisting}[language=EasyCrypt,basicstyle=\ttfamily\scriptsize]
paper_sim_sample_lowbits_carrier
\end{lstlisting}
\noindent\textbf{Checked EasyCrypt statement.}
\begin{lstlisting}[language=EasyCrypt,basicstyle=\ttfamily\scriptsize]
op paper_sim_sample_lowbits_carrier
   (s : paper_sim_signature_sample) : poly =
  s.`2.
\end{lstlisting}
\noindent\textbf{Formal mathematical reading.}
Mathematical definition. This is a total EasyCrypt operation.  The exact displayed statement gives the domain, codomain, and defining equation when present.

\noindent\textbf{Role in the proof.}
Use in this chapter. It contributes a sampler or sampler-derived object to the distributional model.

\end{formaldefinition}

\begin{formaldefinition}[EasyCrypt line 862]
\noindent\textbf{EasyCrypt name.}
\begin{lstlisting}[language=EasyCrypt,basicstyle=\ttfamily\scriptsize]
paper_sim_sample_entropy
\end{lstlisting}
\noindent\textbf{Checked EasyCrypt statement.}
\begin{lstlisting}[language=EasyCrypt,basicstyle=\ttfamily\scriptsize]
op paper_sim_sample_entropy (s : paper_sim_signature_sample) : int =
  s.`3.
\end{lstlisting}
\noindent\textbf{Formal mathematical reading.}
Mathematical definition. This is a total EasyCrypt operation.  The exact displayed statement gives the domain, codomain, and defining equation when present.

\noindent\textbf{Role in the proof.}
Use in this chapter. It contributes a sampler or sampler-derived object to the distributional model.

\end{formaldefinition}

\begin{formaldefinition}[EasyCrypt line 864]
\noindent\textbf{EasyCrypt name.}
\begin{lstlisting}[language=EasyCrypt,basicstyle=\ttfamily\scriptsize]
paper_sim_sample_response
\end{lstlisting}
\noindent\textbf{Checked EasyCrypt statement.}
\begin{lstlisting}[language=EasyCrypt,basicstyle=\ttfamily\scriptsize]
op paper_sim_sample_response (s : paper_sim_signature_sample) : polyvecl =
  (paper_sim_sample_token s).`1.
\end{lstlisting}
\noindent\textbf{Formal mathematical reading.}
Mathematical definition. This is a total EasyCrypt operation.  The exact displayed statement gives the domain, codomain, and defining equation when present.

\noindent\textbf{Role in the proof.}
Use in this chapter. It contributes a sampler or sampler-derived object to the distributional model.

\end{formaldefinition}

\begin{formaldefinition}[EasyCrypt line 866]
\noindent\textbf{EasyCrypt name.}
\begin{lstlisting}[language=EasyCrypt,basicstyle=\ttfamily\scriptsize]
paper_sim_sample_response_aux
\end{lstlisting}
\noindent\textbf{Checked EasyCrypt statement.}
\begin{lstlisting}[language=EasyCrypt,basicstyle=\ttfamily\scriptsize]
op paper_sim_sample_response_aux
   (s : paper_sim_signature_sample) : polyveck =
  (paper_sim_sample_token s).`2.
\end{lstlisting}
\noindent\textbf{Formal mathematical reading.}
Mathematical definition. This is a total EasyCrypt operation.  The exact displayed statement gives the domain, codomain, and defining equation when present.

\noindent\textbf{Role in the proof.}
Use in this chapter. It contributes a sampler or sampler-derived object to the distributional model.

\end{formaldefinition}

\begin{formaldefinition}[EasyCrypt line 869]
\noindent\textbf{EasyCrypt name.}
\begin{lstlisting}[language=EasyCrypt,basicstyle=\ttfamily\scriptsize]
paper_sim_sample_hint
\end{lstlisting}
\noindent\textbf{Checked EasyCrypt statement.}
\begin{lstlisting}[language=EasyCrypt,basicstyle=\ttfamily\scriptsize]
op paper_sim_sample_hint (s : paper_sim_signature_sample) : hint_t =
  (paper_sim_sample_token s).`3.
\end{lstlisting}
\noindent\textbf{Formal mathematical reading.}
Mathematical definition. This is a total EasyCrypt operation.  The exact displayed statement gives the domain, codomain, and defining equation when present.

\noindent\textbf{Role in the proof.}
Use in this chapter. It contributes a sampler or sampler-derived object to the distributional model.

\end{formaldefinition}

\begin{formaldefinition}[EasyCrypt line 872]
\noindent\textbf{EasyCrypt name.}
\begin{lstlisting}[language=EasyCrypt,basicstyle=\ttfamily\scriptsize]
paper_sim_commitment_lowbits
\end{lstlisting}
\noindent\textbf{Checked EasyCrypt statement.}
\begin{lstlisting}[language=EasyCrypt,basicstyle=\ttfamily\scriptsize]
op paper_sim_commitment_lowbits
   (md : mode) (s : paper_sim_signature_sample) : poly =
  mkseq
    (fun i =>
      if i = 0 then paper_sim_sample_entropy s
      else poly_coeff (paper_sim_sample_lowbits_carrier s) i)
    n.
\end{lstlisting}
\noindent\textbf{Formal mathematical reading.}
Mathematical definition. This is a total EasyCrypt operation.  The exact displayed statement gives the domain, codomain, and defining equation when present.

\noindent\textbf{Role in the proof.}
Use in this chapter. It is part of the exact formal vocabulary used by subsequent lemmas and games.

\end{formaldefinition}

\begin{formaldefinition}[EasyCrypt line 880]
\noindent\textbf{EasyCrypt name.}
\begin{lstlisting}[language=EasyCrypt,basicstyle=\ttfamily\scriptsize]
paper_sim_commitment_challenge
\end{lstlisting}
\noindent\textbf{Checked EasyCrypt statement.}
\begin{lstlisting}[language=EasyCrypt,basicstyle=\ttfamily\scriptsize]
op paper_sim_commitment_challenge
   (md : mode) (pk : pkey) (m : message) (ctx : context)
   (s : paper_sim_signature_sample) : challenge =
  challenge_hash md
    (paper_sim_sample_response_aux s)
    (paper_sim_commitment_lowbits md s)
    (message_hash pk ctx m).
\end{lstlisting}
\noindent\textbf{Formal mathematical reading.}
Mathematical definition. This is a total EasyCrypt operation.  The exact displayed statement gives the domain, codomain, and defining equation when present.

\noindent\textbf{Role in the proof.}
Use in this chapter. It is part of the exact formal vocabulary used by subsequent lemmas and games.

\end{formaldefinition}

\begin{formaldefinition}[EasyCrypt line 888]
\noindent\textbf{EasyCrypt name.}
\begin{lstlisting}[language=EasyCrypt,basicstyle=\ttfamily\scriptsize]
paper_sim_commitment_highbits
\end{lstlisting}
\noindent\textbf{Checked EasyCrypt statement.}
\begin{lstlisting}[language=EasyCrypt,basicstyle=\ttfamily\scriptsize]
op paper_sim_commitment_highbits
   (md : mode) (pk : pkey) (m : message) (ctx : context)
   (s : paper_sim_signature_sample) : polyveck =
  let ch = paper_sim_commitment_challenge md pk m ctx s in
  reconstructed_highbits md
    (paper_sim_sample_response_aux s)
    (public_key_challenge_term md pk ch).
\end{lstlisting}
\noindent\textbf{Formal mathematical reading.}
Mathematical definition. This is a total EasyCrypt operation.  The exact displayed statement gives the domain, codomain, and defining equation when present.

\noindent\textbf{Role in the proof.}
Use in this chapter. It is part of the exact formal vocabulary used by subsequent lemmas and games.

\end{formaldefinition}

\begin{formaldefinition}[EasyCrypt line 896]
\noindent\textbf{EasyCrypt name.}
\begin{lstlisting}[language=EasyCrypt,basicstyle=\ttfamily\scriptsize]
paper_sim_signature
\end{lstlisting}
\noindent\textbf{Checked EasyCrypt statement.}
\begin{lstlisting}[language=EasyCrypt,basicstyle=\ttfamily\scriptsize]
op paper_sim_signature
   (md : mode) (pk : pkey) (m : message) (ctx : context)
   (s : paper_sim_signature_sample) : signature =
  let highbits = paper_sim_commitment_highbits md pk m ctx s in
  let lowbits = paper_sim_commitment_lowbits md s in
  let mu = message_hash pk ctx m in
  let ch = paper_sim_commitment_challenge md pk m ctx s in
  (highbits, lowbits, mu, ch, paper_sim_sample_token s, 0, 0).
\end{lstlisting}
\noindent\textbf{Formal mathematical reading.}
Mathematical definition. This is a total EasyCrypt operation.  The exact displayed statement gives the domain, codomain, and defining equation when present.

\noindent\textbf{Role in the proof.}
Use in this chapter. It is part of the exact formal vocabulary used by subsequent lemmas and games.

\end{formaldefinition}

\begin{formaldefinition}[EasyCrypt line 905]
\noindent\textbf{EasyCrypt name.}
\begin{lstlisting}[language=EasyCrypt,basicstyle=\ttfamily\scriptsize]
paper_sim_signature_sample_wf
\end{lstlisting}
\noindent\textbf{Checked EasyCrypt statement.}
\begin{lstlisting}[language=EasyCrypt,basicstyle=\ttfamily\scriptsize]
op paper_sim_signature_sample_wf
   (md : mode) (s : paper_sim_signature_sample) : bool =
  polyvecl_wf md (paper_sim_sample_response s) /\
  polyveck_wf md (paper_sim_sample_response_aux s) /\
  polyveck_wf md (paper_sim_sample_hint s) /\
  poly_wf (paper_sim_sample_lowbits_carrier s).
\end{lstlisting}
\noindent\textbf{Formal mathematical reading.}
Mathematical definition. This is a Boolean predicate.  The exact displayed EasyCrypt statement gives its arguments; mathematically it is an event, invariant, support condition, or well-formedness predicate over those arguments.

\noindent\textbf{Role in the proof.}
Use in this chapter. It names a well-formedness, support, or validity predicate needed by later preservation lemmas.

\end{formaldefinition}

\begin{formaldefinition}[EasyCrypt line 931]
\noindent\textbf{EasyCrypt name.}
\begin{lstlisting}[language=EasyCrypt,basicstyle=\ttfamily\scriptsize]
public_sim_lowbits_carrier
\end{lstlisting}
\noindent\textbf{Checked EasyCrypt statement.}
\begin{lstlisting}[language=EasyCrypt,basicstyle=\ttfamily\scriptsize]
op public_sim_lowbits_carrier
   (md : mode) (pk : pkey) (m : message) (ctx : context)
   (coins : random_coins) : poly =
  nth poly_zero
    (deterministic_polyveck md (public_sim_source pk m ctx coins)) 0.
\end{lstlisting}
\noindent\textbf{Formal mathematical reading.}
Mathematical definition. This is a total EasyCrypt operation.  The exact displayed statement gives the domain, codomain, and defining equation when present.

\noindent\textbf{Role in the proof.}
Use in this chapter. It is part of the exact formal vocabulary used by subsequent lemmas and games.

\end{formaldefinition}

\begin{formaldefinition}[EasyCrypt line 937]
\noindent\textbf{EasyCrypt name.}
\begin{lstlisting}[language=EasyCrypt,basicstyle=\ttfamily\scriptsize]
paper_sim_sample_from_public_coins
\end{lstlisting}
\noindent\textbf{Checked EasyCrypt statement.}
\begin{lstlisting}[language=EasyCrypt,basicstyle=\ttfamily\scriptsize]
op paper_sim_sample_from_public_coins
   (md : mode) (pk : pkey) (m : message) (ctx : context)
   (coins : random_coins) : paper_sim_signature_sample =
  ( public_sim_signature_token md pk m ctx coins,
    public_sim_lowbits_carrier md pk m ctx coins,
    signing_entropy_token_of_coins coins).
\end{lstlisting}
\noindent\textbf{Formal mathematical reading.}
Mathematical definition. This is a total EasyCrypt operation.  The exact displayed statement gives the domain, codomain, and defining equation when present.

\noindent\textbf{Role in the proof.}
Use in this chapter. It contributes a sampler or sampler-derived object to the distributional model.

\end{formaldefinition}

\begin{formaldefinition}[EasyCrypt line 1044]
\noindent\textbf{EasyCrypt name.}
\begin{lstlisting}[language=EasyCrypt,basicstyle=\ttfamily\scriptsize]
real_signing_lowbits_carrier
\end{lstlisting}
\noindent\textbf{Checked EasyCrypt statement.}
\begin{lstlisting}[language=EasyCrypt,basicstyle=\ttfamily\scriptsize]
op real_signing_lowbits_carrier
   (md : mode) (sk : skey) (m : message) (ctx : context)
   (coins : random_coins) : poly =
  nth poly_zero (commitment_raw md sk m ctx coins) 0.
\end{lstlisting}
\noindent\textbf{Formal mathematical reading.}
Mathematical definition. This is a total EasyCrypt operation.  The exact displayed statement gives the domain, codomain, and defining equation when present.

\noindent\textbf{Role in the proof.}
Use in this chapter. It is part of the exact formal vocabulary used by subsequent lemmas and games.

\end{formaldefinition}

\begin{formaldefinition}[EasyCrypt line 1049]
\noindent\textbf{EasyCrypt name.}
\begin{lstlisting}[language=EasyCrypt,basicstyle=\ttfamily\scriptsize]
paper_sim_sample_from_real_signing_coins
\end{lstlisting}
\noindent\textbf{Checked EasyCrypt statement.}
\begin{lstlisting}[language=EasyCrypt,basicstyle=\ttfamily\scriptsize]
op paper_sim_sample_from_real_signing_coins
   (md : mode) (sk : skey) (m : message) (ctx : context)
   (coins : random_coins) : paper_sim_signature_sample =
  ( signature_token md sk m ctx coins,
    real_signing_lowbits_carrier md sk m ctx coins,
    signing_entropy_token_of_coins coins).
\end{lstlisting}
\noindent\textbf{Formal mathematical reading.}
Mathematical definition. This is a total EasyCrypt operation.  The exact displayed statement gives the domain, codomain, and defining equation when present.

\noindent\textbf{Role in the proof.}
Use in this chapter. It contributes a sampler or sampler-derived object to the distributional model.

\end{formaldefinition}

\begin{formaldefinition}[EasyCrypt line 1154]
\noindent\textbf{EasyCrypt name.}
\begin{lstlisting}[language=EasyCrypt,basicstyle=\ttfamily\scriptsize]
paper_sim_sample_from_rejection_attempt
\end{lstlisting}
\noindent\textbf{Checked EasyCrypt statement.}
\begin{lstlisting}[language=EasyCrypt,basicstyle=\ttfamily\scriptsize]
op paper_sim_sample_from_rejection_attempt
   (st : signing_attempt_state) : paper_sim_signature_sample =
  ( signing_attempt_state_token st,
    signing_attempt_state_lowbits_carrier st,
    signing_entropy_token_of_coins (signing_attempt_state_coins st)).
\end{lstlisting}
\noindent\textbf{Formal mathematical reading.}
Mathematical definition. This is a total EasyCrypt operation.  The exact displayed statement gives the domain, codomain, and defining equation when present.

\noindent\textbf{Role in the proof.}
Use in this chapter. It names a bad event or rejection condition whose probability must be zero or explicitly bounded.

\end{formaldefinition}

\begin{formaldefinition}[EasyCrypt line 1160]
\noindent\textbf{EasyCrypt name.}
\begin{lstlisting}[language=EasyCrypt,basicstyle=\ttfamily\scriptsize]
paper_sim_abort_fallback_sample
\end{lstlisting}
\noindent\textbf{Checked EasyCrypt statement.}
\begin{lstlisting}[language=EasyCrypt,basicstyle=\ttfamily\scriptsize]
op paper_sim_abort_fallback_sample
   (md : mode) : paper_sim_signature_sample =
  ((polyvecl_zero md, polyveck_zero md, polyveck_zero md), poly_zero, 0).
\end{lstlisting}
\noindent\textbf{Formal mathematical reading.}
Mathematical definition. This is a total EasyCrypt operation.  The exact displayed statement gives the domain, codomain, and defining equation when present.

\noindent\textbf{Role in the proof.}
Use in this chapter. It contributes a sampler or sampler-derived object to the distributional model.

\end{formaldefinition}

\begin{formaldefinition}[EasyCrypt line 1164]
\noindent\textbf{EasyCrypt name.}
\begin{lstlisting}[language=EasyCrypt,basicstyle=\ttfamily\scriptsize]
haetae_rejection_sample_from_attempt
\end{lstlisting}
\noindent\textbf{Checked EasyCrypt statement.}
\begin{lstlisting}[language=EasyCrypt,basicstyle=\ttfamily\scriptsize]
op haetae_rejection_sample_from_attempt
   (md : mode) (pk : pkey) (m : message) (ctx : context)
   (st : signing_attempt_state) : paper_sim_signature_sample =
  if signing_attempt_state_rejection_accepts md pk m ctx st
  then paper_sim_sample_from_rejection_attempt st
  else paper_sim_abort_fallback_sample md.
\end{lstlisting}
\noindent\textbf{Formal mathematical reading.}
Mathematical definition. This is a total EasyCrypt operation.  The exact displayed statement gives the domain, codomain, and defining equation when present.

\noindent\textbf{Role in the proof.}
Use in this chapter. It names a bad event or rejection condition whose probability must be zero or explicitly bounded.

\end{formaldefinition}

\begin{formaldefinition}[EasyCrypt line 1352]
\noindent\textbf{EasyCrypt name.}
\begin{lstlisting}[language=EasyCrypt,basicstyle=\ttfamily\scriptsize]
PaperSimSigningSampler
\end{lstlisting}
\noindent\textbf{Checked EasyCrypt statement.}
\begin{lstlisting}[language=EasyCrypt,basicstyle=\ttfamily\scriptsize]
module type PaperSimSigningSampler = {
\end{lstlisting}
\noindent\textbf{Formal mathematical reading.}
Mathematical definition. This is an interface for an oracle, adversary, scheme, sampler, solver, or reduction.  A later theorem quantified over this interface is universally quantified over all modules implementing these procedures.

\noindent\textbf{Role in the proof.}
Use in this chapter. It defines the sampling procedure whose distributional behavior is justified by the later loss lemmas.

\end{formaldefinition}

\begin{formaldefinition}[EasyCrypt line 1353]
\noindent\textbf{EasyCrypt name.}
\begin{lstlisting}[language=EasyCrypt,basicstyle=\ttfamily\scriptsize]
init
\end{lstlisting}
\noindent\textbf{Checked EasyCrypt statement.}
\begin{lstlisting}[language=EasyCrypt,basicstyle=\ttfamily\scriptsize]
proc init(pk : pkey) : unit
  proc sample_with_seed(seed_coins : random_coins, m : message, ctx : context)
    : paper_sim_signature_sample
  proc sample(m : message, ctx : context) : paper_sim_signature_sample
}.
\end{lstlisting}
\noindent\textbf{Formal mathematical reading.}
Mathematical definition. This procedure is a command in the probabilistic program.  Its return value, state updates, and oracle calls are the operational content of the game.

\noindent\textbf{Role in the proof.}
Use in this chapter. It supplies an executable component of the formal probabilistic model.

\end{formaldefinition}

\begin{formaldefinition}[EasyCrypt line 1359]
\noindent\textbf{EasyCrypt name.}
\begin{lstlisting}[language=EasyCrypt,basicstyle=\ttfamily\scriptsize]
ROMPaperSimSigningSampler
\end{lstlisting}
\noindent\textbf{Checked EasyCrypt statement.}
\begin{lstlisting}[language=EasyCrypt,basicstyle=\ttfamily\scriptsize]
module ROMPaperSimSigningSampler(H : SIG.POracle) : PaperSimSigningSampler = {
\end{lstlisting}
\noindent\textbf{Formal mathematical reading.}
Mathematical definition. This is a probabilistic program/game or a module adapter.  Its procedures induce the probability space used by the game-based proof.

\noindent\textbf{Role in the proof.}
Use in this chapter. It defines the oracle boundary through which adversaries and simulations interact.

\end{formaldefinition}

\begin{formaldefinition}[EasyCrypt line 1362]
\noindent\textbf{EasyCrypt name.}
\begin{lstlisting}[language=EasyCrypt,basicstyle=\ttfamily\scriptsize]
init
\end{lstlisting}
\noindent\textbf{Checked EasyCrypt statement.}
\begin{lstlisting}[language=EasyCrypt,basicstyle=\ttfamily\scriptsize]
proc init(pk : pkey) : unit = {
\end{lstlisting}
\noindent\textbf{Formal mathematical reading.}
Mathematical definition. This procedure is a command in the probabilistic program.  Its return value, state updates, and oracle calls are the operational content of the game.

\noindent\textbf{Role in the proof.}
Use in this chapter. It supplies an executable component of the formal probabilistic model.

\end{formaldefinition}

\begin{formaldefinition}[EasyCrypt line 1366]
\noindent\textbf{EasyCrypt name.}
\begin{lstlisting}[language=EasyCrypt,basicstyle=\ttfamily\scriptsize]
sample_with_seed
\end{lstlisting}
\noindent\textbf{Checked EasyCrypt statement.}
\begin{lstlisting}[language=EasyCrypt,basicstyle=\ttfamily\scriptsize]
proc sample_with_seed(seed_coins : random_coins, m : message, ctx : context)
    : paper_sim_signature_sample = {
\end{lstlisting}
\noindent\textbf{Formal mathematical reading.}
Mathematical definition. This procedure is a command in the probabilistic program.  Its return value, state updates, and oracle calls are the operational content of the game.

\noindent\textbf{Role in the proof.}
Use in this chapter. It defines the sampling procedure whose distributional behavior is justified by the later loss lemmas.

\end{formaldefinition}

\begin{formaldefinition}[EasyCrypt line 1379]
\noindent\textbf{EasyCrypt name.}
\begin{lstlisting}[language=EasyCrypt,basicstyle=\ttfamily\scriptsize]
sample
\end{lstlisting}
\noindent\textbf{Checked EasyCrypt statement.}
\begin{lstlisting}[language=EasyCrypt,basicstyle=\ttfamily\scriptsize]
proc sample(m : message, ctx : context) : paper_sim_signature_sample = {
\end{lstlisting}
\noindent\textbf{Formal mathematical reading.}
Mathematical definition. This procedure is a command in the probabilistic program.  Its return value, state updates, and oracle calls are the operational content of the game.

\noindent\textbf{Role in the proof.}
Use in this chapter. It defines the sampling procedure whose distributional behavior is justified by the later loss lemmas.

\end{formaldefinition}

\begin{formaldefinition}[EasyCrypt line 1389]
\noindent\textbf{EasyCrypt name.}
\begin{lstlisting}[language=EasyCrypt,basicstyle=\ttfamily\scriptsize]
RealSigningPaperSimSampler
\end{lstlisting}
\noindent\textbf{Checked EasyCrypt statement.}
\begin{lstlisting}[language=EasyCrypt,basicstyle=\ttfamily\scriptsize]
module RealSigningPaperSimSampler(H : SIG.POracle) : PaperSimSigningSampler = {
\end{lstlisting}
\noindent\textbf{Formal mathematical reading.}
Mathematical definition. This is a probabilistic program/game or a module adapter.  Its procedures induce the probability space used by the game-based proof.

\noindent\textbf{Role in the proof.}
Use in this chapter. It defines the sampling procedure whose distributional behavior is justified by the later loss lemmas.

\end{formaldefinition}

\begin{formaldefinition}[EasyCrypt line 1393]
\noindent\textbf{EasyCrypt name.}
\begin{lstlisting}[language=EasyCrypt,basicstyle=\ttfamily\scriptsize]
init
\end{lstlisting}
\noindent\textbf{Checked EasyCrypt statement.}
\begin{lstlisting}[language=EasyCrypt,basicstyle=\ttfamily\scriptsize]
proc init(pk : pkey) : unit = {
\end{lstlisting}
\noindent\textbf{Formal mathematical reading.}
Mathematical definition. This procedure is a command in the probabilistic program.  Its return value, state updates, and oracle calls are the operational content of the game.

\noindent\textbf{Role in the proof.}
Use in this chapter. It supplies an executable component of the formal probabilistic model.

\end{formaldefinition}

\begin{formaldefinition}[EasyCrypt line 1398]
\noindent\textbf{EasyCrypt name.}
\begin{lstlisting}[language=EasyCrypt,basicstyle=\ttfamily\scriptsize]
sample_with_seed
\end{lstlisting}
\noindent\textbf{Checked EasyCrypt statement.}
\begin{lstlisting}[language=EasyCrypt,basicstyle=\ttfamily\scriptsize]
proc sample_with_seed(seed_coins : random_coins, m : message, ctx : context)
    : paper_sim_signature_sample = {
\end{lstlisting}
\noindent\textbf{Formal mathematical reading.}
Mathematical definition. This procedure is a command in the probabilistic program.  Its return value, state updates, and oracle calls are the operational content of the game.

\noindent\textbf{Role in the proof.}
Use in this chapter. It defines the sampling procedure whose distributional behavior is justified by the later loss lemmas.

\end{formaldefinition}

\begin{formaldefinition}[EasyCrypt line 1411]
\noindent\textbf{EasyCrypt name.}
\begin{lstlisting}[language=EasyCrypt,basicstyle=\ttfamily\scriptsize]
sample
\end{lstlisting}
\noindent\textbf{Checked EasyCrypt statement.}
\begin{lstlisting}[language=EasyCrypt,basicstyle=\ttfamily\scriptsize]
proc sample(m : message, ctx : context) : paper_sim_signature_sample = {
\end{lstlisting}
\noindent\textbf{Formal mathematical reading.}
Mathematical definition. This procedure is a command in the probabilistic program.  Its return value, state updates, and oracle calls are the operational content of the game.

\noindent\textbf{Role in the proof.}
Use in this chapter. It defines the sampling procedure whose distributional behavior is justified by the later loss lemmas.

\end{formaldefinition}

\begin{formaldefinition}[EasyCrypt line 1421]
\noindent\textbf{EasyCrypt name.}
\begin{lstlisting}[language=EasyCrypt,basicstyle=\ttfamily\scriptsize]
ROSigningAttemptPaperSimSampler
\end{lstlisting}
\noindent\textbf{Checked EasyCrypt statement.}
\begin{lstlisting}[language=EasyCrypt,basicstyle=\ttfamily\scriptsize]
module ROSigningAttemptPaperSimSampler(H : SIG.POracle)
  : PaperSimSigningSampler = {
\end{lstlisting}
\noindent\textbf{Formal mathematical reading.}
Mathematical definition. This is a probabilistic program/game or a module adapter.  Its procedures induce the probability space used by the game-based proof.

\noindent\textbf{Role in the proof.}
Use in this chapter. It defines the sampling procedure whose distributional behavior is justified by the later loss lemmas.

\end{formaldefinition}

\begin{formaldefinition}[EasyCrypt line 1426]
\noindent\textbf{EasyCrypt name.}
\begin{lstlisting}[language=EasyCrypt,basicstyle=\ttfamily\scriptsize]
init
\end{lstlisting}
\noindent\textbf{Checked EasyCrypt statement.}
\begin{lstlisting}[language=EasyCrypt,basicstyle=\ttfamily\scriptsize]
proc init(pk : pkey) : unit = {
\end{lstlisting}
\noindent\textbf{Formal mathematical reading.}
Mathematical definition. This procedure is a command in the probabilistic program.  Its return value, state updates, and oracle calls are the operational content of the game.

\noindent\textbf{Role in the proof.}
Use in this chapter. It supplies an executable component of the formal probabilistic model.

\end{formaldefinition}

\begin{formaldefinition}[EasyCrypt line 1431]
\noindent\textbf{EasyCrypt name.}
\begin{lstlisting}[language=EasyCrypt,basicstyle=\ttfamily\scriptsize]
sample_with_seed
\end{lstlisting}
\noindent\textbf{Checked EasyCrypt statement.}
\begin{lstlisting}[language=EasyCrypt,basicstyle=\ttfamily\scriptsize]
proc sample_with_seed(seed_coins : random_coins, m : message, ctx : context)
    : paper_sim_signature_sample = {
\end{lstlisting}
\noindent\textbf{Formal mathematical reading.}
Mathematical definition. This procedure is a command in the probabilistic program.  Its return value, state updates, and oracle calls are the operational content of the game.

\noindent\textbf{Role in the proof.}
Use in this chapter. It defines the sampling procedure whose distributional behavior is justified by the later loss lemmas.

\end{formaldefinition}

\begin{formaldefinition}[EasyCrypt line 1446]
\noindent\textbf{EasyCrypt name.}
\begin{lstlisting}[language=EasyCrypt,basicstyle=\ttfamily\scriptsize]
sample
\end{lstlisting}
\noindent\textbf{Checked EasyCrypt statement.}
\begin{lstlisting}[language=EasyCrypt,basicstyle=\ttfamily\scriptsize]
proc sample(m : message, ctx : context) : paper_sim_signature_sample = {
\end{lstlisting}
\noindent\textbf{Formal mathematical reading.}
Mathematical definition. This procedure is a command in the probabilistic program.  Its return value, state updates, and oracle calls are the operational content of the game.

\noindent\textbf{Role in the proof.}
Use in this chapter. It defines the sampling procedure whose distributional behavior is justified by the later loss lemmas.

\end{formaldefinition}

\begin{formaldefinition}[EasyCrypt line 1456]
\noindent\textbf{EasyCrypt name.}
\begin{lstlisting}[language=EasyCrypt,basicstyle=\ttfamily\scriptsize]
HAETAERejectionPaperSimSampler
\end{lstlisting}
\noindent\textbf{Checked EasyCrypt statement.}
\begin{lstlisting}[language=EasyCrypt,basicstyle=\ttfamily\scriptsize]
module HAETAERejectionPaperSimSampler(H : SIG.POracle)
  : PaperSimSigningSampler = {
\end{lstlisting}
\noindent\textbf{Formal mathematical reading.}
Mathematical definition. This is a probabilistic program/game or a module adapter.  Its procedures induce the probability space used by the game-based proof.

\noindent\textbf{Role in the proof.}
Use in this chapter. It defines the sampling procedure whose distributional behavior is justified by the later loss lemmas.

\end{formaldefinition}

\begin{formaldefinition}[EasyCrypt line 1461]
\noindent\textbf{EasyCrypt name.}
\begin{lstlisting}[language=EasyCrypt,basicstyle=\ttfamily\scriptsize]
init
\end{lstlisting}
\noindent\textbf{Checked EasyCrypt statement.}
\begin{lstlisting}[language=EasyCrypt,basicstyle=\ttfamily\scriptsize]
proc init(pk : pkey) : unit = {
\end{lstlisting}
\noindent\textbf{Formal mathematical reading.}
Mathematical definition. This procedure is a command in the probabilistic program.  Its return value, state updates, and oracle calls are the operational content of the game.

\noindent\textbf{Role in the proof.}
Use in this chapter. It supplies an executable component of the formal probabilistic model.

\end{formaldefinition}

\begin{formaldefinition}[EasyCrypt line 1466]
\noindent\textbf{EasyCrypt name.}
\begin{lstlisting}[language=EasyCrypt,basicstyle=\ttfamily\scriptsize]
sample_with_seed
\end{lstlisting}
\noindent\textbf{Checked EasyCrypt statement.}
\begin{lstlisting}[language=EasyCrypt,basicstyle=\ttfamily\scriptsize]
proc sample_with_seed(seed_coins : random_coins, m : message, ctx : context)
    : paper_sim_signature_sample = {
\end{lstlisting}
\noindent\textbf{Formal mathematical reading.}
Mathematical definition. This procedure is a command in the probabilistic program.  Its return value, state updates, and oracle calls are the operational content of the game.

\noindent\textbf{Role in the proof.}
Use in this chapter. It defines the sampling procedure whose distributional behavior is justified by the later loss lemmas.

\end{formaldefinition}

\begin{formaldefinition}[EasyCrypt line 1482]
\noindent\textbf{EasyCrypt name.}
\begin{lstlisting}[language=EasyCrypt,basicstyle=\ttfamily\scriptsize]
sample
\end{lstlisting}
\noindent\textbf{Checked EasyCrypt statement.}
\begin{lstlisting}[language=EasyCrypt,basicstyle=\ttfamily\scriptsize]
proc sample(m : message, ctx : context) : paper_sim_signature_sample = {
\end{lstlisting}
\noindent\textbf{Formal mathematical reading.}
Mathematical definition. This procedure is a command in the probabilistic program.  Its return value, state updates, and oracle calls are the operational content of the game.

\noindent\textbf{Role in the proof.}
Use in this chapter. It defines the sampling procedure whose distributional behavior is justified by the later loss lemmas.

\end{formaldefinition}

\begin{formaldefinition}[EasyCrypt line 1492]
\noindent\textbf{EasyCrypt name.}
\begin{lstlisting}[language=EasyCrypt,basicstyle=\ttfamily\scriptsize]
ExactHyperballHAETAERejectionPaperSimSampler
\end{lstlisting}
\noindent\textbf{Checked EasyCrypt statement.}
\begin{lstlisting}[language=EasyCrypt,basicstyle=\ttfamily\scriptsize]
module ExactHyperballHAETAERejectionPaperSimSampler(H : SIG.POracle)
  : PaperSimSigningSampler = {
\end{lstlisting}
\noindent\textbf{Formal mathematical reading.}
Mathematical definition. This is a probabilistic program/game or a module adapter.  Its procedures induce the probability space used by the game-based proof.

\noindent\textbf{Role in the proof.}
Use in this chapter. It defines the sampling procedure whose distributional behavior is justified by the later loss lemmas.

\end{formaldefinition}

\begin{formaldefinition}[EasyCrypt line 1497]
\noindent\textbf{EasyCrypt name.}
\begin{lstlisting}[language=EasyCrypt,basicstyle=\ttfamily\scriptsize]
init
\end{lstlisting}
\noindent\textbf{Checked EasyCrypt statement.}
\begin{lstlisting}[language=EasyCrypt,basicstyle=\ttfamily\scriptsize]
proc init(pk : pkey) : unit = {
\end{lstlisting}
\noindent\textbf{Formal mathematical reading.}
Mathematical definition. This procedure is a command in the probabilistic program.  Its return value, state updates, and oracle calls are the operational content of the game.

\noindent\textbf{Role in the proof.}
Use in this chapter. It supplies an executable component of the formal probabilistic model.

\end{formaldefinition}

\begin{formaldefinition}[EasyCrypt line 1502]
\noindent\textbf{EasyCrypt name.}
\begin{lstlisting}[language=EasyCrypt,basicstyle=\ttfamily\scriptsize]
sample_with_seed
\end{lstlisting}
\noindent\textbf{Checked EasyCrypt statement.}
\begin{lstlisting}[language=EasyCrypt,basicstyle=\ttfamily\scriptsize]
proc sample_with_seed(seed_coins : random_coins, m : message, ctx : context)
    : paper_sim_signature_sample = {
\end{lstlisting}
\noindent\textbf{Formal mathematical reading.}
Mathematical definition. This procedure is a command in the probabilistic program.  Its return value, state updates, and oracle calls are the operational content of the game.

\noindent\textbf{Role in the proof.}
Use in this chapter. It defines the sampling procedure whose distributional behavior is justified by the later loss lemmas.

\end{formaldefinition}

\begin{formaldefinition}[EasyCrypt line 1513]
\noindent\textbf{EasyCrypt name.}
\begin{lstlisting}[language=EasyCrypt,basicstyle=\ttfamily\scriptsize]
sample
\end{lstlisting}
\noindent\textbf{Checked EasyCrypt statement.}
\begin{lstlisting}[language=EasyCrypt,basicstyle=\ttfamily\scriptsize]
proc sample(m : message, ctx : context) : paper_sim_signature_sample = {
\end{lstlisting}
\noindent\textbf{Formal mathematical reading.}
Mathematical definition. This procedure is a command in the probabilistic program.  Its return value, state updates, and oracle calls are the operational content of the game.

\noindent\textbf{Role in the proof.}
Use in this chapter. It defines the sampling procedure whose distributional behavior is justified by the later loss lemmas.

\end{formaldefinition}

\begin{formaldefinition}[EasyCrypt line 1523]
\noindent\textbf{EasyCrypt name.}
\begin{lstlisting}[language=EasyCrypt,basicstyle=\ttfamily\scriptsize]
ExactHyperballPaperSimSampler
\end{lstlisting}
\noindent\textbf{Checked EasyCrypt statement.}
\begin{lstlisting}[language=EasyCrypt,basicstyle=\ttfamily\scriptsize]
module ExactHyperballPaperSimSampler(H : SIG.POracle)
  : PaperSimSigningSampler = {
\end{lstlisting}
\noindent\textbf{Formal mathematical reading.}
Mathematical definition. This is a probabilistic program/game or a module adapter.  Its procedures induce the probability space used by the game-based proof.

\noindent\textbf{Role in the proof.}
Use in this chapter. It defines the sampling procedure whose distributional behavior is justified by the later loss lemmas.

\end{formaldefinition}

\begin{formaldefinition}[EasyCrypt line 1528]
\noindent\textbf{EasyCrypt name.}
\begin{lstlisting}[language=EasyCrypt,basicstyle=\ttfamily\scriptsize]
init
\end{lstlisting}
\noindent\textbf{Checked EasyCrypt statement.}
\begin{lstlisting}[language=EasyCrypt,basicstyle=\ttfamily\scriptsize]
proc init(pk : pkey) : unit = {
\end{lstlisting}
\noindent\textbf{Formal mathematical reading.}
Mathematical definition. This procedure is a command in the probabilistic program.  Its return value, state updates, and oracle calls are the operational content of the game.

\noindent\textbf{Role in the proof.}
Use in this chapter. It supplies an executable component of the formal probabilistic model.

\end{formaldefinition}

\begin{formaldefinition}[EasyCrypt line 1533]
\noindent\textbf{EasyCrypt name.}
\begin{lstlisting}[language=EasyCrypt,basicstyle=\ttfamily\scriptsize]
sample_with_seed
\end{lstlisting}
\noindent\textbf{Checked EasyCrypt statement.}
\begin{lstlisting}[language=EasyCrypt,basicstyle=\ttfamily\scriptsize]
proc sample_with_seed(seed_coins : random_coins, m : message, ctx : context)
    : paper_sim_signature_sample = {
\end{lstlisting}
\noindent\textbf{Formal mathematical reading.}
Mathematical definition. This procedure is a command in the probabilistic program.  Its return value, state updates, and oracle calls are the operational content of the game.

\noindent\textbf{Role in the proof.}
Use in this chapter. It defines the sampling procedure whose distributional behavior is justified by the later loss lemmas.

\end{formaldefinition}

\begin{formaldefinition}[EasyCrypt line 1544]
\noindent\textbf{EasyCrypt name.}
\begin{lstlisting}[language=EasyCrypt,basicstyle=\ttfamily\scriptsize]
sample
\end{lstlisting}
\noindent\textbf{Checked EasyCrypt statement.}
\begin{lstlisting}[language=EasyCrypt,basicstyle=\ttfamily\scriptsize]
proc sample(m : message, ctx : context) : paper_sim_signature_sample = {
\end{lstlisting}
\noindent\textbf{Formal mathematical reading.}
Mathematical definition. This procedure is a command in the probabilistic program.  Its return value, state updates, and oracle calls are the operational content of the game.

\noindent\textbf{Role in the proof.}
Use in this chapter. It defines the sampling procedure whose distributional behavior is justified by the later loss lemmas.

\end{formaldefinition}

\begin{formaldefinition}[EasyCrypt line 1554]
\noindent\textbf{EasyCrypt name.}
\begin{lstlisting}[language=EasyCrypt,basicstyle=\ttfamily\scriptsize]
ROExactHyperballPaperSimSampler
\end{lstlisting}
\noindent\textbf{Checked EasyCrypt statement.}
\begin{lstlisting}[language=EasyCrypt,basicstyle=\ttfamily\scriptsize]
module ROExactHyperballPaperSimSampler(H : SIG.POracle)
  : PaperSimSigningSampler = {
\end{lstlisting}
\noindent\textbf{Formal mathematical reading.}
Mathematical definition. This is a probabilistic program/game or a module adapter.  Its procedures induce the probability space used by the game-based proof.

\noindent\textbf{Role in the proof.}
Use in this chapter. It defines the sampling procedure whose distributional behavior is justified by the later loss lemmas.

\end{formaldefinition}

\begin{formaldefinition}[EasyCrypt line 1559]
\noindent\textbf{EasyCrypt name.}
\begin{lstlisting}[language=EasyCrypt,basicstyle=\ttfamily\scriptsize]
init
\end{lstlisting}
\noindent\textbf{Checked EasyCrypt statement.}
\begin{lstlisting}[language=EasyCrypt,basicstyle=\ttfamily\scriptsize]
proc init(pk : pkey) : unit = {
\end{lstlisting}
\noindent\textbf{Formal mathematical reading.}
Mathematical definition. This procedure is a command in the probabilistic program.  Its return value, state updates, and oracle calls are the operational content of the game.

\noindent\textbf{Role in the proof.}
Use in this chapter. It supplies an executable component of the formal probabilistic model.

\end{formaldefinition}

\begin{formaldefinition}[EasyCrypt line 1564]
\noindent\textbf{EasyCrypt name.}
\begin{lstlisting}[language=EasyCrypt,basicstyle=\ttfamily\scriptsize]
sample_with_seed
\end{lstlisting}
\noindent\textbf{Checked EasyCrypt statement.}
\begin{lstlisting}[language=EasyCrypt,basicstyle=\ttfamily\scriptsize]
proc sample_with_seed(seed_coins : random_coins, m : message, ctx : context)
    : paper_sim_signature_sample = {
\end{lstlisting}
\noindent\textbf{Formal mathematical reading.}
Mathematical definition. This procedure is a command in the probabilistic program.  Its return value, state updates, and oracle calls are the operational content of the game.

\noindent\textbf{Role in the proof.}
Use in this chapter. It defines the sampling procedure whose distributional behavior is justified by the later loss lemmas.

\end{formaldefinition}

\begin{formaldefinition}[EasyCrypt line 1577]
\noindent\textbf{EasyCrypt name.}
\begin{lstlisting}[language=EasyCrypt,basicstyle=\ttfamily\scriptsize]
sample
\end{lstlisting}
\noindent\textbf{Checked EasyCrypt statement.}
\begin{lstlisting}[language=EasyCrypt,basicstyle=\ttfamily\scriptsize]
proc sample(m : message, ctx : context) : paper_sim_signature_sample = {
\end{lstlisting}
\noindent\textbf{Formal mathematical reading.}
Mathematical definition. This procedure is a command in the probabilistic program.  Its return value, state updates, and oracle calls are the operational content of the game.

\noindent\textbf{Role in the proof.}
Use in this chapter. It defines the sampling procedure whose distributional behavior is justified by the later loss lemmas.

\end{formaldefinition}

\begin{formaldefinition}[EasyCrypt line 1587]
\noindent\textbf{EasyCrypt name.}
\begin{lstlisting}[language=EasyCrypt,basicstyle=\ttfamily\scriptsize]
StructuralAttemptPaperSample
\end{lstlisting}
\noindent\textbf{Checked EasyCrypt statement.}
\begin{lstlisting}[language=EasyCrypt,basicstyle=\ttfamily\scriptsize]
module StructuralAttemptPaperSample = {
\end{lstlisting}
\noindent\textbf{Formal mathematical reading.}
Mathematical definition. This is a probabilistic program/game or a module adapter.  Its procedures induce the probability space used by the game-based proof.

\noindent\textbf{Role in the proof.}
Use in this chapter. It defines the sampling procedure whose distributional behavior is justified by the later loss lemmas.

\end{formaldefinition}

\begin{formaldefinition}[EasyCrypt line 1588]
\noindent\textbf{EasyCrypt name.}
\begin{lstlisting}[language=EasyCrypt,basicstyle=\ttfamily\scriptsize]
sample
\end{lstlisting}
\noindent\textbf{Checked EasyCrypt statement.}
\begin{lstlisting}[language=EasyCrypt,basicstyle=\ttfamily\scriptsize]
proc sample(sk : skey, m : message, ctx : context)
    : paper_sim_signature_sample = {
\end{lstlisting}
\noindent\textbf{Formal mathematical reading.}
Mathematical definition. This procedure is a command in the probabilistic program.  Its return value, state updates, and oracle calls are the operational content of the game.

\noindent\textbf{Role in the proof.}
Use in this chapter. It defines the sampling procedure whose distributional behavior is justified by the later loss lemmas.

\end{formaldefinition}

\begin{formaldefinition}[EasyCrypt line 1599]
\noindent\textbf{EasyCrypt name.}
\begin{lstlisting}[language=EasyCrypt,basicstyle=\ttfamily\scriptsize]
ExactHyperballPaperSample
\end{lstlisting}
\noindent\textbf{Checked EasyCrypt statement.}
\begin{lstlisting}[language=EasyCrypt,basicstyle=\ttfamily\scriptsize]
module ExactHyperballPaperSample = {
\end{lstlisting}
\noindent\textbf{Formal mathematical reading.}
Mathematical definition. This is a probabilistic program/game or a module adapter.  Its procedures induce the probability space used by the game-based proof.

\noindent\textbf{Role in the proof.}
Use in this chapter. It defines the sampling procedure whose distributional behavior is justified by the later loss lemmas.

\end{formaldefinition}

\begin{formaldefinition}[EasyCrypt line 1600]
\noindent\textbf{EasyCrypt name.}
\begin{lstlisting}[language=EasyCrypt,basicstyle=\ttfamily\scriptsize]
sample
\end{lstlisting}
\noindent\textbf{Checked EasyCrypt statement.}
\begin{lstlisting}[language=EasyCrypt,basicstyle=\ttfamily\scriptsize]
proc sample(sk : skey, m : message, ctx : context)
    : paper_sim_signature_sample = {
\end{lstlisting}
\noindent\textbf{Formal mathematical reading.}
Mathematical definition. This procedure is a command in the probabilistic program.  Its return value, state updates, and oracle calls are the operational content of the game.

\noindent\textbf{Role in the proof.}
Use in this chapter. It defines the sampling procedure whose distributional behavior is justified by the later loss lemmas.

\end{formaldefinition}

\begin{formaldefinition}[EasyCrypt line 1611]
\noindent\textbf{EasyCrypt name.}
\begin{lstlisting}[language=EasyCrypt,basicstyle=\ttfamily\scriptsize]
ROMInternalTranscriptPaperSimAsNMA
\end{lstlisting}
\noindent\textbf{Checked EasyCrypt statement.}
\begin{lstlisting}[language=EasyCrypt,basicstyle=\ttfamily\scriptsize]
module ROMInternalTranscriptPaperSimAsNMA(
  A : SIG.Adversary,
  Samp : PaperSimSigningSampler
) (H : SIG.POracle) = {
\end{lstlisting}
\noindent\textbf{Formal mathematical reading.}
Mathematical definition. This is a probabilistic program/game or a module adapter.  Its procedures induce the probability space used by the game-based proof.

\noindent\textbf{Role in the proof.}
Use in this chapter. It defines the game or game interface whose success event is bounded by later theorems.

\end{formaldefinition}

\begin{formaldefinition}[EasyCrypt line 1620]
\noindent\textbf{EasyCrypt name.}
\begin{lstlisting}[language=EasyCrypt,basicstyle=\ttfamily\scriptsize]
O
\end{lstlisting}
\noindent\textbf{Checked EasyCrypt statement.}
\begin{lstlisting}[language=EasyCrypt,basicstyle=\ttfamily\scriptsize]
module O = {
\end{lstlisting}
\noindent\textbf{Formal mathematical reading.}
Mathematical definition. This is a probabilistic program/game or a module adapter.  Its procedures induce the probability space used by the game-based proof.

\noindent\textbf{Role in the proof.}
Use in this chapter. It supplies an executable component of the formal probabilistic model.

\end{formaldefinition}

\begin{formaldefinition}[EasyCrypt line 1621]
\noindent\textbf{EasyCrypt name.}
\begin{lstlisting}[language=EasyCrypt,basicstyle=\ttfamily\scriptsize]
sign
\end{lstlisting}
\noindent\textbf{Checked EasyCrypt statement.}
\begin{lstlisting}[language=EasyCrypt,basicstyle=\ttfamily\scriptsize]
proc sign(m : message, ctx : context) : signature = {
\end{lstlisting}
\noindent\textbf{Formal mathematical reading.}
Mathematical definition. This procedure is a command in the probabilistic program.  Its return value, state updates, and oracle calls are the operational content of the game.

\noindent\textbf{Role in the proof.}
Use in this chapter. It supplies an executable component of the formal probabilistic model.

\end{formaldefinition}

\begin{formaldefinition}[EasyCrypt line 1646]
\noindent\textbf{EasyCrypt name.}
\begin{lstlisting}[language=EasyCrypt,basicstyle=\ttfamily\scriptsize]
A
\end{lstlisting}
\noindent\textbf{Checked EasyCrypt statement.}
\begin{lstlisting}[language=EasyCrypt,basicstyle=\ttfamily\scriptsize]
module A = A(H, O)

  proc forge(pk : pkey) : message * context * signature = {
\end{lstlisting}
\noindent\textbf{Formal mathematical reading.}
Mathematical definition. This is a probabilistic program/game or a module adapter.  Its procedures induce the probability space used by the game-based proof.

\noindent\textbf{Role in the proof.}
Use in this chapter. It supplies an executable component of the formal probabilistic model.

\end{formaldefinition}

\begin{formaldefinition}[EasyCrypt line 1661]
\noindent\textbf{EasyCrypt name.}
\begin{lstlisting}[language=EasyCrypt,basicstyle=\ttfamily\scriptsize]
ROMInternalTranscriptBudgetedPaperSimAsNMA
\end{lstlisting}
\noindent\textbf{Checked EasyCrypt statement.}
\begin{lstlisting}[language=EasyCrypt,basicstyle=\ttfamily\scriptsize]
module ROMInternalTranscriptBudgetedPaperSimAsNMA(
  A : SIG.Adversary,
  Samp : PaperSimSigningSampler
) (H : SIG.POracle) = {
\end{lstlisting}
\noindent\textbf{Formal mathematical reading.}
Mathematical definition. This is a probabilistic program/game or a module adapter.  Its procedures induce the probability space used by the game-based proof.

\noindent\textbf{Role in the proof.}
Use in this chapter. It defines the game or game interface whose success event is bounded by later theorems.

\end{formaldefinition}

\begin{formaldefinition}[EasyCrypt line 1675]
\noindent\textbf{EasyCrypt name.}
\begin{lstlisting}[language=EasyCrypt,basicstyle=\ttfamily\scriptsize]
AH
\end{lstlisting}
\noindent\textbf{Checked EasyCrypt statement.}
\begin{lstlisting}[language=EasyCrypt,basicstyle=\ttfamily\scriptsize]
module AH = {
\end{lstlisting}
\noindent\textbf{Formal mathematical reading.}
Mathematical definition. This is a probabilistic program/game or a module adapter.  Its procedures induce the probability space used by the game-based proof.

\noindent\textbf{Role in the proof.}
Use in this chapter. It supplies an executable component of the formal probabilistic model.

\end{formaldefinition}

\begin{formaldefinition}[EasyCrypt line 1676]
\noindent\textbf{EasyCrypt name.}
\begin{lstlisting}[language=EasyCrypt,basicstyle=\ttfamily\scriptsize]
get
\end{lstlisting}
\noindent\textbf{Checked EasyCrypt statement.}
\begin{lstlisting}[language=EasyCrypt,basicstyle=\ttfamily\scriptsize]
proc get(q : ro_query) : ro_output = {
\end{lstlisting}
\noindent\textbf{Formal mathematical reading.}
Mathematical definition. This procedure is a command in the probabilistic program.  Its return value, state updates, and oracle calls are the operational content of the game.

\noindent\textbf{Role in the proof.}
Use in this chapter. It supplies an executable component of the formal probabilistic model.

\end{formaldefinition}

\begin{formaldefinition}[EasyCrypt line 1695]
\noindent\textbf{EasyCrypt name.}
\begin{lstlisting}[language=EasyCrypt,basicstyle=\ttfamily\scriptsize]
O
\end{lstlisting}
\noindent\textbf{Checked EasyCrypt statement.}
\begin{lstlisting}[language=EasyCrypt,basicstyle=\ttfamily\scriptsize]
module O = {
\end{lstlisting}
\noindent\textbf{Formal mathematical reading.}
Mathematical definition. This is a probabilistic program/game or a module adapter.  Its procedures induce the probability space used by the game-based proof.

\noindent\textbf{Role in the proof.}
Use in this chapter. It supplies an executable component of the formal probabilistic model.

\end{formaldefinition}

\begin{formaldefinition}[EasyCrypt line 1696]
\noindent\textbf{EasyCrypt name.}
\begin{lstlisting}[language=EasyCrypt,basicstyle=\ttfamily\scriptsize]
sign
\end{lstlisting}
\noindent\textbf{Checked EasyCrypt statement.}
\begin{lstlisting}[language=EasyCrypt,basicstyle=\ttfamily\scriptsize]
proc sign(m : message, ctx : context) : signature = {
\end{lstlisting}
\noindent\textbf{Formal mathematical reading.}
Mathematical definition. This procedure is a command in the probabilistic program.  Its return value, state updates, and oracle calls are the operational content of the game.

\noindent\textbf{Role in the proof.}
Use in this chapter. It supplies an executable component of the formal probabilistic model.

\end{formaldefinition}

\begin{formaldefinition}[EasyCrypt line 1746]
\noindent\textbf{EasyCrypt name.}
\begin{lstlisting}[language=EasyCrypt,basicstyle=\ttfamily\scriptsize]
A
\end{lstlisting}
\noindent\textbf{Checked EasyCrypt statement.}
\begin{lstlisting}[language=EasyCrypt,basicstyle=\ttfamily\scriptsize]
module A = A(AH, O)

  proc forge(pk : pkey) : message * context * signature = {
\end{lstlisting}
\noindent\textbf{Formal mathematical reading.}
Mathematical definition. This is a probabilistic program/game or a module adapter.  Its procedures induce the probability space used by the game-based proof.

\noindent\textbf{Role in the proof.}
Use in this chapter. It supplies an executable component of the formal probabilistic model.

\end{formaldefinition}

\begin{formaldefinition}[EasyCrypt line 2624]
\noindent\textbf{EasyCrypt name.}
\begin{lstlisting}[language=EasyCrypt,basicstyle=\ttfamily\scriptsize]
EUF_CMA_ROMInternalTranscriptCountedSign
\end{lstlisting}
\noindent\textbf{Checked EasyCrypt statement.}
\begin{lstlisting}[language=EasyCrypt,basicstyle=\ttfamily\scriptsize]
module EUF_CMA_ROMInternalTranscriptCountedSign(
  H : Oracle,
  A : SIG.Adversary
) = {
\end{lstlisting}
\noindent\textbf{Formal mathematical reading.}
Mathematical definition. This is a probabilistic program/game or a module adapter.  Its procedures induce the probability space used by the game-based proof.

\noindent\textbf{Role in the proof.}
Use in this chapter. It defines the game or game interface whose success event is bounded by later theorems.

\end{formaldefinition}

\begin{formaldefinition}[EasyCrypt line 2628]
\noindent\textbf{EasyCrypt name.}
\begin{lstlisting}[language=EasyCrypt,basicstyle=\ttfamily\scriptsize]
C
\end{lstlisting}
\noindent\textbf{Checked EasyCrypt statement.}
\begin{lstlisting}[language=EasyCrypt,basicstyle=\ttfamily\scriptsize]
module C = CountedLazyROM(H)

  var pk_current : pkey
  var sk_current : skey
  var queries : SIG.query list
  var transcripts : transcript list
  var records : signing_transcript_record list

  module O = {
\end{lstlisting}
\noindent\textbf{Formal mathematical reading.}
Mathematical definition. This is a probabilistic program/game or a module adapter.  Its procedures induce the probability space used by the game-based proof.

\noindent\textbf{Role in the proof.}
Use in this chapter. It supplies an executable component of the formal probabilistic model.

\end{formaldefinition}

\begin{formaldefinition}[EasyCrypt line 2637]
\noindent\textbf{EasyCrypt name.}
\begin{lstlisting}[language=EasyCrypt,basicstyle=\ttfamily\scriptsize]
sign
\end{lstlisting}
\noindent\textbf{Checked EasyCrypt statement.}
\begin{lstlisting}[language=EasyCrypt,basicstyle=\ttfamily\scriptsize]
proc sign(m : message, ctx : context) : signature = {
\end{lstlisting}
\noindent\textbf{Formal mathematical reading.}
Mathematical definition. This procedure is a command in the probabilistic program.  Its return value, state updates, and oracle calls are the operational content of the game.

\noindent\textbf{Role in the proof.}
Use in this chapter. It supplies an executable component of the formal probabilistic model.

\end{formaldefinition}

\begin{formaldefinition}[EasyCrypt line 2667]
\noindent\textbf{EasyCrypt name.}
\begin{lstlisting}[language=EasyCrypt,basicstyle=\ttfamily\scriptsize]
A
\end{lstlisting}
\noindent\textbf{Checked EasyCrypt statement.}
\begin{lstlisting}[language=EasyCrypt,basicstyle=\ttfamily\scriptsize]
module A = A(C, O)

  proc main() : bool = {
\end{lstlisting}
\noindent\textbf{Formal mathematical reading.}
Mathematical definition. This is a probabilistic program/game or a module adapter.  Its procedures induce the probability space used by the game-based proof.

\noindent\textbf{Role in the proof.}
Use in this chapter. It supplies an executable component of the formal probabilistic model.

\end{formaldefinition}

\begin{formaldefinition}[EasyCrypt line 2696]
\noindent\textbf{EasyCrypt name.}
\begin{lstlisting}[language=EasyCrypt,basicstyle=\ttfamily\scriptsize]
EUF_CMA_ROMProgrammedTranscriptCountedSign
\end{lstlisting}
\noindent\textbf{Checked EasyCrypt statement.}
\begin{lstlisting}[language=EasyCrypt,basicstyle=\ttfamily\scriptsize]
module EUF_CMA_ROMProgrammedTranscriptCountedSign(
  H : Oracle,
  A : SIG.Adversary
) = {
\end{lstlisting}
\noindent\textbf{Formal mathematical reading.}
Mathematical definition. This is a probabilistic program/game or a module adapter.  Its procedures induce the probability space used by the game-based proof.

\noindent\textbf{Role in the proof.}
Use in this chapter. It defines the game or game interface whose success event is bounded by later theorems.

\end{formaldefinition}

\begin{formaldefinition}[EasyCrypt line 2700]
\noindent\textbf{EasyCrypt name.}
\begin{lstlisting}[language=EasyCrypt,basicstyle=\ttfamily\scriptsize]
C
\end{lstlisting}
\noindent\textbf{Checked EasyCrypt statement.}
\begin{lstlisting}[language=EasyCrypt,basicstyle=\ttfamily\scriptsize]
module C = CountedLazyROM(H)

  var pk_current : pkey
  var sk_current : skey
  var queries : SIG.query list
  var transcripts : transcript list
  var records : signing_transcript_record list

  module O = {
\end{lstlisting}
\noindent\textbf{Formal mathematical reading.}
Mathematical definition. This is a probabilistic program/game or a module adapter.  Its procedures induce the probability space used by the game-based proof.

\noindent\textbf{Role in the proof.}
Use in this chapter. It supplies an executable component of the formal probabilistic model.

\end{formaldefinition}

\begin{formaldefinition}[EasyCrypt line 2709]
\noindent\textbf{EasyCrypt name.}
\begin{lstlisting}[language=EasyCrypt,basicstyle=\ttfamily\scriptsize]
sign
\end{lstlisting}
\noindent\textbf{Checked EasyCrypt statement.}
\begin{lstlisting}[language=EasyCrypt,basicstyle=\ttfamily\scriptsize]
proc sign(m : message, ctx : context) : signature = {
\end{lstlisting}
\noindent\textbf{Formal mathematical reading.}
Mathematical definition. This procedure is a command in the probabilistic program.  Its return value, state updates, and oracle calls are the operational content of the game.

\noindent\textbf{Role in the proof.}
Use in this chapter. It supplies an executable component of the formal probabilistic model.

\end{formaldefinition}

\begin{formaldefinition}[EasyCrypt line 2741]
\noindent\textbf{EasyCrypt name.}
\begin{lstlisting}[language=EasyCrypt,basicstyle=\ttfamily\scriptsize]
A
\end{lstlisting}
\noindent\textbf{Checked EasyCrypt statement.}
\begin{lstlisting}[language=EasyCrypt,basicstyle=\ttfamily\scriptsize]
module A = A(C, O)

  proc main() : bool = {
\end{lstlisting}
\noindent\textbf{Formal mathematical reading.}
Mathematical definition. This is a probabilistic program/game or a module adapter.  Its procedures induce the probability space used by the game-based proof.

\noindent\textbf{Role in the proof.}
Use in this chapter. It supplies an executable component of the formal probabilistic model.

\end{formaldefinition}

\begin{formaldefinition}[EasyCrypt line 2770]
\noindent\textbf{EasyCrypt name.}
\begin{lstlisting}[language=EasyCrypt,basicstyle=\ttfamily\scriptsize]
EUF_CMA_ROMProgrammedTranscriptBudgetedCountedSign
\end{lstlisting}
\noindent\textbf{Checked EasyCrypt statement.}
\begin{lstlisting}[language=EasyCrypt,basicstyle=\ttfamily\scriptsize]
module EUF_CMA_ROMProgrammedTranscriptBudgetedCountedSign(
  H : Oracle,
  A : SIG.Adversary
) = {
\end{lstlisting}
\noindent\textbf{Formal mathematical reading.}
Mathematical definition. This is a probabilistic program/game or a module adapter.  Its procedures induce the probability space used by the game-based proof.

\noindent\textbf{Role in the proof.}
Use in this chapter. It defines the game or game interface whose success event is bounded by later theorems.

\end{formaldefinition}

\begin{formaldefinition}[EasyCrypt line 2774]
\noindent\textbf{EasyCrypt name.}
\begin{lstlisting}[language=EasyCrypt,basicstyle=\ttfamily\scriptsize]
C
\end{lstlisting}
\noindent\textbf{Checked EasyCrypt statement.}
\begin{lstlisting}[language=EasyCrypt,basicstyle=\ttfamily\scriptsize]
module C = CountedLazyROM(H)

  var pk_current : pkey
  var sk_current : skey
  var queries : SIG.query list
  var transcripts : transcript list
  var records : signing_transcript_record list
  var adversary_hash_count : int
  var signing_count : int

  module AH = {
\end{lstlisting}
\noindent\textbf{Formal mathematical reading.}
Mathematical definition. This is a probabilistic program/game or a module adapter.  Its procedures induce the probability space used by the game-based proof.

\noindent\textbf{Role in the proof.}
Use in this chapter. It supplies an executable component of the formal probabilistic model.

\end{formaldefinition}

\begin{formaldefinition}[EasyCrypt line 2785]
\noindent\textbf{EasyCrypt name.}
\begin{lstlisting}[language=EasyCrypt,basicstyle=\ttfamily\scriptsize]
get
\end{lstlisting}
\noindent\textbf{Checked EasyCrypt statement.}
\begin{lstlisting}[language=EasyCrypt,basicstyle=\ttfamily\scriptsize]
proc get(q : ro_query) : ro_output = {
\end{lstlisting}
\noindent\textbf{Formal mathematical reading.}
Mathematical definition. This procedure is a command in the probabilistic program.  Its return value, state updates, and oracle calls are the operational content of the game.

\noindent\textbf{Role in the proof.}
Use in this chapter. It supplies an executable component of the formal probabilistic model.

\end{formaldefinition}

\begin{formaldefinition}[EasyCrypt line 2803]
\noindent\textbf{EasyCrypt name.}
\begin{lstlisting}[language=EasyCrypt,basicstyle=\ttfamily\scriptsize]
O
\end{lstlisting}
\noindent\textbf{Checked EasyCrypt statement.}
\begin{lstlisting}[language=EasyCrypt,basicstyle=\ttfamily\scriptsize]
module O = {
\end{lstlisting}
\noindent\textbf{Formal mathematical reading.}
Mathematical definition. This is a probabilistic program/game or a module adapter.  Its procedures induce the probability space used by the game-based proof.

\noindent\textbf{Role in the proof.}
Use in this chapter. It supplies an executable component of the formal probabilistic model.

\end{formaldefinition}

\begin{formaldefinition}[EasyCrypt line 2804]
\noindent\textbf{EasyCrypt name.}
\begin{lstlisting}[language=EasyCrypt,basicstyle=\ttfamily\scriptsize]
sign
\end{lstlisting}
\noindent\textbf{Checked EasyCrypt statement.}
\begin{lstlisting}[language=EasyCrypt,basicstyle=\ttfamily\scriptsize]
proc sign(m : message, ctx : context) : signature = {
\end{lstlisting}
\noindent\textbf{Formal mathematical reading.}
Mathematical definition. This procedure is a command in the probabilistic program.  Its return value, state updates, and oracle calls are the operational content of the game.

\noindent\textbf{Role in the proof.}
Use in this chapter. It supplies an executable component of the formal probabilistic model.

\end{formaldefinition}

\begin{formaldefinition}[EasyCrypt line 2843]
\noindent\textbf{EasyCrypt name.}
\begin{lstlisting}[language=EasyCrypt,basicstyle=\ttfamily\scriptsize]
A
\end{lstlisting}
\noindent\textbf{Checked EasyCrypt statement.}
\begin{lstlisting}[language=EasyCrypt,basicstyle=\ttfamily\scriptsize]
module A = A(AH, O)

  proc main() : bool = {
\end{lstlisting}
\noindent\textbf{Formal mathematical reading.}
Mathematical definition. This is a probabilistic program/game or a module adapter.  Its procedures induce the probability space used by the game-based proof.

\noindent\textbf{Role in the proof.}
Use in this chapter. It supplies an executable component of the formal probabilistic model.

\end{formaldefinition}

\begin{formaldefinition}[EasyCrypt line 2874]
\noindent\textbf{EasyCrypt name.}
\begin{lstlisting}[language=EasyCrypt,basicstyle=\ttfamily\scriptsize]
EUF_CMA_ROMProgrammedTranscriptBudgetedSampledSign
\end{lstlisting}
\noindent\textbf{Checked EasyCrypt statement.}
\begin{lstlisting}[language=EasyCrypt,basicstyle=\ttfamily\scriptsize]
module EUF_CMA_ROMProgrammedTranscriptBudgetedSampledSign(
  H : Oracle,
  A : SIG.Adversary,
  Samp : SigningCoinSampler
) = {
\end{lstlisting}
\noindent\textbf{Formal mathematical reading.}
Mathematical definition. This is a probabilistic program/game or a module adapter.  Its procedures induce the probability space used by the game-based proof.

\noindent\textbf{Role in the proof.}
Use in this chapter. It defines the game or game interface whose success event is bounded by later theorems.

\end{formaldefinition}

\begin{formaldefinition}[EasyCrypt line 2879]
\noindent\textbf{EasyCrypt name.}
\begin{lstlisting}[language=EasyCrypt,basicstyle=\ttfamily\scriptsize]
C
\end{lstlisting}
\noindent\textbf{Checked EasyCrypt statement.}
\begin{lstlisting}[language=EasyCrypt,basicstyle=\ttfamily\scriptsize]
module C = CountedLazyROM(H)

  var pk_current : pkey
  var sk_current : skey
  var queries : SIG.query list
  var transcripts : transcript list
  var records : signing_transcript_record list
  var adversary_hash_count : int
  var signing_count : int

  module AH = {
\end{lstlisting}
\noindent\textbf{Formal mathematical reading.}
Mathematical definition. This is a probabilistic program/game or a module adapter.  Its procedures induce the probability space used by the game-based proof.

\noindent\textbf{Role in the proof.}
Use in this chapter. It supplies an executable component of the formal probabilistic model.

\end{formaldefinition}

\begin{formaldefinition}[EasyCrypt line 2890]
\noindent\textbf{EasyCrypt name.}
\begin{lstlisting}[language=EasyCrypt,basicstyle=\ttfamily\scriptsize]
get
\end{lstlisting}
\noindent\textbf{Checked EasyCrypt statement.}
\begin{lstlisting}[language=EasyCrypt,basicstyle=\ttfamily\scriptsize]
proc get(q : ro_query) : ro_output = {
\end{lstlisting}
\noindent\textbf{Formal mathematical reading.}
Mathematical definition. This procedure is a command in the probabilistic program.  Its return value, state updates, and oracle calls are the operational content of the game.

\noindent\textbf{Role in the proof.}
Use in this chapter. It supplies an executable component of the formal probabilistic model.

\end{formaldefinition}

\begin{formaldefinition}[EasyCrypt line 2908]
\noindent\textbf{EasyCrypt name.}
\begin{lstlisting}[language=EasyCrypt,basicstyle=\ttfamily\scriptsize]
O
\end{lstlisting}
\noindent\textbf{Checked EasyCrypt statement.}
\begin{lstlisting}[language=EasyCrypt,basicstyle=\ttfamily\scriptsize]
module O = {
\end{lstlisting}
\noindent\textbf{Formal mathematical reading.}
Mathematical definition. This is a probabilistic program/game or a module adapter.  Its procedures induce the probability space used by the game-based proof.

\noindent\textbf{Role in the proof.}
Use in this chapter. It supplies an executable component of the formal probabilistic model.

\end{formaldefinition}

\begin{formaldefinition}[EasyCrypt line 2909]
\noindent\textbf{EasyCrypt name.}
\begin{lstlisting}[language=EasyCrypt,basicstyle=\ttfamily\scriptsize]
sign
\end{lstlisting}
\noindent\textbf{Checked EasyCrypt statement.}
\begin{lstlisting}[language=EasyCrypt,basicstyle=\ttfamily\scriptsize]
proc sign(m : message, ctx : context) : signature = {
\end{lstlisting}
\noindent\textbf{Formal mathematical reading.}
Mathematical definition. This procedure is a command in the probabilistic program.  Its return value, state updates, and oracle calls are the operational content of the game.

\noindent\textbf{Role in the proof.}
Use in this chapter. It supplies an executable component of the formal probabilistic model.

\end{formaldefinition}

\begin{formaldefinition}[EasyCrypt line 2948]
\noindent\textbf{EasyCrypt name.}
\begin{lstlisting}[language=EasyCrypt,basicstyle=\ttfamily\scriptsize]
A
\end{lstlisting}
\noindent\textbf{Checked EasyCrypt statement.}
\begin{lstlisting}[language=EasyCrypt,basicstyle=\ttfamily\scriptsize]
module A = A(AH, O)

  proc main() : bool = {
\end{lstlisting}
\noindent\textbf{Formal mathematical reading.}
Mathematical definition. This is a probabilistic program/game or a module adapter.  Its procedures induce the probability space used by the game-based proof.

\noindent\textbf{Role in the proof.}
Use in this chapter. It supplies an executable component of the formal probabilistic model.

\end{formaldefinition}

\begin{formaldefinition}[EasyCrypt line 6615]
\noindent\textbf{EasyCrypt name.}
\begin{lstlisting}[language=EasyCrypt,basicstyle=\ttfamily\scriptsize]
structural_to_exact_hyperball_paper_sample_loss_obligation
\end{lstlisting}
\noindent\textbf{Checked EasyCrypt statement.}
\begin{lstlisting}[language=EasyCrypt,basicstyle=\ttfamily\scriptsize]
op structural_to_exact_hyperball_paper_sample_loss_obligation
   (md : mode) : bool =
  forall (sk : skey) (m : message) (ctx : context)
         (p : paper_sim_signature_sample -> bool),
    mu (dsigning_attempt_state md sk m ctx)
      (fun st => p (paper_sim_sample_from_rejection_attempt st)) <=
    mu (dexact_hyperball_signing_attempt_state md sk m ctx)
      (fun st => p (paper_sim_sample_from_rejection_attempt st)) +
      rejection_sampling_loss_term.
\end{lstlisting}
\noindent\textbf{Formal mathematical reading.}
Mathematical definition. This is a Boolean predicate.  The exact displayed EasyCrypt statement gives its arguments; mathematically it is an event, invariant, support condition, or well-formedness predicate over those arguments.

\noindent\textbf{Role in the proof.}
Use in this chapter. It names a numerical term that appears in a game-hop or final security inequality.

\end{formaldefinition}

\begin{formaldefinition}[EasyCrypt line 8424]
\noindent\textbf{EasyCrypt name.}
\begin{lstlisting}[language=EasyCrypt,basicstyle=\ttfamily\scriptsize]
ForkingAdversary
\end{lstlisting}
\noindent\textbf{Checked EasyCrypt statement.}
\begin{lstlisting}[language=EasyCrypt,basicstyle=\ttfamily\scriptsize]
module type ForkingAdversary(H : SIG.POracle) = {
\end{lstlisting}
\noindent\textbf{Formal mathematical reading.}
Mathematical definition. This is an interface for an oracle, adversary, scheme, sampler, solver, or reduction.  A later theorem quantified over this interface is universally quantified over all modules implementing these procedures.

\noindent\textbf{Role in the proof.}
Use in this chapter. It supplies an executable component of the formal probabilistic model.

\end{formaldefinition}

\begin{formaldefinition}[EasyCrypt line 8425]
\noindent\textbf{EasyCrypt name.}
\begin{lstlisting}[language=EasyCrypt,basicstyle=\ttfamily\scriptsize]
fork
\end{lstlisting}
\noindent\textbf{Checked EasyCrypt statement.}
\begin{lstlisting}[language=EasyCrypt,basicstyle=\ttfamily\scriptsize]
proc fork(pk : pkey) : transcript * transcript
}.
\end{lstlisting}
\noindent\textbf{Formal mathematical reading.}
Mathematical definition. This procedure is a command in the probabilistic program.  Its return value, state updates, and oracle calls are the operational content of the game.

\noindent\textbf{Role in the proof.}
Use in this chapter. It supplies an executable component of the formal probabilistic model.

\end{formaldefinition}

\begin{formaldefinition}[EasyCrypt line 8428]
\noindent\textbf{EasyCrypt name.}
\begin{lstlisting}[language=EasyCrypt,basicstyle=\ttfamily\scriptsize]
TranscriptExtractor
\end{lstlisting}
\noindent\textbf{Checked EasyCrypt statement.}
\begin{lstlisting}[language=EasyCrypt,basicstyle=\ttfamily\scriptsize]
module type TranscriptExtractor = {
\end{lstlisting}
\noindent\textbf{Formal mathematical reading.}
Mathematical definition. This is an interface for an oracle, adversary, scheme, sampler, solver, or reduction.  A later theorem quantified over this interface is universally quantified over all modules implementing these procedures.

\noindent\textbf{Role in the proof.}
Use in this chapter. It supplies an executable component of the formal probabilistic model.

\end{formaldefinition}

\begin{formaldefinition}[EasyCrypt line 8429]
\noindent\textbf{EasyCrypt name.}
\begin{lstlisting}[language=EasyCrypt,basicstyle=\ttfamily\scriptsize]
extract
\end{lstlisting}
\noindent\textbf{Checked EasyCrypt statement.}
\begin{lstlisting}[language=EasyCrypt,basicstyle=\ttfamily\scriptsize]
proc extract(tr1 : transcript, tr2 : transcript) :
    module_sis_solution option
}.
\end{lstlisting}
\noindent\textbf{Formal mathematical reading.}
Mathematical definition. This procedure is a command in the probabilistic program.  Its return value, state updates, and oracle calls are the operational content of the game.

\noindent\textbf{Role in the proof.}
Use in this chapter. It supplies an executable component of the formal probabilistic model.

\end{formaldefinition}

\begin{formaldefinition}[EasyCrypt line 8433]
\noindent\textbf{EasyCrypt name.}
\begin{lstlisting}[language=EasyCrypt,basicstyle=\ttfamily\scriptsize]
BimodalTranscriptExtractor
\end{lstlisting}
\noindent\textbf{Checked EasyCrypt statement.}
\begin{lstlisting}[language=EasyCrypt,basicstyle=\ttfamily\scriptsize]
module type BimodalTranscriptExtractor = {
\end{lstlisting}
\noindent\textbf{Formal mathematical reading.}
Mathematical definition. This is an interface for an oracle, adversary, scheme, sampler, solver, or reduction.  A later theorem quantified over this interface is universally quantified over all modules implementing these procedures.

\noindent\textbf{Role in the proof.}
Use in this chapter. It supplies an executable component of the formal probabilistic model.

\end{formaldefinition}

\begin{formaldefinition}[EasyCrypt line 8434]
\noindent\textbf{EasyCrypt name.}
\begin{lstlisting}[language=EasyCrypt,basicstyle=\ttfamily\scriptsize]
extract
\end{lstlisting}
\noindent\textbf{Checked EasyCrypt statement.}
\begin{lstlisting}[language=EasyCrypt,basicstyle=\ttfamily\scriptsize]
proc extract(tr1 : transcript, tr2 : transcript) :
    bimodal_selftarget_msis_solution option
}.
\end{lstlisting}
\noindent\textbf{Formal mathematical reading.}
Mathematical definition. This procedure is a command in the probabilistic program.  Its return value, state updates, and oracle calls are the operational content of the game.

\noindent\textbf{Role in the proof.}
Use in this chapter. It supplies an executable component of the formal probabilistic model.

\end{formaldefinition}

\begin{formaldefinition}[EasyCrypt line 8438]
\noindent\textbf{EasyCrypt name.}
\begin{lstlisting}[language=EasyCrypt,basicstyle=\ttfamily\scriptsize]
BimodalSpecialSoundness_Game
\end{lstlisting}
\noindent\textbf{Checked EasyCrypt statement.}
\begin{lstlisting}[language=EasyCrypt,basicstyle=\ttfamily\scriptsize]
module BimodalSpecialSoundness_Game(
  H : SIG.Oracle,
  S : SIG.Scheme,
  F : ForkingAdversary,
  X : BimodalTranscriptExtractor
) = {
\end{lstlisting}
\noindent\textbf{Formal mathematical reading.}
Mathematical definition. This is a probabilistic program/game or a module adapter.  Its procedures induce the probability space used by the game-based proof.

\noindent\textbf{Role in the proof.}
Use in this chapter. It defines the game or game interface whose success event is bounded by later theorems.

\end{formaldefinition}

\begin{formaldefinition}[EasyCrypt line 8444]
\noindent\textbf{EasyCrypt name.}
\begin{lstlisting}[language=EasyCrypt,basicstyle=\ttfamily\scriptsize]
main
\end{lstlisting}
\noindent\textbf{Checked EasyCrypt statement.}
\begin{lstlisting}[language=EasyCrypt,basicstyle=\ttfamily\scriptsize]
proc main(md : mode) : bool = {
\end{lstlisting}
\noindent\textbf{Formal mathematical reading.}
Mathematical definition. This procedure is a command in the probabilistic program.  Its return value, state updates, and oracle calls are the operational content of the game.

\noindent\textbf{Role in the proof.}
Use in this chapter. It supplies an executable component of the formal probabilistic model.

\end{formaldefinition}

\begin{formaldefinition}[EasyCrypt line 8463]
\noindent\textbf{EasyCrypt name.}
\begin{lstlisting}[language=EasyCrypt,basicstyle=\ttfamily\scriptsize]
SpecialSoundness_Game
\end{lstlisting}
\noindent\textbf{Checked EasyCrypt statement.}
\begin{lstlisting}[language=EasyCrypt,basicstyle=\ttfamily\scriptsize]
module SpecialSoundness_Game(
  H : SIG.Oracle,
  S : SIG.Scheme,
  F : ForkingAdversary,
  X : TranscriptExtractor
) = {
\end{lstlisting}
\noindent\textbf{Formal mathematical reading.}
Mathematical definition. This is a probabilistic program/game or a module adapter.  Its procedures induce the probability space used by the game-based proof.

\noindent\textbf{Role in the proof.}
Use in this chapter. It defines the game or game interface whose success event is bounded by later theorems.

\end{formaldefinition}

\begin{formaldefinition}[EasyCrypt line 8469]
\noindent\textbf{EasyCrypt name.}
\begin{lstlisting}[language=EasyCrypt,basicstyle=\ttfamily\scriptsize]
main
\end{lstlisting}
\noindent\textbf{Checked EasyCrypt statement.}
\begin{lstlisting}[language=EasyCrypt,basicstyle=\ttfamily\scriptsize]
proc main(md : mode) : bool = {
\end{lstlisting}
\noindent\textbf{Formal mathematical reading.}
Mathematical definition. This procedure is a command in the probabilistic program.  Its return value, state updates, and oracle calls are the operational content of the game.

\noindent\textbf{Role in the proof.}
Use in this chapter. It supplies an executable component of the formal probabilistic model.

\end{formaldefinition}

\begin{formaldefinition}[EasyCrypt line 8488]
\noindent\textbf{EasyCrypt name.}
\begin{lstlisting}[language=EasyCrypt,basicstyle=\ttfamily\scriptsize]
BimodalTranscriptExtractor_As_ModuleSIS
\end{lstlisting}
\noindent\textbf{Checked EasyCrypt statement.}
\begin{lstlisting}[language=EasyCrypt,basicstyle=\ttfamily\scriptsize]
module BimodalTranscriptExtractor_As_ModuleSIS(
  X : BimodalTranscriptExtractor
) = {
\end{lstlisting}
\noindent\textbf{Formal mathematical reading.}
Mathematical definition. This is a probabilistic program/game or a module adapter.  Its procedures induce the probability space used by the game-based proof.

\noindent\textbf{Role in the proof.}
Use in this chapter. It supplies an executable component of the formal probabilistic model.

\end{formaldefinition}

\begin{formaldefinition}[EasyCrypt line 8491]
\noindent\textbf{EasyCrypt name.}
\begin{lstlisting}[language=EasyCrypt,basicstyle=\ttfamily\scriptsize]
extract
\end{lstlisting}
\noindent\textbf{Checked EasyCrypt statement.}
\begin{lstlisting}[language=EasyCrypt,basicstyle=\ttfamily\scriptsize]
proc extract(tr1 : transcript, tr2 : transcript) :
    module_sis_solution option = {
\end{lstlisting}
\noindent\textbf{Formal mathematical reading.}
Mathematical definition. This procedure is a command in the probabilistic program.  Its return value, state updates, and oracle calls are the operational content of the game.

\noindent\textbf{Role in the proof.}
Use in this chapter. It supplies an executable component of the formal probabilistic model.

\end{formaldefinition}

\begin{formaldefinition}[EasyCrypt line 8535]
\noindent\textbf{EasyCrypt name.}
\begin{lstlisting}[language=EasyCrypt,basicstyle=\ttfamily\scriptsize]
BimodalToModuleSIS
\end{lstlisting}
\noindent\textbf{Checked EasyCrypt statement.}
\begin{lstlisting}[language=EasyCrypt,basicstyle=\ttfamily\scriptsize]
module type BimodalToModuleSIS = {
\end{lstlisting}
\noindent\textbf{Formal mathematical reading.}
Mathematical definition. This is an interface for an oracle, adversary, scheme, sampler, solver, or reduction.  A later theorem quantified over this interface is universally quantified over all modules implementing these procedures.

\noindent\textbf{Role in the proof.}
Use in this chapter. It supplies an executable component of the formal probabilistic model.

\end{formaldefinition}

\begin{formaldefinition}[EasyCrypt line 8536]
\noindent\textbf{EasyCrypt name.}
\begin{lstlisting}[language=EasyCrypt,basicstyle=\ttfamily\scriptsize]
lift
\end{lstlisting}
\noindent\textbf{Checked EasyCrypt statement.}
\begin{lstlisting}[language=EasyCrypt,basicstyle=\ttfamily\scriptsize]
proc lift(x : bimodal_selftarget_msis_instance,
            sol : bimodal_selftarget_msis_solution) :
    module_sis_solution
}.
\end{lstlisting}
\noindent\textbf{Formal mathematical reading.}
Mathematical definition. This procedure is a command in the probabilistic program.  Its return value, state updates, and oracle calls are the operational content of the game.

\noindent\textbf{Role in the proof.}
Use in this chapter. It supplies an executable component of the formal probabilistic model.

\end{formaldefinition}

\begin{formaldefinition}[EasyCrypt line 8541]
\noindent\textbf{EasyCrypt name.}
\begin{lstlisting}[language=EasyCrypt,basicstyle=\ttfamily\scriptsize]
special_soundness_interfaces_ready
\end{lstlisting}
\noindent\textbf{Checked EasyCrypt statement.}
\begin{lstlisting}[language=EasyCrypt,basicstyle=\ttfamily\scriptsize]
op special_soundness_interfaces_ready (md : mode) : bool =
  fork_public_reconstruction_obligation md /\
  forking_extraction_obligation md /\
  bimodal_to_module_sis_lift_obligation md.
\end{lstlisting}
\noindent\textbf{Formal mathematical reading.}
Mathematical definition. This is a Boolean predicate.  The exact displayed EasyCrypt statement gives its arguments; mathematically it is an event, invariant, support condition, or well-formedness predicate over those arguments.

\noindent\textbf{Role in the proof.}
Use in this chapter. It is part of the exact formal vocabulary used by subsequent lemmas and games.

\end{formaldefinition}

\begin{formaldefinition}[EasyCrypt line 8587]
\noindent\textbf{EasyCrypt name.}
\begin{lstlisting}[language=EasyCrypt,basicstyle=\ttfamily\scriptsize]
signing_simulator_sound
\end{lstlisting}
\noindent\textbf{Checked EasyCrypt statement.}
\begin{lstlisting}[language=EasyCrypt,basicstyle=\ttfamily\scriptsize]
op signing_simulator_sound (md : mode) : bool =
  forall (pk : pkey) (m : message) (ctx : context) (sig : signature),
    signing_simulation_relation md pk m ctx sig =>
    verify_internal md pk m ctx sig.
\end{lstlisting}
\noindent\textbf{Formal mathematical reading.}
Mathematical definition. This is a Boolean predicate.  The exact displayed EasyCrypt statement gives its arguments; mathematically it is an event, invariant, support condition, or well-formedness predicate over those arguments.

\noindent\textbf{Role in the proof.}
Use in this chapter. It is part of the exact formal vocabulary used by subsequent lemmas and games.

\end{formaldefinition}

\begin{formaldefinition}[EasyCrypt line 8615]
\noindent\textbf{EasyCrypt name.}
\begin{lstlisting}[language=EasyCrypt,basicstyle=\ttfamily\scriptsize]
fs_with_aborts_interfaces_ready
\end{lstlisting}
\noindent\textbf{Checked EasyCrypt statement.}
\begin{lstlisting}[language=EasyCrypt,basicstyle=\ttfamily\scriptsize]
op fs_with_aborts_interfaces_ready (md : mode) : bool =
  signing_simulator_sound md /\
  special_soundness_interfaces_ready md /\
  fs_with_aborts_bound_obligation.
\end{lstlisting}
\noindent\textbf{Formal mathematical reading.}
Mathematical definition. This is a Boolean predicate.  The exact displayed EasyCrypt statement gives its arguments; mathematically it is an event, invariant, support condition, or well-formedness predicate over those arguments.

\noindent\textbf{Role in the proof.}
Use in this chapter. It is part of the exact formal vocabulary used by subsequent lemmas and games.

\end{formaldefinition}

\subsubsection*{Lemmas, Theorems, Relational Equivalences, and Their Proof Dependencies}
\begin{formaltheorem}[EasyCrypt line 88]
\noindent\textbf{EasyCrypt name.}
\begin{lstlisting}[language=EasyCrypt,basicstyle=\ttfamily\scriptsize]
transcript_log_consistent_nil
\end{lstlisting}
\noindent\textbf{Checked EasyCrypt statement.}
\begin{lstlisting}[language=EasyCrypt,basicstyle=\ttfamily\scriptsize]
lemma transcript_log_consistent_nil md pk :
  transcript_log_consistent md pk [] [] [].
\end{lstlisting}
\noindent\textbf{Formal mathematical reading.}
Mathematical theorem. This lemma is a universally quantified proposition over the variables shown in the EasyCrypt statement.

\noindent\textbf{Role in the proof.}
Use in this chapter. It is a reusable checked theorem cited by later proof scripts.

\noindent\textbf{Proof-script interpretation.}
No proof block was collected for this item.

\end{formaltheorem}

\begin{formaltheorem}[EasyCrypt line 95]
\noindent\textbf{EasyCrypt name.}
\begin{lstlisting}[language=EasyCrypt,basicstyle=\ttfamily\scriptsize]
transcript_log_consistent_cons
\end{lstlisting}
\noindent\textbf{Checked EasyCrypt statement.}
\begin{lstlisting}[language=EasyCrypt,basicstyle=\ttfamily\scriptsize]
lemma transcript_log_consistent_cons md pk qs trs rs m ctx sig tr :
  transcript_log_consistent md pk qs trs rs =>
  signing_record_relation md pk (m, ctx, sig, tr) =>
  transcript_log_consistent md pk
    ((m, ctx) :: qs)
    (tr :: trs)
    ((m, ctx, sig, tr) :: rs).
\end{lstlisting}
\noindent\textbf{Formal mathematical reading.}
Mathematical theorem. This is an implication, normally turning one invariant, validity predicate, or proof obligation into another.

\noindent\textbf{Role in the proof.}
Use in this chapter. It is a reusable checked theorem cited by later proof scripts.

\noindent\textbf{Proof-script interpretation.}
Proof-script reading. The machine proof decomposes the theorem into rewrites, program-logic obligations, relational equivalences, and arithmetic side conditions.
\emph{Main tactics visible in the proof script:} \ec{rewrite}, \ec{split}, \ec{apply}.
\emph{Named identifiers syntactically cited in the proof block:}
\begin{lstlisting}[language=EasyCrypt,basicstyle=\ttfamily\scriptsize]
qE
rel
signing_record_context
signing_record_log_sound_cons
signing_record_message
signing_record_queries_cons
signing_record_query
signing_record_transcript
signing_record_transcripts_cons
sound
trE
transcript_log_consistent
\end{lstlisting}

\end{formaltheorem}

\begin{formaltheorem}[EasyCrypt line 114]
\noindent\textbf{EasyCrypt name.}
\begin{lstlisting}[language=EasyCrypt,basicstyle=\ttfamily\scriptsize]
transcript_log_consistent_no_min_entropy
\end{lstlisting}
\noindent\textbf{Checked EasyCrypt statement.}
\begin{lstlisting}[language=EasyCrypt,basicstyle=\ttfamily\scriptsize]
lemma transcript_log_consistent_no_min_entropy md pk qs trs rs :
  transcript_log_consistent md pk qs trs rs =>
  transcript_log_min_entropy_clear md trs.
\end{lstlisting}
\noindent\textbf{Formal mathematical reading.}
Mathematical theorem. This is an implication, normally turning one invariant, validity predicate, or proof obligation into another.

\noindent\textbf{Role in the proof.}
Use in this chapter. It is a reusable checked theorem cited by later proof scripts.

\noindent\textbf{Proof-script interpretation.}
Proof-script reading. The machine proof decomposes the theorem into rewrites, program-logic obligations, relational equivalences, and arithmetic side conditions.
\emph{Main tactics visible in the proof script:} \ec{rewrite}, \ec{apply}.
\emph{Named identifiers syntactically cited in the proof block:}
\begin{lstlisting}[language=EasyCrypt,basicstyle=\ttfamily\scriptsize]
md
pk
rs
signing_record_log_sound_transcripts_no_min_entropy
sound
transcript_log_consistent
trsE
\end{lstlisting}

\end{formaltheorem}

\begin{formaltheorem}[EasyCrypt line 124]
\noindent\textbf{EasyCrypt name.}
\begin{lstlisting}[language=EasyCrypt,basicstyle=\ttfamily\scriptsize]
transcript_log_consistent_valid
\end{lstlisting}
\noindent\textbf{Checked EasyCrypt statement.}
\begin{lstlisting}[language=EasyCrypt,basicstyle=\ttfamily\scriptsize]
lemma transcript_log_consistent_valid md pk qs trs rs :
  transcript_log_consistent md pk qs trs rs =>
  transcript_log_valid md trs.
\end{lstlisting}
\noindent\textbf{Formal mathematical reading.}
Mathematical theorem. This is an implication, normally turning one invariant, validity predicate, or proof obligation into another.

\noindent\textbf{Role in the proof.}
Use in this chapter. It proves a structural correctness fact needed by later extraction or verification arguments.

\noindent\textbf{Proof-script interpretation.}
Proof-script reading. The machine proof decomposes the theorem into rewrites, program-logic obligations, relational equivalences, and arithmetic side conditions.
\emph{Main tactics visible in the proof script:} \ec{rewrite}, \ec{apply}.
\emph{Named identifiers syntactically cited in the proof block:}
\begin{lstlisting}[language=EasyCrypt,basicstyle=\ttfamily\scriptsize]
md
pk
rs
signing_record_log_sound_transcripts_valid
sound
transcript_log_consistent
trsE
\end{lstlisting}

\end{formaltheorem}

\begin{formaltheorem}[EasyCrypt line 168]
\noindent\textbf{EasyCrypt name.}
\begin{lstlisting}[language=EasyCrypt,basicstyle=\ttfamily\scriptsize]
sampler_rom_covered_fresh_after_clean_seed
\end{lstlisting}
\noindent\textbf{Checked EasyCrypt statement.}
\begin{lstlisting}[language=EasyCrypt,basicstyle=\ttfamily\scriptsize]
lemma sampler_rom_covered_fresh_after_clean_seed
    seed_coins rom hash_qs sampler_qs bad :
  sampler_rom_covered rom hash_qs sampler_qs =>
  ! (bad \/
     sampler_expand_query seed_coins \in hash_qs \/
     sampler_expand_query seed_coins \in sampler_qs) =>
  sampler_expand_query seed_coins \notin rom.
\end{lstlisting}
\noindent\textbf{Formal mathematical reading.}
Mathematical theorem. This is an implication, normally turning one invariant, validity predicate, or proof obligation into another.

\noindent\textbf{Role in the proof.}
Use in this chapter. It contributes directly to an advantage bound, bad-event bound, or real-arithmetic composition step.

\noindent\textbf{Proof-script interpretation.}
Proof-script reading. The machine proof decomposes the theorem into rewrites, program-logic obligations, relational equivalences, and arithmetic side conditions.
\emph{Main tactics visible in the proof script:} \ec{rewrite}, \ec{smt}.
\emph{Named identifiers syntactically cited in the proof block:}
\begin{lstlisting}[language=EasyCrypt,basicstyle=\ttfamily\scriptsize]
sampler_rom_covered
\end{lstlisting}

\end{formaltheorem}

\begin{formaltheorem}[EasyCrypt line 179]
\noindent\textbf{EasyCrypt name.}
\begin{lstlisting}[language=EasyCrypt,basicstyle=\ttfamily\scriptsize]
sampler_rom_covered_self_log_preserves
\end{lstlisting}
\noindent\textbf{Checked EasyCrypt statement.}
\begin{lstlisting}[language=EasyCrypt,basicstyle=\ttfamily\scriptsize]
lemma sampler_rom_covered_self_log_preserves
    seed_coins rom hash_qs sampler_qs :
  sampler_rom_covered rom hash_qs sampler_qs =>
  sampler_rom_covered rom hash_qs
    (sampler_expand_query seed_coins :: sampler_qs).
\end{lstlisting}
\noindent\textbf{Formal mathematical reading.}
Mathematical theorem. This is an implication, normally turning one invariant, validity predicate, or proof obligation into another.

\noindent\textbf{Role in the proof.}
Use in this chapter. It contributes directly to an advantage bound, bad-event bound, or real-arithmetic composition step.

\noindent\textbf{Proof-script interpretation.}
Proof-script reading. The machine proof decomposes the theorem into rewrites, program-logic obligations, relational equivalences, and arithmetic side conditions.
\emph{Main tactics visible in the proof script:} \ec{rewrite}, \ec{smt}.
\emph{Named identifiers syntactically cited in the proof block:}
\begin{lstlisting}[language=EasyCrypt,basicstyle=\ttfamily\scriptsize]
sampler_rom_covered
\end{lstlisting}

\end{formaltheorem}

\begin{formaltheorem}[EasyCrypt line 188]
\noindent\textbf{EasyCrypt name.}
\begin{lstlisting}[language=EasyCrypt,basicstyle=\ttfamily\scriptsize]
sampler_rom_covered_old_log_clean_boundary
\end{lstlisting}
\noindent\textbf{Checked EasyCrypt statement.}
\begin{lstlisting}[language=EasyCrypt,basicstyle=\ttfamily\scriptsize]
lemma sampler_rom_covered_old_log_clean_boundary
    seed_coins rom hash_qs sampler_qs bad :
  sampler_rom_covered rom hash_qs sampler_qs =>
  ! (bad \/
     sampler_expand_query seed_coins \in hash_qs \/
     sampler_expand_query seed_coins \in sampler_qs) =>
  sampler_expand_query seed_coins \notin rom /\
  sampler_rom_covered rom hash_qs
    (sampler_expand_query seed_coins :: sampler_qs).
\end{lstlisting}
\noindent\textbf{Formal mathematical reading.}
Mathematical theorem. This is an implication, normally turning one invariant, validity predicate, or proof obligation into another.

\noindent\textbf{Role in the proof.}
Use in this chapter. It contributes directly to an advantage bound, bad-event bound, or real-arithmetic composition step.

\noindent\textbf{Proof-script interpretation.}
Proof-script reading. The machine proof decomposes the theorem into rewrites, program-logic obligations, relational equivalences, and arithmetic side conditions.
\emph{Main tactics visible in the proof script:} \ec{split}, \ec{apply}.
\emph{Named identifiers syntactically cited in the proof block:}
\begin{lstlisting}[language=EasyCrypt,basicstyle=\ttfamily\scriptsize]
bad
clean
covered
hash_qs
rom
sampler_qs
sampler_rom_covered_fresh_after_clean_seed
sampler_rom_covered_self_log_preserves
seed_coins
\end{lstlisting}

\end{formaltheorem}

\begin{formaltheorem}[EasyCrypt line 218]
\noindent\textbf{EasyCrypt name.}
\begin{lstlisting}[language=EasyCrypt,basicstyle=\ttfamily\scriptsize]
transcript_log_consistent_ready
\end{lstlisting}
\noindent\textbf{Checked EasyCrypt statement.}
\begin{lstlisting}[language=EasyCrypt,basicstyle=\ttfamily\scriptsize]
lemma transcript_log_consistent_ready md pk qs trs rs :
  transcript_log_consistent md pk qs trs rs =>
  transcript_log_ready md trs.
\end{lstlisting}
\noindent\textbf{Formal mathematical reading.}
Mathematical theorem. This is an implication, normally turning one invariant, validity predicate, or proof obligation into another.

\noindent\textbf{Role in the proof.}
Use in this chapter. It is a reusable checked theorem cited by later proof scripts.

\noindent\textbf{Proof-script interpretation.}
Proof-script reading. The machine proof decomposes the theorem into rewrites, program-logic obligations, relational equivalences, and arithmetic side conditions.
\emph{Main tactics visible in the proof script:} \ec{rewrite}, \ec{split}, \ec{apply}.
\emph{Named identifiers syntactically cited in the proof block:}
\begin{lstlisting}[language=EasyCrypt,basicstyle=\ttfamily\scriptsize]
consistent
md
pk
qs
rs
transcript_log_consistent_no_min_entropy
transcript_log_consistent_valid
transcript_log_ready
trs
\end{lstlisting}

\end{formaltheorem}

\begin{formaltheorem}[EasyCrypt line 229]
\noindent\textbf{EasyCrypt name.}
\begin{lstlisting}[language=EasyCrypt,basicstyle=\ttfamily\scriptsize]
transcript_log_valid_member
\end{lstlisting}
\noindent\textbf{Checked EasyCrypt statement.}
\begin{lstlisting}[language=EasyCrypt,basicstyle=\ttfamily\scriptsize]
lemma transcript_log_valid_member md trs tr :
  transcript_log_valid md trs =>
  tr \in trs =>
  transcript_valid md tr.
\end{lstlisting}
\noindent\textbf{Formal mathematical reading.}
Mathematical theorem. This is an implication, normally turning one invariant, validity predicate, or proof obligation into another.

\noindent\textbf{Role in the proof.}
Use in this chapter. It proves a structural correctness fact needed by later extraction or verification arguments.

\noindent\textbf{Proof-script interpretation.}
Proof-script reading. The machine proof decomposes the theorem into rewrites, program-logic obligations, relational equivalences, and arithmetic side conditions.
\emph{Main tactics visible in the proof script:} \ec{rewrite}, \ec{apply}.
\emph{Named identifiers syntactically cited in the proof block:}
\begin{lstlisting}[language=EasyCrypt,basicstyle=\ttfamily\scriptsize]
allP
mem_tr
transcript_log_valid
valid_all
\end{lstlisting}

\end{formaltheorem}

\begin{formaltheorem}[EasyCrypt line 239]
\noindent\textbf{EasyCrypt name.}
\begin{lstlisting}[language=EasyCrypt,basicstyle=\ttfamily\scriptsize]
transcript_log_min_entropy_clear_member
\end{lstlisting}
\noindent\textbf{Checked EasyCrypt statement.}
\begin{lstlisting}[language=EasyCrypt,basicstyle=\ttfamily\scriptsize]
lemma transcript_log_min_entropy_clear_member md trs tr :
  transcript_log_min_entropy_clear md trs =>
  tr \in trs =>
  ! min_entropy_failure md tr.
\end{lstlisting}
\noindent\textbf{Formal mathematical reading.}
Mathematical theorem. This is an implication, normally turning one invariant, validity predicate, or proof obligation into another.

\noindent\textbf{Role in the proof.}
Use in this chapter. It contributes directly to an advantage bound, bad-event bound, or real-arithmetic composition step.

\noindent\textbf{Proof-script interpretation.}
Proof-script reading. The machine proof decomposes the theorem into rewrites, program-logic obligations, relational equivalences, and arithmetic side conditions.
\emph{Main tactics visible in the proof script:} \ec{rewrite}, \ec{apply}.
\emph{Named identifiers syntactically cited in the proof block:}
\begin{lstlisting}[language=EasyCrypt,basicstyle=\ttfamily\scriptsize]
allP
clear_all
mem_tr
transcript_log_min_entropy_clear
\end{lstlisting}

\end{formaltheorem}

\begin{formaltheorem}[EasyCrypt line 249]
\noindent\textbf{EasyCrypt name.}
\begin{lstlisting}[language=EasyCrypt,basicstyle=\ttfamily\scriptsize]
transcript_log_ready_member_valid
\end{lstlisting}
\noindent\textbf{Checked EasyCrypt statement.}
\begin{lstlisting}[language=EasyCrypt,basicstyle=\ttfamily\scriptsize]
lemma transcript_log_ready_member_valid md trs tr :
  transcript_log_ready md trs =>
  tr \in trs =>
  transcript_valid md tr.
\end{lstlisting}
\noindent\textbf{Formal mathematical reading.}
Mathematical theorem. This is an implication, normally turning one invariant, validity predicate, or proof obligation into another.

\noindent\textbf{Role in the proof.}
Use in this chapter. It proves a structural correctness fact needed by later extraction or verification arguments.

\noindent\textbf{Proof-script interpretation.}
Proof-script reading. The machine proof decomposes the theorem into rewrites, program-logic obligations, relational equivalences, and arithmetic side conditions.
\emph{Main tactics visible in the proof script:} \ec{rewrite}, \ec{apply}.
\emph{Named identifiers syntactically cited in the proof block:}
\begin{lstlisting}[language=EasyCrypt,basicstyle=\ttfamily\scriptsize]
md
mem_tr
tr
transcript_log_ready
transcript_log_valid_member
trs
valid
\end{lstlisting}

\end{formaltheorem}

\begin{formaltheorem}[EasyCrypt line 259]
\noindent\textbf{EasyCrypt name.}
\begin{lstlisting}[language=EasyCrypt,basicstyle=\ttfamily\scriptsize]
transcript_log_ready_member_no_min_entropy
\end{lstlisting}
\noindent\textbf{Checked EasyCrypt statement.}
\begin{lstlisting}[language=EasyCrypt,basicstyle=\ttfamily\scriptsize]
lemma transcript_log_ready_member_no_min_entropy md trs tr :
  transcript_log_ready md trs =>
  tr \in trs =>
  ! min_entropy_failure md tr.
\end{lstlisting}
\noindent\textbf{Formal mathematical reading.}
Mathematical theorem. This is an implication, normally turning one invariant, validity predicate, or proof obligation into another.

\noindent\textbf{Role in the proof.}
Use in this chapter. It is a reusable checked theorem cited by later proof scripts.

\noindent\textbf{Proof-script interpretation.}
Proof-script reading. The machine proof decomposes the theorem into rewrites, program-logic obligations, relational equivalences, and arithmetic side conditions.
\emph{Main tactics visible in the proof script:} \ec{rewrite}, \ec{apply}.
\emph{Named identifiers syntactically cited in the proof block:}
\begin{lstlisting}[language=EasyCrypt,basicstyle=\ttfamily\scriptsize]
clear_log
md
mem_tr
tr
transcript_log_min_entropy_clear_member
transcript_log_ready
trs
\end{lstlisting}

\end{formaltheorem}

\begin{formaltheorem}[EasyCrypt line 269]
\noindent\textbf{EasyCrypt name.}
\begin{lstlisting}[language=EasyCrypt,basicstyle=\ttfamily\scriptsize]
transcript_log_ready_no_failure_for_matching_site
\end{lstlisting}
\noindent\textbf{Checked EasyCrypt statement.}
\begin{lstlisting}[language=EasyCrypt,basicstyle=\ttfamily\scriptsize]
lemma transcript_log_ready_no_failure_for_matching_site md trs tr site :
  transcript_log_ready md trs =>
  tr \in trs =>
  programming_site_matches_transcript md tr site =>
  ! fs_with_aborts_failure md tr site.
\end{lstlisting}
\noindent\textbf{Formal mathematical reading.}
Mathematical theorem. This is an implication, normally turning one invariant, validity predicate, or proof obligation into another.

\noindent\textbf{Role in the proof.}
Use in this chapter. It is a reusable checked theorem cited by later proof scripts.

\noindent\textbf{Proof-script interpretation.}
Proof-script reading. The machine proof decomposes the theorem into rewrites, program-logic obligations, relational equivalences, and arithmetic side conditions.
\emph{Main tactics visible in the proof script:} \ec{rewrite}, \ec{apply}.
\emph{Named identifiers syntactically cited in the proof block:}
\begin{lstlisting}[language=EasyCrypt,basicstyle=\ttfamily\scriptsize]
fs_with_aborts_failure
match_site
md
mem_tr
ready
reprogramming_failure
tr
transcript_log_ready_member_no_min_entropy
trs
\end{lstlisting}

\end{formaltheorem}

\begin{formaltheorem}[EasyCrypt line 280]
\noindent\textbf{EasyCrypt name.}
\begin{lstlisting}[language=EasyCrypt,basicstyle=\ttfamily\scriptsize]
transcript_log_ready_no_failure_for_programming_site
\end{lstlisting}
\noindent\textbf{Checked EasyCrypt statement.}
\begin{lstlisting}[language=EasyCrypt,basicstyle=\ttfamily\scriptsize]
lemma transcript_log_ready_no_failure_for_programming_site md trs tr y :
  transcript_log_ready md trs =>
  tr \in trs =>
  transcript_challenge_matches tr y =>
  ! fs_with_aborts_failure md tr
      (programming_site_of_transcript md tr y).
\end{lstlisting}
\noindent\textbf{Formal mathematical reading.}
Mathematical theorem. This is an implication, normally turning one invariant, validity predicate, or proof obligation into another.

\noindent\textbf{Role in the proof.}
Use in this chapter. It is a reusable checked theorem cited by later proof scripts.

\noindent\textbf{Proof-script interpretation.}
Proof-script reading. The machine proof decomposes the theorem into rewrites, program-logic obligations, relational equivalences, and arithmetic side conditions.
\emph{Main tactics visible in the proof script:} \ec{apply}.
\emph{Named identifiers syntactically cited in the proof block:}
\begin{lstlisting}[language=EasyCrypt,basicstyle=\ttfamily\scriptsize]
match_y
md
mem_tr
programming_site_of_transcript
programming_site_self_matches
ready
tr
transcript_log_ready_no_failure_for_matching_site
trs
\end{lstlisting}

\end{formaltheorem}

\begin{formaltheorem}[EasyCrypt line 293]
\noindent\textbf{EasyCrypt name.}
\begin{lstlisting}[language=EasyCrypt,basicstyle=\ttfamily\scriptsize]
transcript_log_ready_no_failure_for_signature_programming_site
\end{lstlisting}
\noindent\textbf{Checked EasyCrypt statement.}
\begin{lstlisting}[language=EasyCrypt,basicstyle=\ttfamily\scriptsize]
lemma transcript_log_ready_no_failure_for_signature_programming_site
  md trs tr y :
  transcript_log_ready md trs =>
  tr \in trs =>
  transcript_signature_challenge_matches tr y =>
  ! fs_with_aborts_failure md tr
      (signature_programming_site_of_transcript md tr y).
\end{lstlisting}
\noindent\textbf{Formal mathematical reading.}
Mathematical theorem. This is an implication, normally turning one invariant, validity predicate, or proof obligation into another.

\noindent\textbf{Role in the proof.}
Use in this chapter. It is a reusable checked theorem cited by later proof scripts.

\noindent\textbf{Proof-script interpretation.}
Proof-script reading. The machine proof decomposes the theorem into rewrites, program-logic obligations, relational equivalences, and arithmetic side conditions.
\emph{Main tactics visible in the proof script:} \ec{have}, \ec{apply}.
\emph{Named identifiers syntactically cited in the proof block:}
\begin{lstlisting}[language=EasyCrypt,basicstyle=\ttfamily\scriptsize]
match_y
md
mem_tr
ready
signature_programming_site_of_transcript
tr
transcript_log_ready_member_valid
transcript_log_ready_no_failure_for_matching_site
transcript_valid_signature_programming_site_matches
trs
valid
\end{lstlisting}

\end{formaltheorem}

\begin{formaltheorem}[EasyCrypt line 308]
\noindent\textbf{EasyCrypt name.}
\begin{lstlisting}[language=EasyCrypt,basicstyle=\ttfamily\scriptsize]
transcript_log_ready_actual_signature_site_clear_member
\end{lstlisting}
\noindent\textbf{Checked EasyCrypt statement.}
\begin{lstlisting}[language=EasyCrypt,basicstyle=\ttfamily\scriptsize]
lemma transcript_log_ready_actual_signature_site_clear_member md trs tr :
  transcript_log_ready md trs =>
  tr \in trs =>
  transcript_signature_programming_site_clear md tr.
\end{lstlisting}
\noindent\textbf{Formal mathematical reading.}
Mathematical theorem. This is an implication, normally turning one invariant, validity predicate, or proof obligation into another.

\noindent\textbf{Role in the proof.}
Use in this chapter. It contributes directly to an advantage bound, bad-event bound, or real-arithmetic composition step.

\noindent\textbf{Proof-script interpretation.}
Proof-script reading. The machine proof decomposes the theorem into rewrites, program-logic obligations, relational equivalences, and arithmetic side conditions.
\emph{Main tactics visible in the proof script:} \ec{rewrite}, \ec{apply}.
\emph{Named identifiers syntactically cited in the proof block:}
\begin{lstlisting}[language=EasyCrypt,basicstyle=\ttfamily\scriptsize]
actual_signature_programming_site
md
mem_tr
ready
tr
transcript_log_ready_no_failure_for_signature_programming_site
transcript_signature_challenge_output
transcript_signature_challenge_output_matches
transcript_signature_programming_site_clear
trs
\end{lstlisting}

\end{formaltheorem}

\begin{formaltheorem}[EasyCrypt line 322]
\noindent\textbf{EasyCrypt name.}
\begin{lstlisting}[language=EasyCrypt,basicstyle=\ttfamily\scriptsize]
transcript_log_ready_signature_sites_clear
\end{lstlisting}
\noindent\textbf{Checked EasyCrypt statement.}
\begin{lstlisting}[language=EasyCrypt,basicstyle=\ttfamily\scriptsize]
lemma transcript_log_ready_signature_sites_clear md trs :
  transcript_log_ready md trs =>
  transcript_log_signature_programming_sites_clear md trs.
\end{lstlisting}
\noindent\textbf{Formal mathematical reading.}
Mathematical theorem. This is an implication, normally turning one invariant, validity predicate, or proof obligation into another.

\noindent\textbf{Role in the proof.}
Use in this chapter. It contributes directly to an advantage bound, bad-event bound, or real-arithmetic composition step.

\noindent\textbf{Proof-script interpretation.}
Proof-script reading. The machine proof decomposes the theorem into rewrites, program-logic obligations, relational equivalences, and arithmetic side conditions.
\emph{Main tactics visible in the proof script:} \ec{rewrite}, \ec{apply}.
\emph{Named identifiers syntactically cited in the proof block:}
\begin{lstlisting}[language=EasyCrypt,basicstyle=\ttfamily\scriptsize]
allP
md
mem_tr
ready
tr
transcript_log_ready_actual_signature_site_clear_member
transcript_log_signature_programming_sites_clear
trs
\end{lstlisting}

\end{formaltheorem}

\begin{formaltheorem}[EasyCrypt line 338]
\noindent\textbf{EasyCrypt name.}
\begin{lstlisting}[language=EasyCrypt,basicstyle=\ttfamily\scriptsize]
internal_transcript_state_sound_nil
\end{lstlisting}
\noindent\textbf{Checked EasyCrypt statement.}
\begin{lstlisting}[language=EasyCrypt,basicstyle=\ttfamily\scriptsize]
lemma internal_transcript_state_sound_nil md sk :
  internal_transcript_state_sound md
    (public_key_of_secret md sk) sk [] [] [].
\end{lstlisting}
\noindent\textbf{Formal mathematical reading.}
Mathematical theorem. This lemma is a universally quantified proposition over the variables shown in the EasyCrypt statement.

\noindent\textbf{Role in the proof.}
Use in this chapter. It preserves an invariant across an oracle call, adversary call, or game procedure.

\noindent\textbf{Proof-script interpretation.}
Proof-script reading. The machine proof decomposes the theorem into rewrites, program-logic obligations, relational equivalences, and arithmetic side conditions.
\emph{Main tactics visible in the proof script:} \ec{rewrite}, \ec{split}, \ec{apply}.
\emph{Named identifiers syntactically cited in the proof block:}
\begin{lstlisting}[language=EasyCrypt,basicstyle=\ttfamily\scriptsize]
internal_transcript_state_sound
transcript_log_consistent_nil
trivial
\end{lstlisting}

\end{formaltheorem}

\begin{formaltheorem}[EasyCrypt line 348]
\noindent\textbf{EasyCrypt name.}
\begin{lstlisting}[language=EasyCrypt,basicstyle=\ttfamily\scriptsize]
internal_transcript_state_sound_keygen
\end{lstlisting}
\noindent\textbf{Checked EasyCrypt statement.}
\begin{lstlisting}[language=EasyCrypt,basicstyle=\ttfamily\scriptsize]
lemma internal_transcript_state_sound_keygen md sd :
  internal_transcript_state_sound md
    (keygen_internal md sd).`1
    (keygen_internal md sd).`2
    [] [] [].
\end{lstlisting}
\noindent\textbf{Formal mathematical reading.}
Mathematical theorem. This lemma is a universally quantified proposition over the variables shown in the EasyCrypt statement.

\noindent\textbf{Role in the proof.}
Use in this chapter. It preserves an invariant across an oracle call, adversary call, or game procedure.

\noindent\textbf{Proof-script interpretation.}
Proof-script reading. The machine proof decomposes the theorem into rewrites, program-logic obligations, relational equivalences, and arithmetic side conditions.
\emph{Main tactics visible in the proof script:} \ec{rewrite}, \ec{split}, \ec{apply}.
\emph{Named identifiers syntactically cited in the proof block:}
\begin{lstlisting}[language=EasyCrypt,basicstyle=\ttfamily\scriptsize]
internal_transcript_state_sound
keygen_internal
transcript_log_consistent_nil
\end{lstlisting}

\end{formaltheorem}

\begin{formaltheorem}[EasyCrypt line 360]
\noindent\textbf{EasyCrypt name.}
\begin{lstlisting}[language=EasyCrypt,basicstyle=\ttfamily\scriptsize]
internal_transcript_state_sound_no_min_entropy
\end{lstlisting}
\noindent\textbf{Checked EasyCrypt statement.}
\begin{lstlisting}[language=EasyCrypt,basicstyle=\ttfamily\scriptsize]
lemma internal_transcript_state_sound_no_min_entropy
  md pk sk qs trs rs :
  internal_transcript_state_sound md pk sk qs trs rs =>
  transcript_log_min_entropy_clear md trs.
\end{lstlisting}
\noindent\textbf{Formal mathematical reading.}
Mathematical theorem. This is an implication, normally turning one invariant, validity predicate, or proof obligation into another.

\noindent\textbf{Role in the proof.}
Use in this chapter. It preserves an invariant across an oracle call, adversary call, or game procedure.

\noindent\textbf{Proof-script interpretation.}
Proof-script reading. The machine proof decomposes the theorem into rewrites, program-logic obligations, relational equivalences, and arithmetic side conditions.
\emph{Main tactics visible in the proof script:} \ec{rewrite}, \ec{apply}.
\emph{Named identifiers syntactically cited in the proof block:}
\begin{lstlisting}[language=EasyCrypt,basicstyle=\ttfamily\scriptsize]
internal_transcript_state_sound
log_sound
md
pk
qs
rs
transcript_log_consistent_no_min_entropy
trs
\end{lstlisting}

\end{formaltheorem}

\begin{formaltheorem}[EasyCrypt line 370]
\noindent\textbf{EasyCrypt name.}
\begin{lstlisting}[language=EasyCrypt,basicstyle=\ttfamily\scriptsize]
internal_transcript_state_sound_valid
\end{lstlisting}
\noindent\textbf{Checked EasyCrypt statement.}
\begin{lstlisting}[language=EasyCrypt,basicstyle=\ttfamily\scriptsize]
lemma internal_transcript_state_sound_valid md pk sk qs trs rs :
  internal_transcript_state_sound md pk sk qs trs rs =>
  transcript_log_valid md trs.
\end{lstlisting}
\noindent\textbf{Formal mathematical reading.}
Mathematical theorem. This is an implication, normally turning one invariant, validity predicate, or proof obligation into another.

\noindent\textbf{Role in the proof.}
Use in this chapter. It preserves an invariant across an oracle call, adversary call, or game procedure.

\noindent\textbf{Proof-script interpretation.}
Proof-script reading. The machine proof decomposes the theorem into rewrites, program-logic obligations, relational equivalences, and arithmetic side conditions.
\emph{Main tactics visible in the proof script:} \ec{rewrite}, \ec{apply}.
\emph{Named identifiers syntactically cited in the proof block:}
\begin{lstlisting}[language=EasyCrypt,basicstyle=\ttfamily\scriptsize]
internal_transcript_state_sound
log_sound
md
pk
qs
rs
transcript_log_consistent_valid
trs
\end{lstlisting}

\end{formaltheorem}

\begin{formaltheorem}[EasyCrypt line 379]
\noindent\textbf{EasyCrypt name.}
\begin{lstlisting}[language=EasyCrypt,basicstyle=\ttfamily\scriptsize]
internal_transcript_state_sound_log_ready
\end{lstlisting}
\noindent\textbf{Checked EasyCrypt statement.}
\begin{lstlisting}[language=EasyCrypt,basicstyle=\ttfamily\scriptsize]
lemma internal_transcript_state_sound_log_ready md pk sk qs trs rs :
  internal_transcript_state_sound md pk sk qs trs rs =>
  transcript_log_ready md trs.
\end{lstlisting}
\noindent\textbf{Formal mathematical reading.}
Mathematical theorem. This is an implication, normally turning one invariant, validity predicate, or proof obligation into another.

\noindent\textbf{Role in the proof.}
Use in this chapter. It preserves an invariant across an oracle call, adversary call, or game procedure.

\noindent\textbf{Proof-script interpretation.}
Proof-script reading. The machine proof decomposes the theorem into rewrites, program-logic obligations, relational equivalences, and arithmetic side conditions.
\emph{Main tactics visible in the proof script:} \ec{rewrite}, \ec{apply}.
\emph{Named identifiers syntactically cited in the proof block:}
\begin{lstlisting}[language=EasyCrypt,basicstyle=\ttfamily\scriptsize]
internal_transcript_state_sound
log_sound
md
pk
qs
rs
transcript_log_consistent_ready
trs
\end{lstlisting}

\end{formaltheorem}

\begin{formaltheorem}[EasyCrypt line 388]
\noindent\textbf{EasyCrypt name.}
\begin{lstlisting}[language=EasyCrypt,basicstyle=\ttfamily\scriptsize]
transcript_log_consistent_sign_internal_cons
\end{lstlisting}
\noindent\textbf{Checked EasyCrypt statement.}
\begin{lstlisting}[language=EasyCrypt,basicstyle=\ttfamily\scriptsize]
lemma transcript_log_consistent_sign_internal_cons
  md sk qs trs rs m ctx coins :
  transcript_log_consistent md (public_key_of_secret md sk) qs trs rs =>
  transcript_log_consistent md
    (public_key_of_secret md sk)
    ((m, ctx) :: qs)
    (transcript_of_signature md (public_key_of_secret md sk) m ctx
      (sign_internal md sk m ctx coins) :: trs)
    ((m, ctx, sign_internal md sk m ctx coins,
      transcript_of_signature md (public_key_of_secret md sk) m ctx
        (sign_internal md sk m ctx coins)) :: rs).
\end{lstlisting}
\noindent\textbf{Formal mathematical reading.}
Mathematical theorem. This is an implication, normally turning one invariant, validity predicate, or proof obligation into another.

\noindent\textbf{Role in the proof.}
Use in this chapter. It is a reusable checked theorem cited by later proof scripts.

\noindent\textbf{Proof-script interpretation.}
Proof-script reading. The machine proof decomposes the theorem into rewrites, program-logic obligations, relational equivalences, and arithmetic side conditions.
\emph{Main tactics visible in the proof script:} \ec{apply}.
\emph{Named identifiers syntactically cited in the proof block:}
\begin{lstlisting}[language=EasyCrypt,basicstyle=\ttfamily\scriptsize]
log_sound
sign_internal_record_relation
transcript_log_consistent_cons
\end{lstlisting}

\end{formaltheorem}

\begin{formaltheorem}[EasyCrypt line 406]
\noindent\textbf{EasyCrypt name.}
\begin{lstlisting}[language=EasyCrypt,basicstyle=\ttfamily\scriptsize]
internal_transcript_state_sound_sign_internal_cons
\end{lstlisting}
\noindent\textbf{Checked EasyCrypt statement.}
\begin{lstlisting}[language=EasyCrypt,basicstyle=\ttfamily\scriptsize]
lemma internal_transcript_state_sound_sign_internal_cons
  md pk sk qs trs rs m ctx coins :
  internal_transcript_state_sound md pk sk qs trs rs =>
  internal_transcript_state_sound md pk sk
    ((m, ctx) :: qs)
    (transcript_of_signature md pk m ctx
      (sign_internal md sk m ctx coins) :: trs)
    ((m, ctx, sign_internal md sk m ctx coins,
      transcript_of_signature md pk m ctx
        (sign_internal md sk m ctx coins)) :: rs).
\end{lstlisting}
\noindent\textbf{Formal mathematical reading.}
Mathematical theorem. This is an implication, normally turning one invariant, validity predicate, or proof obligation into another.

\noindent\textbf{Role in the proof.}
Use in this chapter. It preserves an invariant across an oracle call, adversary call, or game procedure.

\noindent\textbf{Proof-script interpretation.}
Proof-script reading. The machine proof decomposes the theorem into rewrites, program-logic obligations, relational equivalences, and arithmetic side conditions.
\emph{Main tactics visible in the proof script:} \ec{rewrite}, \ec{split}, \ec{apply}.
\emph{Named identifiers syntactically cited in the proof block:}
\begin{lstlisting}[language=EasyCrypt,basicstyle=\ttfamily\scriptsize]
coins
ctx
internal_transcript_state_sound
log_sound
md
pkE
qs
rs
sk
transcript_log_consistent_sign_internal_cons
trs
\end{lstlisting}

\end{formaltheorem}

\begin{formaltheorem}[EasyCrypt line 428]
\noindent\textbf{EasyCrypt name.}
\begin{lstlisting}[language=EasyCrypt,basicstyle=\ttfamily\scriptsize]
sign_internal_transcript_no_min_entropy
\end{lstlisting}
\noindent\textbf{Checked EasyCrypt statement.}
\begin{lstlisting}[language=EasyCrypt,basicstyle=\ttfamily\scriptsize]
lemma sign_internal_transcript_no_min_entropy md pk sk m ctx coins :
  pk = public_key_of_secret md sk =>
  ! min_entropy_failure md
      (transcript_of_signature md pk m ctx
        (sign_internal md sk m ctx coins)).
\end{lstlisting}
\noindent\textbf{Formal mathematical reading.}
Mathematical theorem. This is an implication, normally turning one invariant, validity predicate, or proof obligation into another.

\noindent\textbf{Role in the proof.}
Use in this chapter. It is a reusable checked theorem cited by later proof scripts.

\noindent\textbf{Proof-script interpretation.}
Proof-script reading. The machine proof decomposes the theorem into rewrites, program-logic obligations, relational equivalences, and arithmetic side conditions.
\emph{Main tactics visible in the proof script:} \ec{rewrite}, \ec{apply}.
\emph{Named identifiers syntactically cited in the proof block:}
\begin{lstlisting}[language=EasyCrypt,basicstyle=\ttfamily\scriptsize]
coins
ctx
md
pkE
public_key_of_secret
sign_internal
sign_internal_transcript_relation
signing_transcript_relation_no_min_entropy_failure
sk
\end{lstlisting}

\end{formaltheorem}

\begin{formaltheorem}[EasyCrypt line 733]
\noindent\textbf{EasyCrypt name.}
\begin{lstlisting}[language=EasyCrypt,basicstyle=\ttfamily\scriptsize]
public_sim_commitment_lowbitsE
\end{lstlisting}
\noindent\textbf{Checked EasyCrypt statement.}
\begin{lstlisting}[language=EasyCrypt,basicstyle=\ttfamily\scriptsize]
lemma public_sim_commitment_lowbitsE md sk m ctx coins :
  public_sim_commitment_lowbits md (public_key_of_secret md sk) m ctx coins =
  commitment_lowbits md sk m ctx coins.
\end{lstlisting}
\noindent\textbf{Formal mathematical reading.}
Mathematical theorem. This is an equality or definitional expansion used for rewriting and normalization.

\noindent\textbf{Role in the proof.}
Use in this chapter. It proves an exact game replacement or simulator equivalence.

\noindent\textbf{Proof-script interpretation.}
Proof-script reading. The machine proof decomposes the theorem into rewrites, program-logic obligations, relational equivalences, and arithmetic side conditions.
\emph{Main tactics visible in the proof script:} \ec{rewrite}.
\emph{Named identifiers syntactically cited in the proof block:}
\begin{lstlisting}[language=EasyCrypt,basicstyle=\ttfamily\scriptsize]
commitment_lowbits
commitment_raw
commitment_source
public_key_of_secret
public_sim_commitment_lowbits
public_sim_source
\end{lstlisting}

\end{formaltheorem}

\begin{formaltheorem}[EasyCrypt line 742]
\noindent\textbf{EasyCrypt name.}
\begin{lstlisting}[language=EasyCrypt,basicstyle=\ttfamily\scriptsize]
public_sim_response_aux_vectorE
\end{lstlisting}
\noindent\textbf{Checked EasyCrypt statement.}
\begin{lstlisting}[language=EasyCrypt,basicstyle=\ttfamily\scriptsize]
lemma public_sim_response_aux_vectorE md sk m ctx coins :
  public_sim_response_aux_vector md (public_key_of_secret md sk) m ctx coins =
  response_aux_vector md sk m ctx coins.
\end{lstlisting}
\noindent\textbf{Formal mathematical reading.}
Mathematical theorem. This is an equality or definitional expansion used for rewriting and normalization.

\noindent\textbf{Role in the proof.}
Use in this chapter. It proves an exact game replacement or simulator equivalence.

\noindent\textbf{Proof-script interpretation.}
Proof-script reading. The machine proof decomposes the theorem into rewrites, program-logic obligations, relational equivalences, and arithmetic side conditions.
\emph{Main tactics visible in the proof script:} \ec{rewrite}.
\emph{Named identifiers syntactically cited in the proof block:}
\begin{lstlisting}[language=EasyCrypt,basicstyle=\ttfamily\scriptsize]
public_key_of_secret
public_sim_response_aux_vector
public_sim_source
response_aux_vector
response_source
\end{lstlisting}

\end{formaltheorem}

\begin{formaltheorem}[EasyCrypt line 750]
\noindent\textbf{EasyCrypt name.}
\begin{lstlisting}[language=EasyCrypt,basicstyle=\ttfamily\scriptsize]
public_sim_commitment_challengeE
\end{lstlisting}
\noindent\textbf{Checked EasyCrypt statement.}
\begin{lstlisting}[language=EasyCrypt,basicstyle=\ttfamily\scriptsize]
lemma public_sim_commitment_challengeE md sk m ctx coins :
  public_sim_commitment_challenge md (public_key_of_secret md sk) m ctx coins =
  commitment_challenge md sk m ctx coins.
\end{lstlisting}
\noindent\textbf{Formal mathematical reading.}
Mathematical theorem. This is an equality or definitional expansion used for rewriting and normalization.

\noindent\textbf{Role in the proof.}
Use in this chapter. It contributes directly to an advantage bound, bad-event bound, or real-arithmetic composition step.

\noindent\textbf{Proof-script interpretation.}
Proof-script reading. The machine proof decomposes the theorem into rewrites, program-logic obligations, relational equivalences, and arithmetic side conditions.
\emph{Main tactics visible in the proof script:} \ec{rewrite}.
\emph{Named identifiers syntactically cited in the proof block:}
\begin{lstlisting}[language=EasyCrypt,basicstyle=\ttfamily\scriptsize]
commitment_challenge
public_sim_commitment_challenge
public_sim_commitment_lowbitsE
public_sim_response_aux_vectorE
\end{lstlisting}

\end{formaltheorem}

\begin{formaltheorem}[EasyCrypt line 758]
\noindent\textbf{EasyCrypt name.}
\begin{lstlisting}[language=EasyCrypt,basicstyle=\ttfamily\scriptsize]
public_sim_commitment_highbitsE
\end{lstlisting}
\noindent\textbf{Checked EasyCrypt statement.}
\begin{lstlisting}[language=EasyCrypt,basicstyle=\ttfamily\scriptsize]
lemma public_sim_commitment_highbitsE md sk m ctx coins :
  public_sim_commitment_highbits md (public_key_of_secret md sk) m ctx coins =
  commitment_highbits md sk m ctx coins.
\end{lstlisting}
\noindent\textbf{Formal mathematical reading.}
Mathematical theorem. This is an equality or definitional expansion used for rewriting and normalization.

\noindent\textbf{Role in the proof.}
Use in this chapter. It proves an exact game replacement or simulator equivalence.

\noindent\textbf{Proof-script interpretation.}
Proof-script reading. The machine proof decomposes the theorem into rewrites, program-logic obligations, relational equivalences, and arithmetic side conditions.
\emph{Main tactics visible in the proof script:} \ec{rewrite}.
\emph{Named identifiers syntactically cited in the proof block:}
\begin{lstlisting}[language=EasyCrypt,basicstyle=\ttfamily\scriptsize]
commitment_highbits
public_sim_commitment_challengeE
public_sim_commitment_highbits
public_sim_response_aux_vectorE
\end{lstlisting}

\end{formaltheorem}

\begin{formaltheorem}[EasyCrypt line 767]
\noindent\textbf{EasyCrypt name.}
\begin{lstlisting}[language=EasyCrypt,basicstyle=\ttfamily\scriptsize]
public_sim_response_vectorE
\end{lstlisting}
\noindent\textbf{Checked EasyCrypt statement.}
\begin{lstlisting}[language=EasyCrypt,basicstyle=\ttfamily\scriptsize]
lemma public_sim_response_vectorE md sk m ctx coins :
  public_sim_response_vector md (public_key_of_secret md sk) m ctx coins =
  response_vector md sk m ctx coins.
\end{lstlisting}
\noindent\textbf{Formal mathematical reading.}
Mathematical theorem. This is an equality or definitional expansion used for rewriting and normalization.

\noindent\textbf{Role in the proof.}
Use in this chapter. It proves an exact game replacement or simulator equivalence.

\noindent\textbf{Proof-script interpretation.}
Proof-script reading. The machine proof decomposes the theorem into rewrites, program-logic obligations, relational equivalences, and arithmetic side conditions.
\emph{Main tactics visible in the proof script:} \ec{rewrite}.
\emph{Named identifiers syntactically cited in the proof block:}
\begin{lstlisting}[language=EasyCrypt,basicstyle=\ttfamily\scriptsize]
public_key_of_secret
public_sim_response_vector
public_sim_source
response_source
response_vector
\end{lstlisting}

\end{formaltheorem}

\begin{formaltheorem}[EasyCrypt line 775]
\noindent\textbf{EasyCrypt name.}
\begin{lstlisting}[language=EasyCrypt,basicstyle=\ttfamily\scriptsize]
public_sim_hint_vectorE
\end{lstlisting}
\noindent\textbf{Checked EasyCrypt statement.}
\begin{lstlisting}[language=EasyCrypt,basicstyle=\ttfamily\scriptsize]
lemma public_sim_hint_vectorE md sk m ctx coins :
  public_sim_hint_vector md (public_key_of_secret md sk) m ctx coins =
  hint_vector md sk m ctx coins.
\end{lstlisting}
\noindent\textbf{Formal mathematical reading.}
Mathematical theorem. This is an equality or definitional expansion used for rewriting and normalization.

\noindent\textbf{Role in the proof.}
Use in this chapter. It proves an exact game replacement or simulator equivalence.

\noindent\textbf{Proof-script interpretation.}
Proof-script reading. The machine proof decomposes the theorem into rewrites, program-logic obligations, relational equivalences, and arithmetic side conditions.
\emph{Main tactics visible in the proof script:} \ec{rewrite}.
\emph{Named identifiers syntactically cited in the proof block:}
\begin{lstlisting}[language=EasyCrypt,basicstyle=\ttfamily\scriptsize]
hint_vector
public_key_of_secret
public_sim_hint_vector
public_sim_source
response_source
\end{lstlisting}

\end{formaltheorem}

\begin{formaltheorem}[EasyCrypt line 783]
\noindent\textbf{EasyCrypt name.}
\begin{lstlisting}[language=EasyCrypt,basicstyle=\ttfamily\scriptsize]
public_sim_signature_tokenE
\end{lstlisting}
\noindent\textbf{Checked EasyCrypt statement.}
\begin{lstlisting}[language=EasyCrypt,basicstyle=\ttfamily\scriptsize]
lemma public_sim_signature_tokenE md sk m ctx coins :
  public_sim_signature_token md (public_key_of_secret md sk) m ctx coins =
  signature_token md sk m ctx coins.
\end{lstlisting}
\noindent\textbf{Formal mathematical reading.}
Mathematical theorem. This is an equality or definitional expansion used for rewriting and normalization.

\noindent\textbf{Role in the proof.}
Use in this chapter. It proves an exact game replacement or simulator equivalence.

\noindent\textbf{Proof-script interpretation.}
Proof-script reading. The machine proof decomposes the theorem into rewrites, program-logic obligations, relational equivalences, and arithmetic side conditions.
\emph{Main tactics visible in the proof script:} \ec{rewrite}.
\emph{Named identifiers syntactically cited in the proof block:}
\begin{lstlisting}[language=EasyCrypt,basicstyle=\ttfamily\scriptsize]
public_sim_hint_vectorE
public_sim_response_aux_vectorE
public_sim_response_vectorE
public_sim_signature_token
signature_token
\end{lstlisting}

\end{formaltheorem}

\begin{formaltheorem}[EasyCrypt line 793]
\noindent\textbf{EasyCrypt name.}
\begin{lstlisting}[language=EasyCrypt,basicstyle=\ttfamily\scriptsize]
public_sim_signatureE
\end{lstlisting}
\noindent\textbf{Checked EasyCrypt statement.}
\begin{lstlisting}[language=EasyCrypt,basicstyle=\ttfamily\scriptsize]
lemma public_sim_signatureE md sk m ctx coins :
  public_sim_signature md (public_key_of_secret md sk) m ctx coins =
  sign_internal md sk m ctx coins.
\end{lstlisting}
\noindent\textbf{Formal mathematical reading.}
Mathematical theorem. This is an equality or definitional expansion used for rewriting and normalization.

\noindent\textbf{Role in the proof.}
Use in this chapter. It proves an exact game replacement or simulator equivalence.

\noindent\textbf{Proof-script interpretation.}
Proof-script reading. The machine proof decomposes the theorem into rewrites, program-logic obligations, relational equivalences, and arithmetic side conditions.
\emph{Main tactics visible in the proof script:} \ec{rewrite}.
\emph{Named identifiers syntactically cited in the proof block:}
\begin{lstlisting}[language=EasyCrypt,basicstyle=\ttfamily\scriptsize]
public_sim_commitment_challengeE
public_sim_commitment_highbitsE
public_sim_commitment_lowbitsE
public_sim_signature
public_sim_signature_tokenE
sign_internal
\end{lstlisting}

\end{formaltheorem}

\begin{formaltheorem}[EasyCrypt line 912]
\noindent\textbf{EasyCrypt name.}
\begin{lstlisting}[language=EasyCrypt,basicstyle=\ttfamily\scriptsize]
paper_sim_commitment_lowbits_wf
\end{lstlisting}
\noindent\textbf{Checked EasyCrypt statement.}
\begin{lstlisting}[language=EasyCrypt,basicstyle=\ttfamily\scriptsize]
lemma paper_sim_commitment_lowbits_wf md s :
  poly_wf (paper_sim_commitment_lowbits md s).
\end{lstlisting}
\noindent\textbf{Formal mathematical reading.}
Mathematical theorem. This lemma is a universally quantified proposition over the variables shown in the EasyCrypt statement.

\noindent\textbf{Role in the proof.}
Use in this chapter. It proves an exact game replacement or simulator equivalence.

\noindent\textbf{Proof-script interpretation.}
Proof-script reading. The machine proof decomposes the theorem into rewrites, program-logic obligations, relational equivalences, and arithmetic side conditions.
\emph{Main tactics visible in the proof script:} \ec{rewrite}.
\emph{Named identifiers syntactically cited in the proof block:}
\begin{lstlisting}[language=EasyCrypt,basicstyle=\ttfamily\scriptsize]
paper_sim_commitment_lowbits
poly_wf
size_mkseq
\end{lstlisting}

\end{formaltheorem}

\begin{formaltheorem}[EasyCrypt line 918]
\noindent\textbf{EasyCrypt name.}
\begin{lstlisting}[language=EasyCrypt,basicstyle=\ttfamily\scriptsize]
paper_sim_signature_valid
\end{lstlisting}
\noindent\textbf{Checked EasyCrypt statement.}
\begin{lstlisting}[language=EasyCrypt,basicstyle=\ttfamily\scriptsize]
lemma paper_sim_signature_valid md pk m ctx s :
  paper_sim_signature_sample_wf md s =>
  valid_signature md (paper_sim_signature md pk m ctx s).
\end{lstlisting}
\noindent\textbf{Formal mathematical reading.}
Mathematical theorem. This is an implication, normally turning one invariant, validity predicate, or proof obligation into another.

\noindent\textbf{Role in the proof.}
Use in this chapter. It proves an exact game replacement or simulator equivalence.

\noindent\textbf{Proof-script interpretation.}
Proof-script reading. The machine proof decomposes the theorem into rewrites, program-logic obligations, relational equivalences, and arithmetic side conditions.
\emph{Main tactics visible in the proof script:} \ec{rewrite}, \ec{smt}.
\emph{Named identifiers syntactically cited in the proof block:}
\begin{lstlisting}[language=EasyCrypt,basicstyle=\ttfamily\scriptsize]
aux_wf
challenge_hash_wf
hint_wf
low_wf
message_hash_size
paper_sim_commitment_challenge
paper_sim_commitment_highbits
paper_sim_commitment_lowbits_wf
paper_sim_signature
paper_sim_signature_sample_wf
polyveck_add_wf
reconstructed_highbits
resp_wf
valid_signature
\end{lstlisting}

\end{formaltheorem}

\begin{formaltheorem}[EasyCrypt line 944]
\noindent\textbf{EasyCrypt name.}
\begin{lstlisting}[language=EasyCrypt,basicstyle=\ttfamily\scriptsize]
public_sim_lowbits_carrier_wf
\end{lstlisting}
\noindent\textbf{Checked EasyCrypt statement.}
\begin{lstlisting}[language=EasyCrypt,basicstyle=\ttfamily\scriptsize]
lemma public_sim_lowbits_carrier_wf md pk m ctx coins :
  poly_wf (public_sim_lowbits_carrier md pk m ctx coins).
\end{lstlisting}
\noindent\textbf{Formal mathematical reading.}
Mathematical theorem. This lemma is a universally quantified proposition over the variables shown in the EasyCrypt statement.

\noindent\textbf{Role in the proof.}
Use in this chapter. It proves an exact game replacement or simulator equivalence.

\noindent\textbf{Proof-script interpretation.}
Proof-script reading. The machine proof decomposes the theorem into rewrites, program-logic obligations, relational equivalences, and arithmetic side conditions.
\emph{Main tactics visible in the proof script:} \ec{rewrite}, \ec{apply}, \ec{smt}.
\emph{Named identifiers syntactically cited in the proof block:}
\begin{lstlisting}[language=EasyCrypt,basicstyle=\ttfamily\scriptsize]
coins
ctx
deterministic_polyveck
deterministic_polyveck_wf
md
mode_k_gt0
pk
polyveck_wf_nth
public_sim_lowbits_carrier
public_sim_source
\end{lstlisting}

\end{formaltheorem}

\begin{formaltheorem}[EasyCrypt line 955]
\noindent\textbf{EasyCrypt name.}
\begin{lstlisting}[language=EasyCrypt,basicstyle=\ttfamily\scriptsize]
paper_sim_sample_from_public_coins_wf
\end{lstlisting}
\noindent\textbf{Checked EasyCrypt statement.}
\begin{lstlisting}[language=EasyCrypt,basicstyle=\ttfamily\scriptsize]
lemma paper_sim_sample_from_public_coins_wf md pk m ctx coins :
  paper_sim_signature_sample_wf md
    (paper_sim_sample_from_public_coins md pk m ctx coins).
\end{lstlisting}
\noindent\textbf{Formal mathematical reading.}
Mathematical theorem. This lemma is a universally quantified proposition over the variables shown in the EasyCrypt statement.

\noindent\textbf{Role in the proof.}
Use in this chapter. It contributes directly to an advantage bound, bad-event bound, or real-arithmetic composition step.

\noindent\textbf{Proof-script interpretation.}
Proof-script reading. The machine proof decomposes the theorem into rewrites, program-logic obligations, relational equivalences, and arithmetic side conditions.
\emph{Main tactics visible in the proof script:} \ec{rewrite}, \ec{smt}.
\emph{Named identifiers syntactically cited in the proof block:}
\begin{lstlisting}[language=EasyCrypt,basicstyle=\ttfamily\scriptsize]
deterministic_unit_polyveck_wf
deterministic_unit_polyvecl_wf
paper_sim_sample_from_public_coins
paper_sim_sample_hint
paper_sim_sample_lowbits_carrier
paper_sim_sample_response
paper_sim_sample_response_aux
paper_sim_sample_token
paper_sim_signature_sample_wf
public_sim_hint_vector
public_sim_lowbits_carrier_wf
public_sim_response_aux_vector
public_sim_response_vector
public_sim_signature_token
\end{lstlisting}

\end{formaltheorem}

\begin{formaltheorem}[EasyCrypt line 970]
\noindent\textbf{EasyCrypt name.}
\begin{lstlisting}[language=EasyCrypt,basicstyle=\ttfamily\scriptsize]
paper_sim_sample_response_public_coinsE
\end{lstlisting}
\noindent\textbf{Checked EasyCrypt statement.}
\begin{lstlisting}[language=EasyCrypt,basicstyle=\ttfamily\scriptsize]
lemma paper_sim_sample_response_public_coinsE md pk m ctx coins :
  paper_sim_sample_response
    (paper_sim_sample_from_public_coins md pk m ctx coins) =
  public_sim_response_vector md pk m ctx coins.
\end{lstlisting}
\noindent\textbf{Formal mathematical reading.}
Mathematical theorem. This is an equality or definitional expansion used for rewriting and normalization.

\noindent\textbf{Role in the proof.}
Use in this chapter. It contributes directly to an advantage bound, bad-event bound, or real-arithmetic composition step.

\noindent\textbf{Proof-script interpretation.}
No proof block was collected for this item.

\end{formaltheorem}

\begin{formaltheorem}[EasyCrypt line 978]
\noindent\textbf{EasyCrypt name.}
\begin{lstlisting}[language=EasyCrypt,basicstyle=\ttfamily\scriptsize]
paper_sim_sample_response_aux_public_coinsE
\end{lstlisting}
\noindent\textbf{Checked EasyCrypt statement.}
\begin{lstlisting}[language=EasyCrypt,basicstyle=\ttfamily\scriptsize]
lemma paper_sim_sample_response_aux_public_coinsE md pk m ctx coins :
  paper_sim_sample_response_aux
    (paper_sim_sample_from_public_coins md pk m ctx coins) =
  public_sim_response_aux_vector md pk m ctx coins.
\end{lstlisting}
\noindent\textbf{Formal mathematical reading.}
Mathematical theorem. This is an equality or definitional expansion used for rewriting and normalization.

\noindent\textbf{Role in the proof.}
Use in this chapter. It contributes directly to an advantage bound, bad-event bound, or real-arithmetic composition step.

\noindent\textbf{Proof-script interpretation.}
No proof block was collected for this item.

\end{formaltheorem}

\begin{formaltheorem}[EasyCrypt line 986]
\noindent\textbf{EasyCrypt name.}
\begin{lstlisting}[language=EasyCrypt,basicstyle=\ttfamily\scriptsize]
paper_sim_sample_hint_public_coinsE
\end{lstlisting}
\noindent\textbf{Checked EasyCrypt statement.}
\begin{lstlisting}[language=EasyCrypt,basicstyle=\ttfamily\scriptsize]
lemma paper_sim_sample_hint_public_coinsE md pk m ctx coins :
  paper_sim_sample_hint
    (paper_sim_sample_from_public_coins md pk m ctx coins) =
  public_sim_hint_vector md pk m ctx coins.
\end{lstlisting}
\noindent\textbf{Formal mathematical reading.}
Mathematical theorem. This is an equality or definitional expansion used for rewriting and normalization.

\noindent\textbf{Role in the proof.}
Use in this chapter. It contributes directly to an advantage bound, bad-event bound, or real-arithmetic composition step.

\noindent\textbf{Proof-script interpretation.}
No proof block was collected for this item.

\end{formaltheorem}

\begin{formaltheorem}[EasyCrypt line 994]
\noindent\textbf{EasyCrypt name.}
\begin{lstlisting}[language=EasyCrypt,basicstyle=\ttfamily\scriptsize]
paper_sim_sample_token_public_coinsE
\end{lstlisting}
\noindent\textbf{Checked EasyCrypt statement.}
\begin{lstlisting}[language=EasyCrypt,basicstyle=\ttfamily\scriptsize]
lemma paper_sim_sample_token_public_coinsE md pk m ctx coins :
  paper_sim_sample_token
    (paper_sim_sample_from_public_coins md pk m ctx coins) =
  public_sim_signature_token md pk m ctx coins.
\end{lstlisting}
\noindent\textbf{Formal mathematical reading.}
Mathematical theorem. This is an equality or definitional expansion used for rewriting and normalization.

\noindent\textbf{Role in the proof.}
Use in this chapter. It contributes directly to an advantage bound, bad-event bound, or real-arithmetic composition step.

\noindent\textbf{Proof-script interpretation.}
No proof block was collected for this item.

\end{formaltheorem}

\begin{formaltheorem}[EasyCrypt line 1001]
\noindent\textbf{EasyCrypt name.}
\begin{lstlisting}[language=EasyCrypt,basicstyle=\ttfamily\scriptsize]
paper_sim_commitment_lowbits_public_coinsE
\end{lstlisting}
\noindent\textbf{Checked EasyCrypt statement.}
\begin{lstlisting}[language=EasyCrypt,basicstyle=\ttfamily\scriptsize]
lemma paper_sim_commitment_lowbits_public_coinsE md pk m ctx coins :
  paper_sim_commitment_lowbits md
    (paper_sim_sample_from_public_coins md pk m ctx coins) =
  public_sim_commitment_lowbits md pk m ctx coins.
\end{lstlisting}
\noindent\textbf{Formal mathematical reading.}
Mathematical theorem. This is an equality or definitional expansion used for rewriting and normalization.

\noindent\textbf{Role in the proof.}
Use in this chapter. It proves an exact game replacement or simulator equivalence.

\noindent\textbf{Proof-script interpretation.}
Proof-script reading. The machine proof decomposes the theorem into rewrites, program-logic obligations, relational equivalences, and arithmetic side conditions.
\emph{Main tactics visible in the proof script:} \ec{rewrite}.
\emph{Named identifiers syntactically cited in the proof block:}
\begin{lstlisting}[language=EasyCrypt,basicstyle=\ttfamily\scriptsize]
paper_sim_commitment_lowbits
paper_sim_sample_entropy
paper_sim_sample_from_public_coins
paper_sim_sample_lowbits_carrier
public_sim_commitment_lowbits
public_sim_lowbits_carrier
\end{lstlisting}

\end{formaltheorem}

\begin{formaltheorem}[EasyCrypt line 1012]
\noindent\textbf{EasyCrypt name.}
\begin{lstlisting}[language=EasyCrypt,basicstyle=\ttfamily\scriptsize]
paper_sim_commitment_challenge_public_coinsE
\end{lstlisting}
\noindent\textbf{Checked EasyCrypt statement.}
\begin{lstlisting}[language=EasyCrypt,basicstyle=\ttfamily\scriptsize]
lemma paper_sim_commitment_challenge_public_coinsE md pk m ctx coins :
  paper_sim_commitment_challenge md pk m ctx
    (paper_sim_sample_from_public_coins md pk m ctx coins) =
  public_sim_commitment_challenge md pk m ctx coins.
\end{lstlisting}
\noindent\textbf{Formal mathematical reading.}
Mathematical theorem. This is an equality or definitional expansion used for rewriting and normalization.

\noindent\textbf{Role in the proof.}
Use in this chapter. It contributes directly to an advantage bound, bad-event bound, or real-arithmetic composition step.

\noindent\textbf{Proof-script interpretation.}
Proof-script reading. The machine proof decomposes the theorem into rewrites, program-logic obligations, relational equivalences, and arithmetic side conditions.
\emph{Main tactics visible in the proof script:} \ec{rewrite}.
\emph{Named identifiers syntactically cited in the proof block:}
\begin{lstlisting}[language=EasyCrypt,basicstyle=\ttfamily\scriptsize]
paper_sim_commitment_challenge
paper_sim_commitment_lowbits_public_coinsE
paper_sim_sample_response_aux_public_coinsE
public_sim_commitment_challenge
\end{lstlisting}

\end{formaltheorem}

\begin{formaltheorem}[EasyCrypt line 1022]
\noindent\textbf{EasyCrypt name.}
\begin{lstlisting}[language=EasyCrypt,basicstyle=\ttfamily\scriptsize]
paper_sim_commitment_highbits_public_coinsE
\end{lstlisting}
\noindent\textbf{Checked EasyCrypt statement.}
\begin{lstlisting}[language=EasyCrypt,basicstyle=\ttfamily\scriptsize]
lemma paper_sim_commitment_highbits_public_coinsE md pk m ctx coins :
  paper_sim_commitment_highbits md pk m ctx
    (paper_sim_sample_from_public_coins md pk m ctx coins) =
  public_sim_commitment_highbits md pk m ctx coins.
\end{lstlisting}
\noindent\textbf{Formal mathematical reading.}
Mathematical theorem. This is an equality or definitional expansion used for rewriting and normalization.

\noindent\textbf{Role in the proof.}
Use in this chapter. It proves an exact game replacement or simulator equivalence.

\noindent\textbf{Proof-script interpretation.}
Proof-script reading. The machine proof decomposes the theorem into rewrites, program-logic obligations, relational equivalences, and arithmetic side conditions.
\emph{Main tactics visible in the proof script:} \ec{rewrite}.
\emph{Named identifiers syntactically cited in the proof block:}
\begin{lstlisting}[language=EasyCrypt,basicstyle=\ttfamily\scriptsize]
paper_sim_commitment_challenge_public_coinsE
paper_sim_commitment_highbits
paper_sim_sample_response_aux_public_coinsE
public_sim_commitment_highbits
\end{lstlisting}

\end{formaltheorem}

\begin{formaltheorem}[EasyCrypt line 1032]
\noindent\textbf{EasyCrypt name.}
\begin{lstlisting}[language=EasyCrypt,basicstyle=\ttfamily\scriptsize]
paper_sim_signature_public_coinsE
\end{lstlisting}
\noindent\textbf{Checked EasyCrypt statement.}
\begin{lstlisting}[language=EasyCrypt,basicstyle=\ttfamily\scriptsize]
lemma paper_sim_signature_public_coinsE md pk m ctx coins :
  paper_sim_signature md pk m ctx
    (paper_sim_sample_from_public_coins md pk m ctx coins) =
  public_sim_signature md pk m ctx coins.
\end{lstlisting}
\noindent\textbf{Formal mathematical reading.}
Mathematical theorem. This is an equality or definitional expansion used for rewriting and normalization.

\noindent\textbf{Role in the proof.}
Use in this chapter. It proves an exact game replacement or simulator equivalence.

\noindent\textbf{Proof-script interpretation.}
Proof-script reading. The machine proof decomposes the theorem into rewrites, program-logic obligations, relational equivalences, and arithmetic side conditions.
\emph{Main tactics visible in the proof script:} \ec{rewrite}.
\emph{Named identifiers syntactically cited in the proof block:}
\begin{lstlisting}[language=EasyCrypt,basicstyle=\ttfamily\scriptsize]
paper_sim_commitment_challenge_public_coinsE
paper_sim_commitment_highbits_public_coinsE
paper_sim_commitment_lowbits_public_coinsE
paper_sim_sample_token_public_coinsE
paper_sim_signature
public_sim_signature
\end{lstlisting}

\end{formaltheorem}

\begin{formaltheorem}[EasyCrypt line 1056]
\noindent\textbf{EasyCrypt name.}
\begin{lstlisting}[language=EasyCrypt,basicstyle=\ttfamily\scriptsize]
real_signing_lowbits_carrier_wf
\end{lstlisting}
\noindent\textbf{Checked EasyCrypt statement.}
\begin{lstlisting}[language=EasyCrypt,basicstyle=\ttfamily\scriptsize]
lemma real_signing_lowbits_carrier_wf md sk m ctx coins :
  poly_wf (real_signing_lowbits_carrier md sk m ctx coins).
\end{lstlisting}
\noindent\textbf{Formal mathematical reading.}
Mathematical theorem. This lemma is a universally quantified proposition over the variables shown in the EasyCrypt statement.

\noindent\textbf{Role in the proof.}
Use in this chapter. It proves a structural correctness fact needed by later extraction or verification arguments.

\noindent\textbf{Proof-script interpretation.}
Proof-script reading. The machine proof decomposes the theorem into rewrites, program-logic obligations, relational equivalences, and arithmetic side conditions.
\emph{Main tactics visible in the proof script:} \ec{rewrite}, \ec{apply}, \ec{smt}.
\emph{Named identifiers syntactically cited in the proof block:}
\begin{lstlisting}[language=EasyCrypt,basicstyle=\ttfamily\scriptsize]
coins
commitment_raw
commitment_raw_wf
ctx
md
mode_k_gt0
polyveck_wf_nth
real_signing_lowbits_carrier
sk
\end{lstlisting}

\end{formaltheorem}

\begin{formaltheorem}[EasyCrypt line 1065]
\noindent\textbf{EasyCrypt name.}
\begin{lstlisting}[language=EasyCrypt,basicstyle=\ttfamily\scriptsize]
paper_sim_sample_from_real_signing_coins_wf
\end{lstlisting}
\noindent\textbf{Checked EasyCrypt statement.}
\begin{lstlisting}[language=EasyCrypt,basicstyle=\ttfamily\scriptsize]
lemma paper_sim_sample_from_real_signing_coins_wf md sk m ctx coins :
  paper_sim_signature_sample_wf md
    (paper_sim_sample_from_real_signing_coins md sk m ctx coins).
\end{lstlisting}
\noindent\textbf{Formal mathematical reading.}
Mathematical theorem. This lemma is a universally quantified proposition over the variables shown in the EasyCrypt statement.

\noindent\textbf{Role in the proof.}
Use in this chapter. It contributes directly to an advantage bound, bad-event bound, or real-arithmetic composition step.

\noindent\textbf{Proof-script interpretation.}
Proof-script reading. The machine proof decomposes the theorem into rewrites, program-logic obligations, relational equivalences, and arithmetic side conditions.
\emph{Main tactics visible in the proof script:} \ec{rewrite}, \ec{smt}.
\emph{Named identifiers syntactically cited in the proof block:}
\begin{lstlisting}[language=EasyCrypt,basicstyle=\ttfamily\scriptsize]
hint_vector_wf
paper_sim_sample_from_real_signing_coins
paper_sim_sample_hint
paper_sim_sample_lowbits_carrier
paper_sim_sample_response
paper_sim_sample_response_aux
paper_sim_sample_token
paper_sim_signature_sample_wf
real_signing_lowbits_carrier_wf
response_aux_vector_wf
response_vector_wf
\end{lstlisting}

\end{formaltheorem}

\begin{formaltheorem}[EasyCrypt line 1078]
\noindent\textbf{EasyCrypt name.}
\begin{lstlisting}[language=EasyCrypt,basicstyle=\ttfamily\scriptsize]
paper_sim_sample_response_real_signing_coinsE
\end{lstlisting}
\noindent\textbf{Checked EasyCrypt statement.}
\begin{lstlisting}[language=EasyCrypt,basicstyle=\ttfamily\scriptsize]
lemma paper_sim_sample_response_real_signing_coinsE md sk m ctx coins :
  paper_sim_sample_response
    (paper_sim_sample_from_real_signing_coins md sk m ctx coins) =
  response_vector md sk m ctx coins.
\end{lstlisting}
\noindent\textbf{Formal mathematical reading.}
Mathematical theorem. This is an equality or definitional expansion used for rewriting and normalization.

\noindent\textbf{Role in the proof.}
Use in this chapter. It contributes directly to an advantage bound, bad-event bound, or real-arithmetic composition step.

\noindent\textbf{Proof-script interpretation.}
No proof block was collected for this item.

\end{formaltheorem}

\begin{formaltheorem}[EasyCrypt line 1086]
\noindent\textbf{EasyCrypt name.}
\begin{lstlisting}[language=EasyCrypt,basicstyle=\ttfamily\scriptsize]
paper_sim_sample_response_aux_real_signing_coinsE
\end{lstlisting}
\noindent\textbf{Checked EasyCrypt statement.}
\begin{lstlisting}[language=EasyCrypt,basicstyle=\ttfamily\scriptsize]
lemma paper_sim_sample_response_aux_real_signing_coinsE md sk m ctx coins :
  paper_sim_sample_response_aux
    (paper_sim_sample_from_real_signing_coins md sk m ctx coins) =
  response_aux_vector md sk m ctx coins.
\end{lstlisting}
\noindent\textbf{Formal mathematical reading.}
Mathematical theorem. This is an equality or definitional expansion used for rewriting and normalization.

\noindent\textbf{Role in the proof.}
Use in this chapter. It contributes directly to an advantage bound, bad-event bound, or real-arithmetic composition step.

\noindent\textbf{Proof-script interpretation.}
No proof block was collected for this item.

\end{formaltheorem}

\begin{formaltheorem}[EasyCrypt line 1094]
\noindent\textbf{EasyCrypt name.}
\begin{lstlisting}[language=EasyCrypt,basicstyle=\ttfamily\scriptsize]
paper_sim_sample_hint_real_signing_coinsE
\end{lstlisting}
\noindent\textbf{Checked EasyCrypt statement.}
\begin{lstlisting}[language=EasyCrypt,basicstyle=\ttfamily\scriptsize]
lemma paper_sim_sample_hint_real_signing_coinsE md sk m ctx coins :
  paper_sim_sample_hint
    (paper_sim_sample_from_real_signing_coins md sk m ctx coins) =
  hint_vector md sk m ctx coins.
\end{lstlisting}
\noindent\textbf{Formal mathematical reading.}
Mathematical theorem. This is an equality or definitional expansion used for rewriting and normalization.

\noindent\textbf{Role in the proof.}
Use in this chapter. It contributes directly to an advantage bound, bad-event bound, or real-arithmetic composition step.

\noindent\textbf{Proof-script interpretation.}
No proof block was collected for this item.

\end{formaltheorem}

\begin{formaltheorem}[EasyCrypt line 1102]
\noindent\textbf{EasyCrypt name.}
\begin{lstlisting}[language=EasyCrypt,basicstyle=\ttfamily\scriptsize]
paper_sim_sample_token_real_signing_coinsE
\end{lstlisting}
\noindent\textbf{Checked EasyCrypt statement.}
\begin{lstlisting}[language=EasyCrypt,basicstyle=\ttfamily\scriptsize]
lemma paper_sim_sample_token_real_signing_coinsE md sk m ctx coins :
  paper_sim_sample_token
    (paper_sim_sample_from_real_signing_coins md sk m ctx coins) =
  signature_token md sk m ctx coins.
\end{lstlisting}
\noindent\textbf{Formal mathematical reading.}
Mathematical theorem. This is an equality or definitional expansion used for rewriting and normalization.

\noindent\textbf{Role in the proof.}
Use in this chapter. It contributes directly to an advantage bound, bad-event bound, or real-arithmetic composition step.

\noindent\textbf{Proof-script interpretation.}
No proof block was collected for this item.

\end{formaltheorem}

\begin{formaltheorem}[EasyCrypt line 1109]
\noindent\textbf{EasyCrypt name.}
\begin{lstlisting}[language=EasyCrypt,basicstyle=\ttfamily\scriptsize]
paper_sim_commitment_lowbits_real_signing_coinsE
\end{lstlisting}
\noindent\textbf{Checked EasyCrypt statement.}
\begin{lstlisting}[language=EasyCrypt,basicstyle=\ttfamily\scriptsize]
lemma paper_sim_commitment_lowbits_real_signing_coinsE md sk m ctx coins :
  paper_sim_commitment_lowbits md
    (paper_sim_sample_from_real_signing_coins md sk m ctx coins) =
  commitment_lowbits md sk m ctx coins.
\end{lstlisting}
\noindent\textbf{Formal mathematical reading.}
Mathematical theorem. This is an equality or definitional expansion used for rewriting and normalization.

\noindent\textbf{Role in the proof.}
Use in this chapter. It proves an exact game replacement or simulator equivalence.

\noindent\textbf{Proof-script interpretation.}
Proof-script reading. The machine proof decomposes the theorem into rewrites, program-logic obligations, relational equivalences, and arithmetic side conditions.
\emph{Main tactics visible in the proof script:} \ec{rewrite}.
\emph{Named identifiers syntactically cited in the proof block:}
\begin{lstlisting}[language=EasyCrypt,basicstyle=\ttfamily\scriptsize]
commitment_lowbits
paper_sim_commitment_lowbits
paper_sim_sample_entropy
paper_sim_sample_from_real_signing_coins
paper_sim_sample_lowbits_carrier
real_signing_lowbits_carrier
\end{lstlisting}

\end{formaltheorem}

\begin{formaltheorem}[EasyCrypt line 1120]
\noindent\textbf{EasyCrypt name.}
\begin{lstlisting}[language=EasyCrypt,basicstyle=\ttfamily\scriptsize]
paper_sim_commitment_challenge_real_signing_coinsE
\end{lstlisting}
\noindent\textbf{Checked EasyCrypt statement.}
\begin{lstlisting}[language=EasyCrypt,basicstyle=\ttfamily\scriptsize]
lemma paper_sim_commitment_challenge_real_signing_coinsE
   md sk m ctx coins :
  paper_sim_commitment_challenge md (public_key_of_secret md sk) m ctx
    (paper_sim_sample_from_real_signing_coins md sk m ctx coins) =
  commitment_challenge md sk m ctx coins.
\end{lstlisting}
\noindent\textbf{Formal mathematical reading.}
Mathematical theorem. This is an equality or definitional expansion used for rewriting and normalization.

\noindent\textbf{Role in the proof.}
Use in this chapter. It contributes directly to an advantage bound, bad-event bound, or real-arithmetic composition step.

\noindent\textbf{Proof-script interpretation.}
Proof-script reading. The machine proof decomposes the theorem into rewrites, program-logic obligations, relational equivalences, and arithmetic side conditions.
\emph{Main tactics visible in the proof script:} \ec{rewrite}.
\emph{Named identifiers syntactically cited in the proof block:}
\begin{lstlisting}[language=EasyCrypt,basicstyle=\ttfamily\scriptsize]
commitment_challenge
paper_sim_commitment_challenge
paper_sim_commitment_lowbits_real_signing_coinsE
paper_sim_sample_response_aux_real_signing_coinsE
\end{lstlisting}

\end{formaltheorem}

\begin{formaltheorem}[EasyCrypt line 1131]
\noindent\textbf{EasyCrypt name.}
\begin{lstlisting}[language=EasyCrypt,basicstyle=\ttfamily\scriptsize]
paper_sim_commitment_highbits_real_signing_coinsE
\end{lstlisting}
\noindent\textbf{Checked EasyCrypt statement.}
\begin{lstlisting}[language=EasyCrypt,basicstyle=\ttfamily\scriptsize]
lemma paper_sim_commitment_highbits_real_signing_coinsE
   md sk m ctx coins :
  paper_sim_commitment_highbits md (public_key_of_secret md sk) m ctx
    (paper_sim_sample_from_real_signing_coins md sk m ctx coins) =
  commitment_highbits md sk m ctx coins.
\end{lstlisting}
\noindent\textbf{Formal mathematical reading.}
Mathematical theorem. This is an equality or definitional expansion used for rewriting and normalization.

\noindent\textbf{Role in the proof.}
Use in this chapter. It proves an exact game replacement or simulator equivalence.

\noindent\textbf{Proof-script interpretation.}
Proof-script reading. The machine proof decomposes the theorem into rewrites, program-logic obligations, relational equivalences, and arithmetic side conditions.
\emph{Main tactics visible in the proof script:} \ec{rewrite}.
\emph{Named identifiers syntactically cited in the proof block:}
\begin{lstlisting}[language=EasyCrypt,basicstyle=\ttfamily\scriptsize]
commitment_highbits
paper_sim_commitment_challenge_real_signing_coinsE
paper_sim_commitment_highbits
paper_sim_sample_response_aux_real_signing_coinsE
\end{lstlisting}

\end{formaltheorem}

\begin{formaltheorem}[EasyCrypt line 1142]
\noindent\textbf{EasyCrypt name.}
\begin{lstlisting}[language=EasyCrypt,basicstyle=\ttfamily\scriptsize]
paper_sim_signature_real_signing_coinsE
\end{lstlisting}
\noindent\textbf{Checked EasyCrypt statement.}
\begin{lstlisting}[language=EasyCrypt,basicstyle=\ttfamily\scriptsize]
lemma paper_sim_signature_real_signing_coinsE md sk m ctx coins :
  paper_sim_signature md (public_key_of_secret md sk) m ctx
    (paper_sim_sample_from_real_signing_coins md sk m ctx coins) =
  sign_internal md sk m ctx coins.
\end{lstlisting}
\noindent\textbf{Formal mathematical reading.}
Mathematical theorem. This is an equality or definitional expansion used for rewriting and normalization.

\noindent\textbf{Role in the proof.}
Use in this chapter. It proves an exact game replacement or simulator equivalence.

\noindent\textbf{Proof-script interpretation.}
Proof-script reading. The machine proof decomposes the theorem into rewrites, program-logic obligations, relational equivalences, and arithmetic side conditions.
\emph{Main tactics visible in the proof script:} \ec{rewrite}.
\emph{Named identifiers syntactically cited in the proof block:}
\begin{lstlisting}[language=EasyCrypt,basicstyle=\ttfamily\scriptsize]
paper_sim_commitment_challenge_real_signing_coinsE
paper_sim_commitment_highbits_real_signing_coinsE
paper_sim_commitment_lowbits_real_signing_coinsE
paper_sim_sample_token_real_signing_coinsE
paper_sim_signature
sign_internal
\end{lstlisting}

\end{formaltheorem}

\begin{formaltheorem}[EasyCrypt line 1171]
\noindent\textbf{EasyCrypt name.}
\begin{lstlisting}[language=EasyCrypt,basicstyle=\ttfamily\scriptsize]
paper_sim_abort_fallback_sample_wf
\end{lstlisting}
\noindent\textbf{Checked EasyCrypt statement.}
\begin{lstlisting}[language=EasyCrypt,basicstyle=\ttfamily\scriptsize]
lemma paper_sim_abort_fallback_sample_wf md :
  paper_sim_signature_sample_wf md (paper_sim_abort_fallback_sample md).
\end{lstlisting}
\noindent\textbf{Formal mathematical reading.}
Mathematical theorem. This lemma is a universally quantified proposition over the variables shown in the EasyCrypt statement.

\noindent\textbf{Role in the proof.}
Use in this chapter. It contributes directly to an advantage bound, bad-event bound, or real-arithmetic composition step.

\noindent\textbf{Proof-script interpretation.}
Proof-script reading. The machine proof decomposes the theorem into rewrites, program-logic obligations, relational equivalences, and arithmetic side conditions.
\emph{Main tactics visible in the proof script:} \ec{rewrite}.
\emph{Named identifiers syntactically cited in the proof block:}
\begin{lstlisting}[language=EasyCrypt,basicstyle=\ttfamily\scriptsize]
paper_sim_abort_fallback_sample
paper_sim_sample_hint
paper_sim_sample_lowbits_carrier
paper_sim_sample_response
paper_sim_sample_response_aux
paper_sim_sample_token
paper_sim_signature_sample_wf
poly_zero_wf
polyveck_zero_wf
polyvecl_zero_wf
\end{lstlisting}

\end{formaltheorem}

\begin{formaltheorem}[EasyCrypt line 1182]
\noindent\textbf{EasyCrypt name.}
\begin{lstlisting}[language=EasyCrypt,basicstyle=\ttfamily\scriptsize]
paper_sim_sample_from_rejection_attempt_of_coinsE
\end{lstlisting}
\noindent\textbf{Checked EasyCrypt statement.}
\begin{lstlisting}[language=EasyCrypt,basicstyle=\ttfamily\scriptsize]
lemma paper_sim_sample_from_rejection_attempt_of_coinsE md sk m ctx coins :
  paper_sim_sample_from_rejection_attempt
    (signing_attempt_state_of_coins md sk m ctx coins) =
  paper_sim_sample_from_real_signing_coins md sk m ctx coins.
\end{lstlisting}
\noindent\textbf{Formal mathematical reading.}
Mathematical theorem. This is an equality or definitional expansion used for rewriting and normalization.

\noindent\textbf{Role in the proof.}
Use in this chapter. It contributes directly to an advantage bound, bad-event bound, or real-arithmetic composition step.

\noindent\textbf{Proof-script interpretation.}
Proof-script reading. The machine proof decomposes the theorem into rewrites, program-logic obligations, relational equivalences, and arithmetic side conditions.
\emph{Main tactics visible in the proof script:} \ec{rewrite}.
\emph{Named identifiers syntactically cited in the proof block:}
\begin{lstlisting}[language=EasyCrypt,basicstyle=\ttfamily\scriptsize]
paper_sim_sample_from_real_signing_coins
paper_sim_sample_from_rejection_attempt
real_signing_lowbits_carrier
signing_attempt_state_coins
signing_attempt_state_of_coins
signing_attempt_state_of_coins_lowbits_carrierE
signing_attempt_state_of_coins_tokenE
\end{lstlisting}

\end{formaltheorem}

\begin{formaltheorem}[EasyCrypt line 1195]
\noindent\textbf{EasyCrypt name.}
\begin{lstlisting}[language=EasyCrypt,basicstyle=\ttfamily\scriptsize]
paper_sim_commitment_lowbits_rejection_attemptE
\end{lstlisting}
\noindent\textbf{Checked EasyCrypt statement.}
\begin{lstlisting}[language=EasyCrypt,basicstyle=\ttfamily\scriptsize]
lemma paper_sim_commitment_lowbits_rejection_attemptE md st :
  paper_sim_commitment_lowbits md
    (paper_sim_sample_from_rejection_attempt st) =
  signing_attempt_state_programmed_lowbits st.
\end{lstlisting}
\noindent\textbf{Formal mathematical reading.}
Mathematical theorem. This is an equality or definitional expansion used for rewriting and normalization.

\noindent\textbf{Role in the proof.}
Use in this chapter. It proves an exact game replacement or simulator equivalence.

\noindent\textbf{Proof-script interpretation.}
Proof-script reading. The machine proof decomposes the theorem into rewrites, program-logic obligations, relational equivalences, and arithmetic side conditions.
\emph{Main tactics visible in the proof script:} \ec{rewrite}.
\emph{Named identifiers syntactically cited in the proof block:}
\begin{lstlisting}[language=EasyCrypt,basicstyle=\ttfamily\scriptsize]
paper_sim_commitment_lowbits
paper_sim_sample_entropy
paper_sim_sample_from_rejection_attempt
paper_sim_sample_lowbits_carrier
signing_attempt_state_programmed_lowbits
\end{lstlisting}

\end{formaltheorem}

\begin{formaltheorem}[EasyCrypt line 1207]
\noindent\textbf{EasyCrypt name.}
\begin{lstlisting}[language=EasyCrypt,basicstyle=\ttfamily\scriptsize]
paper_sim_commitment_challenge_rejection_attemptE
\end{lstlisting}
\noindent\textbf{Checked EasyCrypt statement.}
\begin{lstlisting}[language=EasyCrypt,basicstyle=\ttfamily\scriptsize]
lemma paper_sim_commitment_challenge_rejection_attemptE md pk m ctx st :
  signing_attempt_state_lowbits_consistent st =>
  paper_sim_commitment_challenge md pk m ctx
    (paper_sim_sample_from_rejection_attempt st) =
  signing_attempt_state_programmed_challenge md pk m ctx st.
\end{lstlisting}
\noindent\textbf{Formal mathematical reading.}
Mathematical theorem. This is an implication, normally turning one invariant, validity predicate, or proof obligation into another.

\noindent\textbf{Role in the proof.}
Use in this chapter. It contributes directly to an advantage bound, bad-event bound, or real-arithmetic composition step.

\noindent\textbf{Proof-script interpretation.}
Proof-script reading. The machine proof decomposes the theorem into rewrites, program-logic obligations, relational equivalences, and arithmetic side conditions.
\emph{Main tactics visible in the proof script:} \ec{rewrite}.
\emph{Named identifiers syntactically cited in the proof block:}
\begin{lstlisting}[language=EasyCrypt,basicstyle=\ttfamily\scriptsize]
lowbits_ok
paper_sim_commitment_challenge
paper_sim_commitment_lowbits_rejection_attemptE
paper_sim_sample_from_rejection_attempt
paper_sim_sample_response_aux
paper_sim_sample_token
signing_attempt_state_lowbits_consistent
signing_attempt_state_programmed_challenge
\end{lstlisting}

\end{formaltheorem}

\begin{formaltheorem}[EasyCrypt line 1225]
\noindent\textbf{EasyCrypt name.}
\begin{lstlisting}[language=EasyCrypt,basicstyle=\ttfamily\scriptsize]
paper_sim_commitment_highbits_rejection_attemptE
\end{lstlisting}
\noindent\textbf{Checked EasyCrypt statement.}
\begin{lstlisting}[language=EasyCrypt,basicstyle=\ttfamily\scriptsize]
lemma paper_sim_commitment_highbits_rejection_attemptE md pk m ctx st :
  signing_attempt_state_lowbits_consistent st =>
  signing_attempt_state_challenge_consistent md pk m ctx st =>
  signing_attempt_state_public_equation_consistent md pk st =>
  paper_sim_commitment_highbits md pk m ctx
    (paper_sim_sample_from_rejection_attempt st) =
  signing_attempt_state_highbits st.
\end{lstlisting}
\noindent\textbf{Formal mathematical reading.}
Mathematical theorem. This is an implication, normally turning one invariant, validity predicate, or proof obligation into another.

\noindent\textbf{Role in the proof.}
Use in this chapter. It proves an exact game replacement or simulator equivalence.

\noindent\textbf{Proof-script interpretation.}
Proof-script reading. The machine proof decomposes the theorem into rewrites, program-logic obligations, relational equivalences, and arithmetic side conditions.
\emph{Main tactics visible in the proof script:} \ec{rewrite}, \ec{smt}.
\emph{Named identifiers syntactically cited in the proof block:}
\begin{lstlisting}[language=EasyCrypt,basicstyle=\ttfamily\scriptsize]
challengeE
challenge_ok
ctx
lowbits_ok
md
paper_sim_commitment_challenge_rejection_attemptE
paper_sim_commitment_highbits
paper_sim_sample_from_rejection_attempt
paper_sim_sample_response_aux
paper_sim_sample_token
pk
publicE
public_ok
signing_attempt_state_challenge_consistent
signing_attempt_state_public_equation_consistent
signing_attempt_state_reconstructed_highbits
st
\end{lstlisting}

\end{formaltheorem}

\begin{formaltheorem}[EasyCrypt line 1249]
\noindent\textbf{EasyCrypt name.}
\begin{lstlisting}[language=EasyCrypt,basicstyle=\ttfamily\scriptsize]
paper_sim_signature_rejection_attempt_fieldsE
\end{lstlisting}
\noindent\textbf{Checked EasyCrypt statement.}
\begin{lstlisting}[language=EasyCrypt,basicstyle=\ttfamily\scriptsize]
lemma paper_sim_signature_rejection_attempt_fieldsE md pk m ctx st :
  signing_attempt_state_field_accepts md pk m ctx st =>
  paper_sim_signature md pk m ctx
    (paper_sim_sample_from_rejection_attempt st) =
  signing_attempt_state_signature_of_fields pk m ctx st.
\end{lstlisting}
\noindent\textbf{Formal mathematical reading.}
Mathematical theorem. This is an implication, normally turning one invariant, validity predicate, or proof obligation into another.

\noindent\textbf{Role in the proof.}
Use in this chapter. It proves an exact game replacement or simulator equivalence.

\noindent\textbf{Proof-script interpretation.}
Proof-script reading. The machine proof decomposes the theorem into rewrites, program-logic obligations, relational equivalences, and arithmetic side conditions.
\emph{Main tactics visible in the proof script:} \ec{rewrite}, \ec{have}.
\emph{Named identifiers syntactically cited in the proof block:}
\begin{lstlisting}[language=EasyCrypt,basicstyle=\ttfamily\scriptsize]
challengeE
challenge_ok
ctx
fields_ok
lowbitsE
lowbits_ok
md
paper_sim_commitment_challenge_rejection_attemptE
paper_sim_commitment_highbits_rejection_attemptE
paper_sim_commitment_lowbits_rejection_attemptE
paper_sim_sample_from_rejection_attempt
paper_sim_sample_token
paper_sim_signature
pk
public_ok
signing_attempt_state_challenge_consistent
signing_attempt_state_field_accepts
signing_attempt_state_lowbits_consistent
signing_attempt_state_signature_of_fields
st
\end{lstlisting}

\end{formaltheorem}

\begin{formaltheorem}[EasyCrypt line 1277]
\noindent\textbf{EasyCrypt name.}
\begin{lstlisting}[language=EasyCrypt,basicstyle=\ttfamily\scriptsize]
paper_sim_signature_rejection_attempt_matches_state
\end{lstlisting}
\noindent\textbf{Checked EasyCrypt statement.}
\begin{lstlisting}[language=EasyCrypt,basicstyle=\ttfamily\scriptsize]
lemma paper_sim_signature_rejection_attempt_matches_state md pk m ctx st :
  signing_attempt_state_field_accepts md pk m ctx st =>
  paper_sim_signature md pk m ctx
    (paper_sim_sample_from_rejection_attempt st) =
  signing_attempt_state_signature st.
\end{lstlisting}
\noindent\textbf{Formal mathematical reading.}
Mathematical theorem. This is an implication, normally turning one invariant, validity predicate, or proof obligation into another.

\noindent\textbf{Role in the proof.}
Use in this chapter. It proves an exact game replacement or simulator equivalence.

\noindent\textbf{Proof-script interpretation.}
Proof-script reading. The machine proof decomposes the theorem into rewrites, program-logic obligations, relational equivalences, and arithmetic side conditions.
\emph{Main tactics visible in the proof script:} \ec{rewrite}.
\emph{Named identifiers syntactically cited in the proof block:}
\begin{lstlisting}[language=EasyCrypt,basicstyle=\ttfamily\scriptsize]
ctx
field_ok
match_ok
md
paper_sim_signature_rejection_attempt_fieldsE
pk
signing_attempt_state_field_accepts
signing_attempt_state_fields_match_signature
st
\end{lstlisting}

\end{formaltheorem}

\begin{formaltheorem}[EasyCrypt line 1293]
\noindent\textbf{EasyCrypt name.}
\begin{lstlisting}[language=EasyCrypt,basicstyle=\ttfamily\scriptsize]
paper_sim_signature_rejection_attempt_of_coinsE
\end{lstlisting}
\noindent\textbf{Checked EasyCrypt statement.}
\begin{lstlisting}[language=EasyCrypt,basicstyle=\ttfamily\scriptsize]
lemma paper_sim_signature_rejection_attempt_of_coinsE md sk m ctx coins :
  paper_sim_signature md (public_key_of_secret md sk) m ctx
    (paper_sim_sample_from_rejection_attempt
      (signing_attempt_state_of_coins md sk m ctx coins)) =
  sign_internal md sk m ctx coins.
\end{lstlisting}
\noindent\textbf{Formal mathematical reading.}
Mathematical theorem. This is an equality or definitional expansion used for rewriting and normalization.

\noindent\textbf{Role in the proof.}
Use in this chapter. It proves an exact game replacement or simulator equivalence.

\noindent\textbf{Proof-script interpretation.}
Proof-script reading. The machine proof decomposes the theorem into rewrites, program-logic obligations, relational equivalences, and arithmetic side conditions.
\emph{Main tactics visible in the proof script:} \ec{rewrite}.
\emph{Named identifiers syntactically cited in the proof block:}
\begin{lstlisting}[language=EasyCrypt,basicstyle=\ttfamily\scriptsize]
coins
ctx
md
paper_sim_signature_rejection_attempt_matches_state
public_key_of_secret
signing_attempt_state_of_coins
signing_attempt_state_of_coins_field_accepts
signing_attempt_state_of_coins_signatureE
sk
\end{lstlisting}

\end{formaltheorem}

\begin{formaltheorem}[EasyCrypt line 1306]
\noindent\textbf{EasyCrypt name.}
\begin{lstlisting}[language=EasyCrypt,basicstyle=\ttfamily\scriptsize]
haetae_rejection_sample_from_attempt_signatureE
\end{lstlisting}
\noindent\textbf{Checked EasyCrypt statement.}
\begin{lstlisting}[language=EasyCrypt,basicstyle=\ttfamily\scriptsize]
lemma haetae_rejection_sample_from_attempt_signatureE md pk m ctx st :
  signing_attempt_state_rejection_accepts md pk m ctx st =>
  paper_sim_signature md pk m ctx
    (haetae_rejection_sample_from_attempt md pk m ctx st) =
  signing_attempt_state_signature st.
\end{lstlisting}
\noindent\textbf{Formal mathematical reading.}
Mathematical theorem. This is an implication, normally turning one invariant, validity predicate, or proof obligation into another.

\noindent\textbf{Role in the proof.}
Use in this chapter. It contributes directly to an advantage bound, bad-event bound, or real-arithmetic composition step.

\noindent\textbf{Proof-script interpretation.}
Proof-script reading. The machine proof decomposes the theorem into rewrites, program-logic obligations, relational equivalences, and arithmetic side conditions.
\emph{Main tactics visible in the proof script:} \ec{rewrite}.
\emph{Named identifiers syntactically cited in the proof block:}
\begin{lstlisting}[language=EasyCrypt,basicstyle=\ttfamily\scriptsize]
ctx
field_ok
haetae_rejection_sample_from_attempt
md
paper_sim_signature_rejection_attempt_matches_state
pk
reject_ok
signing_attempt_state_rejection_accepts
signing_attempt_state_verification_accepts
st
\end{lstlisting}

\end{formaltheorem}

\begin{formaltheorem}[EasyCrypt line 1323]
\noindent\textbf{EasyCrypt name.}
\begin{lstlisting}[language=EasyCrypt,basicstyle=\ttfamily\scriptsize]
haetae_rejection_sample_from_attempt_no_fallback
\end{lstlisting}
\noindent\textbf{Checked EasyCrypt statement.}
\begin{lstlisting}[language=EasyCrypt,basicstyle=\ttfamily\scriptsize]
lemma haetae_rejection_sample_from_attempt_no_fallback md pk m ctx st :
  signing_attempt_state_rejection_accepts md pk m ctx st =>
  haetae_rejection_sample_from_attempt md pk m ctx st =
  paper_sim_sample_from_rejection_attempt st.
\end{lstlisting}
\noindent\textbf{Formal mathematical reading.}
Mathematical theorem. This is an implication, normally turning one invariant, validity predicate, or proof obligation into another.

\noindent\textbf{Role in the proof.}
Use in this chapter. It contributes directly to an advantage bound, bad-event bound, or real-arithmetic composition step.

\noindent\textbf{Proof-script interpretation.}
Proof-script reading. The machine proof decomposes the theorem into rewrites, program-logic obligations, relational equivalences, and arithmetic side conditions.
\emph{Main tactics visible in the proof script:} \ec{rewrite}.
\emph{Named identifiers syntactically cited in the proof block:}
\begin{lstlisting}[language=EasyCrypt,basicstyle=\ttfamily\scriptsize]
haetae_rejection_sample_from_attempt
reject_ok
\end{lstlisting}

\end{formaltheorem}

\begin{formaltheorem}[EasyCrypt line 1332]
\noindent\textbf{EasyCrypt name.}
\begin{lstlisting}[language=EasyCrypt,basicstyle=\ttfamily\scriptsize]
haetae_rejection_sample_from_attempt_of_coinsE
\end{lstlisting}
\noindent\textbf{Checked EasyCrypt statement.}
\begin{lstlisting}[language=EasyCrypt,basicstyle=\ttfamily\scriptsize]
lemma haetae_rejection_sample_from_attempt_of_coinsE md sk m ctx coins :
  haetae_rejection_sample_from_attempt md (public_key_of_secret md sk) m ctx
    (signing_attempt_state_of_coins md sk m ctx coins) =
  paper_sim_sample_from_real_signing_coins md sk m ctx coins.
\end{lstlisting}
\noindent\textbf{Formal mathematical reading.}
Mathematical theorem. This is an equality or definitional expansion used for rewriting and normalization.

\noindent\textbf{Role in the proof.}
Use in this chapter. It contributes directly to an advantage bound, bad-event bound, or real-arithmetic composition step.

\noindent\textbf{Proof-script interpretation.}
Proof-script reading. The machine proof decomposes the theorem into rewrites, program-logic obligations, relational equivalences, and arithmetic side conditions.
\emph{Main tactics visible in the proof script:} \ec{rewrite}.
\emph{Named identifiers syntactically cited in the proof block:}
\begin{lstlisting}[language=EasyCrypt,basicstyle=\ttfamily\scriptsize]
haetae_rejection_sample_from_attempt
paper_sim_sample_from_rejection_attempt_of_coinsE
signing_attempt_state_of_coins_rejection_accepts_current
\end{lstlisting}

\end{formaltheorem}

\begin{formaltheorem}[EasyCrypt line 1343]
\noindent\textbf{EasyCrypt name.}
\begin{lstlisting}[language=EasyCrypt,basicstyle=\ttfamily\scriptsize]
haetae_rejection_sample_from_attempt_of_coins_wf
\end{lstlisting}
\noindent\textbf{Checked EasyCrypt statement.}
\begin{lstlisting}[language=EasyCrypt,basicstyle=\ttfamily\scriptsize]
lemma haetae_rejection_sample_from_attempt_of_coins_wf md sk m ctx coins :
  paper_sim_signature_sample_wf md
    (haetae_rejection_sample_from_attempt md (public_key_of_secret md sk) m ctx
       (signing_attempt_state_of_coins md sk m ctx coins)).
\end{lstlisting}
\noindent\textbf{Formal mathematical reading.}
Mathematical theorem. This lemma is a universally quantified proposition over the variables shown in the EasyCrypt statement.

\noindent\textbf{Role in the proof.}
Use in this chapter. It contributes directly to an advantage bound, bad-event bound, or real-arithmetic composition step.

\noindent\textbf{Proof-script interpretation.}
Proof-script reading. The machine proof decomposes the theorem into rewrites, program-logic obligations, relational equivalences, and arithmetic side conditions.
\emph{Main tactics visible in the proof script:} \ec{rewrite}, \ec{apply}.
\emph{Named identifiers syntactically cited in the proof block:}
\begin{lstlisting}[language=EasyCrypt,basicstyle=\ttfamily\scriptsize]
haetae_rejection_sample_from_attempt_of_coinsE
paper_sim_sample_from_real_signing_coins_wf
\end{lstlisting}

\end{formaltheorem}

\begin{formaltheorem}[EasyCrypt line 1774]
\noindent\textbf{EasyCrypt name.}
\begin{lstlisting}[language=EasyCrypt,basicstyle=\ttfamily\scriptsize]
budgeted_paper_sim_sign_oracle_preserves_signing_count
\end{lstlisting}
\noindent\textbf{Checked EasyCrypt statement.}
\begin{lstlisting}[language=EasyCrypt,basicstyle=\ttfamily\scriptsize]
lemma budgeted_paper_sim_sign_oracle_preserves_signing_count :
  hoare[ROMInternalTranscriptBudgetedPaperSimAsNMA(A, Samp, H).O.sign :
    budgeted_paper_sim_signing_count_discipline
      ROMInternalTranscriptBudgetedPaperSimAsNMA.signing_count ==>
    budgeted_paper_sim_signing_count_discipline
      ROMInternalTranscriptBudgetedPaperSimAsNMA.signing_count].
\end{lstlisting}
\noindent\textbf{Formal mathematical reading.}
Mathematical theorem. This is a relational program theorem: two games or procedures are coupled so that a precondition on their initial states implies the stated postcondition on final states and return values.

\noindent\textbf{Role in the proof.}
Use in this chapter. It contributes directly to an advantage bound, bad-event bound, or real-arithmetic composition step.

\noindent\textbf{Proof-script interpretation.}
Proof-script reading. The machine proof decomposes the theorem into rewrites, program-logic obligations, relational equivalences, and arithmetic side conditions.
\emph{Main tactics visible in the proof script:} \ec{proc}, \ec{wp}, \ec{call}, \ec{rewrite}, \ec{smt}.
\emph{Named identifiers syntactically cited in the proof block:}
\begin{lstlisting}[language=EasyCrypt,basicstyle=\ttfamily\scriptsize]
budgeted_paper_sim_signing_count_discipline
signature_query_budget_count
\end{lstlisting}

\end{formaltheorem}

\begin{formaltheorem}[EasyCrypt line 1795]
\noindent\textbf{EasyCrypt name.}
\begin{lstlisting}[language=EasyCrypt,basicstyle=\ttfamily\scriptsize]
adversary_budgeted_paper_sim_signing_count_preserved
\end{lstlisting}
\noindent\textbf{Checked EasyCrypt statement.}
\begin{lstlisting}[language=EasyCrypt,basicstyle=\ttfamily\scriptsize]
lemma adversary_budgeted_paper_sim_signing_count_preserved :
  hoare[
    ROMInternalTranscriptBudgetedPaperSimAsNMA(A, Samp, H).A.forge :
    budgeted_paper_sim_signing_count_discipline
      ROMInternalTranscriptBudgetedPaperSimAsNMA.signing_count ==>
    budgeted_paper_sim_signing_count_discipline
      ROMInternalTranscriptBudgetedPaperSimAsNMA.signing_count].
\end{lstlisting}
\noindent\textbf{Formal mathematical reading.}
Mathematical theorem. This is a relational program theorem: two games or procedures are coupled so that a precondition on their initial states implies the stated postcondition on final states and return values.

\noindent\textbf{Role in the proof.}
Use in this chapter. It proves an exact game replacement or simulator equivalence.

\noindent\textbf{Proof-script interpretation.}
Proof-script reading. The machine proof decomposes the theorem into rewrites, program-logic obligations, relational equivalences, and arithmetic side conditions.
\emph{Main tactics visible in the proof script:} \ec{proc}, \ec{wp}, \ec{call}, \ec{apply}.
\emph{Named identifiers syntactically cited in the proof block:}
\begin{lstlisting}[language=EasyCrypt,basicstyle=\ttfamily\scriptsize]
ROMInternalTranscriptBudgetedPaperSimAsNMA
budgeted_paper_sim_sign_oracle_preserves_signing_count
budgeted_paper_sim_signing_count_discipline
signing_count
\end{lstlisting}

\end{formaltheorem}

\begin{formaltheorem}[EasyCrypt line 1816]
\noindent\textbf{EasyCrypt name.}
\begin{lstlisting}[language=EasyCrypt,basicstyle=\ttfamily\scriptsize]
budgeted_paper_sim_main_signing_count_preserved
\end{lstlisting}
\noindent\textbf{Checked EasyCrypt statement.}
\begin{lstlisting}[language=EasyCrypt,basicstyle=\ttfamily\scriptsize]
lemma budgeted_paper_sim_main_signing_count_preserved :
  hoare[ROMInternalTranscriptBudgetedPaperSimAsNMA(A, Samp, H).forge :
    true ==>
    budgeted_paper_sim_signing_count_discipline
      ROMInternalTranscriptBudgetedPaperSimAsNMA.signing_count].
\end{lstlisting}
\noindent\textbf{Formal mathematical reading.}
Mathematical theorem. This is a relational program theorem: two games or procedures are coupled so that a precondition on their initial states implies the stated postcondition on final states and return values.

\noindent\textbf{Role in the proof.}
Use in this chapter. It proves an exact game replacement or simulator equivalence.

\noindent\textbf{Proof-script interpretation.}
Proof-script reading. The machine proof decomposes the theorem into rewrites, program-logic obligations, relational equivalences, and arithmetic side conditions.
\emph{Main tactics visible in the proof script:} \ec{proc}, \ec{wp}, \ec{call}, \ec{conseq}, \ec{rewrite}.
\emph{Named identifiers syntactically cited in the proof block:}
\begin{lstlisting}[language=EasyCrypt,basicstyle=\ttfamily\scriptsize]
ROMInternalTranscriptBudgetedPaperSimAsNMA
adversary_budgeted_paper_sim_signing_count_preserved
budgeted_paper_sim_signing_count_discipline
signature_query_budget_count
signing_count
\end{lstlisting}

\end{formaltheorem}

\begin{formaltheorem}[EasyCrypt line 1836]
\noindent\textbf{EasyCrypt name.}
\begin{lstlisting}[language=EasyCrypt,basicstyle=\ttfamily\scriptsize]
sampler_expand_query_point_bound
\end{lstlisting}
\noindent\textbf{Checked EasyCrypt statement.}
\begin{lstlisting}[language=EasyCrypt,basicstyle=\ttfamily\scriptsize]
lemma sampler_expand_query_point_bound q :
  mu drandom_coins (fun coins => sampler_expand_query coins = q) <=
  signing_randomness_point_bound.
\end{lstlisting}
\noindent\textbf{Formal mathematical reading.}
Mathematical theorem. This is an inequality, usually an arithmetic loss bound, event inclusion bound, or monotonicity fact used to compose probabilities.

\noindent\textbf{Role in the proof.}
Use in this chapter. It contributes directly to an advantage bound, bad-event bound, or real-arithmetic composition step.

\noindent\textbf{Proof-script interpretation.}
Proof-script reading. The machine proof decomposes the theorem into rewrites, program-logic obligations, relational equivalences, and arithmetic side conditions.
\emph{Main tactics visible in the proof script:} \ec{rewrite}, \ec{case}, \ec{smt}, \ec{apply}.
\emph{Named identifiers syntactically cited in the proof block:}
\begin{lstlisting}[language=EasyCrypt,basicstyle=\ttfamily\scriptsize]
coins
drandom_coins
eq_q
ler_trans
mu
mu0
mu_le
sampler_expand_query
signing_coin_distribution_token_point_bound
signing_entropy_token_cardinality
signing_entropy_token_of_coins
signing_randomness_point_bound
target
\end{lstlisting}

\end{formaltheorem}

\begin{formaltheorem}[EasyCrypt line 1859]
\noindent\textbf{EasyCrypt name.}
\begin{lstlisting}[language=EasyCrypt,basicstyle=\ttfamily\scriptsize]
sampler_expand_query_log_fset_bound
\end{lstlisting}
\noindent\textbf{Checked EasyCrypt statement.}
\begin{lstlisting}[language=EasyCrypt,basicstyle=\ttfamily\scriptsize]
lemma sampler_expand_query_log_fset_bound (qs : ro_query list) :
  mu drandom_coins
    (fun coins => sampler_expand_query coins \in qs) <=
  (card (oflist qs))%r * signing_randomness_point_bound.
\end{lstlisting}
\noindent\textbf{Formal mathematical reading.}
Mathematical theorem. This is an inequality, usually an arithmetic loss bound, event inclusion bound, or monotonicity fact used to compose probabilities.

\noindent\textbf{Role in the proof.}
Use in this chapter. It contributes directly to an advantage bound, bad-event bound, or real-arithmetic composition step.

\noindent\textbf{Proof-script interpretation.}
Proof-script reading. The machine proof decomposes the theorem into rewrites, program-logic obligations, relational equivalences, and arithmetic side conditions.
\emph{Main tactics visible in the proof script:} \ec{rewrite}, \ec{smt}, \ec{apply}.
\emph{Named identifiers syntactically cited in the proof block:}
\begin{lstlisting}[language=EasyCrypt,basicstyle=\ttfamily\scriptsize]
coins
drandom_coins
mem_oflist
mu_eq
mu_mem_le_gen
oflist
qs
ro_query
sampler_expand_query
sampler_expand_query_point_bound
signing_randomness_point_bound
\end{lstlisting}

\end{formaltheorem}

\begin{formaltheorem}[EasyCrypt line 1877]
\noindent\textbf{EasyCrypt name.}
\begin{lstlisting}[language=EasyCrypt,basicstyle=\ttfamily\scriptsize]
sampler_expand_query_log_budget_bound
\end{lstlisting}
\noindent\textbf{Checked EasyCrypt statement.}
\begin{lstlisting}[language=EasyCrypt,basicstyle=\ttfamily\scriptsize]
lemma sampler_expand_query_log_budget_bound (qs : ro_query list) hc :
  0 <= hc =>
  hc <= hash_query_budget_count =>
  size qs <= hc =>
  mu drandom_coins
    (fun coins => sampler_expand_query coins \in qs) <=
  rom_hash_query_budget / challenge_support_cardinality_lower_bound.
\end{lstlisting}
\noindent\textbf{Formal mathematical reading.}
Mathematical theorem. This is an inequality, usually an arithmetic loss bound, event inclusion bound, or monotonicity fact used to compose probabilities.

\noindent\textbf{Role in the proof.}
Use in this chapter. It contributes directly to an advantage bound, bad-event bound, or real-arithmetic composition step.

\noindent\textbf{Proof-script interpretation.}
Proof-script reading. The machine proof decomposes the theorem into rewrites, program-logic obligations, relational equivalences, and arithmetic side conditions.
\emph{Main tactics visible in the proof script:} \ec{apply}, \ec{rewrite}, \ec{smt}.
\emph{Named identifiers syntactically cited in the proof block:}
\begin{lstlisting}[language=EasyCrypt,basicstyle=\ttfamily\scriptsize]
card
challenge_support_cardinality_lower_bound
fcard_oflist
hash_query_budget_count
hash_query_count
hc_budget
hc_ge0
ler_trans
oflist
qs
rom_hash_query_budget
sampler_expand_query_log_fset_bound
signing_entropy_token_cardinality
signing_randomness_point_bound
size_bound
\end{lstlisting}

\end{formaltheorem}

\begin{formaltheorem}[EasyCrypt line 1896]
\noindent\textbf{EasyCrypt name.}
\begin{lstlisting}[language=EasyCrypt,basicstyle=\ttfamily\scriptsize]
sampler_expand_query_joint_log_budget_bound
\end{lstlisting}
\noindent\textbf{Checked EasyCrypt statement.}
\begin{lstlisting}[language=EasyCrypt,basicstyle=\ttfamily\scriptsize]
lemma sampler_expand_query_joint_log_budget_bound
    (hash_qs sampler_qs : ro_query list) hc sc :
  0 <= hc =>
  hc <= hash_query_budget_count =>
  0 <= sc =>
  sc <= signature_query_budget_count =>
  size hash_qs <= hc =>
  size sampler_qs <= sc =>
  mu drandom_coins
    (fun coins =>
      sampler_expand_query coins \in hash_qs \/
      sampler_expand_query coins \in sampler_qs) <=
  (rom_hash_query_budget + rom_signature_query_budget) /
    challenge_support_cardinality_lower_bound.
\end{lstlisting}
\noindent\textbf{Formal mathematical reading.}
Mathematical theorem. This is an inequality, usually an arithmetic loss bound, event inclusion bound, or monotonicity fact used to compose probabilities.

\noindent\textbf{Role in the proof.}
Use in this chapter. It contributes directly to an advantage bound, bad-event bound, or real-arithmetic composition step.

\noindent\textbf{Proof-script interpretation.}
Proof-script reading. The machine proof decomposes the theorem into rewrites, program-logic obligations, relational equivalences, and arithmetic side conditions.
\emph{Main tactics visible in the proof script:} \ec{rewrite}, \ec{apply}, \ec{smt}.
\emph{Named identifiers syntactically cited in the proof block:}
\begin{lstlisting}[language=EasyCrypt,basicstyle=\ttfamily\scriptsize]
card
challenge_support_cardinality_lower_bound
coins
fcard_oflist
hash_qs
hash_query_budget_count
hash_query_count
hash_size
hc_budget
hc_ge0
ler_trans
mem_cat
mu_eq
oflist
rom_hash_query_budget
rom_signature_query_budget
sampler_expand_query
sampler_expand_query_log_fset_bound
sampler_qs
sampler_size
sc_budget
sc_ge0
signature_query_budget_count
signature_query_count
signing_entropy_token_cardinality
signing_randomness_point_bound
size_cat
\end{lstlisting}

\end{formaltheorem}

\begin{formaltheorem}[EasyCrypt line 1929]
\noindent\textbf{EasyCrypt name.}
\begin{lstlisting}[language=EasyCrypt,basicstyle=\ttfamily\scriptsize]
concrete_lazy_rom_sampler_expand_query_fresh_le
\end{lstlisting}
\noindent\textbf{Checked EasyCrypt statement.}
\begin{lstlisting}[language=EasyCrypt,basicstyle=\ttfamily\scriptsize]
lemma concrete_lazy_rom_sampler_expand_query_fresh_le
    seed_coins (p : random_coins -> bool) &m :
  sampler_expand_query seed_coins \notin HAETAE_RO.FRO.m{m} =>
  Pr[HAETAE_RO.FRO.get(sampler_expand_query seed_coins) @ &m :
       p (ro_signing_coins res)] <=
  mu drandom_coins p.
\end{lstlisting}
\noindent\textbf{Formal mathematical reading.}
Mathematical theorem. This theorem has the form \(\Pr[G:E] \le B\) is a game-probability statement.  It is one of the exact hops, bad-event bounds, or final advantage inequalities used in the reduction.

\noindent\textbf{Role in the proof.}
Use in this chapter. It contributes directly to an advantage bound, bad-event bound, or real-arithmetic composition step.

\noindent\textbf{Proof-script interpretation.}
Proof-script reading. The machine proof decomposes the theorem into rewrites, program-logic obligations, relational equivalences, and arithmetic side conditions.
\emph{Main tactics visible in the proof script:} \ec{byphoare}, \ec{proc}, \ec{wp}, \ec{rnd}, \ec{rewrite}.
\emph{Named identifiers syntactically cited in the proof block:}
\begin{lstlisting}[language=EasyCrypt,basicstyle=\ttfamily\scriptsize]
FRO
HAETAE_RO
arg
fresh
hr
lerr
notin
ro_signing_coins
sampler_expand_query
sampler_expand_query_dro_output_ro_signing_coins
seed_coins
\end{lstlisting}

\end{formaltheorem}

\begin{formaltheorem}[EasyCrypt line 1949]
\noindent\textbf{EasyCrypt name.}
\begin{lstlisting}[language=EasyCrypt,basicstyle=\ttfamily\scriptsize]
concrete_lazy_rom_sampler_expand_query_fresh_eq
\end{lstlisting}
\noindent\textbf{Checked EasyCrypt statement.}
\begin{lstlisting}[language=EasyCrypt,basicstyle=\ttfamily\scriptsize]
lemma concrete_lazy_rom_sampler_expand_query_fresh_eq
    seed_coins (p : random_coins -> bool) :
  phoare[HAETAE_RO.FRO.get :
    arg = sampler_expand_query seed_coins /\
    sampler_expand_query seed_coins \notin HAETAE_RO.FRO.m
    ==> p (ro_signing_coins res)] =
  (mu drandom_coins p).
\end{lstlisting}
\noindent\textbf{Formal mathematical reading.}
Mathematical theorem. This is a relational program theorem: two games or procedures are coupled so that a precondition on their initial states implies the stated postcondition on final states and return values.

\noindent\textbf{Role in the proof.}
Use in this chapter. It contributes directly to an advantage bound, bad-event bound, or real-arithmetic composition step.

\noindent\textbf{Proof-script interpretation.}
Proof-script reading. The machine proof decomposes the theorem into rewrites, program-logic obligations, relational equivalences, and arithmetic side conditions.
\emph{Main tactics visible in the proof script:} \ec{proc}, \ec{wp}, \ec{rnd}, \ec{rewrite}.
\emph{Named identifiers syntactically cited in the proof block:}
\begin{lstlisting}[language=EasyCrypt,basicstyle=\ttfamily\scriptsize]
hr
sampler_expand_query_dro_output_ro_signing_coins
\end{lstlisting}

\end{formaltheorem}

\begin{formaltheorem}[EasyCrypt line 1964]
\noindent\textbf{EasyCrypt name.}
\begin{lstlisting}[language=EasyCrypt,basicstyle=\ttfamily\scriptsize]
concrete_lazy_rom_get_preserves_sampler_rom_covered
\end{lstlisting}
\noindent\textbf{Checked EasyCrypt statement.}
\begin{lstlisting}[language=EasyCrypt,basicstyle=\ttfamily\scriptsize]
lemma concrete_lazy_rom_get_preserves_sampler_rom_covered
    q hash_qs sampler_qs :
  hoare[HAETAE_RO.FRO.get :
    arg = q /\
    sampler_rom_covered HAETAE_RO.FRO.m hash_qs sampler_qs /\
    (forall seed_coins,
       q = sampler_expand_query seed_coins =>
       q \in hash_qs \/ q \in sampler_qs) ==>
    sampler_rom_covered HAETAE_RO.FRO.m hash_qs sampler_qs].
\end{lstlisting}
\noindent\textbf{Formal mathematical reading.}
Mathematical theorem. This is a relational program theorem: two games or procedures are coupled so that a precondition on their initial states implies the stated postcondition on final states and return values.

\noindent\textbf{Role in the proof.}
Use in this chapter. It contributes directly to an advantage bound, bad-event bound, or real-arithmetic composition step.

\noindent\textbf{Proof-script interpretation.}
Proof-script reading. The machine proof decomposes the theorem into rewrites, program-logic obligations, relational equivalences, and arithmetic side conditions.
\emph{Main tactics visible in the proof script:} \ec{proc}, \ec{wp}, \ec{rnd}, \ec{rewrite}, \ec{smt}.
\emph{Named identifiers syntactically cited in the proof block:}
\begin{lstlisting}[language=EasyCrypt,basicstyle=\ttfamily\scriptsize]
mem_set
sampler_rom_covered
\end{lstlisting}

\end{formaltheorem}

\begin{formaltheorem}[EasyCrypt line 1982]
\noindent\textbf{EasyCrypt name.}
\begin{lstlisting}[language=EasyCrypt,basicstyle=\ttfamily\scriptsize]
concrete_lazy_rom_get_preserves_sampler_rom_covered_arg
\end{lstlisting}
\noindent\textbf{Checked EasyCrypt statement.}
\begin{lstlisting}[language=EasyCrypt,basicstyle=\ttfamily\scriptsize]
lemma concrete_lazy_rom_get_preserves_sampler_rom_covered_arg
    hash_qs sampler_qs :
  hoare[HAETAE_RO.FRO.get :
    sampler_rom_covered HAETAE_RO.FRO.m hash_qs sampler_qs /\
    (forall seed_coins,
       arg = sampler_expand_query seed_coins =>
       arg \in hash_qs \/ arg \in sampler_qs) ==>
    sampler_rom_covered HAETAE_RO.FRO.m hash_qs sampler_qs].
\end{lstlisting}
\noindent\textbf{Formal mathematical reading.}
Mathematical theorem. This is a relational program theorem: two games or procedures are coupled so that a precondition on their initial states implies the stated postcondition on final states and return values.

\noindent\textbf{Role in the proof.}
Use in this chapter. It contributes directly to an advantage bound, bad-event bound, or real-arithmetic composition step.

\noindent\textbf{Proof-script interpretation.}
Proof-script reading. The machine proof decomposes the theorem into rewrites, program-logic obligations, relational equivalences, and arithmetic side conditions.
\emph{Main tactics visible in the proof script:} \ec{proc}, \ec{wp}, \ec{rnd}, \ec{rewrite}, \ec{smt}.
\emph{Named identifiers syntactically cited in the proof block:}
\begin{lstlisting}[language=EasyCrypt,basicstyle=\ttfamily\scriptsize]
mem_set
sampler_rom_covered
\end{lstlisting}

\end{formaltheorem}

\begin{formaltheorem}[EasyCrypt line 1999]
\noindent\textbf{EasyCrypt name.}
\begin{lstlisting}[language=EasyCrypt,basicstyle=\ttfamily\scriptsize]
concrete_lazy_rom_get_preserves_sampler_rom_covered_budgeted_logs
\end{lstlisting}
\noindent\textbf{Checked EasyCrypt statement.}
\begin{lstlisting}[language=EasyCrypt,basicstyle=\ttfamily\scriptsize]
lemma concrete_lazy_rom_get_preserves_sampler_rom_covered_budgeted_logs :
  hoare[HAETAE_RO.FRO.get :
    sampler_rom_covered
      HAETAE_RO.FRO.m
      ROMInternalTranscriptBudgetedPaperSimAsNMA.adversary_hash_queries
      ROMInternalTranscriptBudgetedPaperSimAsNMA.sampler_expand_queries /\
    (forall seed_coins,
       arg = sampler_expand_query seed_coins =>
       arg \in
         ROMInternalTranscriptBudgetedPaperSimAsNMA.adversary_hash_queries \/
       arg \in
         ROMInternalTranscriptBudgetedPaperSimAsNMA.sampler_expand_queries) ==>
    sampler_rom_covered
      HAETAE_RO.FRO.m
      ROMInternalTranscriptBudgetedPaperSimAsNMA.adversary_hash_queries
      ROMInternalTranscriptBudgetedPaperSimAsNMA.sampler_expand_queries].
\end{lstlisting}
\noindent\textbf{Formal mathematical reading.}
Mathematical theorem. This is a relational program theorem: two games or procedures are coupled so that a precondition on their initial states implies the stated postcondition on final states and return values.

\noindent\textbf{Role in the proof.}
Use in this chapter. It contributes directly to an advantage bound, bad-event bound, or real-arithmetic composition step.

\noindent\textbf{Proof-script interpretation.}
Proof-script reading. The machine proof decomposes the theorem into rewrites, program-logic obligations, relational equivalences, and arithmetic side conditions.
\emph{Main tactics visible in the proof script:} \ec{proc}, \ec{wp}, \ec{rnd}, \ec{rewrite}, \ec{smt}.
\emph{Named identifiers syntactically cited in the proof block:}
\begin{lstlisting}[language=EasyCrypt,basicstyle=\ttfamily\scriptsize]
mem_set
sampler_rom_covered
\end{lstlisting}

\end{formaltheorem}

\begin{formaltheorem}[EasyCrypt line 2024]
\noindent\textbf{EasyCrypt name.}
\begin{lstlisting}[language=EasyCrypt,basicstyle=\ttfamily\scriptsize]
concrete_lazy_rom_get_preserves_sampler_rom_covered_hash_cons
\end{lstlisting}
\noindent\textbf{Checked EasyCrypt statement.}
\begin{lstlisting}[language=EasyCrypt,basicstyle=\ttfamily\scriptsize]
lemma concrete_lazy_rom_get_preserves_sampler_rom_covered_hash_cons
    q hash_qs sampler_qs :
  hoare[HAETAE_RO.FRO.get :
    arg = q /\
    sampler_rom_covered HAETAE_RO.FRO.m hash_qs sampler_qs ==>
    sampler_rom_covered HAETAE_RO.FRO.m (q :: hash_qs) sampler_qs].
\end{lstlisting}
\noindent\textbf{Formal mathematical reading.}
Mathematical theorem. This is a relational program theorem: two games or procedures are coupled so that a precondition on their initial states implies the stated postcondition on final states and return values.

\noindent\textbf{Role in the proof.}
Use in this chapter. It contributes directly to an advantage bound, bad-event bound, or real-arithmetic composition step.

\noindent\textbf{Proof-script interpretation.}
Proof-script reading. The machine proof decomposes the theorem into rewrites, program-logic obligations, relational equivalences, and arithmetic side conditions.
\emph{Main tactics visible in the proof script:} \ec{proc}, \ec{wp}, \ec{rnd}, \ec{rewrite}, \ec{smt}.
\emph{Named identifiers syntactically cited in the proof block:}
\begin{lstlisting}[language=EasyCrypt,basicstyle=\ttfamily\scriptsize]
mem_set
sampler_rom_covered
\end{lstlisting}

\end{formaltheorem}

\begin{formaltheorem}[EasyCrypt line 2039]
\noindent\textbf{EasyCrypt name.}
\begin{lstlisting}[language=EasyCrypt,basicstyle=\ttfamily\scriptsize]
concrete_lazy_rom_get_preserves_sampler_rom_covered_sampler_cons
\end{lstlisting}
\noindent\textbf{Checked EasyCrypt statement.}
\begin{lstlisting}[language=EasyCrypt,basicstyle=\ttfamily\scriptsize]
lemma concrete_lazy_rom_get_preserves_sampler_rom_covered_sampler_cons
    q hash_qs sampler_qs :
  hoare[HAETAE_RO.FRO.get :
    arg = q /\
    sampler_rom_covered HAETAE_RO.FRO.m hash_qs sampler_qs ==>
    sampler_rom_covered HAETAE_RO.FRO.m hash_qs (q :: sampler_qs)].
\end{lstlisting}
\noindent\textbf{Formal mathematical reading.}
Mathematical theorem. This is a relational program theorem: two games or procedures are coupled so that a precondition on their initial states implies the stated postcondition on final states and return values.

\noindent\textbf{Role in the proof.}
Use in this chapter. It contributes directly to an advantage bound, bad-event bound, or real-arithmetic composition step.

\noindent\textbf{Proof-script interpretation.}
Proof-script reading. The machine proof decomposes the theorem into rewrites, program-logic obligations, relational equivalences, and arithmetic side conditions.
\emph{Main tactics visible in the proof script:} \ec{proc}, \ec{wp}, \ec{rnd}, \ec{rewrite}, \ec{smt}.
\emph{Named identifiers syntactically cited in the proof block:}
\begin{lstlisting}[language=EasyCrypt,basicstyle=\ttfamily\scriptsize]
mem_set
sampler_rom_covered
\end{lstlisting}

\end{formaltheorem}

\begin{formaltheorem}[EasyCrypt line 2054]
\noindent\textbf{EasyCrypt name.}
\begin{lstlisting}[language=EasyCrypt,basicstyle=\ttfamily\scriptsize]
sampler_bad_prequery_one_step_bound_nonnegative
\end{lstlisting}
\noindent\textbf{Checked EasyCrypt statement.}
\begin{lstlisting}[language=EasyCrypt,basicstyle=\ttfamily\scriptsize]
lemma sampler_bad_prequery_one_step_bound_nonnegative :
  0%r <= rom_hash_query_budget / challenge_support_cardinality_lower_bound.
\end{lstlisting}
\noindent\textbf{Formal mathematical reading.}
Mathematical theorem. This is an inequality, usually an arithmetic loss bound, event inclusion bound, or monotonicity fact used to compose probabilities.

\noindent\textbf{Role in the proof.}
Use in this chapter. It contributes directly to an advantage bound, bad-event bound, or real-arithmetic composition step.

\noindent\textbf{Proof-script interpretation.}
Proof-script reading. The machine proof decomposes the theorem into rewrites, program-logic obligations, relational equivalences, and arithmetic side conditions.
\emph{Main tactics visible in the proof script:} \ec{apply}.
\emph{Named identifiers syntactically cited in the proof block:}
\begin{lstlisting}[language=EasyCrypt,basicstyle=\ttfamily\scriptsize]
challenge_support_cardinality_lower_bound_nonnegative
divr_ge0
rom_hash_query_budget_nonnegative
\end{lstlisting}

\end{formaltheorem}

\begin{formaltheorem}[EasyCrypt line 2062]
\noindent\textbf{EasyCrypt name.}
\begin{lstlisting}[language=EasyCrypt,basicstyle=\ttfamily\scriptsize]
budgeted_paper_sim_sampler_bad_prequery_forge_fel_bound
\end{lstlisting}
\noindent\textbf{Checked EasyCrypt statement.}
\begin{lstlisting}[language=EasyCrypt,basicstyle=\ttfamily\scriptsize]
lemma budgeted_paper_sim_sampler_bad_prequery_forge_fel_bound pk &m :
  Pr[ROMInternalTranscriptBudgetedPaperSimAsNMA(A, Samp, H).forge(pk) @ &m :
       ROMInternalTranscriptBudgetedPaperSimAsNMA.sampler_bad_prequery] <=
    rom_signature_query_budget *
      ((rom_hash_query_budget + rom_signature_query_budget) /
       challenge_support_cardinality_lower_bound).
\end{lstlisting}
\noindent\textbf{Formal mathematical reading.}
Mathematical theorem. This theorem has the form \(\Pr[G:E] \le B\) is a game-probability statement.  It is one of the exact hops, bad-event bounds, or final advantage inequalities used in the reduction.

\noindent\textbf{Role in the proof.}
Use in this chapter. It contributes directly to an advantage bound, bad-event bound, or real-arithmetic composition step.

\noindent\textbf{Proof-script interpretation.}
Proof-script reading. The machine proof decomposes the theorem into rewrites, program-logic obligations, relational equivalences, and arithmetic side conditions.
\emph{Main tactics visible in the proof script:} \ec{rewrite}, \ec{smt}, \ec{inline}, \ec{wp}, \ec{proc}, \ec{conseq}, \ec{hoare}, \ec{call}, \ec{rnd}, \ec{apply}.
\emph{Named identifiers syntactically cited in the proof block:}
\begin{lstlisting}[language=EasyCrypt,basicstyle=\ttfamily\scriptsize]
AH
BRA
Bigreal
RField
ROMInternalTranscriptBudgetedPaperSimAsNMA
Samp
StdBigop
adversary_hash_count
adversary_hash_queries
budgeted_paper_sim_sampler_freshness_discipline
challenge_support_cardinality_lower_bound
fel
forge
get
hash_query_budget_count
hr
intmulr
old_bad
old_count
rom_hash_query_budget
rom_signature_query_budget
sampler_bad_prequery
sampler_expand_queries
sampler_expand_query
sampler_expand_query_joint_log_budget_bound
seed_coins
sign
signature_query_budget_count
signature_query_count
signing_coin_distribution_lossless
signing_count
sumri_const
trivial
\end{lstlisting}

\end{formaltheorem}

\begin{formaltheorem}[EasyCrypt line 2202]
\noindent\textbf{EasyCrypt name.}
\begin{lstlisting}[language=EasyCrypt,basicstyle=\ttfamily\scriptsize]
budgeted_paper_sim_sampler_bad_prequery_forge_fel_phoare
\end{lstlisting}
\noindent\textbf{Checked EasyCrypt statement.}
\begin{lstlisting}[language=EasyCrypt,basicstyle=\ttfamily\scriptsize]
lemma budgeted_paper_sim_sampler_bad_prequery_forge_fel_phoare :
  phoare[ROMInternalTranscriptBudgetedPaperSimAsNMA(A, Samp, H).forge :
       true ==> ROMInternalTranscriptBudgetedPaperSimAsNMA.sampler_bad_prequery] <=
    (rom_signature_query_budget *
      ((rom_hash_query_budget + rom_signature_query_budget) /
       challenge_support_cardinality_lower_bound)).
\end{lstlisting}
\noindent\textbf{Formal mathematical reading.}
Mathematical theorem. This is a relational program theorem: two games or procedures are coupled so that a precondition on their initial states implies the stated postcondition on final states and return values.

\noindent\textbf{Role in the proof.}
Use in this chapter. It contributes directly to an advantage bound, bad-event bound, or real-arithmetic composition step.

\noindent\textbf{Proof-script interpretation.}
Proof-script reading. The machine proof decomposes the theorem into rewrites, program-logic obligations, relational equivalences, and arithmetic side conditions.
\emph{Named identifiers syntactically cited in the proof block:}
\begin{lstlisting}[language=EasyCrypt,basicstyle=\ttfamily\scriptsize]
arg
budgeted_paper_sim_sampler_bad_prequery_forge_fel_bound
bypr
exact
\end{lstlisting}

\end{formaltheorem}

\begin{formaltheorem}[EasyCrypt line 2214]
\noindent\textbf{EasyCrypt name.}
\begin{lstlisting}[language=EasyCrypt,basicstyle=\ttfamily\scriptsize]
budgeted_paper_sim_sampler_bad_prequery_nma_fel_bound
\end{lstlisting}
\noindent\textbf{Checked EasyCrypt statement.}
\begin{lstlisting}[language=EasyCrypt,basicstyle=\ttfamily\scriptsize]
lemma budgeted_paper_sim_sampler_bad_prequery_nma_fel_bound &m :
  Pr[SIG.UF_NMA(H, HAETAE,
       ROMInternalTranscriptBudgetedPaperSimAsNMA(A, Samp)).main() @ &m :
       ROMInternalTranscriptBudgetedPaperSimAsNMA.sampler_bad_prequery] <=
    rom_signature_query_budget *
      ((rom_hash_query_budget + rom_signature_query_budget) /
       challenge_support_cardinality_lower_bound).
\end{lstlisting}
\noindent\textbf{Formal mathematical reading.}
Mathematical theorem. This theorem has the form \(\Pr[G:E] \le B\) is a game-probability statement.  It is one of the exact hops, bad-event bounds, or final advantage inequalities used in the reduction.

\noindent\textbf{Role in the proof.}
Use in this chapter. It contributes directly to an advantage bound, bad-event bound, or real-arithmetic composition step.

\noindent\textbf{Proof-script interpretation.}
Proof-script reading. The machine proof decomposes the theorem into rewrites, program-logic obligations, relational equivalences, and arithmetic side conditions.
\emph{Main tactics visible in the proof script:} \ec{byphoare}, \ec{proc}, \ec{inline}, \ec{wp}, \ec{call}, \ec{rnd}.
\emph{Named identifiers syntactically cited in the proof block:}
\begin{lstlisting}[language=EasyCrypt,basicstyle=\ttfamily\scriptsize]
HAETAE
ROMInternalTranscriptBudgetedPaperSimAsNMA
budgeted_paper_sim_sampler_bad_prequery_forge_fel_phoare
kg
sampler_bad_prequery
verify
\end{lstlisting}

\end{formaltheorem}

\begin{formaltheorem}[EasyCrypt line 2239]
\noindent\textbf{EasyCrypt name.}
\begin{lstlisting}[language=EasyCrypt,basicstyle=\ttfamily\scriptsize]
concrete_lazy_rom_sampler_expand_query_fresh_phoare
\end{lstlisting}
\noindent\textbf{Checked EasyCrypt statement.}
\begin{lstlisting}[language=EasyCrypt,basicstyle=\ttfamily\scriptsize]
lemma concrete_lazy_rom_sampler_expand_query_fresh_phoare
    seed_coins (p : random_coins -> bool) :
  phoare[HAETAE_RO.FRO.get :
      arg = sampler_expand_query seed_coins /\
      sampler_expand_query seed_coins \notin HAETAE_RO.FRO.m
      ==> p (ro_signing_coins res)] <=
    (mu drandom_coins p).
\end{lstlisting}
\noindent\textbf{Formal mathematical reading.}
Mathematical theorem. This is a relational program theorem: two games or procedures are coupled so that a precondition on their initial states implies the stated postcondition on final states and return values.

\noindent\textbf{Role in the proof.}
Use in this chapter. It contributes directly to an advantage bound, bad-event bound, or real-arithmetic composition step.

\noindent\textbf{Proof-script interpretation.}
Proof-script reading. The machine proof decomposes the theorem into rewrites, program-logic obligations, relational equivalences, and arithmetic side conditions.
\emph{Main tactics visible in the proof script:} \ec{rewrite}, \ec{apply}.
\emph{Named identifiers syntactically cited in the proof block:}
\begin{lstlisting}[language=EasyCrypt,basicstyle=\ttfamily\scriptsize]
arg_eq
bypr
concrete_lazy_rom_sampler_expand_query_fresh_le
fresh
seed_coins
\end{lstlisting}

\end{formaltheorem}

\begin{formaltheorem}[EasyCrypt line 2252]
\noindent\textbf{EasyCrypt name.}
\begin{lstlisting}[language=EasyCrypt,basicstyle=\ttfamily\scriptsize]
concrete_lazy_rom_sampler_expand_query_fresh_arg_phoare
\end{lstlisting}
\noindent\textbf{Checked EasyCrypt statement.}
\begin{lstlisting}[language=EasyCrypt,basicstyle=\ttfamily\scriptsize]
lemma concrete_lazy_rom_sampler_expand_query_fresh_arg_phoare
    (p : random_coins -> bool) :
  phoare[HAETAE_RO.FRO.get :
      (exists seed_coins,
         arg = sampler_expand_query seed_coins /\
         sampler_expand_query seed_coins \notin HAETAE_RO.FRO.m)
      ==> p (ro_signing_coins res)] <=
    (mu drandom_coins p).
\end{lstlisting}
\noindent\textbf{Formal mathematical reading.}
Mathematical theorem. This is a relational program theorem: two games or procedures are coupled so that a precondition on their initial states implies the stated postcondition on final states and return values.

\noindent\textbf{Role in the proof.}
Use in this chapter. It contributes directly to an advantage bound, bad-event bound, or real-arithmetic composition step.

\noindent\textbf{Proof-script interpretation.}
Proof-script reading. The machine proof decomposes the theorem into rewrites, program-logic obligations, relational equivalences, and arithmetic side conditions.
\emph{Main tactics visible in the proof script:} \ec{rewrite}, \ec{apply}.
\emph{Named identifiers syntactically cited in the proof block:}
\begin{lstlisting}[language=EasyCrypt,basicstyle=\ttfamily\scriptsize]
arg_eq
bypr
concrete_lazy_rom_sampler_expand_query_fresh_le
fresh
seed_coins
\end{lstlisting}

\end{formaltheorem}

\begin{formaltheorem}[EasyCrypt line 2266]
\noindent\textbf{EasyCrypt name.}
\begin{lstlisting}[language=EasyCrypt,basicstyle=\ttfamily\scriptsize]
concrete_lazy_rom_sampler_expand_query_fresh_arg_clean_phoare
\end{lstlisting}
\noindent\textbf{Checked EasyCrypt statement.}
\begin{lstlisting}[language=EasyCrypt,basicstyle=\ttfamily\scriptsize]
lemma concrete_lazy_rom_sampler_expand_query_fresh_arg_clean_phoare
    (p : random_coins -> bool) :
  phoare[HAETAE_RO.FRO.get :
      (exists seed_coins,
         arg = sampler_expand_query seed_coins /\
         sampler_expand_query seed_coins \notin HAETAE_RO.FRO.m) /\
      ! ROMInternalTranscriptBudgetedPaperSimAsNMA.sampler_bad_prequery
      ==> p (ro_signing_coins res) /\
          ! ROMInternalTranscriptBudgetedPaperSimAsNMA.sampler_bad_prequery] <=
    (mu drandom_coins p).
\end{lstlisting}
\noindent\textbf{Formal mathematical reading.}
Mathematical theorem. This is a relational program theorem: two games or procedures are coupled so that a precondition on their initial states implies the stated postcondition on final states and return values.

\noindent\textbf{Role in the proof.}
Use in this chapter. It contributes directly to an advantage bound, bad-event bound, or real-arithmetic composition step.

\noindent\textbf{Proof-script interpretation.}
Proof-script reading. The machine proof decomposes the theorem into rewrites, program-logic obligations, relational equivalences, and arithmetic side conditions.
\emph{Main tactics visible in the proof script:} \ec{proc}, \ec{wp}, \ec{rnd}, \ec{rewrite}.
\emph{Named identifiers syntactically cited in the proof block:}
\begin{lstlisting}[language=EasyCrypt,basicstyle=\ttfamily\scriptsize]
arg_eq
fresh
hr
lerr
sampler_expand_query_dro_output_ro_signing_coins
seed_coins
\end{lstlisting}

\end{formaltheorem}

\begin{formaltheorem}[EasyCrypt line 2285]
\noindent\textbf{EasyCrypt name.}
\begin{lstlisting}[language=EasyCrypt,basicstyle=\ttfamily\scriptsize]
concrete_ro_signing_attempt_sample_with_seed_fresh_structural_le
\end{lstlisting}
\noindent\textbf{Checked EasyCrypt statement.}
\begin{lstlisting}[language=EasyCrypt,basicstyle=\ttfamily\scriptsize]
lemma concrete_ro_signing_attempt_sample_with_seed_fresh_structural_le
    seed_coins sk m ctx (p : paper_sim_signature_sample -> bool) :
  phoare[ROSigningAttemptPaperSimSampler(HAETAE_RO.FRO).sample_with_seed :
      arg = (seed_coins, m, ctx) /\
      ROSigningAttemptPaperSimSampler.sk_current = sk /\
      sampler_expand_query seed_coins \notin HAETAE_RO.FRO.m
      ==> p res] <=
    (mu (dsigning_attempt_state haetae_mode
          sk m ctx)
       (fun st => p (paper_sim_sample_from_rejection_attempt st))).
\end{lstlisting}
\noindent\textbf{Formal mathematical reading.}
Mathematical theorem. This is a relational program theorem: two games or procedures are coupled so that a precondition on their initial states implies the stated postcondition on final states and return values.

\noindent\textbf{Role in the proof.}
Use in this chapter. It contributes directly to an advantage bound, bad-event bound, or real-arithmetic composition step.

\noindent\textbf{Proof-script interpretation.}
Proof-script reading. The machine proof decomposes the theorem into rewrites, program-logic obligations, relational equivalences, and arithmetic side conditions.
\emph{Main tactics visible in the proof script:} \ec{rewrite}, \ec{proc}, \ec{wp}, \ec{call}.
\emph{Named identifiers syntactically cited in the proof block:}
\begin{lstlisting}[language=EasyCrypt,basicstyle=\ttfamily\scriptsize]
coins
concrete_lazy_rom_sampler_expand_query_fresh_phoare
ctx
dmapE
dsigning_attempt_state
haetae_mode
paper_sim_sample_from_rejection_attempt
seed_coins
signing_attempt_state_of_coins
sk
\end{lstlisting}

\end{formaltheorem}

\begin{formaltheorem}[EasyCrypt line 2307]
\noindent\textbf{EasyCrypt name.}
\begin{lstlisting}[language=EasyCrypt,basicstyle=\ttfamily\scriptsize]
concrete_ro_signing_attempt_sample_with_seed_fresh_structural_arg_le
\end{lstlisting}
\noindent\textbf{Checked EasyCrypt statement.}
\begin{lstlisting}[language=EasyCrypt,basicstyle=\ttfamily\scriptsize]
lemma concrete_ro_signing_attempt_sample_with_seed_fresh_structural_arg_le
    sk m ctx (p : paper_sim_signature_sample -> bool) :
  phoare[ROSigningAttemptPaperSimSampler(HAETAE_RO.FRO).sample_with_seed :
      arg.`2 = m /\
      arg.`3 = ctx /\
      ROSigningAttemptPaperSimSampler.sk_current = sk /\
      sampler_expand_query arg.`1 \notin HAETAE_RO.FRO.m
      ==> p res] <=
    (mu (dsigning_attempt_state haetae_mode
          sk m ctx)
       (fun st => p (paper_sim_sample_from_rejection_attempt st))).
\end{lstlisting}
\noindent\textbf{Formal mathematical reading.}
Mathematical theorem. This is a relational program theorem: two games or procedures are coupled so that a precondition on their initial states implies the stated postcondition on final states and return values.

\noindent\textbf{Role in the proof.}
Use in this chapter. It contributes directly to an advantage bound, bad-event bound, or real-arithmetic composition step.

\noindent\textbf{Proof-script interpretation.}
Proof-script reading. The machine proof decomposes the theorem into rewrites, program-logic obligations, relational equivalences, and arithmetic side conditions.
\emph{Main tactics visible in the proof script:} \ec{rewrite}, \ec{proc}, \ec{wp}, \ec{call}, \ec{smt}.
\emph{Named identifiers syntactically cited in the proof block:}
\begin{lstlisting}[language=EasyCrypt,basicstyle=\ttfamily\scriptsize]
coins
concrete_lazy_rom_sampler_expand_query_fresh_arg_phoare
ctx
dmapE
dsigning_attempt_state
haetae_mode
paper_sim_sample_from_rejection_attempt
signing_attempt_state_of_coins
sk
\end{lstlisting}

\end{formaltheorem}

\begin{formaltheorem}[EasyCrypt line 2330]
\noindent\textbf{EasyCrypt name.}
\begin{lstlisting}[language=EasyCrypt,basicstyle=\ttfamily\scriptsize]
concrete_ro_signing_attempt_sample_with_seed_fresh_structural_arg_clean_le
\end{lstlisting}
\noindent\textbf{Checked EasyCrypt statement.}
\begin{lstlisting}[language=EasyCrypt,basicstyle=\ttfamily\scriptsize]
lemma concrete_ro_signing_attempt_sample_with_seed_fresh_structural_arg_clean_le
    sk m ctx (p : paper_sim_signature_sample -> bool) :
  phoare[ROSigningAttemptPaperSimSampler(HAETAE_RO.FRO).sample_with_seed :
      arg.`2 = m /\
      arg.`3 = ctx /\
      ROSigningAttemptPaperSimSampler.sk_current = sk /\
      ! ROMInternalTranscriptBudgetedPaperSimAsNMA.sampler_bad_prequery /\
      sampler_expand_query arg.`1 \notin HAETAE_RO.FRO.m
      ==> p res /\
          ! ROMInternalTranscriptBudgetedPaperSimAsNMA.sampler_bad_prequery] <=
    (mu (dsigning_attempt_state haetae_mode
          sk m ctx)
       (fun st => p (paper_sim_sample_from_rejection_attempt st))).
\end{lstlisting}
\noindent\textbf{Formal mathematical reading.}
Mathematical theorem. This is a relational program theorem: two games or procedures are coupled so that a precondition on their initial states implies the stated postcondition on final states and return values.

\noindent\textbf{Role in the proof.}
Use in this chapter. It contributes directly to an advantage bound, bad-event bound, or real-arithmetic composition step.

\noindent\textbf{Proof-script interpretation.}
Proof-script reading. The machine proof decomposes the theorem into rewrites, program-logic obligations, relational equivalences, and arithmetic side conditions.
\emph{Main tactics visible in the proof script:} \ec{rewrite}, \ec{proc}, \ec{wp}, \ec{call}, \ec{smt}.
\emph{Named identifiers syntactically cited in the proof block:}
\begin{lstlisting}[language=EasyCrypt,basicstyle=\ttfamily\scriptsize]
coins
concrete_lazy_rom_sampler_expand_query_fresh_arg_clean_phoare
ctx
dmapE
dsigning_attempt_state
haetae_mode
paper_sim_sample_from_rejection_attempt
signing_attempt_state_of_coins
sk
\end{lstlisting}

\end{formaltheorem}

\begin{formaltheorem}[EasyCrypt line 2355]
\noindent\textbf{EasyCrypt name.}
\begin{lstlisting}[language=EasyCrypt,basicstyle=\ttfamily\scriptsize]
concrete_lazy_rom_sampler_expand_query_guarded_clean_phoare
\end{lstlisting}
\noindent\textbf{Checked EasyCrypt statement.}
\begin{lstlisting}[language=EasyCrypt,basicstyle=\ttfamily\scriptsize]
lemma concrete_lazy_rom_sampler_expand_query_guarded_clean_phoare
    (p : random_coins -> bool) :
  phoare[HAETAE_RO.FRO.get :
      (! ROMInternalTranscriptBudgetedPaperSimAsNMA.sampler_bad_prequery =>
       exists seed_coins,
         arg = sampler_expand_query seed_coins /\
         sampler_expand_query seed_coins \notin HAETAE_RO.FRO.m)
      ==> p (ro_signing_coins res) /\
          ! ROMInternalTranscriptBudgetedPaperSimAsNMA.sampler_bad_prequery] <=
    (mu drandom_coins p).
\end{lstlisting}
\noindent\textbf{Formal mathematical reading.}
Mathematical theorem. This is a relational program theorem: two games or procedures are coupled so that a precondition on their initial states implies the stated postcondition on final states and return values.

\noindent\textbf{Role in the proof.}
Use in this chapter. It contributes directly to an advantage bound, bad-event bound, or real-arithmetic composition step.

\noindent\textbf{Proof-script interpretation.}
Proof-script reading. The machine proof decomposes the theorem into rewrites, program-logic obligations, relational equivalences, and arithmetic side conditions.
\emph{Main tactics visible in the proof script:} \ec{proc}, \ec{wp}, \ec{rnd}, \ec{case}, \ec{rewrite}, \ec{smt}, \ec{elim}.
\emph{Named identifiers syntactically cited in the proof block:}
\begin{lstlisting}[language=EasyCrypt,basicstyle=\ttfamily\scriptsize]
FRO
HAETAE_RO
ROMInternalTranscriptBudgetedPaperSimAsNMA
arg_eq
bad
clean
fresh
guarded
hr
lerr
mu0
mu_bounded
mu_eq
ro_output
sampler_bad_prequery
sampler_expand_query_dro_output_ro_signing_coins
seed_coins
\end{lstlisting}

\end{formaltheorem}

\begin{formaltheorem}[EasyCrypt line 2382]
\noindent\textbf{EasyCrypt name.}
\begin{lstlisting}[language=EasyCrypt,basicstyle=\ttfamily\scriptsize]
concrete_ro_signing_attempt_sample_with_seed_guarded_clean_structural_arg_le
\end{lstlisting}
\noindent\textbf{Checked EasyCrypt statement.}
\begin{lstlisting}[language=EasyCrypt,basicstyle=\ttfamily\scriptsize]
lemma concrete_ro_signing_attempt_sample_with_seed_guarded_clean_structural_arg_le
    sk m ctx (p : paper_sim_signature_sample -> bool) :
  phoare[ROSigningAttemptPaperSimSampler(HAETAE_RO.FRO).sample_with_seed :
      arg.`2 = m /\
      arg.`3 = ctx /\
      ROSigningAttemptPaperSimSampler.sk_current = sk /\
      (! ROMInternalTranscriptBudgetedPaperSimAsNMA.sampler_bad_prequery =>
       sampler_expand_query arg.`1 \notin HAETAE_RO.FRO.m)
      ==> p res /\
          ! ROMInternalTranscriptBudgetedPaperSimAsNMA.sampler_bad_prequery] <=
    (mu (dsigning_attempt_state haetae_mode
          sk m ctx)
       (fun st => p (paper_sim_sample_from_rejection_attempt st))).
\end{lstlisting}
\noindent\textbf{Formal mathematical reading.}
Mathematical theorem. This is a relational program theorem: two games or procedures are coupled so that a precondition on their initial states implies the stated postcondition on final states and return values.

\noindent\textbf{Role in the proof.}
Use in this chapter. It contributes directly to an advantage bound, bad-event bound, or real-arithmetic composition step.

\noindent\textbf{Proof-script interpretation.}
Proof-script reading. The machine proof decomposes the theorem into rewrites, program-logic obligations, relational equivalences, and arithmetic side conditions.
\emph{Main tactics visible in the proof script:} \ec{rewrite}, \ec{proc}, \ec{wp}, \ec{call}, \ec{smt}.
\emph{Named identifiers syntactically cited in the proof block:}
\begin{lstlisting}[language=EasyCrypt,basicstyle=\ttfamily\scriptsize]
coins
concrete_lazy_rom_sampler_expand_query_guarded_clean_phoare
ctx
dmapE
dsigning_attempt_state
haetae_mode
paper_sim_sample_from_rejection_attempt
signing_attempt_state_of_coins
sk
\end{lstlisting}

\end{formaltheorem}

\begin{formaltheorem}[EasyCrypt line 2407]
\noindent\textbf{EasyCrypt name.}
\begin{lstlisting}[language=EasyCrypt,basicstyle=\ttfamily\scriptsize]
concrete_ro_signing_attempt_sample_with_seed_fresh_structural_eq
\end{lstlisting}
\noindent\textbf{Checked EasyCrypt statement.}
\begin{lstlisting}[language=EasyCrypt,basicstyle=\ttfamily\scriptsize]
lemma concrete_ro_signing_attempt_sample_with_seed_fresh_structural_eq
    seed_coins sk m ctx (p : paper_sim_signature_sample -> bool) :
  phoare[ROSigningAttemptPaperSimSampler(HAETAE_RO.FRO).sample_with_seed :
      arg = (seed_coins, m, ctx) /\
      ROSigningAttemptPaperSimSampler.sk_current = sk /\
      sampler_expand_query seed_coins \notin HAETAE_RO.FRO.m
      ==> p res] =
    (mu (dsigning_attempt_state haetae_mode
          sk m ctx)
       (fun st => p (paper_sim_sample_from_rejection_attempt st))).
\end{lstlisting}
\noindent\textbf{Formal mathematical reading.}
Mathematical theorem. This is a relational program theorem: two games or procedures are coupled so that a precondition on their initial states implies the stated postcondition on final states and return values.

\noindent\textbf{Role in the proof.}
Use in this chapter. It contributes directly to an advantage bound, bad-event bound, or real-arithmetic composition step.

\noindent\textbf{Proof-script interpretation.}
Proof-script reading. The machine proof decomposes the theorem into rewrites, program-logic obligations, relational equivalences, and arithmetic side conditions.
\emph{Main tactics visible in the proof script:} \ec{rewrite}, \ec{proc}, \ec{wp}, \ec{call}.
\emph{Named identifiers syntactically cited in the proof block:}
\begin{lstlisting}[language=EasyCrypt,basicstyle=\ttfamily\scriptsize]
coins
concrete_lazy_rom_sampler_expand_query_fresh_eq
ctx
dmapE
dsigning_attempt_state
haetae_mode
paper_sim_sample_from_rejection_attempt
seed_coins
signing_attempt_state_of_coins
sk
\end{lstlisting}

\end{formaltheorem}

\begin{formaltheorem}[EasyCrypt line 2429]
\noindent\textbf{EasyCrypt name.}
\begin{lstlisting}[language=EasyCrypt,basicstyle=\ttfamily\scriptsize]
concrete_ro_exact_hyperball_sample_with_seed_exact
\end{lstlisting}
\noindent\textbf{Checked EasyCrypt statement.}
\begin{lstlisting}[language=EasyCrypt,basicstyle=\ttfamily\scriptsize]
lemma concrete_ro_exact_hyperball_sample_with_seed_exact
    seed_coins sk m ctx (p : paper_sim_signature_sample -> bool) :
  phoare[ROExactHyperballPaperSimSampler(HAETAE_RO.FRO).sample_with_seed :
      arg = (seed_coins, m, ctx) /\
      ROExactHyperballPaperSimSampler.sk_current = sk
      ==> p res] =
    (mu (dexact_hyperball_signing_attempt_state haetae_mode
          sk m ctx)
       (fun st => p (paper_sim_sample_from_rejection_attempt st))).
\end{lstlisting}
\noindent\textbf{Formal mathematical reading.}
Mathematical theorem. This is a relational program theorem: two games or procedures are coupled so that a precondition on their initial states implies the stated postcondition on final states and return values.

\noindent\textbf{Role in the proof.}
Use in this chapter. It contributes directly to an advantage bound, bad-event bound, or real-arithmetic composition step.

\noindent\textbf{Proof-script interpretation.}
Proof-script reading. The machine proof decomposes the theorem into rewrites, program-logic obligations, relational equivalences, and arithmetic side conditions.
\emph{Main tactics visible in the proof script:} \ec{proc}, \ec{wp}, \ec{rnd}, \ec{call}, \ec{smt}.
\emph{Named identifiers syntactically cited in the proof block:}
\begin{lstlisting}[language=EasyCrypt,basicstyle=\ttfamily\scriptsize]
ROExactHyperballPaperSimSampler
ro_output_distribution_lossless
sk
sk_current
\end{lstlisting}

\end{formaltheorem}

\begin{formaltheorem}[EasyCrypt line 2448]
\noindent\textbf{EasyCrypt name.}
\begin{lstlisting}[language=EasyCrypt,basicstyle=\ttfamily\scriptsize]
concrete_ro_exact_hyperball_sample_with_seed_exact_no_seed
\end{lstlisting}
\noindent\textbf{Checked EasyCrypt statement.}
\begin{lstlisting}[language=EasyCrypt,basicstyle=\ttfamily\scriptsize]
lemma concrete_ro_exact_hyperball_sample_with_seed_exact_no_seed
    sk m ctx (p : paper_sim_signature_sample -> bool) :
  phoare[ROExactHyperballPaperSimSampler(HAETAE_RO.FRO).sample_with_seed :
      arg.`2 = m /\
      arg.`3 = ctx /\
      ROExactHyperballPaperSimSampler.sk_current = sk
      ==> p res] =
    (mu (dexact_hyperball_signing_attempt_state haetae_mode
          sk m ctx)
       (fun st => p (paper_sim_sample_from_rejection_attempt st))).
\end{lstlisting}
\noindent\textbf{Formal mathematical reading.}
Mathematical theorem. This is a relational program theorem: two games or procedures are coupled so that a precondition on their initial states implies the stated postcondition on final states and return values.

\noindent\textbf{Role in the proof.}
Use in this chapter. It contributes directly to an advantage bound, bad-event bound, or real-arithmetic composition step.

\noindent\textbf{Proof-script interpretation.}
Proof-script reading. The machine proof decomposes the theorem into rewrites, program-logic obligations, relational equivalences, and arithmetic side conditions.
\emph{Main tactics visible in the proof script:} \ec{proc}, \ec{wp}, \ec{rnd}, \ec{call}, \ec{smt}.
\emph{Named identifiers syntactically cited in the proof block:}
\begin{lstlisting}[language=EasyCrypt,basicstyle=\ttfamily\scriptsize]
ROExactHyperballPaperSimSampler
ro_output_distribution_lossless
sk
sk_current
\end{lstlisting}

\end{formaltheorem}

\begin{formaltheorem}[EasyCrypt line 2468]
\noindent\textbf{EasyCrypt name.}
\begin{lstlisting}[language=EasyCrypt,basicstyle=\ttfamily\scriptsize]
concrete_ro_exact_hyperball_sample_with_seed_exact_pr
\end{lstlisting}
\noindent\textbf{Checked EasyCrypt statement.}
\begin{lstlisting}[language=EasyCrypt,basicstyle=\ttfamily\scriptsize]
lemma concrete_ro_exact_hyperball_sample_with_seed_exact_pr
    seed_coins m ctx (p : paper_sim_signature_sample -> bool) &m :
  Pr[ROExactHyperballPaperSimSampler(HAETAE_RO.FRO).sample_with_seed
       (seed_coins, m, ctx) @ &m : p res] =
    mu (dexact_hyperball_signing_attempt_state haetae_mode
          ROExactHyperballPaperSimSampler.sk_current{m} m ctx)
       (fun st => p (paper_sim_sample_from_rejection_attempt st)).
\end{lstlisting}
\noindent\textbf{Formal mathematical reading.}
Mathematical theorem. This theorem has the form \(\Pr[G:E] = B\) is a game-probability statement.  It is one of the exact hops, bad-event bounds, or final advantage inequalities used in the reduction.

\noindent\textbf{Role in the proof.}
Use in this chapter. It contributes directly to an advantage bound, bad-event bound, or real-arithmetic composition step.

\noindent\textbf{Proof-script interpretation.}
Proof-script reading. The machine proof decomposes the theorem into rewrites, program-logic obligations, relational equivalences, and arithmetic side conditions.
\emph{Main tactics visible in the proof script:} \ec{byphoare}.
\emph{Named identifiers syntactically cited in the proof block:}
\begin{lstlisting}[language=EasyCrypt,basicstyle=\ttfamily\scriptsize]
ROExactHyperballPaperSimSampler
concrete_ro_exact_hyperball_sample_with_seed_exact
ctx
seed_coins
sk_current
\end{lstlisting}

\end{formaltheorem}

\begin{formaltheorem}[EasyCrypt line 2480]
\noindent\textbf{EasyCrypt name.}
\begin{lstlisting}[language=EasyCrypt,basicstyle=\ttfamily\scriptsize]
concrete_ro_signing_attempt_sample_with_seed_preserves_sampler_rom_covered
\end{lstlisting}
\noindent\textbf{Checked EasyCrypt statement.}
\begin{lstlisting}[language=EasyCrypt,basicstyle=\ttfamily\scriptsize]
lemma concrete_ro_signing_attempt_sample_with_seed_preserves_sampler_rom_covered
    hash_qs sampler_qs :
  hoare[ROSigningAttemptPaperSimSampler(HAETAE_RO.FRO).sample_with_seed :
    sampler_rom_covered HAETAE_RO.FRO.m hash_qs sampler_qs /\
    sampler_expand_query arg.`1 \in sampler_qs ==>
    sampler_rom_covered HAETAE_RO.FRO.m hash_qs sampler_qs].
\end{lstlisting}
\noindent\textbf{Formal mathematical reading.}
Mathematical theorem. This is a relational program theorem: two games or procedures are coupled so that a precondition on their initial states implies the stated postcondition on final states and return values.

\noindent\textbf{Role in the proof.}
Use in this chapter. It contributes directly to an advantage bound, bad-event bound, or real-arithmetic composition step.

\noindent\textbf{Proof-script interpretation.}
Proof-script reading. The machine proof decomposes the theorem into rewrites, program-logic obligations, relational equivalences, and arithmetic side conditions.
\emph{Main tactics visible in the proof script:} \ec{proc}, \ec{wp}, \ec{call}, \ec{smt}.
\emph{Named identifiers syntactically cited in the proof block:}
\begin{lstlisting}[language=EasyCrypt,basicstyle=\ttfamily\scriptsize]
concrete_lazy_rom_get_preserves_sampler_rom_covered_arg
hash_qs
sampler_qs
\end{lstlisting}

\end{formaltheorem}

\begin{formaltheorem}[EasyCrypt line 2494]
\noindent\textbf{EasyCrypt name.}
\begin{lstlisting}[language=EasyCrypt,basicstyle=\ttfamily\scriptsize]
concrete_ro_exact_hyperball_sample_with_seed_preserves_sampler_rom_covered
\end{lstlisting}
\noindent\textbf{Checked EasyCrypt statement.}
\begin{lstlisting}[language=EasyCrypt,basicstyle=\ttfamily\scriptsize]
lemma concrete_ro_exact_hyperball_sample_with_seed_preserves_sampler_rom_covered
    hash_qs sampler_qs :
  hoare[ROExactHyperballPaperSimSampler(HAETAE_RO.FRO).sample_with_seed :
    sampler_rom_covered HAETAE_RO.FRO.m hash_qs sampler_qs /\
    sampler_expand_query arg.`1 \in sampler_qs ==>
    sampler_rom_covered HAETAE_RO.FRO.m hash_qs sampler_qs].
\end{lstlisting}
\noindent\textbf{Formal mathematical reading.}
Mathematical theorem. This is a relational program theorem: two games or procedures are coupled so that a precondition on their initial states implies the stated postcondition on final states and return values.

\noindent\textbf{Role in the proof.}
Use in this chapter. It contributes directly to an advantage bound, bad-event bound, or real-arithmetic composition step.

\noindent\textbf{Proof-script interpretation.}
Proof-script reading. The machine proof decomposes the theorem into rewrites, program-logic obligations, relational equivalences, and arithmetic side conditions.
\emph{Main tactics visible in the proof script:} \ec{proc}, \ec{wp}, \ec{rnd}, \ec{call}, \ec{smt}.
\emph{Named identifiers syntactically cited in the proof block:}
\begin{lstlisting}[language=EasyCrypt,basicstyle=\ttfamily\scriptsize]
concrete_lazy_rom_get_preserves_sampler_rom_covered_arg
hash_qs
sampler_qs
\end{lstlisting}

\end{formaltheorem}

\begin{formaltheorem}[EasyCrypt line 2510]
\noindent\textbf{EasyCrypt name.}
\begin{lstlisting}[language=EasyCrypt,basicstyle=\ttfamily\scriptsize]
concrete_ro_signing_attempt_sample_with_seed_preserves_sampler_rom_covered_budgeted_logs
\end{lstlisting}
\noindent\textbf{Checked EasyCrypt statement.}
\begin{lstlisting}[language=EasyCrypt,basicstyle=\ttfamily\scriptsize]
lemma concrete_ro_signing_attempt_sample_with_seed_preserves_sampler_rom_covered_budgeted_logs :
  hoare[ROSigningAttemptPaperSimSampler(HAETAE_RO.FRO).sample_with_seed :
    sampler_rom_covered
      HAETAE_RO.FRO.m
      ROMInternalTranscriptBudgetedPaperSimAsNMA.adversary_hash_queries
      ROMInternalTranscriptBudgetedPaperSimAsNMA.sampler_expand_queries /\
    sampler_expand_query arg.`1 \in
      ROMInternalTranscriptBudgetedPaperSimAsNMA.sampler_expand_queries ==>
    sampler_rom_covered
      HAETAE_RO.FRO.m
      ROMInternalTranscriptBudgetedPaperSimAsNMA.adversary_hash_queries
      ROMInternalTranscriptBudgetedPaperSimAsNMA.sampler_expand_queries].
\end{lstlisting}
\noindent\textbf{Formal mathematical reading.}
Mathematical theorem. This is a relational program theorem: two games or procedures are coupled so that a precondition on their initial states implies the stated postcondition on final states and return values.

\noindent\textbf{Role in the proof.}
Use in this chapter. It contributes directly to an advantage bound, bad-event bound, or real-arithmetic composition step.

\noindent\textbf{Proof-script interpretation.}
Proof-script reading. The machine proof decomposes the theorem into rewrites, program-logic obligations, relational equivalences, and arithmetic side conditions.
\emph{Main tactics visible in the proof script:} \ec{proc}, \ec{wp}, \ec{call}, \ec{smt}.
\emph{Named identifiers syntactically cited in the proof block:}
\begin{lstlisting}[language=EasyCrypt,basicstyle=\ttfamily\scriptsize]
concrete_lazy_rom_get_preserves_sampler_rom_covered_budgeted_logs
\end{lstlisting}

\end{formaltheorem}

\begin{formaltheorem}[EasyCrypt line 2529]
\noindent\textbf{EasyCrypt name.}
\begin{lstlisting}[language=EasyCrypt,basicstyle=\ttfamily\scriptsize]
concrete_ro_exact_hyperball_sample_with_seed_preserves_sampler_rom_covered_budgeted_logs
\end{lstlisting}
\noindent\textbf{Checked EasyCrypt statement.}
\begin{lstlisting}[language=EasyCrypt,basicstyle=\ttfamily\scriptsize]
lemma concrete_ro_exact_hyperball_sample_with_seed_preserves_sampler_rom_covered_budgeted_logs :
  hoare[ROExactHyperballPaperSimSampler(HAETAE_RO.FRO).sample_with_seed :
    sampler_rom_covered
      HAETAE_RO.FRO.m
      ROMInternalTranscriptBudgetedPaperSimAsNMA.adversary_hash_queries
      ROMInternalTranscriptBudgetedPaperSimAsNMA.sampler_expand_queries /\
    sampler_expand_query arg.`1 \in
      ROMInternalTranscriptBudgetedPaperSimAsNMA.sampler_expand_queries ==>
    sampler_rom_covered
      HAETAE_RO.FRO.m
      ROMInternalTranscriptBudgetedPaperSimAsNMA.adversary_hash_queries
      ROMInternalTranscriptBudgetedPaperSimAsNMA.sampler_expand_queries].
\end{lstlisting}
\noindent\textbf{Formal mathematical reading.}
Mathematical theorem. This is a relational program theorem: two games or procedures are coupled so that a precondition on their initial states implies the stated postcondition on final states and return values.

\noindent\textbf{Role in the proof.}
Use in this chapter. It contributes directly to an advantage bound, bad-event bound, or real-arithmetic composition step.

\noindent\textbf{Proof-script interpretation.}
Proof-script reading. The machine proof decomposes the theorem into rewrites, program-logic obligations, relational equivalences, and arithmetic side conditions.
\emph{Main tactics visible in the proof script:} \ec{proc}, \ec{wp}, \ec{rnd}, \ec{call}, \ec{smt}.
\emph{Named identifiers syntactically cited in the proof block:}
\begin{lstlisting}[language=EasyCrypt,basicstyle=\ttfamily\scriptsize]
concrete_lazy_rom_get_preserves_sampler_rom_covered_budgeted_logs
\end{lstlisting}

\end{formaltheorem}

\begin{formaltheorem}[EasyCrypt line 2550]
\noindent\textbf{EasyCrypt name.}
\begin{lstlisting}[language=EasyCrypt,basicstyle=\ttfamily\scriptsize]
concrete_budgeted_ro_signing_attempt_o_sign_preserves_sampler_rom_covered
\end{lstlisting}
\noindent\textbf{Checked EasyCrypt statement.}
\begin{lstlisting}[language=EasyCrypt,basicstyle=\ttfamily\scriptsize]
lemma concrete_budgeted_ro_signing_attempt_o_sign_preserves_sampler_rom_covered :
  hoare[ROMInternalTranscriptBudgetedPaperSimAsNMA
          (A, ROSigningAttemptPaperSimSampler(HAETAE_RO.FRO),
           HAETAE_RO.FRO).O.sign :
    sampler_rom_covered
      HAETAE_RO.FRO.m
      ROMInternalTranscriptBudgetedPaperSimAsNMA.adversary_hash_queries
      ROMInternalTranscriptBudgetedPaperSimAsNMA.sampler_expand_queries ==>
    sampler_rom_covered
      HAETAE_RO.FRO.m
      ROMInternalTranscriptBudgetedPaperSimAsNMA.adversary_hash_queries
      ROMInternalTranscriptBudgetedPaperSimAsNMA.sampler_expand_queries].
\end{lstlisting}
\noindent\textbf{Formal mathematical reading.}
Mathematical theorem. This is a relational program theorem: two games or procedures are coupled so that a precondition on their initial states implies the stated postcondition on final states and return values.

\noindent\textbf{Role in the proof.}
Use in this chapter. It contributes directly to an advantage bound, bad-event bound, or real-arithmetic composition step.

\noindent\textbf{Proof-script interpretation.}
Proof-script reading. The machine proof decomposes the theorem into rewrites, program-logic obligations, relational equivalences, and arithmetic side conditions.
\emph{Main tactics visible in the proof script:} \ec{proc}, \ec{wp}, \ec{call}, \ec{rnd}, \ec{smt}.
\emph{Named identifiers syntactically cited in the proof block:}
\begin{lstlisting}[language=EasyCrypt,basicstyle=\ttfamily\scriptsize]
concrete_lazy_rom_get_preserves_sampler_rom_covered_budgeted_logs
concrete_ro_signing_attempt_sample_with_seed_preserves_sampler_rom_covered_budgeted_logs
\end{lstlisting}

\end{formaltheorem}

\begin{formaltheorem}[EasyCrypt line 2577]
\noindent\textbf{EasyCrypt name.}
\begin{lstlisting}[language=EasyCrypt,basicstyle=\ttfamily\scriptsize]
concrete_budgeted_ro_exact_hyperball_o_sign_preserves_sampler_rom_covered
\end{lstlisting}
\noindent\textbf{Checked EasyCrypt statement.}
\begin{lstlisting}[language=EasyCrypt,basicstyle=\ttfamily\scriptsize]
lemma concrete_budgeted_ro_exact_hyperball_o_sign_preserves_sampler_rom_covered :
  hoare[ROMInternalTranscriptBudgetedPaperSimAsNMA
          (A, ROExactHyperballPaperSimSampler(HAETAE_RO.FRO),
           HAETAE_RO.FRO).O.sign :
    sampler_rom_covered
      HAETAE_RO.FRO.m
      ROMInternalTranscriptBudgetedPaperSimAsNMA.adversary_hash_queries
      ROMInternalTranscriptBudgetedPaperSimAsNMA.sampler_expand_queries ==>
    sampler_rom_covered
      HAETAE_RO.FRO.m
      ROMInternalTranscriptBudgetedPaperSimAsNMA.adversary_hash_queries
      ROMInternalTranscriptBudgetedPaperSimAsNMA.sampler_expand_queries].
\end{lstlisting}
\noindent\textbf{Formal mathematical reading.}
Mathematical theorem. This is a relational program theorem: two games or procedures are coupled so that a precondition on their initial states implies the stated postcondition on final states and return values.

\noindent\textbf{Role in the proof.}
Use in this chapter. It contributes directly to an advantage bound, bad-event bound, or real-arithmetic composition step.

\noindent\textbf{Proof-script interpretation.}
Proof-script reading. The machine proof decomposes the theorem into rewrites, program-logic obligations, relational equivalences, and arithmetic side conditions.
\emph{Main tactics visible in the proof script:} \ec{proc}, \ec{wp}, \ec{call}, \ec{rnd}, \ec{smt}.
\emph{Named identifiers syntactically cited in the proof block:}
\begin{lstlisting}[language=EasyCrypt,basicstyle=\ttfamily\scriptsize]
concrete_lazy_rom_get_preserves_sampler_rom_covered_budgeted_logs
concrete_ro_exact_hyperball_sample_with_seed_preserves_sampler_rom_covered_budgeted_logs
\end{lstlisting}

\end{formaltheorem}

\begin{formaltheorem}[EasyCrypt line 2604]
\noindent\textbf{EasyCrypt name.}
\begin{lstlisting}[language=EasyCrypt,basicstyle=\ttfamily\scriptsize]
concrete_budgeted_ro_signing_attempt_o_sign_clean_sampler_fresh
\end{lstlisting}
\noindent\textbf{Checked EasyCrypt statement.}
\begin{lstlisting}[language=EasyCrypt,basicstyle=\ttfamily\scriptsize]
lemma concrete_budgeted_ro_signing_attempt_o_sign_clean_sampler_fresh :
  hoare[ROMInternalTranscriptBudgetedPaperSimAsNMA
          (A, ROSigningAttemptPaperSimSampler(HAETAE_RO.FRO),
           HAETAE_RO.FRO).O.sign :
    sampler_rom_covered
      HAETAE_RO.FRO.m
      ROMInternalTranscriptBudgetedPaperSimAsNMA.adversary_hash_queries
      ROMInternalTranscriptBudgetedPaperSimAsNMA.sampler_expand_queries /\
    ! ROMInternalTranscriptBudgetedPaperSimAsNMA.sampler_bad_prequery ==>
    sampler_rom_covered
      HAETAE_RO.FRO.m
      ROMInternalTranscriptBudgetedPaperSimAsNMA.adversary_hash_queries
      ROMInternalTranscriptBudgetedPaperSimAsNMA.sampler_expand_queries].
\end{lstlisting}
\noindent\textbf{Formal mathematical reading.}
Mathematical theorem. This is a relational program theorem: two games or procedures are coupled so that a precondition on their initial states implies the stated postcondition on final states and return values.

\noindent\textbf{Role in the proof.}
Use in this chapter. It contributes directly to an advantage bound, bad-event bound, or real-arithmetic composition step.

\noindent\textbf{Proof-script interpretation.}
Proof-script reading. The machine proof decomposes the theorem into rewrites, program-logic obligations, relational equivalences, and arithmetic side conditions.
\emph{Main tactics visible in the proof script:} \ec{conseq}.
\emph{Named identifiers syntactically cited in the proof block:}
\begin{lstlisting}[language=EasyCrypt,basicstyle=\ttfamily\scriptsize]
concrete_budgeted_ro_signing_attempt_o_sign_preserves_sampler_rom_covered
\end{lstlisting}

\end{formaltheorem}

\begin{formaltheorem}[EasyCrypt line 2986]
\noindent\textbf{EasyCrypt name.}
\begin{lstlisting}[language=EasyCrypt,basicstyle=\ttfamily\scriptsize]
internal_transcript_sign_oracle_preserves_state
\end{lstlisting}
\noindent\textbf{Checked EasyCrypt statement.}
\begin{lstlisting}[language=EasyCrypt,basicstyle=\ttfamily\scriptsize]
lemma internal_transcript_sign_oracle_preserves_state :
  hoare[EUF_CMA_InternalTranscriptSign(H, A).O.sign :
    internal_transcript_state_sound haetae_mode
      EUF_CMA_InternalTranscriptSign.pk_current
      EUF_CMA_InternalTranscriptSign.sk_current
      EUF_CMA_InternalTranscriptSign.queries
      EUF_CMA_InternalTranscriptSign.transcripts
      EUF_CMA_InternalTranscriptSign.records ==>
    internal_transcript_state_sound haetae_mode
      EUF_CMA_InternalTranscriptSign.pk_current
      EUF_CMA_InternalTranscriptSign.sk_current
      EUF_CMA_InternalTranscriptSign.queries
      EUF_CMA_InternalTranscriptSign.transcripts
      EUF_CMA_InternalTranscriptSign.records].
\end{lstlisting}
\noindent\textbf{Formal mathematical reading.}
Mathematical theorem. This is a relational program theorem: two games or procedures are coupled so that a precondition on their initial states implies the stated postcondition on final states and return values.

\noindent\textbf{Role in the proof.}
Use in this chapter. It contributes directly to an advantage bound, bad-event bound, or real-arithmetic composition step.

\noindent\textbf{Proof-script interpretation.}
Proof-script reading. The machine proof decomposes the theorem into rewrites, program-logic obligations, relational equivalences, and arithmetic side conditions.
\emph{Main tactics visible in the proof script:} \ec{proc}, \ec{wp}, \ec{rnd}, \ec{apply}.
\emph{Named identifiers syntactically cited in the proof block:}
\begin{lstlisting}[language=EasyCrypt,basicstyle=\ttfamily\scriptsize]
EUF_CMA_InternalTranscriptSign
coins
ctx
haetae_mode
hr
internal_transcript_state_sound_sign_internal_cons
pk_current
queries
records
sk_current
state_sound
transcripts
\end{lstlisting}

\end{formaltheorem}

\begin{formaltheorem}[EasyCrypt line 3016]
\noindent\textbf{EasyCrypt name.}
\begin{lstlisting}[language=EasyCrypt,basicstyle=\ttfamily\scriptsize]
adversary_internal_transcript_state_preserved
\end{lstlisting}
\noindent\textbf{Checked EasyCrypt statement.}
\begin{lstlisting}[language=EasyCrypt,basicstyle=\ttfamily\scriptsize]
lemma adversary_internal_transcript_state_preserved :
  hoare[A(H, EUF_CMA_InternalTranscriptSign(H, A).O).forge :
    internal_transcript_state_sound haetae_mode
      EUF_CMA_InternalTranscriptSign.pk_current
      EUF_CMA_InternalTranscriptSign.sk_current
      EUF_CMA_InternalTranscriptSign.queries
      EUF_CMA_InternalTranscriptSign.transcripts
      EUF_CMA_InternalTranscriptSign.records ==>
    internal_transcript_state_sound haetae_mode
      EUF_CMA_InternalTranscriptSign.pk_current
      EUF_CMA_InternalTranscriptSign.sk_current
      EUF_CMA_InternalTranscriptSign.queries
      EUF_CMA_InternalTranscriptSign.transcripts
      EUF_CMA_InternalTranscriptSign.records].
\end{lstlisting}
\noindent\textbf{Formal mathematical reading.}
Mathematical theorem. This is a relational program theorem: two games or procedures are coupled so that a precondition on their initial states implies the stated postcondition on final states and return values.

\noindent\textbf{Role in the proof.}
Use in this chapter. It preserves an invariant across an oracle call, adversary call, or game procedure.

\noindent\textbf{Proof-script interpretation.}
Proof-script reading. The machine proof decomposes the theorem into rewrites, program-logic obligations, relational equivalences, and arithmetic side conditions.
\emph{Main tactics visible in the proof script:} \ec{proc}, \ec{call}, \ec{apply}.
\emph{Named identifiers syntactically cited in the proof block:}
\begin{lstlisting}[language=EasyCrypt,basicstyle=\ttfamily\scriptsize]
EUF_CMA_InternalTranscriptSign
haetae_mode
internal_transcript_sign_oracle_preserves_state
internal_transcript_state_sound
pk_current
queries
records
sk_current
transcripts
\end{lstlisting}

\end{formaltheorem}

\begin{formaltheorem}[EasyCrypt line 3044]
\noindent\textbf{EasyCrypt name.}
\begin{lstlisting}[language=EasyCrypt,basicstyle=\ttfamily\scriptsize]
internal_transcript_sign_main_state_sound
\end{lstlisting}
\noindent\textbf{Checked EasyCrypt statement.}
\begin{lstlisting}[language=EasyCrypt,basicstyle=\ttfamily\scriptsize]
lemma internal_transcript_sign_main_state_sound :
  hoare[EUF_CMA_InternalTranscriptSign(H, A).main :
    true ==>
    internal_transcript_state_sound haetae_mode
      EUF_CMA_InternalTranscriptSign.pk_current
      EUF_CMA_InternalTranscriptSign.sk_current
      EUF_CMA_InternalTranscriptSign.queries
      EUF_CMA_InternalTranscriptSign.transcripts
      EUF_CMA_InternalTranscriptSign.records].
\end{lstlisting}
\noindent\textbf{Formal mathematical reading.}
Mathematical theorem. This is a relational program theorem: two games or procedures are coupled so that a precondition on their initial states implies the stated postcondition on final states and return values.

\noindent\textbf{Role in the proof.}
Use in this chapter. It preserves an invariant across an oracle call, adversary call, or game procedure.

\noindent\textbf{Proof-script interpretation.}
Proof-script reading. The machine proof decomposes the theorem into rewrites, program-logic obligations, relational equivalences, and arithmetic side conditions.
\emph{Main tactics visible in the proof script:} \ec{proc}, \ec{wp}, \ec{call}, \ec{apply}, \ec{rnd}.
\emph{Named identifiers syntactically cited in the proof block:}
\begin{lstlisting}[language=EasyCrypt,basicstyle=\ttfamily\scriptsize]
EUF_CMA_InternalTranscriptSign
adversary_internal_transcript_state_preserved
haetae_mode
internal_transcript_state_sound
pk_current
queries
records
sk_current
transcripts
\end{lstlisting}

\end{formaltheorem}

\begin{formaltheorem}[EasyCrypt line 3076]
\noindent\textbf{EasyCrypt name.}
\begin{lstlisting}[language=EasyCrypt,basicstyle=\ttfamily\scriptsize]
internal_transcript_sign_main_transcripts_valid
\end{lstlisting}
\noindent\textbf{Checked EasyCrypt statement.}
\begin{lstlisting}[language=EasyCrypt,basicstyle=\ttfamily\scriptsize]
lemma internal_transcript_sign_main_transcripts_valid :
  hoare[EUF_CMA_InternalTranscriptSign(H, A).main :
    true ==>
    transcript_log_valid haetae_mode
      EUF_CMA_InternalTranscriptSign.transcripts].
\end{lstlisting}
\noindent\textbf{Formal mathematical reading.}
Mathematical theorem. This is a relational program theorem: two games or procedures are coupled so that a precondition on their initial states implies the stated postcondition on final states and return values.

\noindent\textbf{Role in the proof.}
Use in this chapter. It proves a structural correctness fact needed by later extraction or verification arguments.

\noindent\textbf{Proof-script interpretation.}
Proof-script reading. The machine proof decomposes the theorem into rewrites, program-logic obligations, relational equivalences, and arithmetic side conditions.
\emph{Main tactics visible in the proof script:} \ec{proc}, \ec{wp}, \ec{call}, \ec{conseq}, \ec{apply}, \ec{rnd}.
\emph{Named identifiers syntactically cited in the proof block:}
\begin{lstlisting}[language=EasyCrypt,basicstyle=\ttfamily\scriptsize]
EUF_CMA_InternalTranscriptSign
adversary_internal_transcript_state_preserved
haetae_mode
hr
internal_transcript_state_sound
internal_transcript_state_sound_valid
pk_current
qs0
queries
records
rs0
sk_current
state_sound
transcript_log_valid
transcripts
trs0
\end{lstlisting}

\end{formaltheorem}

\begin{formaltheorem}[EasyCrypt line 3118]
\noindent\textbf{EasyCrypt name.}
\begin{lstlisting}[language=EasyCrypt,basicstyle=\ttfamily\scriptsize]
internal_transcript_sign_main_log_ready
\end{lstlisting}
\noindent\textbf{Checked EasyCrypt statement.}
\begin{lstlisting}[language=EasyCrypt,basicstyle=\ttfamily\scriptsize]
lemma internal_transcript_sign_main_log_ready :
  hoare[EUF_CMA_InternalTranscriptSign(H, A).main :
    true ==>
    transcript_log_ready haetae_mode
      EUF_CMA_InternalTranscriptSign.transcripts].
\end{lstlisting}
\noindent\textbf{Formal mathematical reading.}
Mathematical theorem. This is a relational program theorem: two games or procedures are coupled so that a precondition on their initial states implies the stated postcondition on final states and return values.

\noindent\textbf{Role in the proof.}
Use in this chapter. It is a reusable checked theorem cited by later proof scripts.

\noindent\textbf{Proof-script interpretation.}
Proof-script reading. The machine proof decomposes the theorem into rewrites, program-logic obligations, relational equivalences, and arithmetic side conditions.
\emph{Main tactics visible in the proof script:} \ec{proc}, \ec{wp}, \ec{call}, \ec{conseq}, \ec{apply}, \ec{rnd}.
\emph{Named identifiers syntactically cited in the proof block:}
\begin{lstlisting}[language=EasyCrypt,basicstyle=\ttfamily\scriptsize]
EUF_CMA_InternalTranscriptSign
adversary_internal_transcript_state_preserved
haetae_mode
hr
internal_transcript_state_sound
internal_transcript_state_sound_log_ready
pk_current
qs0
queries
records
rs0
sk_current
state_sound
transcript_log_ready
transcripts
trs0
\end{lstlisting}

\end{formaltheorem}

\begin{formaltheorem}[EasyCrypt line 3160]
\noindent\textbf{EasyCrypt name.}
\begin{lstlisting}[language=EasyCrypt,basicstyle=\ttfamily\scriptsize]
internal_transcript_sign_main_no_min_entropy
\end{lstlisting}
\noindent\textbf{Checked EasyCrypt statement.}
\begin{lstlisting}[language=EasyCrypt,basicstyle=\ttfamily\scriptsize]
lemma internal_transcript_sign_main_no_min_entropy :
  hoare[EUF_CMA_InternalTranscriptSign(H, A).main :
    true ==>
    transcript_log_min_entropy_clear haetae_mode
      EUF_CMA_InternalTranscriptSign.transcripts].
\end{lstlisting}
\noindent\textbf{Formal mathematical reading.}
Mathematical theorem. This is a relational program theorem: two games or procedures are coupled so that a precondition on their initial states implies the stated postcondition on final states and return values.

\noindent\textbf{Role in the proof.}
Use in this chapter. It is a reusable checked theorem cited by later proof scripts.

\noindent\textbf{Proof-script interpretation.}
Proof-script reading. The machine proof decomposes the theorem into rewrites, program-logic obligations, relational equivalences, and arithmetic side conditions.
\emph{Main tactics visible in the proof script:} \ec{proc}, \ec{wp}, \ec{call}, \ec{conseq}, \ec{apply}, \ec{rnd}.
\emph{Named identifiers syntactically cited in the proof block:}
\begin{lstlisting}[language=EasyCrypt,basicstyle=\ttfamily\scriptsize]
EUF_CMA_InternalTranscriptSign
adversary_internal_transcript_state_preserved
haetae_mode
hr
internal_transcript_state_sound
internal_transcript_state_sound_no_min_entropy
pk_current
qs0
queries
records
rs0
sk_current
state_sound
transcript_log_min_entropy_clear
transcripts
trs0
\end{lstlisting}

\end{formaltheorem}

\begin{formaltheorem}[EasyCrypt line 3202]
\noindent\textbf{EasyCrypt name.}
\begin{lstlisting}[language=EasyCrypt,basicstyle=\ttfamily\scriptsize]
rom_internal_transcript_sign_oracle_preserves_state
\end{lstlisting}
\noindent\textbf{Checked EasyCrypt statement.}
\begin{lstlisting}[language=EasyCrypt,basicstyle=\ttfamily\scriptsize]
lemma rom_internal_transcript_sign_oracle_preserves_state :
  hoare[EUF_CMA_ROMInternalTranscriptSign(H, A).O.sign :
    internal_transcript_state_sound haetae_mode
      EUF_CMA_ROMInternalTranscriptSign.pk_current
      EUF_CMA_ROMInternalTranscriptSign.sk_current
      EUF_CMA_ROMInternalTranscriptSign.queries
      EUF_CMA_ROMInternalTranscriptSign.transcripts
      EUF_CMA_ROMInternalTranscriptSign.records ==>
    internal_transcript_state_sound haetae_mode
      EUF_CMA_ROMInternalTranscriptSign.pk_current
      EUF_CMA_ROMInternalTranscriptSign.sk_current
      EUF_CMA_ROMInternalTranscriptSign.queries
      EUF_CMA_ROMInternalTranscriptSign.transcripts
      EUF_CMA_ROMInternalTranscriptSign.records].
\end{lstlisting}
\noindent\textbf{Formal mathematical reading.}
Mathematical theorem. This is a relational program theorem: two games or procedures are coupled so that a precondition on their initial states implies the stated postcondition on final states and return values.

\noindent\textbf{Role in the proof.}
Use in this chapter. It contributes directly to an advantage bound, bad-event bound, or real-arithmetic composition step.

\noindent\textbf{Proof-script interpretation.}
Proof-script reading. The machine proof decomposes the theorem into rewrites, program-logic obligations, relational equivalences, and arithmetic side conditions.
\emph{Main tactics visible in the proof script:} \ec{proc}, \ec{wp}, \ec{call}, \ec{rnd}, \ec{apply}.
\emph{Named identifiers syntactically cited in the proof block:}
\begin{lstlisting}[language=EasyCrypt,basicstyle=\ttfamily\scriptsize]
EUF_CMA_ROMInternalTranscriptSign
ctx
haetae_mode
hr
internal_transcript_state_sound_sign_internal_cons
pk_current
queries
records
result
ro_signing_coins
seed_coins
sk_current
state_sound
transcripts
\end{lstlisting}

\end{formaltheorem}

\begin{formaltheorem}[EasyCrypt line 3238]
\noindent\textbf{EasyCrypt name.}
\begin{lstlisting}[language=EasyCrypt,basicstyle=\ttfamily\scriptsize]
adversary_rom_internal_transcript_state_preserved
\end{lstlisting}
\noindent\textbf{Checked EasyCrypt statement.}
\begin{lstlisting}[language=EasyCrypt,basicstyle=\ttfamily\scriptsize]
lemma adversary_rom_internal_transcript_state_preserved :
  hoare[A(H, EUF_CMA_ROMInternalTranscriptSign(H, A).O).forge :
    internal_transcript_state_sound haetae_mode
      EUF_CMA_ROMInternalTranscriptSign.pk_current
      EUF_CMA_ROMInternalTranscriptSign.sk_current
      EUF_CMA_ROMInternalTranscriptSign.queries
      EUF_CMA_ROMInternalTranscriptSign.transcripts
      EUF_CMA_ROMInternalTranscriptSign.records ==>
    internal_transcript_state_sound haetae_mode
      EUF_CMA_ROMInternalTranscriptSign.pk_current
      EUF_CMA_ROMInternalTranscriptSign.sk_current
      EUF_CMA_ROMInternalTranscriptSign.queries
      EUF_CMA_ROMInternalTranscriptSign.transcripts
      EUF_CMA_ROMInternalTranscriptSign.records].
\end{lstlisting}
\noindent\textbf{Formal mathematical reading.}
Mathematical theorem. This is a relational program theorem: two games or procedures are coupled so that a precondition on their initial states implies the stated postcondition on final states and return values.

\noindent\textbf{Role in the proof.}
Use in this chapter. It preserves an invariant across an oracle call, adversary call, or game procedure.

\noindent\textbf{Proof-script interpretation.}
Proof-script reading. The machine proof decomposes the theorem into rewrites, program-logic obligations, relational equivalences, and arithmetic side conditions.
\emph{Main tactics visible in the proof script:} \ec{proc}, \ec{call}, \ec{apply}.
\emph{Named identifiers syntactically cited in the proof block:}
\begin{lstlisting}[language=EasyCrypt,basicstyle=\ttfamily\scriptsize]
EUF_CMA_ROMInternalTranscriptSign
haetae_mode
internal_transcript_state_sound
pk_current
queries
records
rom_internal_transcript_sign_oracle_preserves_state
sk_current
transcripts
\end{lstlisting}

\end{formaltheorem}

\begin{formaltheorem}[EasyCrypt line 3266]
\noindent\textbf{EasyCrypt name.}
\begin{lstlisting}[language=EasyCrypt,basicstyle=\ttfamily\scriptsize]
rom_internal_transcript_sign_main_state_sound
\end{lstlisting}
\noindent\textbf{Checked EasyCrypt statement.}
\begin{lstlisting}[language=EasyCrypt,basicstyle=\ttfamily\scriptsize]
lemma rom_internal_transcript_sign_main_state_sound :
  hoare[EUF_CMA_ROMInternalTranscriptSign(H, A).main :
    true ==>
    internal_transcript_state_sound haetae_mode
      EUF_CMA_ROMInternalTranscriptSign.pk_current
      EUF_CMA_ROMInternalTranscriptSign.sk_current
      EUF_CMA_ROMInternalTranscriptSign.queries
      EUF_CMA_ROMInternalTranscriptSign.transcripts
      EUF_CMA_ROMInternalTranscriptSign.records].
\end{lstlisting}
\noindent\textbf{Formal mathematical reading.}
Mathematical theorem. This is a relational program theorem: two games or procedures are coupled so that a precondition on their initial states implies the stated postcondition on final states and return values.

\noindent\textbf{Role in the proof.}
Use in this chapter. It preserves an invariant across an oracle call, adversary call, or game procedure.

\noindent\textbf{Proof-script interpretation.}
Proof-script reading. The machine proof decomposes the theorem into rewrites, program-logic obligations, relational equivalences, and arithmetic side conditions.
\emph{Main tactics visible in the proof script:} \ec{proc}, \ec{wp}, \ec{call}, \ec{apply}, \ec{rnd}.
\emph{Named identifiers syntactically cited in the proof block:}
\begin{lstlisting}[language=EasyCrypt,basicstyle=\ttfamily\scriptsize]
EUF_CMA_ROMInternalTranscriptSign
adversary_rom_internal_transcript_state_preserved
haetae_mode
internal_transcript_state_sound
pk_current
queries
records
sk_current
transcripts
\end{lstlisting}

\end{formaltheorem}

\begin{formaltheorem}[EasyCrypt line 3300]
\noindent\textbf{EasyCrypt name.}
\begin{lstlisting}[language=EasyCrypt,basicstyle=\ttfamily\scriptsize]
rom_internal_transcript_sign_main_transcripts_valid
\end{lstlisting}
\noindent\textbf{Checked EasyCrypt statement.}
\begin{lstlisting}[language=EasyCrypt,basicstyle=\ttfamily\scriptsize]
lemma rom_internal_transcript_sign_main_transcripts_valid :
  hoare[EUF_CMA_ROMInternalTranscriptSign(H, A).main :
    true ==>
    transcript_log_valid haetae_mode
      EUF_CMA_ROMInternalTranscriptSign.transcripts].
\end{lstlisting}
\noindent\textbf{Formal mathematical reading.}
Mathematical theorem. This is a relational program theorem: two games or procedures are coupled so that a precondition on their initial states implies the stated postcondition on final states and return values.

\noindent\textbf{Role in the proof.}
Use in this chapter. It proves a structural correctness fact needed by later extraction or verification arguments.

\noindent\textbf{Proof-script interpretation.}
Proof-script reading. The machine proof decomposes the theorem into rewrites, program-logic obligations, relational equivalences, and arithmetic side conditions.
\emph{Main tactics visible in the proof script:} \ec{proc}, \ec{wp}, \ec{call}, \ec{conseq}, \ec{apply}, \ec{rnd}.
\emph{Named identifiers syntactically cited in the proof block:}
\begin{lstlisting}[language=EasyCrypt,basicstyle=\ttfamily\scriptsize]
EUF_CMA_ROMInternalTranscriptSign
adversary_rom_internal_transcript_state_preserved
haetae_mode
hr
internal_transcript_state_sound
internal_transcript_state_sound_valid
pk_current
qs0
queries
records
rs0
sk_current
state_sound
transcript_log_valid
transcripts
trs0
\end{lstlisting}

\end{formaltheorem}

\begin{formaltheorem}[EasyCrypt line 3344]
\noindent\textbf{EasyCrypt name.}
\begin{lstlisting}[language=EasyCrypt,basicstyle=\ttfamily\scriptsize]
rom_internal_transcript_sign_main_log_ready
\end{lstlisting}
\noindent\textbf{Checked EasyCrypt statement.}
\begin{lstlisting}[language=EasyCrypt,basicstyle=\ttfamily\scriptsize]
lemma rom_internal_transcript_sign_main_log_ready :
  hoare[EUF_CMA_ROMInternalTranscriptSign(H, A).main :
    true ==>
    transcript_log_ready haetae_mode
      EUF_CMA_ROMInternalTranscriptSign.transcripts].
\end{lstlisting}
\noindent\textbf{Formal mathematical reading.}
Mathematical theorem. This is a relational program theorem: two games or procedures are coupled so that a precondition on their initial states implies the stated postcondition on final states and return values.

\noindent\textbf{Role in the proof.}
Use in this chapter. It is a reusable checked theorem cited by later proof scripts.

\noindent\textbf{Proof-script interpretation.}
Proof-script reading. The machine proof decomposes the theorem into rewrites, program-logic obligations, relational equivalences, and arithmetic side conditions.
\emph{Main tactics visible in the proof script:} \ec{proc}, \ec{wp}, \ec{call}, \ec{conseq}, \ec{apply}, \ec{rnd}.
\emph{Named identifiers syntactically cited in the proof block:}
\begin{lstlisting}[language=EasyCrypt,basicstyle=\ttfamily\scriptsize]
EUF_CMA_ROMInternalTranscriptSign
adversary_rom_internal_transcript_state_preserved
haetae_mode
hr
internal_transcript_state_sound
internal_transcript_state_sound_log_ready
pk_current
qs0
queries
records
rs0
sk_current
state_sound
transcript_log_ready
transcripts
trs0
\end{lstlisting}

\end{formaltheorem}

\begin{formaltheorem}[EasyCrypt line 3388]
\noindent\textbf{EasyCrypt name.}
\begin{lstlisting}[language=EasyCrypt,basicstyle=\ttfamily\scriptsize]
rom_internal_transcript_sign_main_no_min_entropy
\end{lstlisting}
\noindent\textbf{Checked EasyCrypt statement.}
\begin{lstlisting}[language=EasyCrypt,basicstyle=\ttfamily\scriptsize]
lemma rom_internal_transcript_sign_main_no_min_entropy :
  hoare[EUF_CMA_ROMInternalTranscriptSign(H, A).main :
    true ==>
    transcript_log_min_entropy_clear haetae_mode
      EUF_CMA_ROMInternalTranscriptSign.transcripts].
\end{lstlisting}
\noindent\textbf{Formal mathematical reading.}
Mathematical theorem. This is a relational program theorem: two games or procedures are coupled so that a precondition on their initial states implies the stated postcondition on final states and return values.

\noindent\textbf{Role in the proof.}
Use in this chapter. It is a reusable checked theorem cited by later proof scripts.

\noindent\textbf{Proof-script interpretation.}
Proof-script reading. The machine proof decomposes the theorem into rewrites, program-logic obligations, relational equivalences, and arithmetic side conditions.
\emph{Main tactics visible in the proof script:} \ec{proc}, \ec{wp}, \ec{call}, \ec{conseq}, \ec{apply}, \ec{rnd}.
\emph{Named identifiers syntactically cited in the proof block:}
\begin{lstlisting}[language=EasyCrypt,basicstyle=\ttfamily\scriptsize]
EUF_CMA_ROMInternalTranscriptSign
adversary_rom_internal_transcript_state_preserved
haetae_mode
hr
internal_transcript_state_sound
internal_transcript_state_sound_no_min_entropy
pk_current
qs0
queries
records
rs0
sk_current
state_sound
transcript_log_min_entropy_clear
transcripts
trs0
\end{lstlisting}

\end{formaltheorem}

\begin{formaltheorem}[EasyCrypt line 3432]
\noindent\textbf{EasyCrypt name.}
\begin{lstlisting}[language=EasyCrypt,basicstyle=\ttfamily\scriptsize]
rom_internal_transcript_sign_main_no_failure_for_signature_site
\end{lstlisting}
\noindent\textbf{Checked EasyCrypt statement.}
\begin{lstlisting}[language=EasyCrypt,basicstyle=\ttfamily\scriptsize]
lemma rom_internal_transcript_sign_main_no_failure_for_signature_site :
  hoare[EUF_CMA_ROMInternalTranscriptSign(H, A).main :
    true ==>
    forall tr y,
      tr \in EUF_CMA_ROMInternalTranscriptSign.transcripts =>
      transcript_signature_challenge_matches tr y =>
      ! fs_with_aborts_failure haetae_mode tr
          (signature_programming_site_of_transcript haetae_mode tr y)].
\end{lstlisting}
\noindent\textbf{Formal mathematical reading.}
Mathematical theorem. This is a relational program theorem: two games or procedures are coupled so that a precondition on their initial states implies the stated postcondition on final states and return values.

\noindent\textbf{Role in the proof.}
Use in this chapter. It is a reusable checked theorem cited by later proof scripts.

\noindent\textbf{Proof-script interpretation.}
Proof-script reading. The machine proof decomposes the theorem into rewrites, program-logic obligations, relational equivalences, and arithmetic side conditions.
\emph{Main tactics visible in the proof script:} \ec{proc}, \ec{wp}, \ec{call}, \ec{conseq}, \ec{have}, \ec{apply}, \ec{rnd}.
\emph{Named identifiers syntactically cited in the proof block:}
\begin{lstlisting}[language=EasyCrypt,basicstyle=\ttfamily\scriptsize]
EUF_CMA_ROMInternalTranscriptSign
adversary_rom_internal_transcript_state_preserved
fs_with_aborts_failure
haetae_mode
hr
internal_transcript_state_sound
internal_transcript_state_sound_log_ready
match_y
mem_tr
pk_current
qs0
queries
ready
records
rs0
signature_programming_site_of_transcript
sk_current
state_sound
tr
transcript_log_ready_no_failure_for_signature_programming_site
transcript_signature_challenge_matches
transcripts
trs0
\end{lstlisting}

\end{formaltheorem}

\begin{formaltheorem}[EasyCrypt line 3485]
\noindent\textbf{EasyCrypt name.}
\begin{lstlisting}[language=EasyCrypt,basicstyle=\ttfamily\scriptsize]
rom_internal_transcript_sign_main_signature_sites_clear
\end{lstlisting}
\noindent\textbf{Checked EasyCrypt statement.}
\begin{lstlisting}[language=EasyCrypt,basicstyle=\ttfamily\scriptsize]
lemma rom_internal_transcript_sign_main_signature_sites_clear :
  hoare[EUF_CMA_ROMInternalTranscriptSign(H, A).main :
    true ==>
    transcript_log_signature_programming_sites_clear haetae_mode
      EUF_CMA_ROMInternalTranscriptSign.transcripts].
\end{lstlisting}
\noindent\textbf{Formal mathematical reading.}
Mathematical theorem. This is a relational program theorem: two games or procedures are coupled so that a precondition on their initial states implies the stated postcondition on final states and return values.

\noindent\textbf{Role in the proof.}
Use in this chapter. It contributes directly to an advantage bound, bad-event bound, or real-arithmetic composition step.

\noindent\textbf{Proof-script interpretation.}
Proof-script reading. The machine proof decomposes the theorem into rewrites, program-logic obligations, relational equivalences, and arithmetic side conditions.
\emph{Main tactics visible in the proof script:} \ec{proc}, \ec{wp}, \ec{call}, \ec{conseq}, \ec{have}, \ec{apply}, \ec{rnd}.
\emph{Named identifiers syntactically cited in the proof block:}
\begin{lstlisting}[language=EasyCrypt,basicstyle=\ttfamily\scriptsize]
EUF_CMA_ROMInternalTranscriptSign
adversary_rom_internal_transcript_state_preserved
haetae_mode
hr
internal_transcript_state_sound
internal_transcript_state_sound_log_ready
pk_current
qs0
queries
ready
records
rs0
sk_current
state_sound
transcript_log_ready_signature_sites_clear
transcript_log_signature_programming_sites_clear
transcripts
trs0
\end{lstlisting}

\end{formaltheorem}

\begin{formaltheorem}[EasyCrypt line 3531]
\noindent\textbf{EasyCrypt name.}
\begin{lstlisting}[language=EasyCrypt,basicstyle=\ttfamily\scriptsize]
rom_internal_transcript_bad_forgery_zero
\end{lstlisting}
\noindent\textbf{Checked EasyCrypt statement.}
\begin{lstlisting}[language=EasyCrypt,basicstyle=\ttfamily\scriptsize]
lemma rom_internal_transcript_bad_forgery_zero &m :
  Pr[EUF_CMA_ROMInternalTranscriptSign(H, A).main() @ &m :
    res /\ ! transcript_log_signature_programming_sites_clear haetae_mode
      EUF_CMA_ROMInternalTranscriptSign.transcripts] = 0%r.
\end{lstlisting}
\noindent\textbf{Formal mathematical reading.}
Mathematical theorem. This theorem has the form \(\Pr[G:E] = B\) is a game-probability statement.  It is one of the exact hops, bad-event bounds, or final advantage inequalities used in the reduction.

\noindent\textbf{Role in the proof.}
Use in this chapter. It is a reusable checked theorem cited by later proof scripts.

\noindent\textbf{Proof-script interpretation.}
Proof-script reading. The machine proof decomposes the theorem into rewrites, program-logic obligations, relational equivalences, and arithmetic side conditions.
\emph{Main tactics visible in the proof script:} \ec{byphoare}, \ec{hoare}, \ec{conseq}, \ec{rewrite}, \ec{apply}.
\emph{Named identifiers syntactically cited in the proof block:}
\begin{lstlisting}[language=EasyCrypt,basicstyle=\ttfamily\scriptsize]
EUF_CMA_ROMInternalTranscriptSign
clear_sites
forge
haetae_mode
hr
rom_internal_transcript_sign_main_signature_sites_clear
transcript_log_signature_programming_sites_clear
transcripts
trs0
\end{lstlisting}

\end{formaltheorem}

\begin{formaltheorem}[EasyCrypt line 3548]
\noindent\textbf{EasyCrypt name.}
\begin{lstlisting}[language=EasyCrypt,basicstyle=\ttfamily\scriptsize]
rom_internal_transcript_bad_forgery_bound
\end{lstlisting}
\noindent\textbf{Checked EasyCrypt statement.}
\begin{lstlisting}[language=EasyCrypt,basicstyle=\ttfamily\scriptsize]
lemma rom_internal_transcript_bad_forgery_bound &m :
  Pr[EUF_CMA_ROMInternalTranscriptSign(H, A).main() @ &m :
    res /\ ! transcript_log_signature_programming_sites_clear haetae_mode
      EUF_CMA_ROMInternalTranscriptSign.transcripts] <=
    fs_with_aborts_reprogramming_term + fs_with_aborts_min_entropy_term.
\end{lstlisting}
\noindent\textbf{Formal mathematical reading.}
Mathematical theorem. This theorem has the form \(\Pr[G:E] \le B\) is a game-probability statement.  It is one of the exact hops, bad-event bounds, or final advantage inequalities used in the reduction.

\noindent\textbf{Role in the proof.}
Use in this chapter. It contributes directly to an advantage bound, bad-event bound, or real-arithmetic composition step.

\noindent\textbf{Proof-script interpretation.}
Proof-script reading. The machine proof decomposes the theorem into rewrites, program-logic obligations, relational equivalences, and arithmetic side conditions.
\emph{Main tactics visible in the proof script:} \ec{rewrite}, \ec{have}, \ec{apply}.
\emph{Named identifiers syntactically cited in the proof block:}
\begin{lstlisting}[language=EasyCrypt,basicstyle=\ttfamily\scriptsize]
fs_with_aborts_min_entropy_term
fs_with_aborts_min_entropy_term_nonnegative
fs_with_aborts_reprogramming_term
fs_with_aborts_reprogramming_term_nonnegative
ge_entropy
ge_reprogram
ge_sum
ler_add
rom_internal_transcript_bad_forgery_zero
\end{lstlisting}

\end{formaltheorem}

\begin{formaltheorem}[EasyCrypt line 3564]
\noindent\textbf{EasyCrypt name.}
\begin{lstlisting}[language=EasyCrypt,basicstyle=\ttfamily\scriptsize]
rom_internal_transcript_bad_forgery_counted_bound
\end{lstlisting}
\noindent\textbf{Checked EasyCrypt statement.}
\begin{lstlisting}[language=EasyCrypt,basicstyle=\ttfamily\scriptsize]
lemma rom_internal_transcript_bad_forgery_counted_bound &m :
  Pr[EUF_CMA_ROMInternalTranscriptSign(H, A).main() @ &m :
    res /\ ! transcript_log_signature_programming_sites_clear haetae_mode
      EUF_CMA_ROMInternalTranscriptSign.transcripts] <=
    counted_rom_programming_loss_term.
\end{lstlisting}
\noindent\textbf{Formal mathematical reading.}
Mathematical theorem. This theorem has the form \(\Pr[G:E] \le B\) is a game-probability statement.  It is one of the exact hops, bad-event bounds, or final advantage inequalities used in the reduction.

\noindent\textbf{Role in the proof.}
Use in this chapter. It contributes directly to an advantage bound, bad-event bound, or real-arithmetic composition step.

\noindent\textbf{Proof-script interpretation.}
Proof-script reading. The machine proof decomposes the theorem into rewrites, program-logic obligations, relational equivalences, and arithmetic side conditions.
\emph{Main tactics visible in the proof script:} \ec{rewrite}, \ec{apply}.
\emph{Named identifiers syntactically cited in the proof block:}
\begin{lstlisting}[language=EasyCrypt,basicstyle=\ttfamily\scriptsize]
counted_rom_programming_loss_term_nonnegative
rom_internal_transcript_bad_forgery_zero
\end{lstlisting}

\end{formaltheorem}

\begin{formaltheorem}[EasyCrypt line 3583]
\noindent\textbf{EasyCrypt name.}
\begin{lstlisting}[language=EasyCrypt,basicstyle=\ttfamily\scriptsize]
counted_rom_internal_transcript_sign_oracle_preserves_budget_state
\end{lstlisting}
\noindent\textbf{Checked EasyCrypt statement.}
\begin{lstlisting}[language=EasyCrypt,basicstyle=\ttfamily\scriptsize]
lemma counted_rom_internal_transcript_sign_oracle_preserves_budget_state :
  hoare[EUF_CMA_ROMInternalTranscriptCountedSign(H, A).O.sign :
    internal_transcript_state_sound haetae_mode
      EUF_CMA_ROMInternalTranscriptCountedSign.pk_current
      EUF_CMA_ROMInternalTranscriptCountedSign.sk_current
      EUF_CMA_ROMInternalTranscriptCountedSign.queries
      EUF_CMA_ROMInternalTranscriptCountedSign.transcripts
      EUF_CMA_ROMInternalTranscriptCountedSign.records /\
    ! EUF_CMA_ROMInternalTranscriptCountedSign.C.bad_prequery /\
    ! EUF_CMA_ROMInternalTranscriptCountedSign.C.bad_reprogram /\
    ! EUF_CMA_ROMInternalTranscriptCountedSign.C.bad_min_entropy /\
    EUF_CMA_ROMInternalTranscriptCountedSign.C.programmed_sites = [] ==>
    internal_transcript_state_sound haetae_mode
      EUF_CMA_ROMInternalTranscriptCountedSign.pk_current
      EUF_CMA_ROMInternalTranscriptCountedSign.sk_current
      EUF_CMA_ROMInternalTranscriptCountedSign.queries
      EUF_CMA_ROMInternalTranscriptCountedSign.transcripts
      EUF_CMA_ROMInternalTranscriptCountedSign.records /\
    ! EUF_CMA_ROMInternalTranscriptCountedSign.C.bad_prequery /\
    ! EUF_CMA_ROMInternalTranscriptCountedSign.C.bad_reprogram /\
    ! EUF_CMA_ROMInternalTranscriptCountedSign.C.bad_min_entropy /\
    EUF_CMA_ROMInternalTranscriptCountedSign.C.programmed_sites = []].
\end{lstlisting}
\noindent\textbf{Formal mathematical reading.}
Mathematical theorem. This is a relational program theorem: two games or procedures are coupled so that a precondition on their initial states implies the stated postcondition on final states and return values.

\noindent\textbf{Role in the proof.}
Use in this chapter. It contributes directly to an advantage bound, bad-event bound, or real-arithmetic composition step.

\noindent\textbf{Proof-script interpretation.}
Proof-script reading. The machine proof decomposes the theorem into rewrites, program-logic obligations, relational equivalences, and arithmetic side conditions.
\emph{Main tactics visible in the proof script:} \ec{proc}, \ec{inline}, \ec{wp}, \ec{call}, \ec{rnd}, \ec{have}, \ec{case}, \ec{apply}, \ec{split}, \ec{rewrite}.
\emph{Named identifiers syntactically cited in the proof block:}
\begin{lstlisting}[language=EasyCrypt,basicstyle=\ttfamily\scriptsize]
EUF_CMA_ROMInternalTranscriptCountedSign
ctx
entropy_ok
get
haetae_mode
hr
internal_transcript_state_sound_sign_internal_cons
min_entropy_failure
no_min
observe_transcript
pkE
pk_current
preq_ok
programmed_ok
public_key_of_secret
queries
records
reprog_ok
result
ro_signing_coins
seed_coins
sign_internal
sign_internal_transcript_no_min_entropy
sk_current
state_sound
transcript_of_signature
transcripts
\end{lstlisting}

\end{formaltheorem}

\begin{formaltheorem}[EasyCrypt line 3653]
\noindent\textbf{EasyCrypt name.}
\begin{lstlisting}[language=EasyCrypt,basicstyle=\ttfamily\scriptsize]
adversary_counted_rom_internal_transcript_budget_state_preserved
\end{lstlisting}
\noindent\textbf{Checked EasyCrypt statement.}
\begin{lstlisting}[language=EasyCrypt,basicstyle=\ttfamily\scriptsize]
lemma adversary_counted_rom_internal_transcript_budget_state_preserved :
  hoare[A(EUF_CMA_ROMInternalTranscriptCountedSign(H, A).C,
          EUF_CMA_ROMInternalTranscriptCountedSign(H, A).O).forge :
    internal_transcript_state_sound haetae_mode
      EUF_CMA_ROMInternalTranscriptCountedSign.pk_current
      EUF_CMA_ROMInternalTranscriptCountedSign.sk_current
      EUF_CMA_ROMInternalTranscriptCountedSign.queries
      EUF_CMA_ROMInternalTranscriptCountedSign.transcripts
      EUF_CMA_ROMInternalTranscriptCountedSign.records /\
    ! EUF_CMA_ROMInternalTranscriptCountedSign.C.bad_prequery /\
    ! EUF_CMA_ROMInternalTranscriptCountedSign.C.bad_reprogram /\
    ! EUF_CMA_ROMInternalTranscriptCountedSign.C.bad_min_entropy /\
    EUF_CMA_ROMInternalTranscriptCountedSign.C.programmed_sites = [] ==>
    internal_transcript_state_sound haetae_mode
      EUF_CMA_ROMInternalTranscriptCountedSign.pk_current
      EUF_CMA_ROMInternalTranscriptCountedSign.sk_current
      EUF_CMA_ROMInternalTranscriptCountedSign.queries
      EUF_CMA_ROMInternalTranscriptCountedSign.transcripts
      EUF_CMA_ROMInternalTranscriptCountedSign.records /\
    ! EUF_CMA_ROMInternalTranscriptCountedSign.C.bad_prequery /\
    ! EUF_CMA_ROMInternalTranscriptCountedSign.C.bad_reprogram /\
    ! EUF_CMA_ROMInternalTranscriptCountedSign.C.bad_min_entropy /\
    EUF_CMA_ROMInternalTranscriptCountedSign.C.programmed_sites = []].
\end{lstlisting}
\noindent\textbf{Formal mathematical reading.}
Mathematical theorem. This is a relational program theorem: two games or procedures are coupled so that a precondition on their initial states implies the stated postcondition on final states and return values.

\noindent\textbf{Role in the proof.}
Use in this chapter. It preserves an invariant across an oracle call, adversary call, or game procedure.

\noindent\textbf{Proof-script interpretation.}
Proof-script reading. The machine proof decomposes the theorem into rewrites, program-logic obligations, relational equivalences, and arithmetic side conditions.
\emph{Main tactics visible in the proof script:} \ec{proc}, \ec{call}, \ec{apply}.
\emph{Named identifiers syntactically cited in the proof block:}
\begin{lstlisting}[language=EasyCrypt,basicstyle=\ttfamily\scriptsize]
EUF_CMA_ROMInternalTranscriptCountedSign
bad_min_entropy
bad_prequery
bad_reprogram
counted_lazy_rom_get_preserves_clear
counted_rom_internal_transcript_sign_oracle_preserves_budget_state
haetae_mode
internal_transcript_state_sound
pk_current
programmed_sites
queries
records
sk_current
transcripts
\end{lstlisting}

\end{formaltheorem}

\begin{formaltheorem}[EasyCrypt line 3694]
\noindent\textbf{EasyCrypt name.}
\begin{lstlisting}[language=EasyCrypt,basicstyle=\ttfamily\scriptsize]
counted_rom_internal_transcript_main_bad_clear
\end{lstlisting}
\noindent\textbf{Checked EasyCrypt statement.}
\begin{lstlisting}[language=EasyCrypt,basicstyle=\ttfamily\scriptsize]
lemma counted_rom_internal_transcript_main_bad_clear :
  hoare[EUF_CMA_ROMInternalTranscriptCountedSign(H, A).main :
    true ==> ! res].
\end{lstlisting}
\noindent\textbf{Formal mathematical reading.}
Mathematical theorem. This is a relational program theorem: two games or procedures are coupled so that a precondition on their initial states implies the stated postcondition on final states and return values.

\noindent\textbf{Role in the proof.}
Use in this chapter. It contributes directly to an advantage bound, bad-event bound, or real-arithmetic composition step.

\noindent\textbf{Proof-script interpretation.}
Proof-script reading. The machine proof decomposes the theorem into rewrites, program-logic obligations, relational equivalences, and arithmetic side conditions.
\emph{Main tactics visible in the proof script:} \ec{proc}, \ec{wp}, \ec{call}, \ec{apply}, \ec{rnd}.
\emph{Named identifiers syntactically cited in the proof block:}
\begin{lstlisting}[language=EasyCrypt,basicstyle=\ttfamily\scriptsize]
EUF_CMA_ROMInternalTranscriptCountedSign
adversary_counted_rom_internal_transcript_budget_state_preserved
bad_min_entropy
bad_prequery
bad_reprogram
counted_lazy_rom_bad_false_when_clear
counted_lazy_rom_get_preserves_clear
counted_lazy_rom_init_clears
haetae_mode
internal_transcript_state_sound
internal_transcript_state_sound_keygen
pk_current
programmed_sites
queries
records
sd
sk_current
transcripts
\end{lstlisting}

\end{formaltheorem}

\begin{formaltheorem}[EasyCrypt line 3733]
\noindent\textbf{EasyCrypt name.}
\begin{lstlisting}[language=EasyCrypt,basicstyle=\ttfamily\scriptsize]
counted_rom_internal_transcript_main_bad_false
\end{lstlisting}
\noindent\textbf{Checked EasyCrypt statement.}
\begin{lstlisting}[language=EasyCrypt,basicstyle=\ttfamily\scriptsize]
lemma counted_rom_internal_transcript_main_bad_false :
  hoare[EUF_CMA_ROMInternalTranscriptCountedSign(H, A).main :
    true ==> res = false].
\end{lstlisting}
\noindent\textbf{Formal mathematical reading.}
Mathematical theorem. This is a relational program theorem: two games or procedures are coupled so that a precondition on their initial states implies the stated postcondition on final states and return values.

\noindent\textbf{Role in the proof.}
Use in this chapter. It is a reusable checked theorem cited by later proof scripts.

\noindent\textbf{Proof-script interpretation.}
Proof-script reading. The machine proof decomposes the theorem into rewrites, program-logic obligations, relational equivalences, and arithmetic side conditions.
\emph{Main tactics visible in the proof script:} \ec{conseq}.
\emph{Named identifiers syntactically cited in the proof block:}
\begin{lstlisting}[language=EasyCrypt,basicstyle=\ttfamily\scriptsize]
counted_rom_internal_transcript_main_bad_clear
\end{lstlisting}

\end{formaltheorem}

\begin{formaltheorem}[EasyCrypt line 3740]
\noindent\textbf{EasyCrypt name.}
\begin{lstlisting}[language=EasyCrypt,basicstyle=\ttfamily\scriptsize]
counted_rom_internal_transcript_counted_bad_zero
\end{lstlisting}
\noindent\textbf{Checked EasyCrypt statement.}
\begin{lstlisting}[language=EasyCrypt,basicstyle=\ttfamily\scriptsize]
lemma counted_rom_internal_transcript_counted_bad_zero &m :
  Pr[EUF_CMA_ROMInternalTranscriptCountedSign(H, A).main() @ &m :
    res] = 0%r.
\end{lstlisting}
\noindent\textbf{Formal mathematical reading.}
Mathematical theorem. This theorem has the form \(\Pr[G:E] = B\) is a game-probability statement.  It is one of the exact hops, bad-event bounds, or final advantage inequalities used in the reduction.

\noindent\textbf{Role in the proof.}
Use in this chapter. It is a reusable checked theorem cited by later proof scripts.

\noindent\textbf{Proof-script interpretation.}
Proof-script reading. The machine proof decomposes the theorem into rewrites, program-logic obligations, relational equivalences, and arithmetic side conditions.
\emph{Main tactics visible in the proof script:} \ec{byphoare}, \ec{hoare}, \ec{conseq}, \ec{apply}.
\emph{Named identifiers syntactically cited in the proof block:}
\begin{lstlisting}[language=EasyCrypt,basicstyle=\ttfamily\scriptsize]
counted_rom_internal_transcript_main_bad_false
\end{lstlisting}

\end{formaltheorem}

\begin{formaltheorem}[EasyCrypt line 3751]
\noindent\textbf{EasyCrypt name.}
\begin{lstlisting}[language=EasyCrypt,basicstyle=\ttfamily\scriptsize]
counted_rom_internal_transcript_budget_sound
\end{lstlisting}
\noindent\textbf{Checked EasyCrypt statement.}
\begin{lstlisting}[language=EasyCrypt,basicstyle=\ttfamily\scriptsize]
lemma counted_rom_internal_transcript_budget_sound &m :
  counted_lazy_rom_bad_flags_budget_sound
    (Pr[EUF_CMA_ROMInternalTranscriptCountedSign(H, A).main() @ &m :
      res]).
\end{lstlisting}
\noindent\textbf{Formal mathematical reading.}
Mathematical theorem. This theorem has the form \(\Pr[G:E] \bowtie B\) is a game-probability statement.  It is one of the exact hops, bad-event bounds, or final advantage inequalities used in the reduction.

\noindent\textbf{Role in the proof.}
Use in this chapter. It proves a structural correctness fact needed by later extraction or verification arguments.

\noindent\textbf{Proof-script interpretation.}
Proof-script reading. The machine proof decomposes the theorem into rewrites, program-logic obligations, relational equivalences, and arithmetic side conditions.
\emph{Main tactics visible in the proof script:} \ec{rewrite}, \ec{apply}.
\emph{Named identifiers syntactically cited in the proof block:}
\begin{lstlisting}[language=EasyCrypt,basicstyle=\ttfamily\scriptsize]
counted_lazy_rom_bad_flags_budget_sound
counted_rom_internal_transcript_counted_bad_zero
counted_rom_programming_loss_term_nonnegative
\end{lstlisting}

\end{formaltheorem}

\begin{formaltheorem}[EasyCrypt line 3761]
\noindent\textbf{EasyCrypt name.}
\begin{lstlisting}[language=EasyCrypt,basicstyle=\ttfamily\scriptsize]
counted_rom_internal_transcript_counted_bad_subunit
\end{lstlisting}
\noindent\textbf{Checked EasyCrypt statement.}
\begin{lstlisting}[language=EasyCrypt,basicstyle=\ttfamily\scriptsize]
lemma counted_rom_internal_transcript_counted_bad_subunit &m :
  Pr[EUF_CMA_ROMInternalTranscriptCountedSign(H, A).main() @ &m :
    res] < 1%r.
\end{lstlisting}
\noindent\textbf{Formal mathematical reading.}
Mathematical theorem. This theorem has the form \(\Pr[G:E] \bowtie B\) is a game-probability statement.  It is one of the exact hops, bad-event bounds, or final advantage inequalities used in the reduction.

\noindent\textbf{Role in the proof.}
Use in this chapter. It contributes directly to an advantage bound, bad-event bound, or real-arithmetic composition step.

\noindent\textbf{Proof-script interpretation.}
Proof-script reading. The machine proof decomposes the theorem into rewrites, program-logic obligations, relational equivalences, and arithmetic side conditions.
\emph{Main tactics visible in the proof script:} \ec{rewrite}.
\emph{Named identifiers syntactically cited in the proof block:}
\begin{lstlisting}[language=EasyCrypt,basicstyle=\ttfamily\scriptsize]
counted_rom_internal_transcript_counted_bad_zero
trivial
\end{lstlisting}

\end{formaltheorem}

\begin{formaltheorem}[EasyCrypt line 3778]
\noindent\textbf{EasyCrypt name.}
\begin{lstlisting}[language=EasyCrypt,basicstyle=\ttfamily\scriptsize]
programmed_counted_sign_oracle_preserves_min_entropy_clear
\end{lstlisting}
\noindent\textbf{Checked EasyCrypt statement.}
\begin{lstlisting}[language=EasyCrypt,basicstyle=\ttfamily\scriptsize]
lemma programmed_counted_sign_oracle_preserves_min_entropy_clear :
  hoare[EUF_CMA_ROMProgrammedTranscriptCountedSign(H, A).O.sign :
    EUF_CMA_ROMProgrammedTranscriptCountedSign.pk_current =
      public_key_of_secret haetae_mode
        EUF_CMA_ROMProgrammedTranscriptCountedSign.sk_current /\
    ! EUF_CMA_ROMProgrammedTranscriptCountedSign.C.bad_min_entropy ==>
    EUF_CMA_ROMProgrammedTranscriptCountedSign.pk_current =
      public_key_of_secret haetae_mode
        EUF_CMA_ROMProgrammedTranscriptCountedSign.sk_current /\
    ! EUF_CMA_ROMProgrammedTranscriptCountedSign.C.bad_min_entropy].
\end{lstlisting}
\noindent\textbf{Formal mathematical reading.}
Mathematical theorem. This is a relational program theorem: two games or procedures are coupled so that a precondition on their initial states implies the stated postcondition on final states and return values.

\noindent\textbf{Role in the proof.}
Use in this chapter. It contributes directly to an advantage bound, bad-event bound, or real-arithmetic composition step.

\noindent\textbf{Proof-script interpretation.}
Proof-script reading. The machine proof decomposes the theorem into rewrites, program-logic obligations, relational equivalences, and arithmetic side conditions.
\emph{Main tactics visible in the proof script:} \ec{proc}, \ec{inline}, \ec{wp}, \ec{call}, \ec{rnd}, \ec{have}, \ec{apply}, \ec{rewrite}.
\emph{Named identifiers syntactically cited in the proof block:}
\begin{lstlisting}[language=EasyCrypt,basicstyle=\ttfamily\scriptsize]
EUF_CMA_ROMProgrammedTranscriptCountedSign
ctx
entropy_ok
get
haetae_mode
hr
min_entropy_failure
no_min
observe_transcript
pkE
pk_current
program
result
ro_signing_coins
seed_coins
sign_internal
sign_internal_transcript_no_min_entropy
sk_current
transcript_of_signature
\end{lstlisting}

\end{formaltheorem}

\begin{formaltheorem}[EasyCrypt line 3818]
\noindent\textbf{EasyCrypt name.}
\begin{lstlisting}[language=EasyCrypt,basicstyle=\ttfamily\scriptsize]
adversary_programmed_counted_min_entropy_preserved
\end{lstlisting}
\noindent\textbf{Checked EasyCrypt statement.}
\begin{lstlisting}[language=EasyCrypt,basicstyle=\ttfamily\scriptsize]
lemma adversary_programmed_counted_min_entropy_preserved :
  hoare[A(EUF_CMA_ROMProgrammedTranscriptCountedSign(H, A).C,
          EUF_CMA_ROMProgrammedTranscriptCountedSign(H, A).O).forge :
    EUF_CMA_ROMProgrammedTranscriptCountedSign.pk_current =
      public_key_of_secret haetae_mode
        EUF_CMA_ROMProgrammedTranscriptCountedSign.sk_current /\
    ! EUF_CMA_ROMProgrammedTranscriptCountedSign.C.bad_min_entropy ==>
    EUF_CMA_ROMProgrammedTranscriptCountedSign.pk_current =
      public_key_of_secret haetae_mode
        EUF_CMA_ROMProgrammedTranscriptCountedSign.sk_current /\
    ! EUF_CMA_ROMProgrammedTranscriptCountedSign.C.bad_min_entropy].
\end{lstlisting}
\noindent\textbf{Formal mathematical reading.}
Mathematical theorem. This is a relational program theorem: two games or procedures are coupled so that a precondition on their initial states implies the stated postcondition on final states and return values.

\noindent\textbf{Role in the proof.}
Use in this chapter. It preserves an invariant across an oracle call, adversary call, or game procedure.

\noindent\textbf{Proof-script interpretation.}
Proof-script reading. The machine proof decomposes the theorem into rewrites, program-logic obligations, relational equivalences, and arithmetic side conditions.
\emph{Main tactics visible in the proof script:} \ec{proc}, \ec{call}, \ec{apply}.
\emph{Named identifiers syntactically cited in the proof block:}
\begin{lstlisting}[language=EasyCrypt,basicstyle=\ttfamily\scriptsize]
EUF_CMA_ROMProgrammedTranscriptCountedSign
bad_min_entropy
counted_lazy_rom_get_preserves_min_entropy_clear
haetae_mode
pk_current
programmed_counted_sign_oracle_preserves_min_entropy_clear
public_key_of_secret
sk_current
\end{lstlisting}

\end{formaltheorem}

\begin{formaltheorem}[EasyCrypt line 3841]
\noindent\textbf{EasyCrypt name.}
\begin{lstlisting}[language=EasyCrypt,basicstyle=\ttfamily\scriptsize]
programmed_counted_main_min_entropy_clear
\end{lstlisting}
\noindent\textbf{Checked EasyCrypt statement.}
\begin{lstlisting}[language=EasyCrypt,basicstyle=\ttfamily\scriptsize]
lemma programmed_counted_main_min_entropy_clear :
  hoare[EUF_CMA_ROMProgrammedTranscriptCountedSign(H, A).main :
    true ==>
    ! EUF_CMA_ROMProgrammedTranscriptCountedSign.C.bad_min_entropy].
\end{lstlisting}
\noindent\textbf{Formal mathematical reading.}
Mathematical theorem. This is a relational program theorem: two games or procedures are coupled so that a precondition on their initial states implies the stated postcondition on final states and return values.

\noindent\textbf{Role in the proof.}
Use in this chapter. It contributes directly to an advantage bound, bad-event bound, or real-arithmetic composition step.

\noindent\textbf{Proof-script interpretation.}
Proof-script reading. The machine proof decomposes the theorem into rewrites, program-logic obligations, relational equivalences, and arithmetic side conditions.
\emph{Main tactics visible in the proof script:} \ec{proc}, \ec{wp}, \ec{call}, \ec{conseq}, \ec{rnd}, \ec{apply}.
\emph{Named identifiers syntactically cited in the proof block:}
\begin{lstlisting}[language=EasyCrypt,basicstyle=\ttfamily\scriptsize]
EUF_CMA_ROMProgrammedTranscriptCountedSign
adversary_programmed_counted_min_entropy_preserved
bad_min_entropy
counted_lazy_rom_bad_preserves_min_entropy_clear
counted_lazy_rom_get_preserves_min_entropy_clear
counted_lazy_rom_init_clears
haetae_mode
pk_current
public_key_of_keygen
public_key_of_secret
sd
sk_current
\end{lstlisting}

\end{formaltheorem}

\begin{formaltheorem}[EasyCrypt line 3869]
\noindent\textbf{EasyCrypt name.}
\begin{lstlisting}[language=EasyCrypt,basicstyle=\ttfamily\scriptsize]
programmed_counted_min_entropy_bad_zero
\end{lstlisting}
\noindent\textbf{Checked EasyCrypt statement.}
\begin{lstlisting}[language=EasyCrypt,basicstyle=\ttfamily\scriptsize]
lemma programmed_counted_min_entropy_bad_zero &m :
  Pr[EUF_CMA_ROMProgrammedTranscriptCountedSign(H, A).main() @ &m :
    EUF_CMA_ROMProgrammedTranscriptCountedSign.C.bad_min_entropy] = 0%r.
\end{lstlisting}
\noindent\textbf{Formal mathematical reading.}
Mathematical theorem. This theorem has the form \(\Pr[G:E] = B\) is a game-probability statement.  It is one of the exact hops, bad-event bounds, or final advantage inequalities used in the reduction.

\noindent\textbf{Role in the proof.}
Use in this chapter. It is a reusable checked theorem cited by later proof scripts.

\noindent\textbf{Proof-script interpretation.}
Proof-script reading. The machine proof decomposes the theorem into rewrites, program-logic obligations, relational equivalences, and arithmetic side conditions.
\emph{Main tactics visible in the proof script:} \ec{byphoare}, \ec{hoare}, \ec{conseq}.
\emph{Named identifiers syntactically cited in the proof block:}
\begin{lstlisting}[language=EasyCrypt,basicstyle=\ttfamily\scriptsize]
EUF_CMA_ROMProgrammedTranscriptCountedSign
bad_min_entropy
programmed_counted_main_min_entropy_clear
\end{lstlisting}

\end{formaltheorem}

\begin{formaltheorem}[EasyCrypt line 3879]
\noindent\textbf{EasyCrypt name.}
\begin{lstlisting}[language=EasyCrypt,basicstyle=\ttfamily\scriptsize]
programmed_counted_min_entropy_bad_bound
\end{lstlisting}
\noindent\textbf{Checked EasyCrypt statement.}
\begin{lstlisting}[language=EasyCrypt,basicstyle=\ttfamily\scriptsize]
lemma programmed_counted_min_entropy_bad_bound &m :
  Pr[EUF_CMA_ROMProgrammedTranscriptCountedSign(H, A).main() @ &m :
    EUF_CMA_ROMProgrammedTranscriptCountedSign.C.bad_min_entropy] <=
  counted_rom_min_entropy_term.
\end{lstlisting}
\noindent\textbf{Formal mathematical reading.}
Mathematical theorem. This theorem has the form \(\Pr[G:E] \le B\) is a game-probability statement.  It is one of the exact hops, bad-event bounds, or final advantage inequalities used in the reduction.

\noindent\textbf{Role in the proof.}
Use in this chapter. It contributes directly to an advantage bound, bad-event bound, or real-arithmetic composition step.

\noindent\textbf{Proof-script interpretation.}
Proof-script reading. The machine proof decomposes the theorem into rewrites, program-logic obligations, relational equivalences, and arithmetic side conditions.
\emph{Main tactics visible in the proof script:} \ec{rewrite}, \ec{apply}.
\emph{Named identifiers syntactically cited in the proof block:}
\begin{lstlisting}[language=EasyCrypt,basicstyle=\ttfamily\scriptsize]
counted_rom_min_entropy_term_nonnegative
programmed_counted_min_entropy_bad_zero
\end{lstlisting}

\end{formaltheorem}

\begin{formaltheorem}[EasyCrypt line 3888]
\noindent\textbf{EasyCrypt name.}
\begin{lstlisting}[language=EasyCrypt,basicstyle=\ttfamily\scriptsize]
programmed_counted_prequery_reprogramming_bad_bound_from_budget_expr
\end{lstlisting}
\noindent\textbf{Checked EasyCrypt statement.}
\begin{lstlisting}[language=EasyCrypt,basicstyle=\ttfamily\scriptsize]
lemma programmed_counted_prequery_reprogramming_bad_bound_from_budget_expr &m :
  Pr[EUF_CMA_ROMProgrammedTranscriptCountedSign(H, A).main() @ &m :
    EUF_CMA_ROMProgrammedTranscriptCountedSign.C.bad_prequery \/
    EUF_CMA_ROMProgrammedTranscriptCountedSign.C.bad_reprogram] <=
    ((rom_total_query_budget * rom_total_query_budget) +
     (rom_signature_query_budget * (rom_hash_query_budget + 1%r))) /
    challenge_support_cardinality_lower_bound =>
  Pr[EUF_CMA_ROMProgrammedTranscriptCountedSign(H, A).main() @ &m :
    EUF_CMA_ROMProgrammedTranscriptCountedSign.C.bad_prequery \/
    EUF_CMA_ROMProgrammedTranscriptCountedSign.C.bad_reprogram] <=
  counted_rom_prequery_reprogramming_term.
\end{lstlisting}
\noindent\textbf{Formal mathematical reading.}
Mathematical theorem. This theorem has the form \(\Pr[G:E] \le B\) is a game-probability statement.  It is one of the exact hops, bad-event bounds, or final advantage inequalities used in the reduction.

\noindent\textbf{Role in the proof.}
Use in this chapter. It contributes directly to an advantage bound, bad-event bound, or real-arithmetic composition step.

\noindent\textbf{Proof-script interpretation.}
Proof-script reading. The machine proof decomposes the theorem into rewrites, program-logic obligations, relational equivalences, and arithmetic side conditions.
\emph{Main tactics visible in the proof script:} \ec{rewrite}.
\emph{Named identifiers syntactically cited in the proof block:}
\begin{lstlisting}[language=EasyCrypt,basicstyle=\ttfamily\scriptsize]
counted_rom_prequery_reprogramming_term_cardinalityE
\end{lstlisting}

\end{formaltheorem}

\begin{formaltheorem}[EasyCrypt line 3915]
\noindent\textbf{EasyCrypt name.}
\begin{lstlisting}[language=EasyCrypt,basicstyle=\ttfamily\scriptsize]
budgeted_counted_lazy_rom_get_true
\end{lstlisting}
\noindent\textbf{Checked EasyCrypt statement.}
\begin{lstlisting}[language=EasyCrypt,basicstyle=\ttfamily\scriptsize]
lemma budgeted_counted_lazy_rom_get_true :
  hoare[EUF_CMA_ROMProgrammedTranscriptBudgetedCountedSign(H, A).C.get :
    true ==> true].
\end{lstlisting}
\noindent\textbf{Formal mathematical reading.}
Mathematical theorem. This is a relational program theorem: two games or procedures are coupled so that a precondition on their initial states implies the stated postcondition on final states and return values.

\noindent\textbf{Role in the proof.}
Use in this chapter. It is a reusable checked theorem cited by later proof scripts.

\noindent\textbf{Proof-script interpretation.}
Proof-script reading. The machine proof decomposes the theorem into rewrites, program-logic obligations, relational equivalences, and arithmetic side conditions.
\emph{Main tactics visible in the proof script:} \ec{proc}, \ec{wp}, \ec{call}, \ec{rewrite}, \ec{smt}.
\emph{Named identifiers syntactically cited in the proof block:}
\begin{lstlisting}[language=EasyCrypt,basicstyle=\ttfamily\scriptsize]
budgeted_programmed_challenge_query_discipline
\end{lstlisting}

\end{formaltheorem}

\begin{formaltheorem}[EasyCrypt line 3925]
\noindent\textbf{EasyCrypt name.}
\begin{lstlisting}[language=EasyCrypt,basicstyle=\ttfamily\scriptsize]
budgeted_counted_lazy_rom_init_true
\end{lstlisting}
\noindent\textbf{Checked EasyCrypt statement.}
\begin{lstlisting}[language=EasyCrypt,basicstyle=\ttfamily\scriptsize]
lemma budgeted_counted_lazy_rom_init_true :
  hoare[EUF_CMA_ROMProgrammedTranscriptBudgetedCountedSign(H, A).C.init :
    true ==> true].
\end{lstlisting}
\noindent\textbf{Formal mathematical reading.}
Mathematical theorem. This is a relational program theorem: two games or procedures are coupled so that a precondition on their initial states implies the stated postcondition on final states and return values.

\noindent\textbf{Role in the proof.}
Use in this chapter. It is a reusable checked theorem cited by later proof scripts.

\noindent\textbf{Proof-script interpretation.}
Proof-script reading. The machine proof decomposes the theorem into rewrites, program-logic obligations, relational equivalences, and arithmetic side conditions.
\emph{Main tactics visible in the proof script:} \ec{proc}, \ec{wp}, \ec{call}.

\end{formaltheorem}

\begin{formaltheorem}[EasyCrypt line 3935]
\noindent\textbf{EasyCrypt name.}
\begin{lstlisting}[language=EasyCrypt,basicstyle=\ttfamily\scriptsize]
budgeted_counted_lazy_rom_bad_true
\end{lstlisting}
\noindent\textbf{Checked EasyCrypt statement.}
\begin{lstlisting}[language=EasyCrypt,basicstyle=\ttfamily\scriptsize]
lemma budgeted_counted_lazy_rom_bad_true :
  hoare[EUF_CMA_ROMProgrammedTranscriptBudgetedCountedSign(H, A).C.bad :
    true ==> true].
\end{lstlisting}
\noindent\textbf{Formal mathematical reading.}
Mathematical theorem. This is a relational program theorem: two games or procedures are coupled so that a precondition on their initial states implies the stated postcondition on final states and return values.

\noindent\textbf{Role in the proof.}
Use in this chapter. It is a reusable checked theorem cited by later proof scripts.

\noindent\textbf{Proof-script interpretation.}
Proof-script reading. The machine proof decomposes the theorem into rewrites, program-logic obligations, relational equivalences, and arithmetic side conditions.
\emph{Main tactics visible in the proof script:} \ec{proc}.

\end{formaltheorem}

\begin{formaltheorem}[EasyCrypt line 3942]
\noindent\textbf{EasyCrypt name.}
\begin{lstlisting}[language=EasyCrypt,basicstyle=\ttfamily\scriptsize]
budgeted_programmed_hash_get_preserves_challenge_query_discipline
\end{lstlisting}
\noindent\textbf{Checked EasyCrypt statement.}
\begin{lstlisting}[language=EasyCrypt,basicstyle=\ttfamily\scriptsize]
lemma budgeted_programmed_hash_get_preserves_challenge_query_discipline :
  hoare[EUF_CMA_ROMProgrammedTranscriptBudgetedCountedSign(H, A).AH.get :
    budgeted_programmed_challenge_query_discipline
      EUF_CMA_ROMProgrammedTranscriptBudgetedCountedSign.adversary_hash_count
      EUF_CMA_ROMProgrammedTranscriptBudgetedCountedSign.C.hash_queries ==>
    budgeted_programmed_challenge_query_discipline
      EUF_CMA_ROMProgrammedTranscriptBudgetedCountedSign.adversary_hash_count
      EUF_CMA_ROMProgrammedTranscriptBudgetedCountedSign.C.hash_queries].
\end{lstlisting}
\noindent\textbf{Formal mathematical reading.}
Mathematical theorem. This is a relational program theorem: two games or procedures are coupled so that a precondition on their initial states implies the stated postcondition on final states and return values.

\noindent\textbf{Role in the proof.}
Use in this chapter. It contributes directly to an advantage bound, bad-event bound, or real-arithmetic composition step.

\noindent\textbf{Proof-script interpretation.}
Proof-script reading. The machine proof decomposes the theorem into rewrites, program-logic obligations, relational equivalences, and arithmetic side conditions.
\emph{Main tactics visible in the proof script:} \ec{proc}, \ec{inline}, \ec{wp}, \ec{call}, \ec{rewrite}, \ec{smt}.
\emph{Named identifiers syntactically cited in the proof block:}
\begin{lstlisting}[language=EasyCrypt,basicstyle=\ttfamily\scriptsize]
EUF_CMA_ROMProgrammedTranscriptBudgetedCountedSign
budgeted_programmed_challenge_query_discipline
challenge_query_log_count_cons_le
get
hash_query_budget_count
\end{lstlisting}

\end{formaltheorem}

\begin{formaltheorem}[EasyCrypt line 3963]
\noindent\textbf{EasyCrypt name.}
\begin{lstlisting}[language=EasyCrypt,basicstyle=\ttfamily\scriptsize]
budgeted_programmed_sign_oracle_preserves_challenge_query_discipline
\end{lstlisting}
\noindent\textbf{Checked EasyCrypt statement.}
\begin{lstlisting}[language=EasyCrypt,basicstyle=\ttfamily\scriptsize]
lemma budgeted_programmed_sign_oracle_preserves_challenge_query_discipline :
  hoare[EUF_CMA_ROMProgrammedTranscriptBudgetedCountedSign(H, A).O.sign :
    budgeted_programmed_challenge_query_discipline
      EUF_CMA_ROMProgrammedTranscriptBudgetedCountedSign.adversary_hash_count
      EUF_CMA_ROMProgrammedTranscriptBudgetedCountedSign.C.hash_queries ==>
    budgeted_programmed_challenge_query_discipline
      EUF_CMA_ROMProgrammedTranscriptBudgetedCountedSign.adversary_hash_count
      EUF_CMA_ROMProgrammedTranscriptBudgetedCountedSign.C.hash_queries].
\end{lstlisting}
\noindent\textbf{Formal mathematical reading.}
Mathematical theorem. This is a relational program theorem: two games or procedures are coupled so that a precondition on their initial states implies the stated postcondition on final states and return values.

\noindent\textbf{Role in the proof.}
Use in this chapter. It contributes directly to an advantage bound, bad-event bound, or real-arithmetic composition step.

\noindent\textbf{Proof-script interpretation.}
Proof-script reading. The machine proof decomposes the theorem into rewrites, program-logic obligations, relational equivalences, and arithmetic side conditions.
\emph{Main tactics visible in the proof script:} \ec{proc}, \ec{inline}, \ec{wp}, \ec{call}, \ec{rnd}, \ec{rewrite}, \ec{smt}.
\emph{Named identifiers syntactically cited in the proof block:}
\begin{lstlisting}[language=EasyCrypt,basicstyle=\ttfamily\scriptsize]
DirectSigningCoinSampler
EUF_CMA_ROMProgrammedTranscriptBudgetedCountedSign
budgeted_programmed_challenge_query_discipline
challenge_query_log_count_message_cons
get
observe_transcript
program
sample
\end{lstlisting}

\end{formaltheorem}

\begin{formaltheorem}[EasyCrypt line 3989]
\noindent\textbf{EasyCrypt name.}
\begin{lstlisting}[language=EasyCrypt,basicstyle=\ttfamily\scriptsize]
adversary_budgeted_programmed_challenge_query_discipline_preserved
\end{lstlisting}
\noindent\textbf{Checked EasyCrypt statement.}
\begin{lstlisting}[language=EasyCrypt,basicstyle=\ttfamily\scriptsize]
lemma adversary_budgeted_programmed_challenge_query_discipline_preserved :
  hoare[
    A(EUF_CMA_ROMProgrammedTranscriptBudgetedCountedSign(H, A).AH,
      EUF_CMA_ROMProgrammedTranscriptBudgetedCountedSign(H, A).O).forge :
    budgeted_programmed_challenge_query_discipline
      EUF_CMA_ROMProgrammedTranscriptBudgetedCountedSign.adversary_hash_count
      EUF_CMA_ROMProgrammedTranscriptBudgetedCountedSign.C.hash_queries ==>
    budgeted_programmed_challenge_query_discipline
      EUF_CMA_ROMProgrammedTranscriptBudgetedCountedSign.adversary_hash_count
      EUF_CMA_ROMProgrammedTranscriptBudgetedCountedSign.C.hash_queries].
\end{lstlisting}
\noindent\textbf{Formal mathematical reading.}
Mathematical theorem. This is a relational program theorem: two games or procedures are coupled so that a precondition on their initial states implies the stated postcondition on final states and return values.

\noindent\textbf{Role in the proof.}
Use in this chapter. It contributes directly to an advantage bound, bad-event bound, or real-arithmetic composition step.

\noindent\textbf{Proof-script interpretation.}
Proof-script reading. The machine proof decomposes the theorem into rewrites, program-logic obligations, relational equivalences, and arithmetic side conditions.
\emph{Main tactics visible in the proof script:} \ec{proc}, \ec{call}, \ec{apply}.
\emph{Named identifiers syntactically cited in the proof block:}
\begin{lstlisting}[language=EasyCrypt,basicstyle=\ttfamily\scriptsize]
EUF_CMA_ROMProgrammedTranscriptBudgetedCountedSign
adversary_hash_count
budgeted_programmed_challenge_query_discipline
budgeted_programmed_hash_get_preserves_challenge_query_discipline
budgeted_programmed_sign_oracle_preserves_challenge_query_discipline
hash_queries
\end{lstlisting}

\end{formaltheorem}

\begin{formaltheorem}[EasyCrypt line 4011]
\noindent\textbf{EasyCrypt name.}
\begin{lstlisting}[language=EasyCrypt,basicstyle=\ttfamily\scriptsize]
budgeted_programmed_main_challenge_query_discipline
\end{lstlisting}
\noindent\textbf{Checked EasyCrypt statement.}
\begin{lstlisting}[language=EasyCrypt,basicstyle=\ttfamily\scriptsize]
lemma budgeted_programmed_main_challenge_query_discipline :
  hoare[EUF_CMA_ROMProgrammedTranscriptBudgetedCountedSign(H, A).main :
    true ==>
    budgeted_programmed_challenge_query_discipline
      EUF_CMA_ROMProgrammedTranscriptBudgetedCountedSign.adversary_hash_count
      EUF_CMA_ROMProgrammedTranscriptBudgetedCountedSign.C.hash_queries].
\end{lstlisting}
\noindent\textbf{Formal mathematical reading.}
Mathematical theorem. This is a relational program theorem: two games or procedures are coupled so that a precondition on their initial states implies the stated postcondition on final states and return values.

\noindent\textbf{Role in the proof.}
Use in this chapter. It contributes directly to an advantage bound, bad-event bound, or real-arithmetic composition step.

\noindent\textbf{Proof-script interpretation.}
Proof-script reading. The machine proof decomposes the theorem into rewrites, program-logic obligations, relational equivalences, and arithmetic side conditions.
\emph{Main tactics visible in the proof script:} \ec{proc}, \ec{wp}, \ec{call}, \ec{conseq}, \ec{rnd}, \ec{rewrite}, \ec{smt}.
\emph{Named identifiers syntactically cited in the proof block:}
\begin{lstlisting}[language=EasyCrypt,basicstyle=\ttfamily\scriptsize]
EUF_CMA_ROMProgrammedTranscriptBudgetedCountedSign
adversary_budgeted_programmed_challenge_query_discipline_preserved
adversary_hash_count
budgeted_counted_lazy_rom_bad_true
budgeted_programmed_challenge_query_discipline
challenge_query_log_count
counted_lazy_rom_get_nonchallenge_preserves_challenge_query_count
counted_lazy_rom_init_clears_hash_queries
hash_queries
hash_query_budget_count
matrix_expand_query
ro_query_is_challenge
sd
\end{lstlisting}

\end{formaltheorem}

\begin{formaltheorem}[EasyCrypt line 4042]
\noindent\textbf{EasyCrypt name.}
\begin{lstlisting}[language=EasyCrypt,basicstyle=\ttfamily\scriptsize]
budgeted_programmed_main_challenge_query_count_bound
\end{lstlisting}
\noindent\textbf{Checked EasyCrypt statement.}
\begin{lstlisting}[language=EasyCrypt,basicstyle=\ttfamily\scriptsize]
lemma budgeted_programmed_main_challenge_query_count_bound :
  hoare[EUF_CMA_ROMProgrammedTranscriptBudgetedCountedSign(H, A).main :
    true ==>
    challenge_query_log_count
      EUF_CMA_ROMProgrammedTranscriptBudgetedCountedSign.C.hash_queries <=
      hash_query_budget_count + 1].
\end{lstlisting}
\noindent\textbf{Formal mathematical reading.}
Mathematical theorem. This is a relational program theorem: two games or procedures are coupled so that a precondition on their initial states implies the stated postcondition on final states and return values.

\noindent\textbf{Role in the proof.}
Use in this chapter. It contributes directly to an advantage bound, bad-event bound, or real-arithmetic composition step.

\noindent\textbf{Proof-script interpretation.}
Proof-script reading. The machine proof decomposes the theorem into rewrites, program-logic obligations, relational equivalences, and arithmetic side conditions.
\emph{Main tactics visible in the proof script:} \ec{conseq}, \ec{rewrite}, \ec{smt}.
\emph{Named identifiers syntactically cited in the proof block:}
\begin{lstlisting}[language=EasyCrypt,basicstyle=\ttfamily\scriptsize]
budgeted_programmed_challenge_query_discipline
budgeted_programmed_main_challenge_query_discipline
\end{lstlisting}

\end{formaltheorem}

\begin{formaltheorem}[EasyCrypt line 4054]
\noindent\textbf{EasyCrypt name.}
\begin{lstlisting}[language=EasyCrypt,basicstyle=\ttfamily\scriptsize]
budgeted_programmed_adaptive_prequery_lift
\end{lstlisting}
\noindent\textbf{Checked EasyCrypt statement.}
\begin{lstlisting}[language=EasyCrypt,basicstyle=\ttfamily\scriptsize]
lemma budgeted_programmed_adaptive_prequery_lift
  md sk m ctx qs hc :
  budgeted_programmed_challenge_query_discipline hc qs =>
  mu drandom_coins
    (fun coins =>
      programming_site_prequeried qs
        (honest_signing_programming_site md sk m ctx coins)) <=
  (rom_hash_query_budget + 1%r) /
    challenge_support_cardinality_lower_bound.
\end{lstlisting}
\noindent\textbf{Formal mathematical reading.}
Mathematical theorem. This is an inequality, usually an arithmetic loss bound, event inclusion bound, or monotonicity fact used to compose probabilities.

\noindent\textbf{Role in the proof.}
Use in this chapter. It is a reusable checked theorem cited by later proof scripts.

\noindent\textbf{Proof-script interpretation.}
Proof-script reading. The machine proof decomposes the theorem into rewrites, program-logic obligations, relational equivalences, and arithmetic side conditions.
\emph{Main tactics visible in the proof script:} \ec{rewrite}, \ec{apply}.
\emph{Named identifiers syntactically cited in the proof block:}
\begin{lstlisting}[language=EasyCrypt,basicstyle=\ttfamily\scriptsize]
budgeted_programmed_challenge_query_discipline
ctx
hc
hc_budget
hc_ge0
md
programmed_site_prequery_one_step_challenge_counter_bound
qs
query_count
sk
\end{lstlisting}

\end{formaltheorem}

\begin{formaltheorem}[EasyCrypt line 4070]
\noindent\textbf{EasyCrypt name.}
\begin{lstlisting}[language=EasyCrypt,basicstyle=\ttfamily\scriptsize]
budgeted_programmed_adaptive_prequery_lift_checked_31
\end{lstlisting}
\noindent\textbf{Checked EasyCrypt statement.}
\begin{lstlisting}[language=EasyCrypt,basicstyle=\ttfamily\scriptsize]
lemma budgeted_programmed_adaptive_prequery_lift_checked_31
  md sk m ctx qs hc :
  budgeted_programmed_challenge_query_discipline hc qs =>
  mu drandom_coins
    (fun coins =>
      programming_site_prequeried qs
        (honest_signing_programming_site md sk m ctx coins)) <=
  31%r / 288230376151711744%r.
\end{lstlisting}
\noindent\textbf{Formal mathematical reading.}
Mathematical theorem. This is an inequality, usually an arithmetic loss bound, event inclusion bound, or monotonicity fact used to compose probabilities.

\noindent\textbf{Role in the proof.}
Use in this chapter. It is a reusable checked theorem cited by later proof scripts.

\noindent\textbf{Proof-script interpretation.}
Proof-script reading. The machine proof decomposes the theorem into rewrites, program-logic obligations, relational equivalences, and arithmetic side conditions.
\emph{Main tactics visible in the proof script:} \ec{apply}, \ec{rewrite}, \ec{smt}.
\emph{Named identifiers syntactically cited in the proof block:}
\begin{lstlisting}[language=EasyCrypt,basicstyle=\ttfamily\scriptsize]
budgeted_programmed_adaptive_prequery_lift
challenge_support_cardinality_lower_bound
ctx
discipline
hash_query_count
hc
ler_trans
md
qs
rom_hash_query_budget
sk
\end{lstlisting}

\end{formaltheorem}

\begin{formaltheorem}[EasyCrypt line 4089]
\noindent\textbf{EasyCrypt name.}
\begin{lstlisting}[language=EasyCrypt,basicstyle=\ttfamily\scriptsize]
budgeted_programmed_hash_get_preserves_discipline
\end{lstlisting}
\noindent\textbf{Checked EasyCrypt statement.}
\begin{lstlisting}[language=EasyCrypt,basicstyle=\ttfamily\scriptsize]
lemma budgeted_programmed_hash_get_preserves_discipline :
  hoare[EUF_CMA_ROMProgrammedTranscriptBudgetedCountedSign(H, A).AH.get :
    budgeted_programmed_query_count_discipline
      EUF_CMA_ROMProgrammedTranscriptBudgetedCountedSign.adversary_hash_count
      EUF_CMA_ROMProgrammedTranscriptBudgetedCountedSign.signing_count /\
    budgeted_programmed_reprogram_discipline
      EUF_CMA_ROMProgrammedTranscriptBudgetedCountedSign.pk_current
      EUF_CMA_ROMProgrammedTranscriptBudgetedCountedSign.sk_current
      EUF_CMA_ROMProgrammedTranscriptBudgetedCountedSign.C.programmed_sites
      EUF_CMA_ROMProgrammedTranscriptBudgetedCountedSign.C.bad_reprogram ==>
    budgeted_programmed_query_count_discipline
      EUF_CMA_ROMProgrammedTranscriptBudgetedCountedSign.adversary_hash_count
      EUF_CMA_ROMProgrammedTranscriptBudgetedCountedSign.signing_count /\
    budgeted_programmed_reprogram_discipline
      EUF_CMA_ROMProgrammedTranscriptBudgetedCountedSign.pk_current
      EUF_CMA_ROMProgrammedTranscriptBudgetedCountedSign.sk_current
      EUF_CMA_ROMProgrammedTranscriptBudgetedCountedSign.C.programmed_sites
      EUF_CMA_ROMProgrammedTranscriptBudgetedCountedSign.C.bad_reprogram].
\end{lstlisting}
\noindent\textbf{Formal mathematical reading.}
Mathematical theorem. This is a relational program theorem: two games or procedures are coupled so that a precondition on their initial states implies the stated postcondition on final states and return values.

\noindent\textbf{Role in the proof.}
Use in this chapter. It preserves an invariant across an oracle call, adversary call, or game procedure.

\noindent\textbf{Proof-script interpretation.}
Proof-script reading. The machine proof decomposes the theorem into rewrites, program-logic obligations, relational equivalences, and arithmetic side conditions.
\emph{Main tactics visible in the proof script:} \ec{proc}, \ec{call}, \ec{wp}, \ec{rewrite}, \ec{smt}.
\emph{Named identifiers syntactically cited in the proof block:}
\begin{lstlisting}[language=EasyCrypt,basicstyle=\ttfamily\scriptsize]
budgeted_counted_lazy_rom_get_true
budgeted_programmed_query_count_discipline
hash_query_budget_count
signature_query_budget_count
\end{lstlisting}

\end{formaltheorem}

\begin{formaltheorem}[EasyCrypt line 4119]
\noindent\textbf{EasyCrypt name.}
\begin{lstlisting}[language=EasyCrypt,basicstyle=\ttfamily\scriptsize]
budgeted_programmed_reprogram_discipline_after_actual_signature
\end{lstlisting}
\noindent\textbf{Checked EasyCrypt statement.}
\begin{lstlisting}[language=EasyCrypt,basicstyle=\ttfamily\scriptsize]
lemma budgeted_programmed_reprogram_discipline_after_actual_signature
  pk sk m ctx coins sites bad_reprogram :
  budgeted_programmed_reprogram_discipline pk sk sites bad_reprogram =>
  pk = public_key_of_secret haetae_mode sk =>
  budgeted_programmed_reprogram_discipline
    pk sk
    (actual_signature_programming_site haetae_mode
      (transcript_of_signature haetae_mode pk m ctx
        (sign_internal haetae_mode sk m ctx coins)) :: sites)
    (bad_reprogram \/
      programming_site_reprograms sites
        (actual_signature_programming_site haetae_mode
          (transcript_of_signature haetae_mode pk m ctx
            (sign_internal haetae_mode sk m ctx coins)))).
\end{lstlisting}
\noindent\textbf{Formal mathematical reading.}
Mathematical theorem. This is an implication, normally turning one invariant, validity predicate, or proof obligation into another.

\noindent\textbf{Role in the proof.}
Use in this chapter. It preserves an invariant across an oracle call, adversary call, or game procedure.

\noindent\textbf{Proof-script interpretation.}
Proof-script reading. The machine proof decomposes the theorem into rewrites, program-logic obligations, relational equivalences, and arithmetic side conditions.
\emph{Main tactics visible in the proof script:} \ec{rewrite}, \ec{have}, \ec{apply}, \ec{split}.
\emph{Named identifiers syntactically cited in the proof block:}
\begin{lstlisting}[language=EasyCrypt,basicstyle=\ttfamily\scriptsize]
actual_signature_programming_site
actual_signature_programming_site_sign_internal_conflict_step
budgeted_programmed_reprogram_discipline
coins
ctx
exact
haetae_mode
hclear
hconf
hconf_next
hkey
hno_reprogram
hsite
pk
programming_site_log_conflict_free_for_honest
programming_site_reprograms
sign_internal
sites
sk
transcript_of_signature
\end{lstlisting}

\end{formaltheorem}

\begin{formaltheorem}[EasyCrypt line 4157]
\noindent\textbf{EasyCrypt name.}
\begin{lstlisting}[language=EasyCrypt,basicstyle=\ttfamily\scriptsize]
budgeted_programmed_sign_oracle_preserves_discipline
\end{lstlisting}
\noindent\textbf{Checked EasyCrypt statement.}
\begin{lstlisting}[language=EasyCrypt,basicstyle=\ttfamily\scriptsize]
lemma budgeted_programmed_sign_oracle_preserves_discipline :
  hoare[EUF_CMA_ROMProgrammedTranscriptBudgetedCountedSign(H, A).O.sign :
    budgeted_programmed_query_count_discipline
      EUF_CMA_ROMProgrammedTranscriptBudgetedCountedSign.adversary_hash_count
      EUF_CMA_ROMProgrammedTranscriptBudgetedCountedSign.signing_count /\
    budgeted_programmed_reprogram_discipline
      EUF_CMA_ROMProgrammedTranscriptBudgetedCountedSign.pk_current
      EUF_CMA_ROMProgrammedTranscriptBudgetedCountedSign.sk_current
      EUF_CMA_ROMProgrammedTranscriptBudgetedCountedSign.C.programmed_sites
      EUF_CMA_ROMProgrammedTranscriptBudgetedCountedSign.C.bad_reprogram ==>
    budgeted_programmed_query_count_discipline
      EUF_CMA_ROMProgrammedTranscriptBudgetedCountedSign.adversary_hash_count
      EUF_CMA_ROMProgrammedTranscriptBudgetedCountedSign.signing_count /\
    budgeted_programmed_reprogram_discipline
      EUF_CMA_ROMProgrammedTranscriptBudgetedCountedSign.pk_current
      EUF_CMA_ROMProgrammedTranscriptBudgetedCountedSign.sk_current
      EUF_CMA_ROMProgrammedTranscriptBudgetedCountedSign.C.programmed_sites
      EUF_CMA_ROMProgrammedTranscriptBudgetedCountedSign.C.bad_reprogram].
\end{lstlisting}
\noindent\textbf{Formal mathematical reading.}
Mathematical theorem. This is a relational program theorem: two games or procedures are coupled so that a precondition on their initial states implies the stated postcondition on final states and return values.

\noindent\textbf{Role in the proof.}
Use in this chapter. It contributes directly to an advantage bound, bad-event bound, or real-arithmetic composition step.

\noindent\textbf{Proof-script interpretation.}
Proof-script reading. The machine proof decomposes the theorem into rewrites, program-logic obligations, relational equivalences, and arithmetic side conditions.
\emph{Main tactics visible in the proof script:} \ec{proc}, \ec{inline}, \ec{wp}, \ec{call}, \ec{rnd}, \ec{case}, \ec{split}, \ec{rewrite}, \ec{smt}, \ec{apply}.
\emph{Named identifiers syntactically cited in the proof block:}
\begin{lstlisting}[language=EasyCrypt,basicstyle=\ttfamily\scriptsize]
CountedLazyROM
DirectSigningCoinSampler
EUF_CMA_ROMProgrammedTranscriptBudgetedCountedSign
bad_reprogram
budgeted_programmed_query_count_discipline
budgeted_programmed_reprogram_discipline
budgeted_programmed_reprogram_discipline_after_actual_signature
ctx
exact
get
hash_query_budget_count
hcount
hdis
hr
hsign_case
hsign_lt
observe_transcript
pk_current
program
programmed_sites
sample
sampled_coins
signature_query_budget_count
signing_count
sk_current
\end{lstlisting}

\end{formaltheorem}

\begin{formaltheorem}[EasyCrypt line 4223]
\noindent\textbf{EasyCrypt name.}
\begin{lstlisting}[language=EasyCrypt,basicstyle=\ttfamily\scriptsize]
adversary_budgeted_programmed_discipline_preserved
\end{lstlisting}
\noindent\textbf{Checked EasyCrypt statement.}
\begin{lstlisting}[language=EasyCrypt,basicstyle=\ttfamily\scriptsize]
lemma adversary_budgeted_programmed_discipline_preserved :
  hoare[
    A(EUF_CMA_ROMProgrammedTranscriptBudgetedCountedSign(H, A).AH,
      EUF_CMA_ROMProgrammedTranscriptBudgetedCountedSign(H, A).O).forge :
    budgeted_programmed_query_count_discipline
      EUF_CMA_ROMProgrammedTranscriptBudgetedCountedSign.adversary_hash_count
      EUF_CMA_ROMProgrammedTranscriptBudgetedCountedSign.signing_count /\
    budgeted_programmed_reprogram_discipline
      EUF_CMA_ROMProgrammedTranscriptBudgetedCountedSign.pk_current
      EUF_CMA_ROMProgrammedTranscriptBudgetedCountedSign.sk_current
      EUF_CMA_ROMProgrammedTranscriptBudgetedCountedSign.C.programmed_sites
      EUF_CMA_ROMProgrammedTranscriptBudgetedCountedSign.C.bad_reprogram ==>
    budgeted_programmed_query_count_discipline
      EUF_CMA_ROMProgrammedTranscriptBudgetedCountedSign.adversary_hash_count
      EUF_CMA_ROMProgrammedTranscriptBudgetedCountedSign.signing_count /\
    budgeted_programmed_reprogram_discipline
      EUF_CMA_ROMProgrammedTranscriptBudgetedCountedSign.pk_current
      EUF_CMA_ROMProgrammedTranscriptBudgetedCountedSign.sk_current
      EUF_CMA_ROMProgrammedTranscriptBudgetedCountedSign.C.programmed_sites
      EUF_CMA_ROMProgrammedTranscriptBudgetedCountedSign.C.bad_reprogram].
\end{lstlisting}
\noindent\textbf{Formal mathematical reading.}
Mathematical theorem. This is a relational program theorem: two games or procedures are coupled so that a precondition on their initial states implies the stated postcondition on final states and return values.

\noindent\textbf{Role in the proof.}
Use in this chapter. It preserves an invariant across an oracle call, adversary call, or game procedure.

\noindent\textbf{Proof-script interpretation.}
Proof-script reading. The machine proof decomposes the theorem into rewrites, program-logic obligations, relational equivalences, and arithmetic side conditions.
\emph{Main tactics visible in the proof script:} \ec{proc}, \ec{call}, \ec{apply}.
\emph{Named identifiers syntactically cited in the proof block:}
\begin{lstlisting}[language=EasyCrypt,basicstyle=\ttfamily\scriptsize]
EUF_CMA_ROMProgrammedTranscriptBudgetedCountedSign
adversary_hash_count
bad_reprogram
budgeted_programmed_hash_get_preserves_discipline
budgeted_programmed_query_count_discipline
budgeted_programmed_reprogram_discipline
budgeted_programmed_sign_oracle_preserves_discipline
pk_current
programmed_sites
signing_count
sk_current
\end{lstlisting}

\end{formaltheorem}

\begin{formaltheorem}[EasyCrypt line 4260]
\noindent\textbf{EasyCrypt name.}
\begin{lstlisting}[language=EasyCrypt,basicstyle=\ttfamily\scriptsize]
budgeted_programmed_main_reprogram_clear
\end{lstlisting}
\noindent\textbf{Checked EasyCrypt statement.}
\begin{lstlisting}[language=EasyCrypt,basicstyle=\ttfamily\scriptsize]
lemma budgeted_programmed_main_reprogram_clear :
  hoare[EUF_CMA_ROMProgrammedTranscriptBudgetedCountedSign(H, A).main :
    true ==>
    ! EUF_CMA_ROMProgrammedTranscriptBudgetedCountedSign.C.bad_reprogram].
\end{lstlisting}
\noindent\textbf{Formal mathematical reading.}
Mathematical theorem. This is a relational program theorem: two games or procedures are coupled so that a precondition on their initial states implies the stated postcondition on final states and return values.

\noindent\textbf{Role in the proof.}
Use in this chapter. It contributes directly to an advantage bound, bad-event bound, or real-arithmetic composition step.

\noindent\textbf{Proof-script interpretation.}
Proof-script reading. The machine proof decomposes the theorem into rewrites, program-logic obligations, relational equivalences, and arithmetic side conditions.
\emph{Main tactics visible in the proof script:} \ec{proc}, \ec{wp}, \ec{call}, \ec{conseq}, \ec{rnd}, \ec{rewrite}, \ec{smt}.
\emph{Named identifiers syntactically cited in the proof block:}
\begin{lstlisting}[language=EasyCrypt,basicstyle=\ttfamily\scriptsize]
EUF_CMA_ROMProgrammedTranscriptBudgetedCountedSign
adversary_budgeted_programmed_discipline_preserved
adversary_hash_count
bad_reprogram
budgeted_counted_lazy_rom_get_true
budgeted_programmed_query_count_discipline
budgeted_programmed_reprogram_discipline
counted_lazy_rom_init_clears
hash_query_budget_count
pk_current
programmed_sites
sd
signature_query_budget_count
signing_count
sk_current
\end{lstlisting}

\end{formaltheorem}

\begin{formaltheorem}[EasyCrypt line 4310]
\noindent\textbf{EasyCrypt name.}
\begin{lstlisting}[language=EasyCrypt,basicstyle=\ttfamily\scriptsize]
budgeted_programmed_bad_reprogram_zero
\end{lstlisting}
\noindent\textbf{Checked EasyCrypt statement.}
\begin{lstlisting}[language=EasyCrypt,basicstyle=\ttfamily\scriptsize]
lemma budgeted_programmed_bad_reprogram_zero &m :
  Pr[EUF_CMA_ROMProgrammedTranscriptBudgetedCountedSign(H, A).main() @ &m :
    EUF_CMA_ROMProgrammedTranscriptBudgetedCountedSign.C.bad_reprogram] =
  0%r.
\end{lstlisting}
\noindent\textbf{Formal mathematical reading.}
Mathematical theorem. This theorem has the form \(\Pr[G:E] = B\) is a game-probability statement.  It is one of the exact hops, bad-event bounds, or final advantage inequalities used in the reduction.

\noindent\textbf{Role in the proof.}
Use in this chapter. It is a reusable checked theorem cited by later proof scripts.

\noindent\textbf{Proof-script interpretation.}
Proof-script reading. The machine proof decomposes the theorem into rewrites, program-logic obligations, relational equivalences, and arithmetic side conditions.
\emph{Main tactics visible in the proof script:} \ec{byphoare}, \ec{hoare}, \ec{conseq}.
\emph{Named identifiers syntactically cited in the proof block:}
\begin{lstlisting}[language=EasyCrypt,basicstyle=\ttfamily\scriptsize]
EUF_CMA_ROMProgrammedTranscriptBudgetedCountedSign
bad_reprogram
budgeted_programmed_main_reprogram_clear
\end{lstlisting}

\end{formaltheorem}

\begin{formaltheorem}[EasyCrypt line 4321]
\noindent\textbf{EasyCrypt name.}
\begin{lstlisting}[language=EasyCrypt,basicstyle=\ttfamily\scriptsize]
budgeted_programmed_prequery_reprogram_bad_bound_from_prequery
\end{lstlisting}
\noindent\textbf{Checked EasyCrypt statement.}
\begin{lstlisting}[language=EasyCrypt,basicstyle=\ttfamily\scriptsize]
lemma budgeted_programmed_prequery_reprogram_bad_bound_from_prequery &m :
  Pr[EUF_CMA_ROMProgrammedTranscriptBudgetedCountedSign(H, A).main() @ &m :
    EUF_CMA_ROMProgrammedTranscriptBudgetedCountedSign.C.bad_prequery] <=
    ((rom_total_query_budget * rom_total_query_budget) +
     (rom_signature_query_budget * (rom_hash_query_budget + 1%r))) /
    challenge_support_cardinality_lower_bound =>
  Pr[EUF_CMA_ROMProgrammedTranscriptBudgetedCountedSign(H, A).main() @ &m :
    EUF_CMA_ROMProgrammedTranscriptBudgetedCountedSign.C.bad_prequery \/
    EUF_CMA_ROMProgrammedTranscriptBudgetedCountedSign.C.bad_reprogram] <=
    ((rom_total_query_budget * rom_total_query_budget) +
     (rom_signature_query_budget * (rom_hash_query_budget + 1%r))) /
    challenge_support_cardinality_lower_bound.
\end{lstlisting}
\noindent\textbf{Formal mathematical reading.}
Mathematical theorem. This theorem has the form \(\Pr[G:E] \le B\) is a game-probability statement.  It is one of the exact hops, bad-event bounds, or final advantage inequalities used in the reduction.

\noindent\textbf{Role in the proof.}
Use in this chapter. It contributes directly to an advantage bound, bad-event bound, or real-arithmetic composition step.

\noindent\textbf{Proof-script interpretation.}
Proof-script reading. The machine proof decomposes the theorem into rewrites, program-logic obligations, relational equivalences, and arithmetic side conditions.
\emph{Main tactics visible in the proof script:} \ec{apply}, \ec{rewrite}, \ec{smt}.
\emph{Named identifiers syntactically cited in the proof block:}
\begin{lstlisting}[language=EasyCrypt,basicstyle=\ttfamily\scriptsize]
EUF_CMA_ROMProgrammedTranscriptBudgetedCountedSign
Pr
bad_prequery
bad_reprogram
budgeted_programmed_bad_reprogram_zero
ler_trans
main
mu_or
prequery_bound
\end{lstlisting}

\end{formaltheorem}

\begin{formaltheorem}[EasyCrypt line 4345]
\noindent\textbf{EasyCrypt name.}
\begin{lstlisting}[language=EasyCrypt,basicstyle=\ttfamily\scriptsize]
budgeted_programmed_hash_get_preserves_query_budget
\end{lstlisting}
\noindent\textbf{Checked EasyCrypt statement.}
\begin{lstlisting}[language=EasyCrypt,basicstyle=\ttfamily\scriptsize]
lemma budgeted_programmed_hash_get_preserves_query_budget :
  hoare[EUF_CMA_ROMProgrammedTranscriptBudgetedCountedSign(H, A).AH.get :
    0 <= EUF_CMA_ROMProgrammedTranscriptBudgetedCountedSign.adversary_hash_count /\
    EUF_CMA_ROMProgrammedTranscriptBudgetedCountedSign.adversary_hash_count <=
      hash_query_budget_count /\
    0 <= EUF_CMA_ROMProgrammedTranscriptBudgetedCountedSign.signing_count /\
    EUF_CMA_ROMProgrammedTranscriptBudgetedCountedSign.signing_count <=
      signature_query_budget_count ==>
    0 <= EUF_CMA_ROMProgrammedTranscriptBudgetedCountedSign.adversary_hash_count /\
    EUF_CMA_ROMProgrammedTranscriptBudgetedCountedSign.adversary_hash_count <=
      hash_query_budget_count /\
    0 <= EUF_CMA_ROMProgrammedTranscriptBudgetedCountedSign.signing_count /\
    EUF_CMA_ROMProgrammedTranscriptBudgetedCountedSign.signing_count <=
      signature_query_budget_count].
\end{lstlisting}
\noindent\textbf{Formal mathematical reading.}
Mathematical theorem. This is a relational program theorem: two games or procedures are coupled so that a precondition on their initial states implies the stated postcondition on final states and return values.

\noindent\textbf{Role in the proof.}
Use in this chapter. It preserves an invariant across an oracle call, adversary call, or game procedure.

\noindent\textbf{Proof-script interpretation.}
Proof-script reading. The machine proof decomposes the theorem into rewrites, program-logic obligations, relational equivalences, and arithmetic side conditions.
\emph{Main tactics visible in the proof script:} \ec{proc}, \ec{call}, \ec{wp}, \ec{rewrite}, \ec{smt}.
\emph{Named identifiers syntactically cited in the proof block:}
\begin{lstlisting}[language=EasyCrypt,basicstyle=\ttfamily\scriptsize]
budgeted_counted_lazy_rom_get_true
hash_query_budget_count
signature_query_budget_count
\end{lstlisting}

\end{formaltheorem}

\begin{formaltheorem}[EasyCrypt line 4368]
\noindent\textbf{EasyCrypt name.}
\begin{lstlisting}[language=EasyCrypt,basicstyle=\ttfamily\scriptsize]
budgeted_programmed_sign_oracle_preserves_query_budget
\end{lstlisting}
\noindent\textbf{Checked EasyCrypt statement.}
\begin{lstlisting}[language=EasyCrypt,basicstyle=\ttfamily\scriptsize]
lemma budgeted_programmed_sign_oracle_preserves_query_budget :
  hoare[EUF_CMA_ROMProgrammedTranscriptBudgetedCountedSign(H, A).O.sign :
    0 <= EUF_CMA_ROMProgrammedTranscriptBudgetedCountedSign.adversary_hash_count /\
    EUF_CMA_ROMProgrammedTranscriptBudgetedCountedSign.adversary_hash_count <=
      hash_query_budget_count /\
    0 <= EUF_CMA_ROMProgrammedTranscriptBudgetedCountedSign.signing_count /\
    EUF_CMA_ROMProgrammedTranscriptBudgetedCountedSign.signing_count <=
      signature_query_budget_count ==>
    0 <= EUF_CMA_ROMProgrammedTranscriptBudgetedCountedSign.adversary_hash_count /\
    EUF_CMA_ROMProgrammedTranscriptBudgetedCountedSign.adversary_hash_count <=
      hash_query_budget_count /\
    0 <= EUF_CMA_ROMProgrammedTranscriptBudgetedCountedSign.signing_count /\
    EUF_CMA_ROMProgrammedTranscriptBudgetedCountedSign.signing_count <=
      signature_query_budget_count].
\end{lstlisting}
\noindent\textbf{Formal mathematical reading.}
Mathematical theorem. This is a relational program theorem: two games or procedures are coupled so that a precondition on their initial states implies the stated postcondition on final states and return values.

\noindent\textbf{Role in the proof.}
Use in this chapter. It contributes directly to an advantage bound, bad-event bound, or real-arithmetic composition step.

\noindent\textbf{Proof-script interpretation.}
Proof-script reading. The machine proof decomposes the theorem into rewrites, program-logic obligations, relational equivalences, and arithmetic side conditions.
\emph{Main tactics visible in the proof script:} \ec{proc}, \ec{inline}, \ec{wp}, \ec{call}, \ec{rnd}, \ec{rewrite}.
\emph{Named identifiers syntactically cited in the proof block:}
\begin{lstlisting}[language=EasyCrypt,basicstyle=\ttfamily\scriptsize]
DirectSigningCoinSampler
EUF_CMA_ROMProgrammedTranscriptBudgetedCountedSign
get
observe_transcript
program
sample
signature_query_budget_count
\end{lstlisting}

\end{formaltheorem}

\begin{formaltheorem}[EasyCrypt line 4399]
\noindent\textbf{EasyCrypt name.}
\begin{lstlisting}[language=EasyCrypt,basicstyle=\ttfamily\scriptsize]
adversary_budgeted_programmed_query_budget_preserved
\end{lstlisting}
\noindent\textbf{Checked EasyCrypt statement.}
\begin{lstlisting}[language=EasyCrypt,basicstyle=\ttfamily\scriptsize]
lemma adversary_budgeted_programmed_query_budget_preserved :
  hoare[
    A(EUF_CMA_ROMProgrammedTranscriptBudgetedCountedSign(H, A).AH,
      EUF_CMA_ROMProgrammedTranscriptBudgetedCountedSign(H, A).O).forge :
    0 <= EUF_CMA_ROMProgrammedTranscriptBudgetedCountedSign.adversary_hash_count /\
    EUF_CMA_ROMProgrammedTranscriptBudgetedCountedSign.adversary_hash_count <=
      hash_query_budget_count /\
    0 <= EUF_CMA_ROMProgrammedTranscriptBudgetedCountedSign.signing_count /\
    EUF_CMA_ROMProgrammedTranscriptBudgetedCountedSign.signing_count <=
      signature_query_budget_count ==>
    0 <= EUF_CMA_ROMProgrammedTranscriptBudgetedCountedSign.adversary_hash_count /\
    EUF_CMA_ROMProgrammedTranscriptBudgetedCountedSign.adversary_hash_count <=
      hash_query_budget_count /\
    0 <= EUF_CMA_ROMProgrammedTranscriptBudgetedCountedSign.signing_count /\
    EUF_CMA_ROMProgrammedTranscriptBudgetedCountedSign.signing_count <=
      signature_query_budget_count].
\end{lstlisting}
\noindent\textbf{Formal mathematical reading.}
Mathematical theorem. This is a relational program theorem: two games or procedures are coupled so that a precondition on their initial states implies the stated postcondition on final states and return values.

\noindent\textbf{Role in the proof.}
Use in this chapter. It preserves an invariant across an oracle call, adversary call, or game procedure.

\noindent\textbf{Proof-script interpretation.}
Proof-script reading. The machine proof decomposes the theorem into rewrites, program-logic obligations, relational equivalences, and arithmetic side conditions.
\emph{Main tactics visible in the proof script:} \ec{proc}, \ec{call}, \ec{apply}.
\emph{Named identifiers syntactically cited in the proof block:}
\begin{lstlisting}[language=EasyCrypt,basicstyle=\ttfamily\scriptsize]
EUF_CMA_ROMProgrammedTranscriptBudgetedCountedSign
adversary_hash_count
budgeted_programmed_hash_get_preserves_query_budget
budgeted_programmed_sign_oracle_preserves_query_budget
hash_query_budget_count
signature_query_budget_count
signing_count
\end{lstlisting}

\end{formaltheorem}

\begin{formaltheorem}[EasyCrypt line 4430]
\noindent\textbf{EasyCrypt name.}
\begin{lstlisting}[language=EasyCrypt,basicstyle=\ttfamily\scriptsize]
budgeted_programmed_main_query_budget_preserved
\end{lstlisting}
\noindent\textbf{Checked EasyCrypt statement.}
\begin{lstlisting}[language=EasyCrypt,basicstyle=\ttfamily\scriptsize]
lemma budgeted_programmed_main_query_budget_preserved :
  hoare[EUF_CMA_ROMProgrammedTranscriptBudgetedCountedSign(H, A).main :
    true ==>
    0 <= EUF_CMA_ROMProgrammedTranscriptBudgetedCountedSign.adversary_hash_count /\
    EUF_CMA_ROMProgrammedTranscriptBudgetedCountedSign.adversary_hash_count <=
      hash_query_budget_count /\
    0 <= EUF_CMA_ROMProgrammedTranscriptBudgetedCountedSign.signing_count /\
    EUF_CMA_ROMProgrammedTranscriptBudgetedCountedSign.signing_count <=
      signature_query_budget_count].
\end{lstlisting}
\noindent\textbf{Formal mathematical reading.}
Mathematical theorem. This is a relational program theorem: two games or procedures are coupled so that a precondition on their initial states implies the stated postcondition on final states and return values.

\noindent\textbf{Role in the proof.}
Use in this chapter. It preserves an invariant across an oracle call, adversary call, or game procedure.

\noindent\textbf{Proof-script interpretation.}
Proof-script reading. The machine proof decomposes the theorem into rewrites, program-logic obligations, relational equivalences, and arithmetic side conditions.
\emph{Main tactics visible in the proof script:} \ec{proc}, \ec{wp}, \ec{call}, \ec{conseq}, \ec{rnd}, \ec{rewrite}.
\emph{Named identifiers syntactically cited in the proof block:}
\begin{lstlisting}[language=EasyCrypt,basicstyle=\ttfamily\scriptsize]
EUF_CMA_ROMProgrammedTranscriptBudgetedCountedSign
adversary_budgeted_programmed_query_budget_preserved
adversary_hash_count
budgeted_counted_lazy_rom_bad_true
budgeted_counted_lazy_rom_get_true
budgeted_counted_lazy_rom_init_true
hash_query_budget_count
signature_query_budget_count
signing_count
\end{lstlisting}

\end{formaltheorem}

\begin{formaltheorem}[EasyCrypt line 4467]
\noindent\textbf{EasyCrypt name.}
\begin{lstlisting}[language=EasyCrypt,basicstyle=\ttfamily\scriptsize]
budgeted_programmed_query_budget_certifies_branch_bound
\end{lstlisting}
\noindent\textbf{Checked EasyCrypt statement.}
\begin{lstlisting}[language=EasyCrypt,basicstyle=\ttfamily\scriptsize]
lemma budgeted_programmed_query_budget_certifies_branch_bound &m :
  Pr[EUF_CMA_ROMProgrammedTranscriptBudgetedCountedSign(H, A).main() @ &m :
    EUF_CMA_ROMProgrammedTranscriptBudgetedCountedSign.C.bad_prequery \/
    EUF_CMA_ROMProgrammedTranscriptBudgetedCountedSign.C.bad_reprogram] <=
    ((rom_total_query_budget * rom_total_query_budget) +
     (rom_signature_query_budget * (rom_hash_query_budget + 1%r))) /
    challenge_support_cardinality_lower_bound =>
  Pr[EUF_CMA_ROMProgrammedTranscriptBudgetedCountedSign(H, A).main() @ &m :
    EUF_CMA_ROMProgrammedTranscriptBudgetedCountedSign.C.bad_prequery \/
    EUF_CMA_ROMProgrammedTranscriptBudgetedCountedSign.C.bad_reprogram] <=
  counted_rom_prequery_reprogramming_term.
\end{lstlisting}
\noindent\textbf{Formal mathematical reading.}
Mathematical theorem. This theorem has the form \(\Pr[G:E] \le B\) is a game-probability statement.  It is one of the exact hops, bad-event bounds, or final advantage inequalities used in the reduction.

\noindent\textbf{Role in the proof.}
Use in this chapter. It contributes directly to an advantage bound, bad-event bound, or real-arithmetic composition step.

\noindent\textbf{Proof-script interpretation.}
Proof-script reading. The machine proof decomposes the theorem into rewrites, program-logic obligations, relational equivalences, and arithmetic side conditions.
\emph{Main tactics visible in the proof script:} \ec{rewrite}.
\emph{Named identifiers syntactically cited in the proof block:}
\begin{lstlisting}[language=EasyCrypt,basicstyle=\ttfamily\scriptsize]
counted_rom_prequery_reprogramming_term_cardinalityE
\end{lstlisting}

\end{formaltheorem}

\begin{formaltheorem}[EasyCrypt line 4483]
\noindent\textbf{EasyCrypt name.}
\begin{lstlisting}[language=EasyCrypt,basicstyle=\ttfamily\scriptsize]
budgeted_programmed_query_budget_certifies_branch_bound_from_prequery
\end{lstlisting}
\noindent\textbf{Checked EasyCrypt statement.}
\begin{lstlisting}[language=EasyCrypt,basicstyle=\ttfamily\scriptsize]
lemma budgeted_programmed_query_budget_certifies_branch_bound_from_prequery
  &m :
  Pr[EUF_CMA_ROMProgrammedTranscriptBudgetedCountedSign(H, A).main() @ &m :
    EUF_CMA_ROMProgrammedTranscriptBudgetedCountedSign.C.bad_prequery] <=
    ((rom_total_query_budget * rom_total_query_budget) +
     (rom_signature_query_budget * (rom_hash_query_budget + 1%r))) /
    challenge_support_cardinality_lower_bound =>
  Pr[EUF_CMA_ROMProgrammedTranscriptBudgetedCountedSign(H, A).main() @ &m :
    EUF_CMA_ROMProgrammedTranscriptBudgetedCountedSign.C.bad_prequery \/
    EUF_CMA_ROMProgrammedTranscriptBudgetedCountedSign.C.bad_reprogram] <=
  counted_rom_prequery_reprogramming_term.
\end{lstlisting}
\noindent\textbf{Formal mathematical reading.}
Mathematical theorem. This theorem has the form \(\Pr[G:E] \le B\) is a game-probability statement.  It is one of the exact hops, bad-event bounds, or final advantage inequalities used in the reduction.

\noindent\textbf{Role in the proof.}
Use in this chapter. It contributes directly to an advantage bound, bad-event bound, or real-arithmetic composition step.

\noindent\textbf{Proof-script interpretation.}
Proof-script reading. The machine proof decomposes the theorem into rewrites, program-logic obligations, relational equivalences, and arithmetic side conditions.
\emph{Main tactics visible in the proof script:} \ec{apply}.
\emph{Named identifiers syntactically cited in the proof block:}
\begin{lstlisting}[language=EasyCrypt,basicstyle=\ttfamily\scriptsize]
budgeted_programmed_prequery_reprogram_bad_bound_from_prequery
budgeted_programmed_query_budget_certifies_branch_bound
prequery_bound
\end{lstlisting}

\end{formaltheorem}

\begin{formaltheorem}[EasyCrypt line 4500]
\noindent\textbf{EasyCrypt name.}
\begin{lstlisting}[language=EasyCrypt,basicstyle=\ttfamily\scriptsize]
budgeted_programmed_query_budget_certifies_checked_31_bound
\end{lstlisting}
\noindent\textbf{Checked EasyCrypt statement.}
\begin{lstlisting}[language=EasyCrypt,basicstyle=\ttfamily\scriptsize]
lemma budgeted_programmed_query_budget_certifies_checked_31_bound
  &m :
  Pr[EUF_CMA_ROMProgrammedTranscriptBudgetedCountedSign(H, A).main() @ &m :
    EUF_CMA_ROMProgrammedTranscriptBudgetedCountedSign.C.bad_prequery] <=
    ((rom_total_query_budget * rom_total_query_budget) +
     (rom_signature_query_budget * (rom_hash_query_budget + 1%r))) /
    challenge_support_cardinality_lower_bound =>
  Pr[EUF_CMA_ROMProgrammedTranscriptBudgetedCountedSign(H, A).main() @ &m :
    EUF_CMA_ROMProgrammedTranscriptBudgetedCountedSign.C.bad_prequery \/
    EUF_CMA_ROMProgrammedTranscriptBudgetedCountedSign.C.bad_reprogram] <=
  31%r / 288230376151711744%r.
\end{lstlisting}
\noindent\textbf{Formal mathematical reading.}
Mathematical theorem. This theorem has the form \(\Pr[G:E] \le B\) is a game-probability statement.  It is one of the exact hops, bad-event bounds, or final advantage inequalities used in the reduction.

\noindent\textbf{Role in the proof.}
Use in this chapter. It contributes directly to an advantage bound, bad-event bound, or real-arithmetic composition step.

\noindent\textbf{Proof-script interpretation.}
Proof-script reading. The machine proof decomposes the theorem into rewrites, program-logic obligations, relational equivalences, and arithmetic side conditions.
\emph{Main tactics visible in the proof script:} \ec{apply}, \ec{rewrite}.
\emph{Named identifiers syntactically cited in the proof block:}
\begin{lstlisting}[language=EasyCrypt,basicstyle=\ttfamily\scriptsize]
budgeted_programmed_query_budget_certifies_branch_bound_from_prequery
counted_rom_prequery_reprogramming_term
counted_rom_prequery_reprogramming_term_concreteE
ler_trans
prequery_bound
\end{lstlisting}

\end{formaltheorem}

\begin{formaltheorem}[EasyCrypt line 4529]
\noindent\textbf{EasyCrypt name.}
\begin{lstlisting}[language=EasyCrypt,basicstyle=\ttfamily\scriptsize]
budgeted_sampled_direct_counted_lazy_rom_get_true
\end{lstlisting}
\noindent\textbf{Checked EasyCrypt statement.}
\begin{lstlisting}[language=EasyCrypt,basicstyle=\ttfamily\scriptsize]
lemma budgeted_sampled_direct_counted_lazy_rom_get_true :
  hoare[EUF_CMA_ROMProgrammedTranscriptBudgetedSampledSign
          (H, A, DirectSigningCoinSampler).C.get :
    true ==> true].
\end{lstlisting}
\noindent\textbf{Formal mathematical reading.}
Mathematical theorem. This is a relational program theorem: two games or procedures are coupled so that a precondition on their initial states implies the stated postcondition on final states and return values.

\noindent\textbf{Role in the proof.}
Use in this chapter. It contributes directly to an advantage bound, bad-event bound, or real-arithmetic composition step.

\noindent\textbf{Proof-script interpretation.}
Proof-script reading. The machine proof decomposes the theorem into rewrites, program-logic obligations, relational equivalences, and arithmetic side conditions.
\emph{Main tactics visible in the proof script:} \ec{proc}, \ec{wp}, \ec{call}.

\end{formaltheorem}

\begin{formaltheorem}[EasyCrypt line 4540]
\noindent\textbf{EasyCrypt name.}
\begin{lstlisting}[language=EasyCrypt,basicstyle=\ttfamily\scriptsize]
budgeted_sampled_direct_counted_lazy_rom_init_true
\end{lstlisting}
\noindent\textbf{Checked EasyCrypt statement.}
\begin{lstlisting}[language=EasyCrypt,basicstyle=\ttfamily\scriptsize]
lemma budgeted_sampled_direct_counted_lazy_rom_init_true :
  hoare[EUF_CMA_ROMProgrammedTranscriptBudgetedSampledSign
          (H, A, DirectSigningCoinSampler).C.init :
    true ==> true].
\end{lstlisting}
\noindent\textbf{Formal mathematical reading.}
Mathematical theorem. This is a relational program theorem: two games or procedures are coupled so that a precondition on their initial states implies the stated postcondition on final states and return values.

\noindent\textbf{Role in the proof.}
Use in this chapter. It contributes directly to an advantage bound, bad-event bound, or real-arithmetic composition step.

\noindent\textbf{Proof-script interpretation.}
Proof-script reading. The machine proof decomposes the theorem into rewrites, program-logic obligations, relational equivalences, and arithmetic side conditions.
\emph{Main tactics visible in the proof script:} \ec{proc}, \ec{wp}, \ec{call}.

\end{formaltheorem}

\begin{formaltheorem}[EasyCrypt line 4551]
\noindent\textbf{EasyCrypt name.}
\begin{lstlisting}[language=EasyCrypt,basicstyle=\ttfamily\scriptsize]
budgeted_sampled_direct_counted_lazy_rom_bad_true
\end{lstlisting}
\noindent\textbf{Checked EasyCrypt statement.}
\begin{lstlisting}[language=EasyCrypt,basicstyle=\ttfamily\scriptsize]
lemma budgeted_sampled_direct_counted_lazy_rom_bad_true :
  hoare[EUF_CMA_ROMProgrammedTranscriptBudgetedSampledSign
          (H, A, DirectSigningCoinSampler).C.bad :
    true ==> true].
\end{lstlisting}
\noindent\textbf{Formal mathematical reading.}
Mathematical theorem. This is a relational program theorem: two games or procedures are coupled so that a precondition on their initial states implies the stated postcondition on final states and return values.

\noindent\textbf{Role in the proof.}
Use in this chapter. It contributes directly to an advantage bound, bad-event bound, or real-arithmetic composition step.

\noindent\textbf{Proof-script interpretation.}
Proof-script reading. The machine proof decomposes the theorem into rewrites, program-logic obligations, relational equivalences, and arithmetic side conditions.
\emph{Main tactics visible in the proof script:} \ec{proc}.

\end{formaltheorem}

\begin{formaltheorem}[EasyCrypt line 4559]
\noindent\textbf{EasyCrypt name.}
\begin{lstlisting}[language=EasyCrypt,basicstyle=\ttfamily\scriptsize]
budgeted_sampled_direct_hash_get_preserves_challenge_query_discipline
\end{lstlisting}
\noindent\textbf{Checked EasyCrypt statement.}
\begin{lstlisting}[language=EasyCrypt,basicstyle=\ttfamily\scriptsize]
lemma budgeted_sampled_direct_hash_get_preserves_challenge_query_discipline :
  hoare[EUF_CMA_ROMProgrammedTranscriptBudgetedSampledSign
          (H, A, DirectSigningCoinSampler).AH.get :
    budgeted_programmed_challenge_query_discipline
      EUF_CMA_ROMProgrammedTranscriptBudgetedSampledSign.adversary_hash_count
      EUF_CMA_ROMProgrammedTranscriptBudgetedSampledSign.C.hash_queries ==>
    budgeted_programmed_challenge_query_discipline
      EUF_CMA_ROMProgrammedTranscriptBudgetedSampledSign.adversary_hash_count
      EUF_CMA_ROMProgrammedTranscriptBudgetedSampledSign.C.hash_queries].
\end{lstlisting}
\noindent\textbf{Formal mathematical reading.}
Mathematical theorem. This is a relational program theorem: two games or procedures are coupled so that a precondition on their initial states implies the stated postcondition on final states and return values.

\noindent\textbf{Role in the proof.}
Use in this chapter. It contributes directly to an advantage bound, bad-event bound, or real-arithmetic composition step.

\noindent\textbf{Proof-script interpretation.}
Proof-script reading. The machine proof decomposes the theorem into rewrites, program-logic obligations, relational equivalences, and arithmetic side conditions.
\emph{Main tactics visible in the proof script:} \ec{proc}, \ec{inline}, \ec{wp}, \ec{call}, \ec{rewrite}, \ec{smt}.
\emph{Named identifiers syntactically cited in the proof block:}
\begin{lstlisting}[language=EasyCrypt,basicstyle=\ttfamily\scriptsize]
DirectSigningCoinSampler
EUF_CMA_ROMProgrammedTranscriptBudgetedSampledSign
budgeted_programmed_challenge_query_discipline
challenge_query_log_count_cons_le
get
hash_query_budget_count
\end{lstlisting}

\end{formaltheorem}

\begin{formaltheorem}[EasyCrypt line 4582]
\noindent\textbf{EasyCrypt name.}
\begin{lstlisting}[language=EasyCrypt,basicstyle=\ttfamily\scriptsize]
budgeted_sampled_direct_sign_oracle_preserves_challenge_query_discipline
\end{lstlisting}
\noindent\textbf{Checked EasyCrypt statement.}
\begin{lstlisting}[language=EasyCrypt,basicstyle=\ttfamily\scriptsize]
lemma budgeted_sampled_direct_sign_oracle_preserves_challenge_query_discipline :
  hoare[EUF_CMA_ROMProgrammedTranscriptBudgetedSampledSign
          (H, A, DirectSigningCoinSampler).O.sign :
    budgeted_programmed_challenge_query_discipline
      EUF_CMA_ROMProgrammedTranscriptBudgetedSampledSign.adversary_hash_count
      EUF_CMA_ROMProgrammedTranscriptBudgetedSampledSign.C.hash_queries ==>
    budgeted_programmed_challenge_query_discipline
      EUF_CMA_ROMProgrammedTranscriptBudgetedSampledSign.adversary_hash_count
      EUF_CMA_ROMProgrammedTranscriptBudgetedSampledSign.C.hash_queries].
\end{lstlisting}
\noindent\textbf{Formal mathematical reading.}
Mathematical theorem. This is a relational program theorem: two games or procedures are coupled so that a precondition on their initial states implies the stated postcondition on final states and return values.

\noindent\textbf{Role in the proof.}
Use in this chapter. It contributes directly to an advantage bound, bad-event bound, or real-arithmetic composition step.

\noindent\textbf{Proof-script interpretation.}
Proof-script reading. The machine proof decomposes the theorem into rewrites, program-logic obligations, relational equivalences, and arithmetic side conditions.
\emph{Main tactics visible in the proof script:} \ec{proc}, \ec{inline}, \ec{wp}, \ec{call}, \ec{rnd}, \ec{rewrite}, \ec{smt}.
\emph{Named identifiers syntactically cited in the proof block:}
\begin{lstlisting}[language=EasyCrypt,basicstyle=\ttfamily\scriptsize]
DirectSigningCoinSampler
EUF_CMA_ROMProgrammedTranscriptBudgetedSampledSign
budgeted_programmed_challenge_query_discipline
challenge_query_log_count_message_cons
get
observe_transcript
program
sample
\end{lstlisting}

\end{formaltheorem}

\begin{formaltheorem}[EasyCrypt line 4612]
\noindent\textbf{EasyCrypt name.}
\begin{lstlisting}[language=EasyCrypt,basicstyle=\ttfamily\scriptsize]
adversary_budgeted_sampled_direct_challenge_query_discipline_preserved
\end{lstlisting}
\noindent\textbf{Checked EasyCrypt statement.}
\begin{lstlisting}[language=EasyCrypt,basicstyle=\ttfamily\scriptsize]
lemma adversary_budgeted_sampled_direct_challenge_query_discipline_preserved :
  hoare[
    A(EUF_CMA_ROMProgrammedTranscriptBudgetedSampledSign
        (H, A, DirectSigningCoinSampler).AH,
      EUF_CMA_ROMProgrammedTranscriptBudgetedSampledSign
        (H, A, DirectSigningCoinSampler).O).forge :
    budgeted_programmed_challenge_query_discipline
      EUF_CMA_ROMProgrammedTranscriptBudgetedSampledSign.adversary_hash_count
      EUF_CMA_ROMProgrammedTranscriptBudgetedSampledSign.C.hash_queries ==>
    budgeted_programmed_challenge_query_discipline
      EUF_CMA_ROMProgrammedTranscriptBudgetedSampledSign.adversary_hash_count
      EUF_CMA_ROMProgrammedTranscriptBudgetedSampledSign.C.hash_queries].
\end{lstlisting}
\noindent\textbf{Formal mathematical reading.}
Mathematical theorem. This is a relational program theorem: two games or procedures are coupled so that a precondition on their initial states implies the stated postcondition on final states and return values.

\noindent\textbf{Role in the proof.}
Use in this chapter. It contributes directly to an advantage bound, bad-event bound, or real-arithmetic composition step.

\noindent\textbf{Proof-script interpretation.}
Proof-script reading. The machine proof decomposes the theorem into rewrites, program-logic obligations, relational equivalences, and arithmetic side conditions.
\emph{Main tactics visible in the proof script:} \ec{proc}, \ec{call}, \ec{apply}.
\emph{Named identifiers syntactically cited in the proof block:}
\begin{lstlisting}[language=EasyCrypt,basicstyle=\ttfamily\scriptsize]
EUF_CMA_ROMProgrammedTranscriptBudgetedSampledSign
adversary_hash_count
budgeted_programmed_challenge_query_discipline
budgeted_sampled_direct_hash_get_preserves_challenge_query_discipline
budgeted_sampled_direct_sign_oracle_preserves_challenge_query_discipline
hash_queries
\end{lstlisting}

\end{formaltheorem}

\begin{formaltheorem}[EasyCrypt line 4636]
\noindent\textbf{EasyCrypt name.}
\begin{lstlisting}[language=EasyCrypt,basicstyle=\ttfamily\scriptsize]
budgeted_sampled_direct_main_challenge_query_discipline
\end{lstlisting}
\noindent\textbf{Checked EasyCrypt statement.}
\begin{lstlisting}[language=EasyCrypt,basicstyle=\ttfamily\scriptsize]
lemma budgeted_sampled_direct_main_challenge_query_discipline :
  hoare[EUF_CMA_ROMProgrammedTranscriptBudgetedSampledSign
          (H, A, DirectSigningCoinSampler).main :
    true ==>
    budgeted_programmed_challenge_query_discipline
      EUF_CMA_ROMProgrammedTranscriptBudgetedSampledSign.adversary_hash_count
      EUF_CMA_ROMProgrammedTranscriptBudgetedSampledSign.C.hash_queries].
\end{lstlisting}
\noindent\textbf{Formal mathematical reading.}
Mathematical theorem. This is a relational program theorem: two games or procedures are coupled so that a precondition on their initial states implies the stated postcondition on final states and return values.

\noindent\textbf{Role in the proof.}
Use in this chapter. It contributes directly to an advantage bound, bad-event bound, or real-arithmetic composition step.

\noindent\textbf{Proof-script interpretation.}
Proof-script reading. The machine proof decomposes the theorem into rewrites, program-logic obligations, relational equivalences, and arithmetic side conditions.
\emph{Main tactics visible in the proof script:} \ec{proc}, \ec{wp}, \ec{call}, \ec{conseq}, \ec{rnd}, \ec{rewrite}, \ec{smt}.
\emph{Named identifiers syntactically cited in the proof block:}
\begin{lstlisting}[language=EasyCrypt,basicstyle=\ttfamily\scriptsize]
EUF_CMA_ROMProgrammedTranscriptBudgetedSampledSign
adversary_budgeted_sampled_direct_challenge_query_discipline_preserved
adversary_hash_count
budgeted_programmed_challenge_query_discipline
budgeted_sampled_direct_counted_lazy_rom_bad_true
challenge_query_log_count
counted_lazy_rom_get_nonchallenge_preserves_challenge_query_count
counted_lazy_rom_init_clears_hash_queries
hash_queries
hash_query_budget_count
matrix_expand_query
ro_query_is_challenge
sd
\end{lstlisting}

\end{formaltheorem}

\begin{formaltheorem}[EasyCrypt line 4668]
\noindent\textbf{EasyCrypt name.}
\begin{lstlisting}[language=EasyCrypt,basicstyle=\ttfamily\scriptsize]
budgeted_sampled_direct_hash_get_preserves_discipline
\end{lstlisting}
\noindent\textbf{Checked EasyCrypt statement.}
\begin{lstlisting}[language=EasyCrypt,basicstyle=\ttfamily\scriptsize]
lemma budgeted_sampled_direct_hash_get_preserves_discipline :
  hoare[EUF_CMA_ROMProgrammedTranscriptBudgetedSampledSign
          (H, A, DirectSigningCoinSampler).AH.get :
    budgeted_programmed_query_count_discipline
      EUF_CMA_ROMProgrammedTranscriptBudgetedSampledSign.adversary_hash_count
      EUF_CMA_ROMProgrammedTranscriptBudgetedSampledSign.signing_count /\
    budgeted_programmed_reprogram_discipline
      EUF_CMA_ROMProgrammedTranscriptBudgetedSampledSign.pk_current
      EUF_CMA_ROMProgrammedTranscriptBudgetedSampledSign.sk_current
      EUF_CMA_ROMProgrammedTranscriptBudgetedSampledSign.C.programmed_sites
      EUF_CMA_ROMProgrammedTranscriptBudgetedSampledSign.C.bad_reprogram ==>
    budgeted_programmed_query_count_discipline
      EUF_CMA_ROMProgrammedTranscriptBudgetedSampledSign.adversary_hash_count
      EUF_CMA_ROMProgrammedTranscriptBudgetedSampledSign.signing_count /\
    budgeted_programmed_reprogram_discipline
      EUF_CMA_ROMProgrammedTranscriptBudgetedSampledSign.pk_current
      EUF_CMA_ROMProgrammedTranscriptBudgetedSampledSign.sk_current
      EUF_CMA_ROMProgrammedTranscriptBudgetedSampledSign.C.programmed_sites
      EUF_CMA_ROMProgrammedTranscriptBudgetedSampledSign.C.bad_reprogram].
\end{lstlisting}
\noindent\textbf{Formal mathematical reading.}
Mathematical theorem. This is a relational program theorem: two games or procedures are coupled so that a precondition on their initial states implies the stated postcondition on final states and return values.

\noindent\textbf{Role in the proof.}
Use in this chapter. It contributes directly to an advantage bound, bad-event bound, or real-arithmetic composition step.

\noindent\textbf{Proof-script interpretation.}
Proof-script reading. The machine proof decomposes the theorem into rewrites, program-logic obligations, relational equivalences, and arithmetic side conditions.
\emph{Main tactics visible in the proof script:} \ec{proc}, \ec{call}, \ec{wp}, \ec{rewrite}, \ec{smt}.
\emph{Named identifiers syntactically cited in the proof block:}
\begin{lstlisting}[language=EasyCrypt,basicstyle=\ttfamily\scriptsize]
budgeted_programmed_query_count_discipline
budgeted_sampled_direct_counted_lazy_rom_get_true
hash_query_budget_count
signature_query_budget_count
\end{lstlisting}

\end{formaltheorem}

\begin{formaltheorem}[EasyCrypt line 4699]
\noindent\textbf{EasyCrypt name.}
\begin{lstlisting}[language=EasyCrypt,basicstyle=\ttfamily\scriptsize]
budgeted_sampled_direct_sign_oracle_preserves_discipline
\end{lstlisting}
\noindent\textbf{Checked EasyCrypt statement.}
\begin{lstlisting}[language=EasyCrypt,basicstyle=\ttfamily\scriptsize]
lemma budgeted_sampled_direct_sign_oracle_preserves_discipline :
  hoare[EUF_CMA_ROMProgrammedTranscriptBudgetedSampledSign
          (H, A, DirectSigningCoinSampler).O.sign :
    budgeted_programmed_query_count_discipline
      EUF_CMA_ROMProgrammedTranscriptBudgetedSampledSign.adversary_hash_count
      EUF_CMA_ROMProgrammedTranscriptBudgetedSampledSign.signing_count /\
    budgeted_programmed_reprogram_discipline
      EUF_CMA_ROMProgrammedTranscriptBudgetedSampledSign.pk_current
      EUF_CMA_ROMProgrammedTranscriptBudgetedSampledSign.sk_current
      EUF_CMA_ROMProgrammedTranscriptBudgetedSampledSign.C.programmed_sites
      EUF_CMA_ROMProgrammedTranscriptBudgetedSampledSign.C.bad_reprogram ==>
    budgeted_programmed_query_count_discipline
      EUF_CMA_ROMProgrammedTranscriptBudgetedSampledSign.adversary_hash_count
      EUF_CMA_ROMProgrammedTranscriptBudgetedSampledSign.signing_count /\
    budgeted_programmed_reprogram_discipline
      EUF_CMA_ROMProgrammedTranscriptBudgetedSampledSign.pk_current
      EUF_CMA_ROMProgrammedTranscriptBudgetedSampledSign.sk_current
      EUF_CMA_ROMProgrammedTranscriptBudgetedSampledSign.C.programmed_sites
      EUF_CMA_ROMProgrammedTranscriptBudgetedSampledSign.C.bad_reprogram].
\end{lstlisting}
\noindent\textbf{Formal mathematical reading.}
Mathematical theorem. This is a relational program theorem: two games or procedures are coupled so that a precondition on their initial states implies the stated postcondition on final states and return values.

\noindent\textbf{Role in the proof.}
Use in this chapter. It contributes directly to an advantage bound, bad-event bound, or real-arithmetic composition step.

\noindent\textbf{Proof-script interpretation.}
Proof-script reading. The machine proof decomposes the theorem into rewrites, program-logic obligations, relational equivalences, and arithmetic side conditions.
\emph{Main tactics visible in the proof script:} \ec{proc}, \ec{inline}, \ec{wp}, \ec{call}, \ec{rnd}, \ec{case}, \ec{split}, \ec{rewrite}, \ec{smt}, \ec{apply}.
\emph{Named identifiers syntactically cited in the proof block:}
\begin{lstlisting}[language=EasyCrypt,basicstyle=\ttfamily\scriptsize]
CountedLazyROM
DirectSigningCoinSampler
EUF_CMA_ROMProgrammedTranscriptBudgetedSampledSign
bad_reprogram
budgeted_programmed_query_count_discipline
budgeted_programmed_reprogram_discipline
budgeted_programmed_reprogram_discipline_after_actual_signature
ctx
exact
get
hash_query_budget_count
hcount
hdis
hr
hsign_case
hsign_lt
observe_transcript
pk_current
program
programmed_sites
sample
sampled_coins
signature_query_budget_count
signing_count
sk_current
\end{lstlisting}

\end{formaltheorem}

\begin{formaltheorem}[EasyCrypt line 4769]
\noindent\textbf{EasyCrypt name.}
\begin{lstlisting}[language=EasyCrypt,basicstyle=\ttfamily\scriptsize]
adversary_budgeted_sampled_direct_discipline_preserved
\end{lstlisting}
\noindent\textbf{Checked EasyCrypt statement.}
\begin{lstlisting}[language=EasyCrypt,basicstyle=\ttfamily\scriptsize]
lemma adversary_budgeted_sampled_direct_discipline_preserved :
  hoare[
    A(EUF_CMA_ROMProgrammedTranscriptBudgetedSampledSign
        (H, A, DirectSigningCoinSampler).AH,
      EUF_CMA_ROMProgrammedTranscriptBudgetedSampledSign
        (H, A, DirectSigningCoinSampler).O).forge :
    budgeted_programmed_query_count_discipline
      EUF_CMA_ROMProgrammedTranscriptBudgetedSampledSign.adversary_hash_count
      EUF_CMA_ROMProgrammedTranscriptBudgetedSampledSign.signing_count /\
    budgeted_programmed_reprogram_discipline
      EUF_CMA_ROMProgrammedTranscriptBudgetedSampledSign.pk_current
      EUF_CMA_ROMProgrammedTranscriptBudgetedSampledSign.sk_current
      EUF_CMA_ROMProgrammedTranscriptBudgetedSampledSign.C.programmed_sites
      EUF_CMA_ROMProgrammedTranscriptBudgetedSampledSign.C.bad_reprogram ==>
    budgeted_programmed_query_count_discipline
      EUF_CMA_ROMProgrammedTranscriptBudgetedSampledSign.adversary_hash_count
      EUF_CMA_ROMProgrammedTranscriptBudgetedSampledSign.signing_count /\
    budgeted_programmed_reprogram_discipline
      EUF_CMA_ROMProgrammedTranscriptBudgetedSampledSign.pk_current
      EUF_CMA_ROMProgrammedTranscriptBudgetedSampledSign.sk_current
      EUF_CMA_ROMProgrammedTranscriptBudgetedSampledSign.C.programmed_sites
      EUF_CMA_ROMProgrammedTranscriptBudgetedSampledSign.C.bad_reprogram].
\end{lstlisting}
\noindent\textbf{Formal mathematical reading.}
Mathematical theorem. This is a relational program theorem: two games or procedures are coupled so that a precondition on their initial states implies the stated postcondition on final states and return values.

\noindent\textbf{Role in the proof.}
Use in this chapter. It contributes directly to an advantage bound, bad-event bound, or real-arithmetic composition step.

\noindent\textbf{Proof-script interpretation.}
Proof-script reading. The machine proof decomposes the theorem into rewrites, program-logic obligations, relational equivalences, and arithmetic side conditions.
\emph{Main tactics visible in the proof script:} \ec{proc}, \ec{call}, \ec{apply}.
\emph{Named identifiers syntactically cited in the proof block:}
\begin{lstlisting}[language=EasyCrypt,basicstyle=\ttfamily\scriptsize]
EUF_CMA_ROMProgrammedTranscriptBudgetedSampledSign
adversary_hash_count
bad_reprogram
budgeted_programmed_query_count_discipline
budgeted_programmed_reprogram_discipline
budgeted_sampled_direct_hash_get_preserves_discipline
budgeted_sampled_direct_sign_oracle_preserves_discipline
pk_current
programmed_sites
signing_count
sk_current
\end{lstlisting}

\end{formaltheorem}

\begin{formaltheorem}[EasyCrypt line 4808]
\noindent\textbf{EasyCrypt name.}
\begin{lstlisting}[language=EasyCrypt,basicstyle=\ttfamily\scriptsize]
budgeted_sampled_direct_main_reprogram_clear
\end{lstlisting}
\noindent\textbf{Checked EasyCrypt statement.}
\begin{lstlisting}[language=EasyCrypt,basicstyle=\ttfamily\scriptsize]
lemma budgeted_sampled_direct_main_reprogram_clear :
  hoare[EUF_CMA_ROMProgrammedTranscriptBudgetedSampledSign
          (H, A, DirectSigningCoinSampler).main :
    true ==>
    ! EUF_CMA_ROMProgrammedTranscriptBudgetedSampledSign.C.bad_reprogram].
\end{lstlisting}
\noindent\textbf{Formal mathematical reading.}
Mathematical theorem. This is a relational program theorem: two games or procedures are coupled so that a precondition on their initial states implies the stated postcondition on final states and return values.

\noindent\textbf{Role in the proof.}
Use in this chapter. It contributes directly to an advantage bound, bad-event bound, or real-arithmetic composition step.

\noindent\textbf{Proof-script interpretation.}
Proof-script reading. The machine proof decomposes the theorem into rewrites, program-logic obligations, relational equivalences, and arithmetic side conditions.
\emph{Main tactics visible in the proof script:} \ec{proc}, \ec{wp}, \ec{call}, \ec{conseq}, \ec{rnd}, \ec{rewrite}, \ec{smt}.
\emph{Named identifiers syntactically cited in the proof block:}
\begin{lstlisting}[language=EasyCrypt,basicstyle=\ttfamily\scriptsize]
EUF_CMA_ROMProgrammedTranscriptBudgetedSampledSign
adversary_budgeted_sampled_direct_discipline_preserved
adversary_hash_count
bad_reprogram
budgeted_programmed_query_count_discipline
budgeted_programmed_reprogram_discipline
budgeted_sampled_direct_counted_lazy_rom_get_true
counted_lazy_rom_init_clears
hash_query_budget_count
pk_current
programmed_sites
sd
signature_query_budget_count
signing_count
sk_current
\end{lstlisting}

\end{formaltheorem}

\begin{formaltheorem}[EasyCrypt line 4859]
\noindent\textbf{EasyCrypt name.}
\begin{lstlisting}[language=EasyCrypt,basicstyle=\ttfamily\scriptsize]
budgeted_sampled_direct_bad_reprogram_zero
\end{lstlisting}
\noindent\textbf{Checked EasyCrypt statement.}
\begin{lstlisting}[language=EasyCrypt,basicstyle=\ttfamily\scriptsize]
lemma budgeted_sampled_direct_bad_reprogram_zero &m :
  Pr[EUF_CMA_ROMProgrammedTranscriptBudgetedSampledSign
        (H, A, DirectSigningCoinSampler).main() @ &m :
    EUF_CMA_ROMProgrammedTranscriptBudgetedSampledSign.C.bad_reprogram] =
  0%r.
\end{lstlisting}
\noindent\textbf{Formal mathematical reading.}
Mathematical theorem. This theorem has the form \(\Pr[G:E] = B\) is a game-probability statement.  It is one of the exact hops, bad-event bounds, or final advantage inequalities used in the reduction.

\noindent\textbf{Role in the proof.}
Use in this chapter. It contributes directly to an advantage bound, bad-event bound, or real-arithmetic composition step.

\noindent\textbf{Proof-script interpretation.}
Proof-script reading. The machine proof decomposes the theorem into rewrites, program-logic obligations, relational equivalences, and arithmetic side conditions.
\emph{Main tactics visible in the proof script:} \ec{byphoare}, \ec{hoare}, \ec{conseq}.
\emph{Named identifiers syntactically cited in the proof block:}
\begin{lstlisting}[language=EasyCrypt,basicstyle=\ttfamily\scriptsize]
EUF_CMA_ROMProgrammedTranscriptBudgetedSampledSign
bad_reprogram
budgeted_sampled_direct_main_reprogram_clear
\end{lstlisting}

\end{formaltheorem}

\begin{formaltheorem}[EasyCrypt line 4871]
\noindent\textbf{EasyCrypt name.}
\begin{lstlisting}[language=EasyCrypt,basicstyle=\ttfamily\scriptsize]
budgeted_sampled_direct_one_step_bound_nonnegative
\end{lstlisting}
\noindent\textbf{Checked EasyCrypt statement.}
\begin{lstlisting}[language=EasyCrypt,basicstyle=\ttfamily\scriptsize]
lemma budgeted_sampled_direct_one_step_bound_nonnegative :
  0%r <=
  (rom_hash_query_budget + 1%r) /
    challenge_support_cardinality_lower_bound.
\end{lstlisting}
\noindent\textbf{Formal mathematical reading.}
Mathematical theorem. This is an inequality, usually an arithmetic loss bound, event inclusion bound, or monotonicity fact used to compose probabilities.

\noindent\textbf{Role in the proof.}
Use in this chapter. It contributes directly to an advantage bound, bad-event bound, or real-arithmetic composition step.

\noindent\textbf{Proof-script interpretation.}
Proof-script reading. The machine proof decomposes the theorem into rewrites, program-logic obligations, relational equivalences, and arithmetic side conditions.
\emph{Main tactics visible in the proof script:} \ec{apply}, \ec{rewrite}.
\emph{Named identifiers syntactically cited in the proof block:}
\begin{lstlisting}[language=EasyCrypt,basicstyle=\ttfamily\scriptsize]
addr_ge0
challenge_support_cardinality_lower_bound_nonnegative
divr_ge0
rom_hash_query_budget_nonnegative
\end{lstlisting}

\end{formaltheorem}

\begin{formaltheorem}[EasyCrypt line 4881]
\noindent\textbf{EasyCrypt name.}
\begin{lstlisting}[language=EasyCrypt,basicstyle=\ttfamily\scriptsize]
budgeted_sampled_direct_sampler_message_hash_prequery_bound
\end{lstlisting}
\noindent\textbf{Checked EasyCrypt statement.}
\begin{lstlisting}[language=EasyCrypt,basicstyle=\ttfamily\scriptsize]
lemma budgeted_sampled_direct_sampler_message_hash_prequery_bound
  pk sk qs m ctx :
  phoare[DirectSigningCoinSampler.sample :
    budgeted_programmed_challenge_query_discipline
      EUF_CMA_ROMProgrammedTranscriptBudgetedSampledSign.adversary_hash_count
      qs /\
    pk = public_key_of_secret haetae_mode sk /\
    EUF_CMA_ROMProgrammedTranscriptBudgetedSampledSign.pk_current = pk /\
    EUF_CMA_ROMProgrammedTranscriptBudgetedSampledSign.sk_current = sk /\
    EUF_CMA_ROMProgrammedTranscriptBudgetedSampledSign.C.hash_queries = qs ==>
    programming_site_prequeried
      (message_hash_query
         (public_key_of_secret haetae_mode sk) ctx m :: qs)
      (actual_signature_programming_site haetae_mode
         (transcript_of_signature haetae_mode
            pk
            m ctx
            (sign_internal haetae_mode
               sk m ctx res)))] <=
  ((rom_hash_query_budget + 1%r) /
    challenge_support_cardinality_lower_bound).
\end{lstlisting}
\noindent\textbf{Formal mathematical reading.}
Mathematical theorem. This is a relational program theorem: two games or procedures are coupled so that a precondition on their initial states implies the stated postcondition on final states and return values.

\noindent\textbf{Role in the proof.}
Use in this chapter. It contributes directly to an advantage bound, bad-event bound, or real-arithmetic composition step.

\noindent\textbf{Proof-script interpretation.}
Proof-script reading. The machine proof decomposes the theorem into rewrites, program-logic obligations, relational equivalences, and arithmetic side conditions.
\emph{Main tactics visible in the proof script:} \ec{rewrite}, \ec{have}, \ec{smt}, \ec{byphoare}, \ec{conseq}.
\emph{Named identifiers syntactically cited in the proof block:}
\begin{lstlisting}[language=EasyCrypt,basicstyle=\ttfamily\scriptsize]
actual_signature_programming_site
budgeted_programmed_challenge_query_discipline
bypr
challenge_query_log_count
ctx
direct_signing_coin_sampler_message_hash_budgeted_prequery_bound
discipline
haetae_mode
hash_query_budget_count
hc_budget
m0
message_hash_query
pk
programming_site_prequeried
public_key_of_secret
qs
query_budget
query_count
sign_internal
sk
transcript_of_signature
\end{lstlisting}

\end{formaltheorem}

\begin{formaltheorem}[EasyCrypt line 4920]
\noindent\textbf{EasyCrypt name.}
\begin{lstlisting}[language=EasyCrypt,basicstyle=\ttfamily\scriptsize]
budgeted_sampled_direct_bad_prequery_bound
\end{lstlisting}
\noindent\textbf{Checked EasyCrypt statement.}
\begin{lstlisting}[language=EasyCrypt,basicstyle=\ttfamily\scriptsize]
lemma budgeted_sampled_direct_bad_prequery_bound &m :
  islossless H.get =>
  islossless H.set =>
  Pr[EUF_CMA_ROMProgrammedTranscriptBudgetedSampledSign
        (H, A, DirectSigningCoinSampler).main() @ &m :
    EUF_CMA_ROMProgrammedTranscriptBudgetedSampledSign.C.bad_prequery] <=
  ((rom_hash_query_budget + rom_signature_query_budget) *
    (rom_hash_query_budget + 1%r)) /
    challenge_support_cardinality_lower_bound.
\end{lstlisting}
\noindent\textbf{Formal mathematical reading.}
Mathematical theorem. This theorem has the form \(\Pr[G:E] \le B\) is a game-probability statement.  It is one of the exact hops, bad-event bounds, or final advantage inequalities used in the reduction.

\noindent\textbf{Role in the proof.}
Use in this chapter. It contributes directly to an advantage bound, bad-event bound, or real-arithmetic composition step.

\noindent\textbf{Proof-script interpretation.}
Proof-script reading. The machine proof decomposes the theorem into rewrites, program-logic obligations, relational equivalences, and arithmetic side conditions.
\emph{Main tactics visible in the proof script:} \ec{have}, \ec{conseq}, \ec{rewrite}, \ec{smt}, \ec{wp}, \ec{inline}, \ec{call}, \ec{rnd}, \ec{proc}, \ec{hoare}, \ec{apply}.
\emph{Named identifiers syntactically cited in the proof block:}
\begin{lstlisting}[language=EasyCrypt,basicstyle=\ttfamily\scriptsize]
AH
BRA
Bigreal
DirectSigningCoinSampler
EUF_CMA_ROMProgrammedTranscriptBudgetedSampledSign
H_get_ll
H_get_true
H_set_ll
H_set_true
StdBigop
actual_signature_programming_site
adversary_hash_count
bad_prequery
budgeted_programmed_challenge_query_discipline
budgeted_sampled_direct_one_step_bound_nonnegative
challenge_query_log_count
challenge_query_log_count_cons_le
challenge_support_cardinality_lower_bound
coins
ctx
fel
get
h1
h2
h3
h4
h5
h6
h7
h8
haetae_mode
hash_queries
hash_query_budget_count
hash_query_count
hr
init
keygen_internal
ler_trans
matrix_expand_query
message_hash_query
observe_transcript
old_bad
old_count
phoare
pk_current
program
programmed_site_prequery_one_step_message_hash_counter_bound
programming_site_prequeried
public_key_of_secret
ro_query_is_challenge
rom_hash_query_budget
rom_signature_query_budget
sample
sampled_coins
sd
set
sign
sign_internal
signature_query_budget_count
signature_query_count
signing_coin_distribution_lossless
signing_count
sk_current
sumri_const
transcript_of_signature
trivial
\end{lstlisting}

\end{formaltheorem}

\begin{formaltheorem}[EasyCrypt line 5182]
\noindent\textbf{EasyCrypt name.}
\begin{lstlisting}[language=EasyCrypt,basicstyle=\ttfamily\scriptsize]
budgeted_sampled_direct_prequery_reprogram_bad_bound
\end{lstlisting}
\noindent\textbf{Checked EasyCrypt statement.}
\begin{lstlisting}[language=EasyCrypt,basicstyle=\ttfamily\scriptsize]
lemma budgeted_sampled_direct_prequery_reprogram_bad_bound &m :
  islossless H.get =>
  islossless H.set =>
  Pr[EUF_CMA_ROMProgrammedTranscriptBudgetedSampledSign
        (H, A, DirectSigningCoinSampler).main() @ &m :
    EUF_CMA_ROMProgrammedTranscriptBudgetedSampledSign.C.bad_prequery \/
    EUF_CMA_ROMProgrammedTranscriptBudgetedSampledSign.C.bad_reprogram] <=
  counted_rom_prequery_reprogramming_term.
\end{lstlisting}
\noindent\textbf{Formal mathematical reading.}
Mathematical theorem. This theorem has the form \(\Pr[G:E] \le B\) is a game-probability statement.  It is one of the exact hops, bad-event bounds, or final advantage inequalities used in the reduction.

\noindent\textbf{Role in the proof.}
Use in this chapter. It contributes directly to an advantage bound, bad-event bound, or real-arithmetic composition step.

\noindent\textbf{Proof-script interpretation.}
Proof-script reading. The machine proof decomposes the theorem into rewrites, program-logic obligations, relational equivalences, and arithmetic side conditions.
\emph{Main tactics visible in the proof script:} \ec{apply}, \ec{rewrite}, \ec{smt}.
\emph{Named identifiers syntactically cited in the proof block:}
\begin{lstlisting}[language=EasyCrypt,basicstyle=\ttfamily\scriptsize]
DirectSigningCoinSampler
EUF_CMA_ROMProgrammedTranscriptBudgetedSampledSign
H_get_ll
H_set_ll
Pr
bad_prequery
bad_reprogram
budgeted_sampled_direct_bad_prequery_bound
budgeted_sampled_direct_bad_reprogram_zero
challenge_support_cardinality_lower_bound
counted_rom_prequery_reprogramming_term_cardinalityE
hash_query_count
ler_trans
main
mu_or
rom_hash_query_budget
rom_signature_query_budget
rom_total_query_budget
signature_query_count
\end{lstlisting}

\end{formaltheorem}

\begin{formaltheorem}[EasyCrypt line 5213]
\noindent\textbf{EasyCrypt name.}
\begin{lstlisting}[language=EasyCrypt,basicstyle=\ttfamily\scriptsize]
budgeted_sampled_direct_prequery_reprogram_bad_checked_31
\end{lstlisting}
\noindent\textbf{Checked EasyCrypt statement.}
\begin{lstlisting}[language=EasyCrypt,basicstyle=\ttfamily\scriptsize]
lemma budgeted_sampled_direct_prequery_reprogram_bad_checked_31 &m :
  islossless H.get =>
  islossless H.set =>
  Pr[EUF_CMA_ROMProgrammedTranscriptBudgetedSampledSign
        (H, A, DirectSigningCoinSampler).main() @ &m :
    EUF_CMA_ROMProgrammedTranscriptBudgetedSampledSign.C.bad_prequery \/
    EUF_CMA_ROMProgrammedTranscriptBudgetedSampledSign.C.bad_reprogram] <=
  31%r / 288230376151711744%r.
\end{lstlisting}
\noindent\textbf{Formal mathematical reading.}
Mathematical theorem. This theorem has the form \(\Pr[G:E] \le B\) is a game-probability statement.  It is one of the exact hops, bad-event bounds, or final advantage inequalities used in the reduction.

\noindent\textbf{Role in the proof.}
Use in this chapter. It contributes directly to an advantage bound, bad-event bound, or real-arithmetic composition step.

\noindent\textbf{Proof-script interpretation.}
Proof-script reading. The machine proof decomposes the theorem into rewrites, program-logic obligations, relational equivalences, and arithmetic side conditions.
\emph{Main tactics visible in the proof script:} \ec{apply}, \ec{rewrite}.
\emph{Named identifiers syntactically cited in the proof block:}
\begin{lstlisting}[language=EasyCrypt,basicstyle=\ttfamily\scriptsize]
H_get_ll
H_set_ll
budgeted_sampled_direct_prequery_reprogram_bad_bound
counted_rom_prequery_reprogramming_term
counted_rom_prequery_reprogramming_term_concreteE
ler_trans
\end{lstlisting}

\end{formaltheorem}

\begin{formaltheorem}[EasyCrypt line 5241]
\noindent\textbf{EasyCrypt name.}
\begin{lstlisting}[language=EasyCrypt,basicstyle=\ttfamily\scriptsize]
budgeted_programmed_counted_sampled_direct_hash_get_equiv
\end{lstlisting}
\noindent\textbf{Checked EasyCrypt statement.}
\begin{lstlisting}[language=EasyCrypt,basicstyle=\ttfamily\scriptsize]
equiv budgeted_programmed_counted_sampled_direct_hash_get_equiv :
  EUF_CMA_ROMProgrammedTranscriptBudgetedCountedSign(H, A).AH.get ~
  EUF_CMA_ROMProgrammedTranscriptBudgetedSampledSign
    (H, A, DirectSigningCoinSampler).AH.get :
    ={glob H, arg} /\
    EUF_CMA_ROMProgrammedTranscriptBudgetedCountedSign.adversary_hash_count{1} =
      EUF_CMA_ROMProgrammedTranscriptBudgetedSampledSign.adversary_hash_count{2} /\
    EUF_CMA_ROMProgrammedTranscriptBudgetedCountedSign.signing_count{1} =
      EUF_CMA_ROMProgrammedTranscriptBudgetedSampledSign.signing_count{2} /\
    EUF_CMA_ROMProgrammedTranscriptBudgetedCountedSign.C.hash_queries{1} =
      EUF_CMA_ROMProgrammedTranscriptBudgetedSampledSign.C.hash_queries{2} /\
    EUF_CMA_ROMProgrammedTranscriptBudgetedCountedSign.C.programmed_sites{1} =
      EUF_CMA_ROMProgrammedTranscriptBudgetedSampledSign.C.programmed_sites{2} /\
    EUF_CMA_ROMProgrammedTranscriptBudgetedCountedSign.C.signature_sites{1} =
      EUF_CMA_ROMProgrammedTranscriptBudgetedSampledSign.C.signature_sites{2} /\
    EUF_CMA_ROMProgrammedTranscriptBudgetedCountedSign.C.bad_prequery{1} =
      EUF_CMA_ROMProgrammedTranscriptBudgetedSampledSign.C.bad_prequery{2} /\
    EUF_CMA_ROMProgrammedTranscriptBudgetedCountedSign.C.bad_reprogram{1} =
      EUF_CMA_ROMProgrammedTranscriptBudgetedSampledSign.C.bad_reprogram{2} /\
    EUF_CMA_ROMProgrammedTranscriptBudgetedCountedSign.C.bad_min_entropy{1} =
      EUF_CMA_ROMProgrammedTranscriptBudgetedSampledSign.C.bad_min_entropy{2}
    ==>
    ={glob H, res} /\
    EUF_CMA_ROMProgrammedTranscriptBudgetedCountedSign.adversary_hash_count{1} =
      EUF_CMA_ROMProgrammedTranscriptBudgetedSampledSign.adversary_hash_count{2} /\
    EUF_CMA_ROMProgrammedTranscriptBudgetedCountedSign.signing_count{1} =
      EUF_CMA_ROMProgrammedTranscriptBudgetedSampledSign.signing_count{2} /\
    EUF_CMA_ROMProgrammedTranscriptBudgetedCountedSign.C.hash_queries{1} =
      EUF_CMA_ROMProgrammedTranscriptBudgetedSampledSign.C.hash_queries{2} /\
    EUF_CMA_ROMProgrammedTranscriptBudgetedCountedSign.C.programmed_sites{1} =
      EUF_CMA_ROMProgrammedTranscriptBudgetedSampledSign.C.programmed_sites{2} /\
    EUF_CMA_ROMProgrammedTranscriptBudgetedCountedSign.C.signature_sites{1} =
      EUF_CMA_ROMProgrammedTranscriptBudgetedSampledSign.C.signature_sites{2} /\
    EUF_CMA_ROMProgrammedTranscriptBudgetedCountedSign.C.bad_prequery{1} =
      EUF_CMA_ROMProgrammedTranscriptBudgetedSampledSign.C.bad_prequery{2} /\
    EUF_CMA_ROMProgrammedTranscriptBudgetedCountedSign.C.bad_reprogram{1} =
      EUF_CMA_ROMProgrammedTranscriptBudgetedSampledSign.C.bad_reprogram{2} /\
    EUF_CMA_ROMProgrammedTranscriptBudgetedCountedSign.C.bad_min_entropy{1} =
      EUF_CMA_ROMProgrammedTranscriptBudgetedSampledSign.C.bad_min_entropy{2}.
\end{lstlisting}
\noindent\textbf{Formal mathematical reading.}
Mathematical theorem. This is a relational program theorem: two games or procedures are coupled so that a precondition on their initial states implies the stated postcondition on final states and return values.

\noindent\textbf{Role in the proof.}
Use in this chapter. It contributes directly to an advantage bound, bad-event bound, or real-arithmetic composition step.

\noindent\textbf{Proof-script interpretation.}
Proof-script reading. The machine proof decomposes the theorem into rewrites, program-logic obligations, relational equivalences, and arithmetic side conditions.
\emph{Main tactics visible in the proof script:} \ec{proc}, \ec{wp}, \ec{inline}, \ec{call}.
\emph{Named identifiers syntactically cited in the proof block:}
\begin{lstlisting}[language=EasyCrypt,basicstyle=\ttfamily\scriptsize]
DirectSigningCoinSampler
EUF_CMA_ROMProgrammedTranscriptBudgetedCountedSign
EUF_CMA_ROMProgrammedTranscriptBudgetedSampledSign
arg
get
glob
\end{lstlisting}

\end{formaltheorem}

\begin{formaltheorem}[EasyCrypt line 5295]
\noindent\textbf{EasyCrypt name.}
\begin{lstlisting}[language=EasyCrypt,basicstyle=\ttfamily\scriptsize]
budgeted_programmed_counted_sampled_direct_sign_equiv
\end{lstlisting}
\noindent\textbf{Checked EasyCrypt statement.}
\begin{lstlisting}[language=EasyCrypt,basicstyle=\ttfamily\scriptsize]
equiv budgeted_programmed_counted_sampled_direct_sign_equiv :
  EUF_CMA_ROMProgrammedTranscriptBudgetedCountedSign(H, A).O.sign ~
  EUF_CMA_ROMProgrammedTranscriptBudgetedSampledSign
    (H, A, DirectSigningCoinSampler).O.sign :
    ={glob H, arg} /\
    EUF_CMA_ROMProgrammedTranscriptBudgetedCountedSign.pk_current{1} =
      EUF_CMA_ROMProgrammedTranscriptBudgetedSampledSign.pk_current{2} /\
    EUF_CMA_ROMProgrammedTranscriptBudgetedCountedSign.sk_current{1} =
      EUF_CMA_ROMProgrammedTranscriptBudgetedSampledSign.sk_current{2} /\
    EUF_CMA_ROMProgrammedTranscriptBudgetedCountedSign.queries{1} =
      EUF_CMA_ROMProgrammedTranscriptBudgetedSampledSign.queries{2} /\
    EUF_CMA_ROMProgrammedTranscriptBudgetedCountedSign.transcripts{1} =
      EUF_CMA_ROMProgrammedTranscriptBudgetedSampledSign.transcripts{2} /\
    EUF_CMA_ROMProgrammedTranscriptBudgetedCountedSign.records{1} =
      EUF_CMA_ROMProgrammedTranscriptBudgetedSampledSign.records{2} /\
    EUF_CMA_ROMProgrammedTranscriptBudgetedCountedSign.adversary_hash_count{1} =
      EUF_CMA_ROMProgrammedTranscriptBudgetedSampledSign.adversary_hash_count{2} /\
    EUF_CMA_ROMProgrammedTranscriptBudgetedCountedSign.signing_count{1} =
      EUF_CMA_ROMProgrammedTranscriptBudgetedSampledSign.signing_count{2} /\
    EUF_CMA_ROMProgrammedTranscriptBudgetedCountedSign.C.hash_queries{1} =
      EUF_CMA_ROMProgrammedTranscriptBudgetedSampledSign.C.hash_queries{2} /\
    EUF_CMA_ROMProgrammedTranscriptBudgetedCountedSign.C.programmed_sites{1} =
      EUF_CMA_ROMProgrammedTranscriptBudgetedSampledSign.C.programmed_sites{2} /\
    EUF_CMA_ROMProgrammedTranscriptBudgetedCountedSign.C.signature_sites{1} =
      EUF_CMA_ROMProgrammedTranscriptBudgetedSampledSign.C.signature_sites{2} /\
    EUF_CMA_ROMProgrammedTranscriptBudgetedCountedSign.C.bad_prequery{1} =
      EUF_CMA_ROMProgrammedTranscriptBudgetedSampledSign.C.bad_prequery{2} /\
    EUF_CMA_ROMProgrammedTranscriptBudgetedCountedSign.C.bad_reprogram{1} =
      EUF_CMA_ROMProgrammedTranscriptBudgetedSampledSign.C.bad_reprogram{2} /\
    EUF_CMA_ROMProgrammedTranscriptBudgetedCountedSign.C.bad_min_entropy{1} =
      EUF_CMA_ROMProgrammedTranscriptBudgetedSampledSign.C.bad_min_entropy{2}
    ==>
    ={glob H, res} /\
    EUF_CMA_ROMProgrammedTranscriptBudgetedCountedSign.pk_current{1} =
      EUF_CMA_ROMProgrammedTranscriptBudgetedSampledSign.pk_current{2} /\
    EUF_CMA_ROMProgrammedTranscriptBudgetedCountedSign.sk_current{1} =
      EUF_CMA_ROMProgrammedTranscriptBudgetedSampledSign.sk_current{2} /\
    EUF_CMA_ROMProgrammedTranscriptBudgetedCountedSign.queries{1} =
      EUF_CMA_ROMProgrammedTranscriptBudgetedSampledSign.queries{2} /\
    EUF_CMA_ROMProgrammedTranscriptBudgetedCountedSign.transcripts{1} =
      EUF_CMA_ROMProgrammedTranscriptBudgetedSampledSign.transcripts{2} /\
    EUF_CMA_ROMProgrammedTranscriptBudgetedCountedSign.records{1} =
      EUF_CMA_ROMProgrammedTranscriptBudgetedSampledSign.records{2} /\
    EUF_CMA_ROMProgrammedTranscriptBudgetedCountedSign.adversary_hash_count{1} =
      EUF_CMA_ROMProgrammedTranscriptBudgetedSampledSign.adversary_hash_count{2} /\
    EUF_CMA_ROMProgrammedTranscriptBudgetedCountedSign.signing_count{1} =
      EUF_CMA_ROMProgrammedTranscriptBudgetedSampledSign.signing_count{2} /\
    EUF_CMA_ROMProgrammedTranscriptBudgetedCountedSign.C.hash_queries{1} =
      EUF_CMA_ROMProgrammedTranscriptBudgetedSampledSign.C.hash_queries{2} /\
    EUF_CMA_ROMProgrammedTranscriptBudgetedCountedSign.C.programmed_sites{1} =
      EUF_CMA_ROMProgrammedTranscriptBudgetedSampledSign.C.programmed_sites{2} /\
    EUF_CMA_ROMProgrammedTranscriptBudgetedCountedSign.C.signature_sites{1} =
      EUF_CMA_ROMProgrammedTranscriptBudgetedSampledSign.C.signature_sites{2} /\
    EUF_CMA_ROMProgrammedTranscriptBudgetedCountedSign.C.bad_prequery{1} =
      EUF_CMA_ROMProgrammedTranscriptBudgetedSampledSign.C.bad_prequery{2} /\
    EUF_CMA_ROMProgrammedTranscriptBudgetedCountedSign.C.bad_reprogram{1} =
      EUF_CMA_ROMProgrammedTranscriptBudgetedSampledSign.C.bad_reprogram{2} /\
    EUF_CMA_ROMProgrammedTranscriptBudgetedCountedSign.C.bad_min_entropy{1} =
      EUF_CMA_ROMProgrammedTranscriptBudgetedSampledSign.C.bad_min_entropy{2}.
\end{lstlisting}
\noindent\textbf{Formal mathematical reading.}
Mathematical theorem. This is a relational program theorem: two games or procedures are coupled so that a precondition on their initial states implies the stated postcondition on final states and return values.

\noindent\textbf{Role in the proof.}
Use in this chapter. It contributes directly to an advantage bound, bad-event bound, or real-arithmetic composition step.

\noindent\textbf{Proof-script interpretation.}
Proof-script reading. The machine proof decomposes the theorem into rewrites, program-logic obligations, relational equivalences, and arithmetic side conditions.
\emph{Main tactics visible in the proof script:} \ec{proc}, \ec{inline}, \ec{wp}, \ec{call}, \ec{rnd}.
\emph{Named identifiers syntactically cited in the proof block:}
\begin{lstlisting}[language=EasyCrypt,basicstyle=\ttfamily\scriptsize]
DirectSigningCoinSampler
EUF_CMA_ROMProgrammedTranscriptBudgetedCountedSign
EUF_CMA_ROMProgrammedTranscriptBudgetedSampledSign
arg
get
glob
observe_transcript
program
sample
\end{lstlisting}

\end{formaltheorem}

\begin{formaltheorem}[EasyCrypt line 5380]
\noindent\textbf{EasyCrypt name.}
\begin{lstlisting}[language=EasyCrypt,basicstyle=\ttfamily\scriptsize]
adversary_budgeted_programmed_counted_sampled_direct_equiv
\end{lstlisting}
\noindent\textbf{Checked EasyCrypt statement.}
\begin{lstlisting}[language=EasyCrypt,basicstyle=\ttfamily\scriptsize]
equiv adversary_budgeted_programmed_counted_sampled_direct_equiv :
  A(EUF_CMA_ROMProgrammedTranscriptBudgetedCountedSign(H, A).AH,
    EUF_CMA_ROMProgrammedTranscriptBudgetedCountedSign(H, A).O).forge ~
  A(EUF_CMA_ROMProgrammedTranscriptBudgetedSampledSign
      (H, A, DirectSigningCoinSampler).AH,
    EUF_CMA_ROMProgrammedTranscriptBudgetedSampledSign
      (H, A, DirectSigningCoinSampler).O).forge :
    ={glob H, glob A, arg} /\
    EUF_CMA_ROMProgrammedTranscriptBudgetedCountedSign.pk_current{1} =
      EUF_CMA_ROMProgrammedTranscriptBudgetedSampledSign.pk_current{2} /\
    EUF_CMA_ROMProgrammedTranscriptBudgetedCountedSign.sk_current{1} =
      EUF_CMA_ROMProgrammedTranscriptBudgetedSampledSign.sk_current{2} /\
    EUF_CMA_ROMProgrammedTranscriptBudgetedCountedSign.queries{1} =
      EUF_CMA_ROMProgrammedTranscriptBudgetedSampledSign.queries{2} /\
    EUF_CMA_ROMProgrammedTranscriptBudgetedCountedSign.transcripts{1} =
      EUF_CMA_ROMProgrammedTranscriptBudgetedSampledSign.transcripts{2} /\
    EUF_CMA_ROMProgrammedTranscriptBudgetedCountedSign.records{1} =
      EUF_CMA_ROMProgrammedTranscriptBudgetedSampledSign.records{2} /\
    EUF_CMA_ROMProgrammedTranscriptBudgetedCountedSign.adversary_hash_count{1} =
      EUF_CMA_ROMProgrammedTranscriptBudgetedSampledSign.adversary_hash_count{2} /\
    EUF_CMA_ROMProgrammedTranscriptBudgetedCountedSign.signing_count{1} =
      EUF_CMA_ROMProgrammedTranscriptBudgetedSampledSign.signing_count{2} /\
    EUF_CMA_ROMProgrammedTranscriptBudgetedCountedSign.C.hash_queries{1} =
      EUF_CMA_ROMProgrammedTranscriptBudgetedSampledSign.C.hash_queries{2} /\
    EUF_CMA_ROMProgrammedTranscriptBudgetedCountedSign.C.programmed_sites{1} =
      EUF_CMA_ROMProgrammedTranscriptBudgetedSampledSign.C.programmed_sites{2} /\
    EUF_CMA_ROMProgrammedTranscriptBudgetedCountedSign.C.signature_sites{1} =
      EUF_CMA_ROMProgrammedTranscriptBudgetedSampledSign.C.signature_sites{2} /\
    EUF_CMA_ROMProgrammedTranscriptBudgetedCountedSign.C.bad_prequery{1} =
      EUF_CMA_ROMProgrammedTranscriptBudgetedSampledSign.C.bad_prequery{2} /\
    EUF_CMA_ROMProgrammedTranscriptBudgetedCountedSign.C.bad_reprogram{1} =
      EUF_CMA_ROMProgrammedTranscriptBudgetedSampledSign.C.bad_reprogram{2} /\
    EUF_CMA_ROMProgrammedTranscriptBudgetedCountedSign.C.bad_min_entropy{1} =
      EUF_CMA_ROMProgrammedTranscriptBudgetedSampledSign.C.bad_min_entropy{2}
    ==>
    ={glob H, res} /\
    EUF_CMA_ROMProgrammedTranscriptBudgetedCountedSign.pk_current{1} =
      EUF_CMA_ROMProgrammedTranscriptBudgetedSampledSign.pk_current{2} /\
    EUF_CMA_ROMProgrammedTranscriptBudgetedCountedSign.sk_current{1} =
      EUF_CMA_ROMProgrammedTranscriptBudgetedSampledSign.sk_current{2} /\
    EUF_CMA_ROMProgrammedTranscriptBudgetedCountedSign.queries{1} =
      EUF_CMA_ROMProgrammedTranscriptBudgetedSampledSign.queries{2} /\
    EUF_CMA_ROMProgrammedTranscriptBudgetedCountedSign.transcripts{1} =
      EUF_CMA_ROMProgrammedTranscriptBudgetedSampledSign.transcripts{2} /\
    EUF_CMA_ROMProgrammedTranscriptBudgetedCountedSign.records{1} =
      EUF_CMA_ROMProgrammedTranscriptBudgetedSampledSign.records{2} /\
    EUF_CMA_ROMProgrammedTranscriptBudgetedCountedSign.adversary_hash_count{1} =
      EUF_CMA_ROMProgrammedTranscriptBudgetedSampledSign.adversary_hash_count{2} /\
    EUF_CMA_ROMProgrammedTranscriptBudgetedCountedSign.signing_count{1} =
      EUF_CMA_ROMProgrammedTranscriptBudgetedSampledSign.signing_count{2} /\
    EUF_CMA_ROMProgrammedTranscriptBudgetedCountedSign.C.hash_queries{1} =
      EUF_CMA_ROMProgrammedTranscriptBudgetedSampledSign.C.hash_queries{2} /\
    EUF_CMA_ROMProgrammedTranscriptBudgetedCountedSign.C.programmed_sites{1} =
      EUF_CMA_ROMProgrammedTranscriptBudgetedSampledSign.C.programmed_sites{2} /\
    EUF_CMA_ROMProgrammedTranscriptBudgetedCountedSign.C.signature_sites{1} =
      EUF_CMA_ROMProgrammedTranscriptBudgetedSampledSign.C.signature_sites{2} /\
    EUF_CMA_ROMProgrammedTranscriptBudgetedCountedSign.C.bad_prequery{1} =
      EUF_CMA_ROMProgrammedTranscriptBudgetedSampledSign.C.bad_prequery{2} /\
    EUF_CMA_ROMProgrammedTranscriptBudgetedCountedSign.C.bad_reprogram{1} =
      EUF_CMA_ROMProgrammedTranscriptBudgetedSampledSign.C.bad_reprogram{2} /\
    EUF_CMA_ROMProgrammedTranscriptBudgetedCountedSign.C.bad_min_entropy{1} =
      EUF_CMA_ROMProgrammedTranscriptBudgetedSampledSign.C.bad_min_entropy{2}.
\end{lstlisting}
\noindent\textbf{Formal mathematical reading.}
Mathematical theorem. This is a relational program theorem: two games or procedures are coupled so that a precondition on their initial states implies the stated postcondition on final states and return values.

\noindent\textbf{Role in the proof.}
Use in this chapter. It contributes directly to an advantage bound, bad-event bound, or real-arithmetic composition step.

\noindent\textbf{Proof-script interpretation.}
Proof-script reading. The machine proof decomposes the theorem into rewrites, program-logic obligations, relational equivalences, and arithmetic side conditions.
\emph{Main tactics visible in the proof script:} \ec{proc}, \ec{call}.
\emph{Named identifiers syntactically cited in the proof block:}
\begin{lstlisting}[language=EasyCrypt,basicstyle=\ttfamily\scriptsize]
EUF_CMA_ROMProgrammedTranscriptBudgetedCountedSign
EUF_CMA_ROMProgrammedTranscriptBudgetedSampledSign
adversary_hash_count
bad_min_entropy
bad_prequery
bad_reprogram
budgeted_programmed_counted_sampled_direct_hash_get_equiv
budgeted_programmed_counted_sampled_direct_sign_equiv
glob
hash_queries
pk_current
programmed_sites
queries
records
signature_sites
signing_count
sk_current
transcripts
\end{lstlisting}

\end{formaltheorem}

\begin{formaltheorem}[EasyCrypt line 5480]
\noindent\textbf{EasyCrypt name.}
\begin{lstlisting}[language=EasyCrypt,basicstyle=\ttfamily\scriptsize]
budgeted_programmed_counted_sampled_direct_main_equiv
\end{lstlisting}
\noindent\textbf{Checked EasyCrypt statement.}
\begin{lstlisting}[language=EasyCrypt,basicstyle=\ttfamily\scriptsize]
equiv budgeted_programmed_counted_sampled_direct_main_equiv :
  EUF_CMA_ROMProgrammedTranscriptBudgetedCountedSign(H, A).main ~
  EUF_CMA_ROMProgrammedTranscriptBudgetedSampledSign
    (H, A, DirectSigningCoinSampler).main :
    ={glob H, glob A} ==>
    ={res} /\
    EUF_CMA_ROMProgrammedTranscriptBudgetedCountedSign.C.bad_prequery{1} =
      EUF_CMA_ROMProgrammedTranscriptBudgetedSampledSign.C.bad_prequery{2} /\
    EUF_CMA_ROMProgrammedTranscriptBudgetedCountedSign.C.bad_reprogram{1} =
      EUF_CMA_ROMProgrammedTranscriptBudgetedSampledSign.C.bad_reprogram{2}.
\end{lstlisting}
\noindent\textbf{Formal mathematical reading.}
Mathematical theorem. This is a relational program theorem: two games or procedures are coupled so that a precondition on their initial states implies the stated postcondition on final states and return values.

\noindent\textbf{Role in the proof.}
Use in this chapter. It contributes directly to an advantage bound, bad-event bound, or real-arithmetic composition step.

\noindent\textbf{Proof-script interpretation.}
Proof-script reading. The machine proof decomposes the theorem into rewrites, program-logic obligations, relational equivalences, and arithmetic side conditions.
\emph{Main tactics visible in the proof script:} \ec{proc}, \ec{inline}, \ec{wp}, \ec{call}, \ec{apply}, \ec{rnd}.
\emph{Named identifiers syntactically cited in the proof block:}
\begin{lstlisting}[language=EasyCrypt,basicstyle=\ttfamily\scriptsize]
DirectSigningCoinSampler
EUF_CMA_ROMProgrammedTranscriptBudgetedCountedSign
EUF_CMA_ROMProgrammedTranscriptBudgetedSampledSign
adversary_budgeted_programmed_counted_sampled_direct_equiv
adversary_hash_count
arg
bad
bad_min_entropy
bad_prequery
bad_reprogram
get
glob
hash_queries
init
pk_current
programmed_sites
queries
records
signature_sites
signing_count
sk_current
transcripts
\end{lstlisting}

\end{formaltheorem}

\begin{formaltheorem}[EasyCrypt line 5572]
\noindent\textbf{EasyCrypt name.}
\begin{lstlisting}[language=EasyCrypt,basicstyle=\ttfamily\scriptsize]
budgeted_programmed_counted_sampled_direct_prequery_reprogram_bad_exact
\end{lstlisting}
\noindent\textbf{Checked EasyCrypt statement.}
\begin{lstlisting}[language=EasyCrypt,basicstyle=\ttfamily\scriptsize]
lemma budgeted_programmed_counted_sampled_direct_prequery_reprogram_bad_exact
  &m :
  Pr[EUF_CMA_ROMProgrammedTranscriptBudgetedCountedSign(H, A).main() @ &m :
    EUF_CMA_ROMProgrammedTranscriptBudgetedCountedSign.C.bad_prequery \/
    EUF_CMA_ROMProgrammedTranscriptBudgetedCountedSign.C.bad_reprogram] =
  Pr[EUF_CMA_ROMProgrammedTranscriptBudgetedSampledSign
        (H, A, DirectSigningCoinSampler).main() @ &m :
    EUF_CMA_ROMProgrammedTranscriptBudgetedSampledSign.C.bad_prequery \/
    EUF_CMA_ROMProgrammedTranscriptBudgetedSampledSign.C.bad_reprogram].
\end{lstlisting}
\noindent\textbf{Formal mathematical reading.}
Mathematical theorem. This theorem has the form \(\Pr[G:E] = B\) is a game-probability statement.  It is one of the exact hops, bad-event bounds, or final advantage inequalities used in the reduction.

\noindent\textbf{Role in the proof.}
Use in this chapter. It contributes directly to an advantage bound, bad-event bound, or real-arithmetic composition step.

\noindent\textbf{Proof-script interpretation.}
Proof-script reading. The machine proof decomposes the theorem into rewrites, program-logic obligations, relational equivalences, and arithmetic side conditions.
\emph{Main tactics visible in the proof script:} \ec{byequiv}.
\emph{Named identifiers syntactically cited in the proof block:}
\begin{lstlisting}[language=EasyCrypt,basicstyle=\ttfamily\scriptsize]
budgeted_programmed_counted_sampled_direct_main_equiv
\end{lstlisting}

\end{formaltheorem}

\begin{formaltheorem}[EasyCrypt line 5585]
\noindent\textbf{EasyCrypt name.}
\begin{lstlisting}[language=EasyCrypt,basicstyle=\ttfamily\scriptsize]
budgeted_programmed_prequery_reprogram_bad_checked_31_via_direct_sampler
\end{lstlisting}
\noindent\textbf{Checked EasyCrypt statement.}
\begin{lstlisting}[language=EasyCrypt,basicstyle=\ttfamily\scriptsize]
lemma budgeted_programmed_prequery_reprogram_bad_checked_31_via_direct_sampler
  &m :
  islossless H.get =>
  islossless H.set =>
  Pr[EUF_CMA_ROMProgrammedTranscriptBudgetedCountedSign(H, A).main() @ &m :
    EUF_CMA_ROMProgrammedTranscriptBudgetedCountedSign.C.bad_prequery \/
    EUF_CMA_ROMProgrammedTranscriptBudgetedCountedSign.C.bad_reprogram] <=
  31%r / 288230376151711744%r.
\end{lstlisting}
\noindent\textbf{Formal mathematical reading.}
Mathematical theorem. This theorem has the form \(\Pr[G:E] \le B\) is a game-probability statement.  It is one of the exact hops, bad-event bounds, or final advantage inequalities used in the reduction.

\noindent\textbf{Role in the proof.}
Use in this chapter. It contributes directly to an advantage bound, bad-event bound, or real-arithmetic composition step.

\noindent\textbf{Proof-script interpretation.}
Proof-script reading. The machine proof decomposes the theorem into rewrites, program-logic obligations, relational equivalences, and arithmetic side conditions.
\emph{Main tactics visible in the proof script:} \ec{rewrite}, \ec{apply}.
\emph{Named identifiers syntactically cited in the proof block:}
\begin{lstlisting}[language=EasyCrypt,basicstyle=\ttfamily\scriptsize]
H_get_ll
H_set_ll
budgeted_programmed_counted_sampled_direct_prequery_reprogram_bad_exact
budgeted_sampled_direct_prequery_reprogram_bad_checked_31
\end{lstlisting}

\end{formaltheorem}

\begin{formaltheorem}[EasyCrypt line 5608]
\noindent\textbf{EasyCrypt name.}
\begin{lstlisting}[language=EasyCrypt,basicstyle=\ttfamily\scriptsize]
haetae_rom_internal_oracle_erasure
\end{lstlisting}
\noindent\textbf{Checked EasyCrypt statement.}
\begin{lstlisting}[language=EasyCrypt,basicstyle=\ttfamily\scriptsize]
equiv haetae_rom_internal_oracle_erasure :
  SIG.EUF_CMA(H, HAETAE, A).O.sign ~
  EUF_CMA_ROMInternalTranscriptSign(H, A).O.sign :
    ={glob H, arg} /\
    SIG.EUF_CMA.sk{1} = EUF_CMA_ROMInternalTranscriptSign.sk_current{2} /\
    SIG.EUF_CMA.queries{1} =
      EUF_CMA_ROMInternalTranscriptSign.queries{2}
    ==>
    ={glob H, res} /\
    SIG.EUF_CMA.sk{1} = EUF_CMA_ROMInternalTranscriptSign.sk_current{2} /\
    SIG.EUF_CMA.queries{1} =
      EUF_CMA_ROMInternalTranscriptSign.queries{2}.
\end{lstlisting}
\noindent\textbf{Formal mathematical reading.}
Mathematical theorem. This is a relational program theorem: two games or procedures are coupled so that a precondition on their initial states implies the stated postcondition on final states and return values.

\noindent\textbf{Role in the proof.}
Use in this chapter. It contributes directly to an advantage bound, bad-event bound, or real-arithmetic composition step.

\noindent\textbf{Proof-script interpretation.}
Proof-script reading. The machine proof decomposes the theorem into rewrites, program-logic obligations, relational equivalences, and arithmetic side conditions.
\emph{Main tactics visible in the proof script:} \ec{proc}, \ec{inline}, \ec{wp}, \ec{call}, \ec{rnd}.
\emph{Named identifiers syntactically cited in the proof block:}
\begin{lstlisting}[language=EasyCrypt,basicstyle=\ttfamily\scriptsize]
HAETAE
arg
glob
sign
\end{lstlisting}

\end{formaltheorem}

\begin{formaltheorem}[EasyCrypt line 5637]
\noindent\textbf{EasyCrypt name.}
\begin{lstlisting}[language=EasyCrypt,basicstyle=\ttfamily\scriptsize]
adversary_haetae_rom_internal_erasure
\end{lstlisting}
\noindent\textbf{Checked EasyCrypt statement.}
\begin{lstlisting}[language=EasyCrypt,basicstyle=\ttfamily\scriptsize]
equiv adversary_haetae_rom_internal_erasure :
  A(H, SIG.EUF_CMA(H, HAETAE, A).O).forge ~
  A(H, EUF_CMA_ROMInternalTranscriptSign(H, A).O).forge :
    ={glob H, glob A, arg} /\
    SIG.EUF_CMA.sk{1} = EUF_CMA_ROMInternalTranscriptSign.sk_current{2} /\
    SIG.EUF_CMA.queries{1} =
      EUF_CMA_ROMInternalTranscriptSign.queries{2}
    ==>
    ={glob H, res} /\
    SIG.EUF_CMA.sk{1} = EUF_CMA_ROMInternalTranscriptSign.sk_current{2} /\
    SIG.EUF_CMA.queries{1} =
      EUF_CMA_ROMInternalTranscriptSign.queries{2}.
\end{lstlisting}
\noindent\textbf{Formal mathematical reading.}
Mathematical theorem. This is a relational program theorem: two games or procedures are coupled so that a precondition on their initial states implies the stated postcondition on final states and return values.

\noindent\textbf{Role in the proof.}
Use in this chapter. It proves an exact game replacement or simulator equivalence.

\noindent\textbf{Proof-script interpretation.}
Proof-script reading. The machine proof decomposes the theorem into rewrites, program-logic obligations, relational equivalences, and arithmetic side conditions.
\emph{Main tactics visible in the proof script:} \ec{proc}, \ec{inline}, \ec{wp}, \ec{call}, \ec{rnd}.
\emph{Named identifiers syntactically cited in the proof block:}
\begin{lstlisting}[language=EasyCrypt,basicstyle=\ttfamily\scriptsize]
EUF_CMA
EUF_CMA_ROMInternalTranscriptSign
HAETAE
SIG
arg
glob
queries
sign
sk
sk_current
\end{lstlisting}

\end{formaltheorem}

\begin{formaltheorem}[EasyCrypt line 5672]
\noindent\textbf{EasyCrypt name.}
\begin{lstlisting}[language=EasyCrypt,basicstyle=\ttfamily\scriptsize]
haetae_rom_internal_transcript_erasure_exact
\end{lstlisting}
\noindent\textbf{Checked EasyCrypt statement.}
\begin{lstlisting}[language=EasyCrypt,basicstyle=\ttfamily\scriptsize]
lemma haetae_rom_internal_transcript_erasure_exact &m :
  Pr[SIG.EUF_CMA(H, HAETAE, A).main() @ &m : res] =
  Pr[EUF_CMA_ROMInternalTranscriptSign(H, A).main() @ &m : res].
\end{lstlisting}
\noindent\textbf{Formal mathematical reading.}
Mathematical theorem. This theorem has the form \(\Pr[G:E] = B\) is a game-probability statement.  It is one of the exact hops, bad-event bounds, or final advantage inequalities used in the reduction.

\noindent\textbf{Role in the proof.}
Use in this chapter. It proves an exact game replacement or simulator equivalence.

\noindent\textbf{Proof-script interpretation.}
Proof-script reading. The machine proof decomposes the theorem into rewrites, program-logic obligations, relational equivalences, and arithmetic side conditions.
\emph{Main tactics visible in the proof script:} \ec{byequiv}, \ec{proc}, \ec{inline}, \ec{wp}, \ec{call}, \ec{apply}, \ec{rnd}.
\emph{Named identifiers syntactically cited in the proof block:}
\begin{lstlisting}[language=EasyCrypt,basicstyle=\ttfamily\scriptsize]
EUF_CMA
EUF_CMA_ROMInternalTranscriptSign
HAETAE
SIG
adversary_haetae_rom_internal_erasure
arg
glob
kg
queries
sk
sk_current
verify
\end{lstlisting}

\end{formaltheorem}

\begin{formaltheorem}[EasyCrypt line 5713]
\noindent\textbf{EasyCrypt name.}
\begin{lstlisting}[language=EasyCrypt,basicstyle=\ttfamily\scriptsize]
rom_internal_oracle_structural_nma
\end{lstlisting}
\noindent\textbf{Checked EasyCrypt statement.}
\begin{lstlisting}[language=EasyCrypt,basicstyle=\ttfamily\scriptsize]
equiv rom_internal_oracle_structural_nma :
  EUF_CMA_ROMInternalTranscriptSign(H, A).O.sign ~
  ROMInternalTranscriptAsNMA(A, H).O.sign :
    ={glob H, arg} /\
    EUF_CMA_ROMInternalTranscriptSign.pk_current{1} =
      ROMInternalTranscriptAsNMA.pk_current{2} /\
    EUF_CMA_ROMInternalTranscriptSign.sk_current{1} =
      ROMInternalTranscriptAsNMA.sk_current{2} /\
    EUF_CMA_ROMInternalTranscriptSign.queries{1} =
      ROMInternalTranscriptAsNMA.queries{2} /\
    EUF_CMA_ROMInternalTranscriptSign.transcripts{1} =
      ROMInternalTranscriptAsNMA.transcripts{2} /\
    EUF_CMA_ROMInternalTranscriptSign.records{1} =
      ROMInternalTranscriptAsNMA.records{2}
    ==>
    ={glob H, res} /\
    EUF_CMA_ROMInternalTranscriptSign.pk_current{1} =
      ROMInternalTranscriptAsNMA.pk_current{2} /\
    EUF_CMA_ROMInternalTranscriptSign.sk_current{1} =
      ROMInternalTranscriptAsNMA.sk_current{2} /\
    EUF_CMA_ROMInternalTranscriptSign.queries{1} =
      ROMInternalTranscriptAsNMA.queries{2} /\
    EUF_CMA_ROMInternalTranscriptSign.transcripts{1} =
      ROMInternalTranscriptAsNMA.transcripts{2} /\
    EUF_CMA_ROMInternalTranscriptSign.records{1} =
      ROMInternalTranscriptAsNMA.records{2}.
\end{lstlisting}
\noindent\textbf{Formal mathematical reading.}
Mathematical theorem. This is a relational program theorem: two games or procedures are coupled so that a precondition on their initial states implies the stated postcondition on final states and return values.

\noindent\textbf{Role in the proof.}
Use in this chapter. It contributes directly to an advantage bound, bad-event bound, or real-arithmetic composition step.

\noindent\textbf{Proof-script interpretation.}
Proof-script reading. The machine proof decomposes the theorem into rewrites, program-logic obligations, relational equivalences, and arithmetic side conditions.
\emph{Main tactics visible in the proof script:} \ec{proc}, \ec{wp}, \ec{call}, \ec{rnd}.
\emph{Named identifiers syntactically cited in the proof block:}
\begin{lstlisting}[language=EasyCrypt,basicstyle=\ttfamily\scriptsize]
arg
glob
\end{lstlisting}

\end{formaltheorem}

\begin{formaltheorem}[EasyCrypt line 5755]
\noindent\textbf{EasyCrypt name.}
\begin{lstlisting}[language=EasyCrypt,basicstyle=\ttfamily\scriptsize]
adversary_rom_internal_structural_nma
\end{lstlisting}
\noindent\textbf{Checked EasyCrypt statement.}
\begin{lstlisting}[language=EasyCrypt,basicstyle=\ttfamily\scriptsize]
equiv adversary_rom_internal_structural_nma :
  A(H, EUF_CMA_ROMInternalTranscriptSign(H, A).O).forge ~
  A(H, ROMInternalTranscriptAsNMA(A, H).O).forge :
    ={glob H, glob A, arg} /\
    EUF_CMA_ROMInternalTranscriptSign.pk_current{1} =
      ROMInternalTranscriptAsNMA.pk_current{2} /\
    EUF_CMA_ROMInternalTranscriptSign.sk_current{1} =
      ROMInternalTranscriptAsNMA.sk_current{2} /\
    EUF_CMA_ROMInternalTranscriptSign.queries{1} =
      ROMInternalTranscriptAsNMA.queries{2} /\
    EUF_CMA_ROMInternalTranscriptSign.transcripts{1} =
      ROMInternalTranscriptAsNMA.transcripts{2} /\
    EUF_CMA_ROMInternalTranscriptSign.records{1} =
      ROMInternalTranscriptAsNMA.records{2}
    ==>
    ={glob H, res} /\
    EUF_CMA_ROMInternalTranscriptSign.pk_current{1} =
      ROMInternalTranscriptAsNMA.pk_current{2} /\
    EUF_CMA_ROMInternalTranscriptSign.sk_current{1} =
      ROMInternalTranscriptAsNMA.sk_current{2} /\
    EUF_CMA_ROMInternalTranscriptSign.queries{1} =
      ROMInternalTranscriptAsNMA.queries{2} /\
    EUF_CMA_ROMInternalTranscriptSign.transcripts{1} =
      ROMInternalTranscriptAsNMA.transcripts{2} /\
    EUF_CMA_ROMInternalTranscriptSign.records{1} =
      ROMInternalTranscriptAsNMA.records{2}.
\end{lstlisting}
\noindent\textbf{Formal mathematical reading.}
Mathematical theorem. This is a relational program theorem: two games or procedures are coupled so that a precondition on their initial states implies the stated postcondition on final states and return values.

\noindent\textbf{Role in the proof.}
Use in this chapter. It is a reusable checked theorem cited by later proof scripts.

\noindent\textbf{Proof-script interpretation.}
Proof-script reading. The machine proof decomposes the theorem into rewrites, program-logic obligations, relational equivalences, and arithmetic side conditions.
\emph{Main tactics visible in the proof script:} \ec{proc}, \ec{wp}, \ec{call}, \ec{rnd}.
\emph{Named identifiers syntactically cited in the proof block:}
\begin{lstlisting}[language=EasyCrypt,basicstyle=\ttfamily\scriptsize]
EUF_CMA_ROMInternalTranscriptSign
ROMInternalTranscriptAsNMA
arg
glob
pk_current
queries
records
sk_current
transcripts
\end{lstlisting}

\end{formaltheorem}

\begin{formaltheorem}[EasyCrypt line 5811]
\noindent\textbf{EasyCrypt name.}
\begin{lstlisting}[language=EasyCrypt,basicstyle=\ttfamily\scriptsize]
rom_internal_transcript_structural_nma_bound
\end{lstlisting}
\noindent\textbf{Checked EasyCrypt statement.}
\begin{lstlisting}[language=EasyCrypt,basicstyle=\ttfamily\scriptsize]
lemma rom_internal_transcript_structural_nma_bound &m :
  Pr[EUF_CMA_ROMInternalTranscriptSign(H, A).main() @ &m : res] <=
  Pr[SIG.UF_NMA(H, HAETAE, ROMInternalTranscriptAsNMA(A)).main()
       @ &m : res].
\end{lstlisting}
\noindent\textbf{Formal mathematical reading.}
Mathematical theorem. This theorem has the form \(\Pr[G:E] \le B\) is a game-probability statement.  It is one of the exact hops, bad-event bounds, or final advantage inequalities used in the reduction.

\noindent\textbf{Role in the proof.}
Use in this chapter. It contributes directly to an advantage bound, bad-event bound, or real-arithmetic composition step.

\noindent\textbf{Proof-script interpretation.}
Proof-script reading. The machine proof decomposes the theorem into rewrites, program-logic obligations, relational equivalences, and arithmetic side conditions.
\emph{Main tactics visible in the proof script:} \ec{byequiv}, \ec{proc}, \ec{inline}, \ec{wp}, \ec{call}, \ec{apply}, \ec{rnd}, \ec{rewrite}.
\emph{Named identifiers syntactically cited in the proof block:}
\begin{lstlisting}[language=EasyCrypt,basicstyle=\ttfamily\scriptsize]
EUF_CMA_ROMInternalTranscriptSign
HAETAE
ROMInternalTranscriptAsNMA
adversary_rom_internal_structural_nma
arg
forge
glob
keygen_internal
kg
pk_current
public_key_of_secret
queries
records
secret_key_of_seed
sk_current
transcripts
verify
\end{lstlisting}

\end{formaltheorem}

\begin{formaltheorem}[EasyCrypt line 5858]
\noindent\textbf{EasyCrypt name.}
\begin{lstlisting}[language=EasyCrypt,basicstyle=\ttfamily\scriptsize]
rom_internal_transcript_clear_site_structural_nma_bound
\end{lstlisting}
\noindent\textbf{Checked EasyCrypt statement.}
\begin{lstlisting}[language=EasyCrypt,basicstyle=\ttfamily\scriptsize]
lemma rom_internal_transcript_clear_site_structural_nma_bound &m :
  Pr[EUF_CMA_ROMInternalTranscriptSign(H, A).main() @ &m :
      res /\
      transcript_log_signature_programming_sites_clear haetae_mode
        EUF_CMA_ROMInternalTranscriptSign.transcripts] <=
  Pr[SIG.UF_NMA(H, HAETAE, ROMInternalTranscriptAsNMA(A)).main()
       @ &m : res].
\end{lstlisting}
\noindent\textbf{Formal mathematical reading.}
Mathematical theorem. This theorem has the form \(\Pr[G:E] \le B\) is a game-probability statement.  It is one of the exact hops, bad-event bounds, or final advantage inequalities used in the reduction.

\noindent\textbf{Role in the proof.}
Use in this chapter. It contributes directly to an advantage bound, bad-event bound, or real-arithmetic composition step.

\noindent\textbf{Proof-script interpretation.}
Proof-script reading. The machine proof decomposes the theorem into rewrites, program-logic obligations, relational equivalences, and arithmetic side conditions.
\emph{Main tactics visible in the proof script:} \ec{apply}, \ec{rewrite}.
\emph{Named identifiers syntactically cited in the proof block:}
\begin{lstlisting}[language=EasyCrypt,basicstyle=\ttfamily\scriptsize]
EUF_CMA_ROMInternalTranscriptSign
Pr
ler_trans
main
mu_sub
rom_internal_transcript_structural_nma_bound
\end{lstlisting}

\end{formaltheorem}

\begin{formaltheorem}[EasyCrypt line 5881]
\noindent\textbf{EasyCrypt name.}
\begin{lstlisting}[language=EasyCrypt,basicstyle=\ttfamily\scriptsize]
rom_internal_public_sim_nma_oracle
\end{lstlisting}
\noindent\textbf{Checked EasyCrypt statement.}
\begin{lstlisting}[language=EasyCrypt,basicstyle=\ttfamily\scriptsize]
equiv rom_internal_public_sim_nma_oracle :
  ROMInternalTranscriptAsNMA(A, H).O.sign ~
  ROMInternalTranscriptPublicSimAsNMA(A, H).O.sign :
    ={glob H, arg} /\
    ROMInternalTranscriptAsNMA.pk_current{1} =
      ROMInternalTranscriptPublicSimAsNMA.pk_current{2} /\
    public_key_of_secret haetae_mode
      ROMInternalTranscriptAsNMA.sk_current{1} =
      ROMInternalTranscriptPublicSimAsNMA.pk_current{2} /\
    ROMInternalTranscriptAsNMA.queries{1} =
      ROMInternalTranscriptPublicSimAsNMA.queries{2} /\
    ROMInternalTranscriptAsNMA.transcripts{1} =
      ROMInternalTranscriptPublicSimAsNMA.transcripts{2} /\
    ROMInternalTranscriptAsNMA.records{1} =
      ROMInternalTranscriptPublicSimAsNMA.records{2}
    ==>
    ={glob H, res} /\
    ROMInternalTranscriptAsNMA.pk_current{1} =
      ROMInternalTranscriptPublicSimAsNMA.pk_current{2} /\
    public_key_of_secret haetae_mode
      ROMInternalTranscriptAsNMA.sk_current{1} =
      ROMInternalTranscriptPublicSimAsNMA.pk_current{2} /\
    ROMInternalTranscriptAsNMA.queries{1} =
      ROMInternalTranscriptPublicSimAsNMA.queries{2} /\
    ROMInternalTranscriptAsNMA.transcripts{1} =
      ROMInternalTranscriptPublicSimAsNMA.transcripts{2} /\
    ROMInternalTranscriptAsNMA.records{1} =
      ROMInternalTranscriptPublicSimAsNMA.records{2}.
\end{lstlisting}
\noindent\textbf{Formal mathematical reading.}
Mathematical theorem. This is a relational program theorem: two games or procedures are coupled so that a precondition on their initial states implies the stated postcondition on final states and return values.

\noindent\textbf{Role in the proof.}
Use in this chapter. It contributes directly to an advantage bound, bad-event bound, or real-arithmetic composition step.

\noindent\textbf{Proof-script interpretation.}
Proof-script reading. The machine proof decomposes the theorem into rewrites, program-logic obligations, relational equivalences, and arithmetic side conditions.
\emph{Main tactics visible in the proof script:} \ec{proc}, \ec{wp}, \ec{call}, \ec{rnd}, \ec{rewrite}.
\emph{Named identifiers syntactically cited in the proof block:}
\begin{lstlisting}[language=EasyCrypt,basicstyle=\ttfamily\scriptsize]
arg
glob
public_sim_commitment_highbitsE
public_sim_commitment_lowbitsE
public_sim_signatureE
\end{lstlisting}

\end{formaltheorem}

\begin{formaltheorem}[EasyCrypt line 5927]
\noindent\textbf{EasyCrypt name.}
\begin{lstlisting}[language=EasyCrypt,basicstyle=\ttfamily\scriptsize]
adversary_rom_internal_public_sim_nma
\end{lstlisting}
\noindent\textbf{Checked EasyCrypt statement.}
\begin{lstlisting}[language=EasyCrypt,basicstyle=\ttfamily\scriptsize]
equiv adversary_rom_internal_public_sim_nma :
  A(H, ROMInternalTranscriptAsNMA(A, H).O).forge ~
  A(H, ROMInternalTranscriptPublicSimAsNMA(A, H).O).forge :
    ={glob H, glob A, arg} /\
    ROMInternalTranscriptAsNMA.pk_current{1} =
      ROMInternalTranscriptPublicSimAsNMA.pk_current{2} /\
    public_key_of_secret haetae_mode
      ROMInternalTranscriptAsNMA.sk_current{1} =
      ROMInternalTranscriptPublicSimAsNMA.pk_current{2} /\
    ROMInternalTranscriptAsNMA.queries{1} =
      ROMInternalTranscriptPublicSimAsNMA.queries{2} /\
    ROMInternalTranscriptAsNMA.transcripts{1} =
      ROMInternalTranscriptPublicSimAsNMA.transcripts{2} /\
    ROMInternalTranscriptAsNMA.records{1} =
      ROMInternalTranscriptPublicSimAsNMA.records{2}
    ==>
    ={glob H, res} /\
    ROMInternalTranscriptAsNMA.pk_current{1} =
      ROMInternalTranscriptPublicSimAsNMA.pk_current{2} /\
    public_key_of_secret haetae_mode
      ROMInternalTranscriptAsNMA.sk_current{1} =
      ROMInternalTranscriptPublicSimAsNMA.pk_current{2} /\
    ROMInternalTranscriptAsNMA.queries{1} =
      ROMInternalTranscriptPublicSimAsNMA.queries{2} /\
    ROMInternalTranscriptAsNMA.transcripts{1} =
      ROMInternalTranscriptPublicSimAsNMA.transcripts{2} /\
    ROMInternalTranscriptAsNMA.records{1} =
      ROMInternalTranscriptPublicSimAsNMA.records{2}.
\end{lstlisting}
\noindent\textbf{Formal mathematical reading.}
Mathematical theorem. This is a relational program theorem: two games or procedures are coupled so that a precondition on their initial states implies the stated postcondition on final states and return values.

\noindent\textbf{Role in the proof.}
Use in this chapter. It proves an exact game replacement or simulator equivalence.

\noindent\textbf{Proof-script interpretation.}
Proof-script reading. The machine proof decomposes the theorem into rewrites, program-logic obligations, relational equivalences, and arithmetic side conditions.
\emph{Main tactics visible in the proof script:} \ec{proc}, \ec{wp}, \ec{call}, \ec{rnd}, \ec{rewrite}.
\emph{Named identifiers syntactically cited in the proof block:}
\begin{lstlisting}[language=EasyCrypt,basicstyle=\ttfamily\scriptsize]
ROMInternalTranscriptAsNMA
ROMInternalTranscriptPublicSimAsNMA
arg
glob
haetae_mode
pk_current
public_key_of_secret
public_sim_commitment_highbitsE
public_sim_commitment_lowbitsE
public_sim_signatureE
queries
records
sk_current
transcripts
\end{lstlisting}

\end{formaltheorem}

\begin{formaltheorem}[EasyCrypt line 5988]
\noindent\textbf{EasyCrypt name.}
\begin{lstlisting}[language=EasyCrypt,basicstyle=\ttfamily\scriptsize]
rom_internal_nma_public_sim_exact
\end{lstlisting}
\noindent\textbf{Checked EasyCrypt statement.}
\begin{lstlisting}[language=EasyCrypt,basicstyle=\ttfamily\scriptsize]
lemma rom_internal_nma_public_sim_exact &m :
  Pr[SIG.UF_NMA(H, HAETAE, ROMInternalTranscriptAsNMA(A)).main()
       @ &m : res] =
  Pr[SIG.UF_NMA(H, HAETAE, ROMInternalTranscriptPublicSimAsNMA(A)).main()
       @ &m : res].
\end{lstlisting}
\noindent\textbf{Formal mathematical reading.}
Mathematical theorem. This theorem has the form \(\Pr[G:E] = B\) is a game-probability statement.  It is one of the exact hops, bad-event bounds, or final advantage inequalities used in the reduction.

\noindent\textbf{Role in the proof.}
Use in this chapter. It proves an exact game replacement or simulator equivalence.

\noindent\textbf{Proof-script interpretation.}
Proof-script reading. The machine proof decomposes the theorem into rewrites, program-logic obligations, relational equivalences, and arithmetic side conditions.
\emph{Main tactics visible in the proof script:} \ec{byequiv}, \ec{proc}, \ec{inline}, \ec{wp}, \ec{call}, \ec{apply}, \ec{rnd}, \ec{rewrite}.
\emph{Named identifiers syntactically cited in the proof block:}
\begin{lstlisting}[language=EasyCrypt,basicstyle=\ttfamily\scriptsize]
HAETAE
ROMInternalTranscriptAsNMA
ROMInternalTranscriptPublicSimAsNMA
adversary_rom_internal_public_sim_nma
arg
forge
glob
haetae_mode
keygen_internal
kg
pk_current
public_key_of_secret
queries
records
secret_key_of_seed
sk_current
transcripts
verify
\end{lstlisting}

\end{formaltheorem}

\begin{formaltheorem}[EasyCrypt line 6050]
\noindent\textbf{EasyCrypt name.}
\begin{lstlisting}[language=EasyCrypt,basicstyle=\ttfamily\scriptsize]
rom_internal_public_sim_rom_paper_sim_nma_oracle
\end{lstlisting}
\noindent\textbf{Checked EasyCrypt statement.}
\begin{lstlisting}[language=EasyCrypt,basicstyle=\ttfamily\scriptsize]
equiv rom_internal_public_sim_rom_paper_sim_nma_oracle :
  ROMInternalTranscriptPublicSimAsNMA(A, H).O.sign ~
  ROMInternalTranscriptPaperSimAsNMA
    (A, ROMPaperSimSigningSampler(H), H).O.sign :
    ={glob H, arg} /\
    ROMInternalTranscriptPublicSimAsNMA.pk_current{1} =
      ROMInternalTranscriptPaperSimAsNMA.pk_current{2} /\
    ROMPaperSimSigningSampler.pk_current{2} =
      ROMInternalTranscriptPublicSimAsNMA.pk_current{1} /\
    ROMInternalTranscriptPublicSimAsNMA.queries{1} =
      ROMInternalTranscriptPaperSimAsNMA.queries{2} /\
    ROMInternalTranscriptPublicSimAsNMA.transcripts{1} =
      ROMInternalTranscriptPaperSimAsNMA.transcripts{2} /\
    ROMInternalTranscriptPublicSimAsNMA.records{1} =
      ROMInternalTranscriptPaperSimAsNMA.records{2}
    ==>
    ={glob H, res} /\
    ROMInternalTranscriptPublicSimAsNMA.pk_current{1} =
      ROMInternalTranscriptPaperSimAsNMA.pk_current{2} /\
    ROMPaperSimSigningSampler.pk_current{2} =
      ROMInternalTranscriptPublicSimAsNMA.pk_current{1} /\
    ROMInternalTranscriptPublicSimAsNMA.queries{1} =
      ROMInternalTranscriptPaperSimAsNMA.queries{2} /\
    ROMInternalTranscriptPublicSimAsNMA.transcripts{1} =
      ROMInternalTranscriptPaperSimAsNMA.transcripts{2} /\
    ROMInternalTranscriptPublicSimAsNMA.records{1} =
      ROMInternalTranscriptPaperSimAsNMA.records{2}.
\end{lstlisting}
\noindent\textbf{Formal mathematical reading.}
Mathematical theorem. This is a relational program theorem: two games or procedures are coupled so that a precondition on their initial states implies the stated postcondition on final states and return values.

\noindent\textbf{Role in the proof.}
Use in this chapter. It contributes directly to an advantage bound, bad-event bound, or real-arithmetic composition step.

\noindent\textbf{Proof-script interpretation.}
Proof-script reading. The machine proof decomposes the theorem into rewrites, program-logic obligations, relational equivalences, and arithmetic side conditions.
\emph{Main tactics visible in the proof script:} \ec{proc}, \ec{inline}, \ec{wp}, \ec{call}, \ec{rnd}, \ec{rewrite}.
\emph{Named identifiers syntactically cited in the proof block:}
\begin{lstlisting}[language=EasyCrypt,basicstyle=\ttfamily\scriptsize]
ROMPaperSimSigningSampler
arg
glob
paper_sim_commitment_highbits_public_coinsE
paper_sim_commitment_lowbits_public_coinsE
paper_sim_signature_public_coinsE
sample
sample_with_seed
\end{lstlisting}

\end{formaltheorem}

\begin{formaltheorem}[EasyCrypt line 6097]
\noindent\textbf{EasyCrypt name.}
\begin{lstlisting}[language=EasyCrypt,basicstyle=\ttfamily\scriptsize]
adversary_rom_internal_public_sim_rom_paper_sim_nma
\end{lstlisting}
\noindent\textbf{Checked EasyCrypt statement.}
\begin{lstlisting}[language=EasyCrypt,basicstyle=\ttfamily\scriptsize]
equiv adversary_rom_internal_public_sim_rom_paper_sim_nma :
  A(H, ROMInternalTranscriptPublicSimAsNMA(A, H).O).forge ~
  A(H, ROMInternalTranscriptPaperSimAsNMA
      (A, ROMPaperSimSigningSampler(H), H).O).forge :
    ={glob H, glob A, arg} /\
    ROMInternalTranscriptPublicSimAsNMA.pk_current{1} =
      ROMInternalTranscriptPaperSimAsNMA.pk_current{2} /\
    ROMPaperSimSigningSampler.pk_current{2} =
      ROMInternalTranscriptPublicSimAsNMA.pk_current{1} /\
    ROMInternalTranscriptPublicSimAsNMA.queries{1} =
      ROMInternalTranscriptPaperSimAsNMA.queries{2} /\
    ROMInternalTranscriptPublicSimAsNMA.transcripts{1} =
      ROMInternalTranscriptPaperSimAsNMA.transcripts{2} /\
    ROMInternalTranscriptPublicSimAsNMA.records{1} =
      ROMInternalTranscriptPaperSimAsNMA.records{2}
    ==>
    ={glob H, res} /\
    ROMInternalTranscriptPublicSimAsNMA.pk_current{1} =
      ROMInternalTranscriptPaperSimAsNMA.pk_current{2} /\
    ROMPaperSimSigningSampler.pk_current{2} =
      ROMInternalTranscriptPublicSimAsNMA.pk_current{1} /\
    ROMInternalTranscriptPublicSimAsNMA.queries{1} =
      ROMInternalTranscriptPaperSimAsNMA.queries{2} /\
    ROMInternalTranscriptPublicSimAsNMA.transcripts{1} =
      ROMInternalTranscriptPaperSimAsNMA.transcripts{2} /\
    ROMInternalTranscriptPublicSimAsNMA.records{1} =
      ROMInternalTranscriptPaperSimAsNMA.records{2}.
\end{lstlisting}
\noindent\textbf{Formal mathematical reading.}
Mathematical theorem. This is a relational program theorem: two games or procedures are coupled so that a precondition on their initial states implies the stated postcondition on final states and return values.

\noindent\textbf{Role in the proof.}
Use in this chapter. It proves an exact game replacement or simulator equivalence.

\noindent\textbf{Proof-script interpretation.}
Proof-script reading. The machine proof decomposes the theorem into rewrites, program-logic obligations, relational equivalences, and arithmetic side conditions.
\emph{Main tactics visible in the proof script:} \ec{proc}, \ec{inline}, \ec{wp}, \ec{call}, \ec{rnd}, \ec{rewrite}.
\emph{Named identifiers syntactically cited in the proof block:}
\begin{lstlisting}[language=EasyCrypt,basicstyle=\ttfamily\scriptsize]
ROMInternalTranscriptPaperSimAsNMA
ROMInternalTranscriptPublicSimAsNMA
ROMPaperSimSigningSampler
arg
glob
paper_sim_commitment_highbits_public_coinsE
paper_sim_commitment_lowbits_public_coinsE
paper_sim_signature_public_coinsE
pk_current
queries
records
sample
sample_with_seed
transcripts
\end{lstlisting}

\end{formaltheorem}

\begin{formaltheorem}[EasyCrypt line 6158]
\noindent\textbf{EasyCrypt name.}
\begin{lstlisting}[language=EasyCrypt,basicstyle=\ttfamily\scriptsize]
rom_internal_nma_public_sim_rom_paper_sim_exact
\end{lstlisting}
\noindent\textbf{Checked EasyCrypt statement.}
\begin{lstlisting}[language=EasyCrypt,basicstyle=\ttfamily\scriptsize]
lemma rom_internal_nma_public_sim_rom_paper_sim_exact &m :
  Pr[SIG.UF_NMA(H, HAETAE, ROMInternalTranscriptPublicSimAsNMA(A)).main()
       @ &m : res] =
  Pr[SIG.UF_NMA(H, HAETAE,
       ROMInternalTranscriptPaperSimAsNMA
         (A, ROMPaperSimSigningSampler(H))).main() @ &m : res].
\end{lstlisting}
\noindent\textbf{Formal mathematical reading.}
Mathematical theorem. This theorem has the form \(\Pr[G:E] = B\) is a game-probability statement.  It is one of the exact hops, bad-event bounds, or final advantage inequalities used in the reduction.

\noindent\textbf{Role in the proof.}
Use in this chapter. It proves an exact game replacement or simulator equivalence.

\noindent\textbf{Proof-script interpretation.}
Proof-script reading. The machine proof decomposes the theorem into rewrites, program-logic obligations, relational equivalences, and arithmetic side conditions.
\emph{Main tactics visible in the proof script:} \ec{byequiv}, \ec{proc}, \ec{inline}, \ec{wp}, \ec{call}, \ec{apply}, \ec{rnd}, \ec{rewrite}.
\emph{Named identifiers syntactically cited in the proof block:}
\begin{lstlisting}[language=EasyCrypt,basicstyle=\ttfamily\scriptsize]
HAETAE
ROMInternalTranscriptPaperSimAsNMA
ROMInternalTranscriptPublicSimAsNMA
ROMPaperSimSigningSampler
adversary_rom_internal_public_sim_rom_paper_sim_nma
arg
forge
glob
init
keygen_internal
kg
pk_current
public_key_of_secret
queries
records
secret_key_of_seed
transcripts
verify
\end{lstlisting}

\end{formaltheorem}

\begin{formaltheorem}[EasyCrypt line 6221]
\noindent\textbf{EasyCrypt name.}
\begin{lstlisting}[language=EasyCrypt,basicstyle=\ttfamily\scriptsize]
rom_internal_real_signing_paper_sim_nma_oracle
\end{lstlisting}
\noindent\textbf{Checked EasyCrypt statement.}
\begin{lstlisting}[language=EasyCrypt,basicstyle=\ttfamily\scriptsize]
equiv rom_internal_real_signing_paper_sim_nma_oracle :
  ROMInternalTranscriptAsNMA(A, H).O.sign ~
  ROMInternalTranscriptPaperSimAsNMA
    (A, RealSigningPaperSimSampler(H), H).O.sign :
    ={glob H, arg} /\
    ROMInternalTranscriptAsNMA.pk_current{1} =
      ROMInternalTranscriptPaperSimAsNMA.pk_current{2} /\
    public_key_of_secret haetae_mode
      ROMInternalTranscriptAsNMA.sk_current{1} =
      ROMInternalTranscriptPaperSimAsNMA.pk_current{2} /\
    RealSigningPaperSimSampler.pk_current{2} =
      ROMInternalTranscriptAsNMA.pk_current{1} /\
    RealSigningPaperSimSampler.sk_current{2} =
      ROMInternalTranscriptAsNMA.sk_current{1} /\
    ROMInternalTranscriptAsNMA.queries{1} =
      ROMInternalTranscriptPaperSimAsNMA.queries{2} /\
    ROMInternalTranscriptAsNMA.transcripts{1} =
      ROMInternalTranscriptPaperSimAsNMA.transcripts{2} /\
    ROMInternalTranscriptAsNMA.records{1} =
      ROMInternalTranscriptPaperSimAsNMA.records{2}
    ==>
    ={glob H, res} /\
    ROMInternalTranscriptAsNMA.pk_current{1} =
      ROMInternalTranscriptPaperSimAsNMA.pk_current{2} /\
    public_key_of_secret haetae_mode
      ROMInternalTranscriptAsNMA.sk_current{1} =
      ROMInternalTranscriptPaperSimAsNMA.pk_current{2} /\
    RealSigningPaperSimSampler.pk_current{2} =
      ROMInternalTranscriptAsNMA.pk_current{1} /\
    RealSigningPaperSimSampler.sk_current{2} =
      ROMInternalTranscriptAsNMA.sk_current{1} /\
    ROMInternalTranscriptAsNMA.queries{1} =
      ROMInternalTranscriptPaperSimAsNMA.queries{2} /\
    ROMInternalTranscriptAsNMA.transcripts{1} =
      ROMInternalTranscriptPaperSimAsNMA.transcripts{2} /\
    ROMInternalTranscriptAsNMA.records{1} =
      ROMInternalTranscriptPaperSimAsNMA.records{2}.
\end{lstlisting}
\noindent\textbf{Formal mathematical reading.}
Mathematical theorem. This is a relational program theorem: two games or procedures are coupled so that a precondition on their initial states implies the stated postcondition on final states and return values.

\noindent\textbf{Role in the proof.}
Use in this chapter. It contributes directly to an advantage bound, bad-event bound, or real-arithmetic composition step.

\noindent\textbf{Proof-script interpretation.}
Proof-script reading. The machine proof decomposes the theorem into rewrites, program-logic obligations, relational equivalences, and arithmetic side conditions.
\emph{Main tactics visible in the proof script:} \ec{proc}, \ec{inline}, \ec{wp}, \ec{call}, \ec{rnd}, \ec{rewrite}.
\emph{Named identifiers syntactically cited in the proof block:}
\begin{lstlisting}[language=EasyCrypt,basicstyle=\ttfamily\scriptsize]
RealSigningPaperSimSampler
arg
glob
paper_sim_commitment_highbits_real_signing_coinsE
paper_sim_commitment_lowbits_real_signing_coinsE
paper_sim_signature_real_signing_coinsE
sample
sample_with_seed
\end{lstlisting}

\end{formaltheorem}

\begin{formaltheorem}[EasyCrypt line 6278]
\noindent\textbf{EasyCrypt name.}
\begin{lstlisting}[language=EasyCrypt,basicstyle=\ttfamily\scriptsize]
adversary_rom_internal_real_signing_paper_sim_nma
\end{lstlisting}
\noindent\textbf{Checked EasyCrypt statement.}
\begin{lstlisting}[language=EasyCrypt,basicstyle=\ttfamily\scriptsize]
equiv adversary_rom_internal_real_signing_paper_sim_nma :
  A(H, ROMInternalTranscriptAsNMA(A, H).O).forge ~
  A(H, ROMInternalTranscriptPaperSimAsNMA
      (A, RealSigningPaperSimSampler(H), H).O).forge :
    ={glob H, glob A, arg} /\
    ROMInternalTranscriptAsNMA.pk_current{1} =
      ROMInternalTranscriptPaperSimAsNMA.pk_current{2} /\
    public_key_of_secret haetae_mode
      ROMInternalTranscriptAsNMA.sk_current{1} =
      ROMInternalTranscriptPaperSimAsNMA.pk_current{2} /\
    RealSigningPaperSimSampler.pk_current{2} =
      ROMInternalTranscriptAsNMA.pk_current{1} /\
    RealSigningPaperSimSampler.sk_current{2} =
      ROMInternalTranscriptAsNMA.sk_current{1} /\
    ROMInternalTranscriptAsNMA.queries{1} =
      ROMInternalTranscriptPaperSimAsNMA.queries{2} /\
    ROMInternalTranscriptAsNMA.transcripts{1} =
      ROMInternalTranscriptPaperSimAsNMA.transcripts{2} /\
    ROMInternalTranscriptAsNMA.records{1} =
      ROMInternalTranscriptPaperSimAsNMA.records{2}
    ==>
    ={glob H, res} /\
    ROMInternalTranscriptAsNMA.pk_current{1} =
      ROMInternalTranscriptPaperSimAsNMA.pk_current{2} /\
    public_key_of_secret haetae_mode
      ROMInternalTranscriptAsNMA.sk_current{1} =
      ROMInternalTranscriptPaperSimAsNMA.pk_current{2} /\
    RealSigningPaperSimSampler.pk_current{2} =
      ROMInternalTranscriptAsNMA.pk_current{1} /\
    RealSigningPaperSimSampler.sk_current{2} =
      ROMInternalTranscriptAsNMA.sk_current{1} /\
    ROMInternalTranscriptAsNMA.queries{1} =
      ROMInternalTranscriptPaperSimAsNMA.queries{2} /\
    ROMInternalTranscriptAsNMA.transcripts{1} =
      ROMInternalTranscriptPaperSimAsNMA.transcripts{2} /\
    ROMInternalTranscriptAsNMA.records{1} =
      ROMInternalTranscriptPaperSimAsNMA.records{2}.
\end{lstlisting}
\noindent\textbf{Formal mathematical reading.}
Mathematical theorem. This is a relational program theorem: two games or procedures are coupled so that a precondition on their initial states implies the stated postcondition on final states and return values.

\noindent\textbf{Role in the proof.}
Use in this chapter. It proves an exact game replacement or simulator equivalence.

\noindent\textbf{Proof-script interpretation.}
Proof-script reading. The machine proof decomposes the theorem into rewrites, program-logic obligations, relational equivalences, and arithmetic side conditions.
\emph{Main tactics visible in the proof script:} \ec{proc}, \ec{inline}, \ec{wp}, \ec{call}, \ec{rnd}, \ec{rewrite}.
\emph{Named identifiers syntactically cited in the proof block:}
\begin{lstlisting}[language=EasyCrypt,basicstyle=\ttfamily\scriptsize]
ROMInternalTranscriptAsNMA
ROMInternalTranscriptPaperSimAsNMA
RealSigningPaperSimSampler
arg
glob
haetae_mode
paper_sim_commitment_highbits_real_signing_coinsE
paper_sim_commitment_lowbits_real_signing_coinsE
paper_sim_signature_real_signing_coinsE
pk_current
public_key_of_secret
queries
records
sample
sample_with_seed
sk_current
transcripts
\end{lstlisting}

\end{formaltheorem}

\begin{formaltheorem}[EasyCrypt line 6354]
\noindent\textbf{EasyCrypt name.}
\begin{lstlisting}[language=EasyCrypt,basicstyle=\ttfamily\scriptsize]
rom_internal_nma_real_signing_paper_sim_exact
\end{lstlisting}
\noindent\textbf{Checked EasyCrypt statement.}
\begin{lstlisting}[language=EasyCrypt,basicstyle=\ttfamily\scriptsize]
lemma rom_internal_nma_real_signing_paper_sim_exact &m :
  Pr[SIG.UF_NMA(H, HAETAE, ROMInternalTranscriptAsNMA(A)).main()
       @ &m : res] =
  Pr[SIG.UF_NMA(H, HAETAE,
       ROMInternalTranscriptPaperSimAsNMA
         (A, RealSigningPaperSimSampler(H))).main() @ &m : res].
\end{lstlisting}
\noindent\textbf{Formal mathematical reading.}
Mathematical theorem. This theorem has the form \(\Pr[G:E] = B\) is a game-probability statement.  It is one of the exact hops, bad-event bounds, or final advantage inequalities used in the reduction.

\noindent\textbf{Role in the proof.}
Use in this chapter. It proves an exact game replacement or simulator equivalence.

\noindent\textbf{Proof-script interpretation.}
Proof-script reading. The machine proof decomposes the theorem into rewrites, program-logic obligations, relational equivalences, and arithmetic side conditions.
\emph{Main tactics visible in the proof script:} \ec{byequiv}, \ec{proc}, \ec{inline}, \ec{wp}, \ec{call}, \ec{apply}, \ec{rnd}, \ec{rewrite}.
\emph{Named identifiers syntactically cited in the proof block:}
\begin{lstlisting}[language=EasyCrypt,basicstyle=\ttfamily\scriptsize]
HAETAE
ROMInternalTranscriptAsNMA
ROMInternalTranscriptPaperSimAsNMA
RealSigningPaperSimSampler
adversary_rom_internal_real_signing_paper_sim_nma
arg
forge
glob
haetae_mode
init
keygen_internal
kg
pk_current
public_key_of_secret
queries
records
secret_key_of_seed
sk_current
transcripts
verify
\end{lstlisting}

\end{formaltheorem}

\begin{formaltheorem}[EasyCrypt line 6427]
\noindent\textbf{EasyCrypt name.}
\begin{lstlisting}[language=EasyCrypt,basicstyle=\ttfamily\scriptsize]
real_signing_ro_attempt_paper_sim_nma_oracle
\end{lstlisting}
\noindent\textbf{Checked EasyCrypt statement.}
\begin{lstlisting}[language=EasyCrypt,basicstyle=\ttfamily\scriptsize]
equiv real_signing_ro_attempt_paper_sim_nma_oracle :
  ROMInternalTranscriptPaperSimAsNMA
    (A, RealSigningPaperSimSampler(H), H).O.sign ~
  ROMInternalTranscriptPaperSimAsNMA
    (A, ROSigningAttemptPaperSimSampler(H), H).O.sign :
    ={glob H, arg} /\
    ROMInternalTranscriptPaperSimAsNMA.pk_current{1} =
      ROMInternalTranscriptPaperSimAsNMA.pk_current{2} /\
    RealSigningPaperSimSampler.pk_current{1} =
      ROSigningAttemptPaperSimSampler.pk_current{2} /\
    RealSigningPaperSimSampler.sk_current{1} =
      ROSigningAttemptPaperSimSampler.sk_current{2} /\
    RealSigningPaperSimSampler.pk_current{1} =
      public_key_of_secret haetae_mode
        RealSigningPaperSimSampler.sk_current{1} /\
    ROMInternalTranscriptPaperSimAsNMA.queries{1} =
      ROMInternalTranscriptPaperSimAsNMA.queries{2} /\
    ROMInternalTranscriptPaperSimAsNMA.transcripts{1} =
      ROMInternalTranscriptPaperSimAsNMA.transcripts{2} /\
    ROMInternalTranscriptPaperSimAsNMA.records{1} =
      ROMInternalTranscriptPaperSimAsNMA.records{2}
    ==>
    ={glob H, res} /\
    ROMInternalTranscriptPaperSimAsNMA.pk_current{1} =
      ROMInternalTranscriptPaperSimAsNMA.pk_current{2} /\
    RealSigningPaperSimSampler.pk_current{1} =
      ROSigningAttemptPaperSimSampler.pk_current{2} /\
    RealSigningPaperSimSampler.sk_current{1} =
      ROSigningAttemptPaperSimSampler.sk_current{2} /\
    RealSigningPaperSimSampler.pk_current{1} =
      public_key_of_secret haetae_mode
        RealSigningPaperSimSampler.sk_current{1} /\
    ROMInternalTranscriptPaperSimAsNMA.queries{1} =
      ROMInternalTranscriptPaperSimAsNMA.queries{2} /\
    ROMInternalTranscriptPaperSimAsNMA.transcripts{1} =
      ROMInternalTranscriptPaperSimAsNMA.transcripts{2} /\
    ROMInternalTranscriptPaperSimAsNMA.records{1} =
      ROMInternalTranscriptPaperSimAsNMA.records{2}.
\end{lstlisting}
\noindent\textbf{Formal mathematical reading.}
Mathematical theorem. This is a relational program theorem: two games or procedures are coupled so that a precondition on their initial states implies the stated postcondition on final states and return values.

\noindent\textbf{Role in the proof.}
Use in this chapter. It contributes directly to an advantage bound, bad-event bound, or real-arithmetic composition step.

\noindent\textbf{Proof-script interpretation.}
Proof-script reading. The machine proof decomposes the theorem into rewrites, program-logic obligations, relational equivalences, and arithmetic side conditions.
\emph{Main tactics visible in the proof script:} \ec{proc}, \ec{inline}, \ec{wp}, \ec{call}, \ec{rnd}, \ec{smt}.
\emph{Named identifiers syntactically cited in the proof block:}
\begin{lstlisting}[language=EasyCrypt,basicstyle=\ttfamily\scriptsize]
ROSigningAttemptPaperSimSampler
RealSigningPaperSimSampler
arg
glob
paper_sim_sample_from_rejection_attempt_of_coinsE
sample
sample_with_seed
\end{lstlisting}

\end{formaltheorem}

\begin{formaltheorem}[EasyCrypt line 6485]
\noindent\textbf{EasyCrypt name.}
\begin{lstlisting}[language=EasyCrypt,basicstyle=\ttfamily\scriptsize]
adversary_real_signing_ro_attempt_paper_sim_nma
\end{lstlisting}
\noindent\textbf{Checked EasyCrypt statement.}
\begin{lstlisting}[language=EasyCrypt,basicstyle=\ttfamily\scriptsize]
equiv adversary_real_signing_ro_attempt_paper_sim_nma :
  A(H, ROMInternalTranscriptPaperSimAsNMA
      (A, RealSigningPaperSimSampler(H), H).O).forge ~
  A(H, ROMInternalTranscriptPaperSimAsNMA
      (A, ROSigningAttemptPaperSimSampler(H), H).O).forge :
    ={glob H, glob A, arg} /\
    ROMInternalTranscriptPaperSimAsNMA.pk_current{1} =
      ROMInternalTranscriptPaperSimAsNMA.pk_current{2} /\
    RealSigningPaperSimSampler.pk_current{1} =
      ROSigningAttemptPaperSimSampler.pk_current{2} /\
    RealSigningPaperSimSampler.sk_current{1} =
      ROSigningAttemptPaperSimSampler.sk_current{2} /\
    RealSigningPaperSimSampler.pk_current{1} =
      public_key_of_secret haetae_mode
        RealSigningPaperSimSampler.sk_current{1} /\
    ROMInternalTranscriptPaperSimAsNMA.queries{1} =
      ROMInternalTranscriptPaperSimAsNMA.queries{2} /\
    ROMInternalTranscriptPaperSimAsNMA.transcripts{1} =
      ROMInternalTranscriptPaperSimAsNMA.transcripts{2} /\
    ROMInternalTranscriptPaperSimAsNMA.records{1} =
      ROMInternalTranscriptPaperSimAsNMA.records{2}
    ==>
    ={glob H, res} /\
    ROMInternalTranscriptPaperSimAsNMA.pk_current{1} =
      ROMInternalTranscriptPaperSimAsNMA.pk_current{2} /\
    RealSigningPaperSimSampler.pk_current{1} =
      ROSigningAttemptPaperSimSampler.pk_current{2} /\
    RealSigningPaperSimSampler.sk_current{1} =
      ROSigningAttemptPaperSimSampler.sk_current{2} /\
    RealSigningPaperSimSampler.pk_current{1} =
      public_key_of_secret haetae_mode
        RealSigningPaperSimSampler.sk_current{1} /\
    ROMInternalTranscriptPaperSimAsNMA.queries{1} =
      ROMInternalTranscriptPaperSimAsNMA.queries{2} /\
    ROMInternalTranscriptPaperSimAsNMA.transcripts{1} =
      ROMInternalTranscriptPaperSimAsNMA.transcripts{2} /\
    ROMInternalTranscriptPaperSimAsNMA.records{1} =
      ROMInternalTranscriptPaperSimAsNMA.records{2}.
\end{lstlisting}
\noindent\textbf{Formal mathematical reading.}
Mathematical theorem. This is a relational program theorem: two games or procedures are coupled so that a precondition on their initial states implies the stated postcondition on final states and return values.

\noindent\textbf{Role in the proof.}
Use in this chapter. It proves an exact game replacement or simulator equivalence.

\noindent\textbf{Proof-script interpretation.}
Proof-script reading. The machine proof decomposes the theorem into rewrites, program-logic obligations, relational equivalences, and arithmetic side conditions.
\emph{Main tactics visible in the proof script:} \ec{proc}, \ec{call}.
\emph{Named identifiers syntactically cited in the proof block:}
\begin{lstlisting}[language=EasyCrypt,basicstyle=\ttfamily\scriptsize]
ROMInternalTranscriptPaperSimAsNMA
ROSigningAttemptPaperSimSampler
RealSigningPaperSimSampler
glob
haetae_mode
pk_current
public_key_of_secret
queries
real_signing_ro_attempt_paper_sim_nma_oracle
records
sk_current
transcripts
\end{lstlisting}

\end{formaltheorem}

\begin{formaltheorem}[EasyCrypt line 6548]
\noindent\textbf{EasyCrypt name.}
\begin{lstlisting}[language=EasyCrypt,basicstyle=\ttfamily\scriptsize]
rom_internal_nma_real_signing_ro_attempt_paper_sim_exact
\end{lstlisting}
\noindent\textbf{Checked EasyCrypt statement.}
\begin{lstlisting}[language=EasyCrypt,basicstyle=\ttfamily\scriptsize]
lemma rom_internal_nma_real_signing_ro_attempt_paper_sim_exact &m :
  Pr[SIG.UF_NMA(H, HAETAE,
       ROMInternalTranscriptPaperSimAsNMA
         (A, RealSigningPaperSimSampler(H))).main() @ &m : res] =
  Pr[SIG.UF_NMA(H, HAETAE,
       ROMInternalTranscriptPaperSimAsNMA
         (A, ROSigningAttemptPaperSimSampler(H))).main() @ &m : res].
\end{lstlisting}
\noindent\textbf{Formal mathematical reading.}
Mathematical theorem. This theorem has the form \(\Pr[G:E] = B\) is a game-probability statement.  It is one of the exact hops, bad-event bounds, or final advantage inequalities used in the reduction.

\noindent\textbf{Role in the proof.}
Use in this chapter. It proves an exact game replacement or simulator equivalence.

\noindent\textbf{Proof-script interpretation.}
Proof-script reading. The machine proof decomposes the theorem into rewrites, program-logic obligations, relational equivalences, and arithmetic side conditions.
\emph{Main tactics visible in the proof script:} \ec{byequiv}, \ec{proc}, \ec{inline}, \ec{wp}, \ec{call}, \ec{apply}, \ec{rnd}, \ec{rewrite}.
\emph{Named identifiers syntactically cited in the proof block:}
\begin{lstlisting}[language=EasyCrypt,basicstyle=\ttfamily\scriptsize]
HAETAE
ROMInternalTranscriptPaperSimAsNMA
ROSigningAttemptPaperSimSampler
RealSigningPaperSimSampler
adversary_real_signing_ro_attempt_paper_sim_nma
arg
forge
glob
haetae_mode
init
keygen_internal
kg
pk_current
public_key_of_secret
queries
records
secret_key_of_seed
sk_current
transcripts
verify
\end{lstlisting}

\end{formaltheorem}

\begin{formaltheorem}[EasyCrypt line 6625]
\noindent\textbf{EasyCrypt name.}
\begin{lstlisting}[language=EasyCrypt,basicstyle=\ttfamily\scriptsize]
concrete_ro_signing_attempt_to_exact_hyperball_sample_with_seed_fresh_loss
\end{lstlisting}
\noindent\textbf{Checked EasyCrypt statement.}
\begin{lstlisting}[language=EasyCrypt,basicstyle=\ttfamily\scriptsize]
lemma concrete_ro_signing_attempt_to_exact_hyperball_sample_with_seed_fresh_loss
    seed_coins m ctx (p : paper_sim_signature_sample -> bool) &m :
  structural_to_exact_hyperball_paper_sample_loss_obligation haetae_mode =>
  sampler_expand_query seed_coins \notin HAETAE_RO.FRO.m{m} =>
  ROSigningAttemptPaperSimSampler.sk_current{m} =
    ROExactHyperballPaperSimSampler.sk_current{m} =>
  Pr[ROSigningAttemptPaperSimSampler(HAETAE_RO.FRO).sample_with_seed
       (seed_coins, m, ctx) @ &m : p res] <=
  Pr[ROExactHyperballPaperSimSampler(HAETAE_RO.FRO).sample_with_seed
       (seed_coins, m, ctx) @ &m : p res] +
    rejection_sampling_loss_term.
\end{lstlisting}
\noindent\textbf{Formal mathematical reading.}
Mathematical theorem. This theorem has the form \(\Pr[G:E] \le B\) is a game-probability statement.  It is one of the exact hops, bad-event bounds, or final advantage inequalities used in the reduction.

\noindent\textbf{Role in the proof.}
Use in this chapter. It contributes directly to an advantage bound, bad-event bound, or real-arithmetic composition step.

\noindent\textbf{Proof-script interpretation.}
Proof-script reading. The machine proof decomposes the theorem into rewrites, program-logic obligations, relational equivalences, and arithmetic side conditions.
\emph{Main tactics visible in the proof script:} \ec{have}, \ec{byphoare}, \ec{smt}.
\emph{Named identifiers syntactically cited in the proof block:}
\begin{lstlisting}[language=EasyCrypt,basicstyle=\ttfamily\scriptsize]
FRO
HAETAE_RO
Pr
ROSigningAttemptPaperSimSampler
concrete_ro_exact_hyperball_sample_with_seed_exact_pr
concrete_ro_signing_attempt_sample_with_seed_fresh_structural_le
ctx
dsigning_attempt_state
exact_eq
fresh
haetae_mode
mu
paper_sim_sample_from_rejection_attempt
sample_loss
sample_step
sample_with_seed
seed_coins
signing_le
sk_current
sk_eq
st
\end{lstlisting}

\end{formaltheorem}

\begin{formaltheorem}[EasyCrypt line 6655]
\noindent\textbf{EasyCrypt name.}
\begin{lstlisting}[language=EasyCrypt,basicstyle=\ttfamily\scriptsize]
concrete_ro_signing_attempt_to_exact_hyperball_sample_with_seed_clean_loss
\end{lstlisting}
\noindent\textbf{Checked EasyCrypt statement.}
\begin{lstlisting}[language=EasyCrypt,basicstyle=\ttfamily\scriptsize]
lemma concrete_ro_signing_attempt_to_exact_hyperball_sample_with_seed_clean_loss
    seed_coins m ctx (p : paper_sim_signature_sample -> bool)
    hash_qs sampler_qs bad &m :
  structural_to_exact_hyperball_paper_sample_loss_obligation haetae_mode =>
  sampler_rom_covered HAETAE_RO.FRO.m{m} hash_qs sampler_qs =>
  ! (bad \/
     sampler_expand_query seed_coins \in hash_qs \/
     sampler_expand_query seed_coins \in sampler_qs) =>
  ROSigningAttemptPaperSimSampler.sk_current{m} =
    ROExactHyperballPaperSimSampler.sk_current{m} =>
  Pr[ROSigningAttemptPaperSimSampler(HAETAE_RO.FRO).sample_with_seed
       (seed_coins, m, ctx) @ &m : p res] <=
  Pr[ROExactHyperballPaperSimSampler(HAETAE_RO.FRO).sample_with_seed
       (seed_coins, m, ctx) @ &m : p res] +
    rejection_sampling_loss_term.
\end{lstlisting}
\noindent\textbf{Formal mathematical reading.}
Mathematical theorem. This theorem has the form \(\Pr[G:E] \le B\) is a game-probability statement.  It is one of the exact hops, bad-event bounds, or final advantage inequalities used in the reduction.

\noindent\textbf{Role in the proof.}
Use in this chapter. It contributes directly to an advantage bound, bad-event bound, or real-arithmetic composition step.

\noindent\textbf{Proof-script interpretation.}
Proof-script reading. The machine proof decomposes the theorem into rewrites, program-logic obligations, relational equivalences, and arithmetic side conditions.
\emph{Main tactics visible in the proof script:} \ec{apply}.
\emph{Named identifiers syntactically cited in the proof block:}
\begin{lstlisting}[language=EasyCrypt,basicstyle=\ttfamily\scriptsize]
FRO
HAETAE_RO
bad
clean
concrete_ro_signing_attempt_to_exact_hyperball_sample_with_seed_fresh_loss
covered
ctx
hash_qs
sample_loss
sampler_qs
sampler_rom_covered_fresh_after_clean_seed
seed_coins
sk_eq
\end{lstlisting}

\end{formaltheorem}

\begin{formaltheorem}[EasyCrypt line 6678]
\noindent\textbf{EasyCrypt name.}
\begin{lstlisting}[language=EasyCrypt,basicstyle=\ttfamily\scriptsize]
concrete_budgeted_o_sign_sample_with_seed_clean_loss_surface
\end{lstlisting}
\noindent\textbf{Checked EasyCrypt statement.}
\begin{lstlisting}[language=EasyCrypt,basicstyle=\ttfamily\scriptsize]
lemma concrete_budgeted_o_sign_sample_with_seed_clean_loss_surface
    seed_coins m ctx (p : paper_sim_signature_sample -> bool) &m :
  structural_to_exact_hyperball_paper_sample_loss_obligation haetae_mode =>
  sampler_rom_covered
    HAETAE_RO.FRO.m{m}
    ROMInternalTranscriptBudgetedPaperSimAsNMA.adversary_hash_queries{m}
    ROMInternalTranscriptBudgetedPaperSimAsNMA.sampler_expand_queries{m} =>
  ! (ROMInternalTranscriptBudgetedPaperSimAsNMA.sampler_bad_prequery{m} \/
     sampler_expand_query seed_coins \in
       ROMInternalTranscriptBudgetedPaperSimAsNMA.adversary_hash_queries{m} \/
     sampler_expand_query seed_coins \in
       ROMInternalTranscriptBudgetedPaperSimAsNMA.sampler_expand_queries{m}) =>
  ROSigningAttemptPaperSimSampler.sk_current{m} =
    ROExactHyperballPaperSimSampler.sk_current{m} =>
  Pr[ROSigningAttemptPaperSimSampler(HAETAE_RO.FRO).sample_with_seed
       (seed_coins, m, ctx) @ &m : p res] <=
  Pr[ROExactHyperballPaperSimSampler(HAETAE_RO.FRO).sample_with_seed
       (seed_coins, m, ctx) @ &m : p res] +
    rejection_sampling_loss_term.
\end{lstlisting}
\noindent\textbf{Formal mathematical reading.}
Mathematical theorem. This theorem has the form \(\Pr[G:E] \le B\) is a game-probability statement.  It is one of the exact hops, bad-event bounds, or final advantage inequalities used in the reduction.

\noindent\textbf{Role in the proof.}
Use in this chapter. It contributes directly to an advantage bound, bad-event bound, or real-arithmetic composition step.

\noindent\textbf{Proof-script interpretation.}
Proof-script reading. The machine proof decomposes the theorem into rewrites, program-logic obligations, relational equivalences, and arithmetic side conditions.
\emph{Main tactics visible in the proof script:} \ec{apply}.
\emph{Named identifiers syntactically cited in the proof block:}
\begin{lstlisting}[language=EasyCrypt,basicstyle=\ttfamily\scriptsize]
ROMInternalTranscriptBudgetedPaperSimAsNMA
adversary_hash_queries
clean
concrete_ro_signing_attempt_to_exact_hyperball_sample_with_seed_clean_loss
covered
ctx
sample_loss
sampler_bad_prequery
sampler_expand_queries
seed_coins
sk_eq
\end{lstlisting}

\end{formaltheorem}

\begin{formaltheorem}[EasyCrypt line 6708]
\noindent\textbf{EasyCrypt name.}
\begin{lstlisting}[language=EasyCrypt,basicstyle=\ttfamily\scriptsize]
concrete_budgeted_o_sign_old_log_sample_with_seed_clean_loss_surface
\end{lstlisting}
\noindent\textbf{Checked EasyCrypt statement.}
\begin{lstlisting}[language=EasyCrypt,basicstyle=\ttfamily\scriptsize]
lemma concrete_budgeted_o_sign_old_log_sample_with_seed_clean_loss_surface
    seed_coins m ctx (p : paper_sim_signature_sample -> bool)
    old_hash_qs old_sampler_qs old_bad &m :
  structural_to_exact_hyperball_paper_sample_loss_obligation haetae_mode =>
  old_hash_qs =
    ROMInternalTranscriptBudgetedPaperSimAsNMA.adversary_hash_queries{m} =>
  old_sampler_qs =
    ROMInternalTranscriptBudgetedPaperSimAsNMA.sampler_expand_queries{m} =>
  old_bad =
    ROMInternalTranscriptBudgetedPaperSimAsNMA.sampler_bad_prequery{m} =>
  sampler_rom_covered HAETAE_RO.FRO.m{m} old_hash_qs old_sampler_qs =>
  ! (old_bad \/
     sampler_expand_query seed_coins \in old_hash_qs \/
     sampler_expand_query seed_coins \in old_sampler_qs) =>
  ROSigningAttemptPaperSimSampler.sk_current{m} =
    ROExactHyperballPaperSimSampler.sk_current{m} =>
  Pr[ROSigningAttemptPaperSimSampler(HAETAE_RO.FRO).sample_with_seed
       (seed_coins, m, ctx) @ &m : p res] <=
  Pr[ROExactHyperballPaperSimSampler(HAETAE_RO.FRO).sample_with_seed
       (seed_coins, m, ctx) @ &m : p res] +
    rejection_sampling_loss_term.
\end{lstlisting}
\noindent\textbf{Formal mathematical reading.}
Mathematical theorem. This theorem has the form \(\Pr[G:E] \le B\) is a game-probability statement.  It is one of the exact hops, bad-event bounds, or final advantage inequalities used in the reduction.

\noindent\textbf{Role in the proof.}
Use in this chapter. It contributes directly to an advantage bound, bad-event bound, or real-arithmetic composition step.

\noindent\textbf{Proof-script interpretation.}
Proof-script reading. The machine proof decomposes the theorem into rewrites, program-logic obligations, relational equivalences, and arithmetic side conditions.
\emph{Main tactics visible in the proof script:} \ec{apply}.
\emph{Named identifiers syntactically cited in the proof block:}
\begin{lstlisting}[language=EasyCrypt,basicstyle=\ttfamily\scriptsize]
clean
concrete_budgeted_o_sign_sample_with_seed_clean_loss_surface
covered
ctx
old_bad
old_bad_eq
old_hash_eq
old_hash_qs
old_sampler_eq
old_sampler_qs
sample_loss
seed_coins
sk_eq
subst
\end{lstlisting}

\end{formaltheorem}

\begin{formaltheorem}[EasyCrypt line 6740]
\noindent\textbf{EasyCrypt name.}
\begin{lstlisting}[language=EasyCrypt,basicstyle=\ttfamily\scriptsize]
concrete_budgeted_o_sign_old_log_signature_from_sample_clean_loss_surface
\end{lstlisting}
\noindent\textbf{Checked EasyCrypt statement.}
\begin{lstlisting}[language=EasyCrypt,basicstyle=\ttfamily\scriptsize]
lemma concrete_budgeted_o_sign_old_log_signature_from_sample_clean_loss_surface
    seed_coins m ctx (p : signature -> bool)
    old_hash_qs old_sampler_qs old_bad &m :
  structural_to_exact_hyperball_paper_sample_loss_obligation haetae_mode =>
  old_hash_qs =
    ROMInternalTranscriptBudgetedPaperSimAsNMA.adversary_hash_queries{m} =>
  old_sampler_qs =
    ROMInternalTranscriptBudgetedPaperSimAsNMA.sampler_expand_queries{m} =>
  old_bad =
    ROMInternalTranscriptBudgetedPaperSimAsNMA.sampler_bad_prequery{m} =>
  sampler_rom_covered HAETAE_RO.FRO.m{m} old_hash_qs old_sampler_qs =>
  ! (old_bad \/
     sampler_expand_query seed_coins \in old_hash_qs \/
     sampler_expand_query seed_coins \in old_sampler_qs) =>
  ROSigningAttemptPaperSimSampler.sk_current{m} =
    ROExactHyperballPaperSimSampler.sk_current{m} =>
  Pr[ROSigningAttemptPaperSimSampler(HAETAE_RO.FRO).sample_with_seed
       (seed_coins, m, ctx) @ &m :
       p (paper_sim_signature haetae_mode
            ROMInternalTranscriptBudgetedPaperSimAsNMA.pk_current{m}
            m ctx res)] <=
  Pr[ROExactHyperballPaperSimSampler(HAETAE_RO.FRO).sample_with_seed
       (seed_coins, m, ctx) @ &m :
       p (paper_sim_signature haetae_mode
            ROMInternalTranscriptBudgetedPaperSimAsNMA.pk_current{m}
            m ctx res)] +
    rejection_sampling_loss_term.
\end{lstlisting}
\noindent\textbf{Formal mathematical reading.}
Mathematical theorem. This theorem has the form \(\Pr[G:E] \le B\) is a game-probability statement.  It is one of the exact hops, bad-event bounds, or final advantage inequalities used in the reduction.

\noindent\textbf{Role in the proof.}
Use in this chapter. It contributes directly to an advantage bound, bad-event bound, or real-arithmetic composition step.

\noindent\textbf{Proof-script interpretation.}
Proof-script reading. The machine proof decomposes the theorem into rewrites, program-logic obligations, relational equivalences, and arithmetic side conditions.
\emph{Main tactics visible in the proof script:} \ec{apply}.
\emph{Named identifiers syntactically cited in the proof block:}
\begin{lstlisting}[language=EasyCrypt,basicstyle=\ttfamily\scriptsize]
ROMInternalTranscriptBudgetedPaperSimAsNMA
clean
concrete_budgeted_o_sign_old_log_sample_with_seed_clean_loss_surface
covered
ctx
haetae_mode
old_bad
old_bad_eq
old_hash_eq
old_hash_qs
old_sampler_eq
old_sampler_qs
paper_sim_signature
pk_current
sample_loss
seed_coins
sk_eq
smp
\end{lstlisting}

\end{formaltheorem}

\begin{formaltheorem}[EasyCrypt line 6779]
\noindent\textbf{EasyCrypt name.}
\begin{lstlisting}[language=EasyCrypt,basicstyle=\ttfamily\scriptsize]
concrete_budgeted_o_sign_old_log_signature_clean_event_loss_surface
\end{lstlisting}
\noindent\textbf{Checked EasyCrypt statement.}
\begin{lstlisting}[language=EasyCrypt,basicstyle=\ttfamily\scriptsize]
lemma concrete_budgeted_o_sign_old_log_signature_clean_event_loss_surface
    seed_coins m ctx (p : signature -> bool)
    old_hash_qs old_sampler_qs old_bad &m :
  structural_to_exact_hyperball_paper_sample_loss_obligation haetae_mode =>
  old_hash_qs =
    ROMInternalTranscriptBudgetedPaperSimAsNMA.adversary_hash_queries{m} =>
  old_sampler_qs =
    ROMInternalTranscriptBudgetedPaperSimAsNMA.sampler_expand_queries{m} =>
  old_bad =
    ROMInternalTranscriptBudgetedPaperSimAsNMA.sampler_bad_prequery{m} =>
  sampler_rom_covered HAETAE_RO.FRO.m{m} old_hash_qs old_sampler_qs =>
  ! (old_bad \/
     sampler_expand_query seed_coins \in old_hash_qs \/
     sampler_expand_query seed_coins \in old_sampler_qs) =>
  ROSigningAttemptPaperSimSampler.sk_current{m} =
    ROExactHyperballPaperSimSampler.sk_current{m} =>
  Pr[ROSigningAttemptPaperSimSampler(HAETAE_RO.FRO).sample_with_seed
       (seed_coins, m, ctx) @ &m :
       p (paper_sim_signature haetae_mode
            ROMInternalTranscriptBudgetedPaperSimAsNMA.pk_current{m}
            m ctx res) /\
       ! (old_bad \/
          sampler_expand_query seed_coins \in old_hash_qs \/
          sampler_expand_query seed_coins \in old_sampler_qs)] <=
  Pr[ROExactHyperballPaperSimSampler(HAETAE_RO.FRO).sample_with_seed
       (seed_coins, m, ctx) @ &m :
       p (paper_sim_signature haetae_mode
            ROMInternalTranscriptBudgetedPaperSimAsNMA.pk_current{m}
            m ctx res)] +
    rejection_sampling_loss_term.
\end{lstlisting}
\noindent\textbf{Formal mathematical reading.}
Mathematical theorem. This theorem has the form \(\Pr[G:E] \le B\) is a game-probability statement.  It is one of the exact hops, bad-event bounds, or final advantage inequalities used in the reduction.

\noindent\textbf{Role in the proof.}
Use in this chapter. It contributes directly to an advantage bound, bad-event bound, or real-arithmetic composition step.

\noindent\textbf{Proof-script interpretation.}
Proof-script reading. The machine proof decomposes the theorem into rewrites, program-logic obligations, relational equivalences, and arithmetic side conditions.
\emph{Main tactics visible in the proof script:} \ec{rewrite}, \ec{apply}.
\emph{Named identifiers syntactically cited in the proof block:}
\begin{lstlisting}[language=EasyCrypt,basicstyle=\ttfamily\scriptsize]
clean
concrete_budgeted_o_sign_old_log_signature_from_sample_clean_loss_surface
covered
ctx
old_bad
old_bad_eq
old_hash_eq
old_hash_qs
old_sampler_eq
old_sampler_qs
sample_loss
seed_coins
sk_eq
\end{lstlisting}

\end{formaltheorem}

\begin{formaltheorem}[EasyCrypt line 6818]
\noindent\textbf{EasyCrypt name.}
\begin{lstlisting}[language=EasyCrypt,basicstyle=\ttfamily\scriptsize]
concrete_frozen_old_log_sample_clean_event_loss_surface
\end{lstlisting}
\noindent\textbf{Checked EasyCrypt statement.}
\begin{lstlisting}[language=EasyCrypt,basicstyle=\ttfamily\scriptsize]
lemma concrete_frozen_old_log_sample_clean_event_loss_surface
    seed_coins m ctx (p : paper_sim_signature_sample -> bool)
    old_hash_qs old_sampler_qs old_bad &m :
  structural_to_exact_hyperball_paper_sample_loss_obligation haetae_mode =>
  sampler_rom_covered HAETAE_RO.FRO.m{m} old_hash_qs old_sampler_qs =>
  ! (old_bad \/
     sampler_expand_query seed_coins \in old_hash_qs \/
     sampler_expand_query seed_coins \in old_sampler_qs) =>
  ROSigningAttemptPaperSimSampler.sk_current{m} =
    ROExactHyperballPaperSimSampler.sk_current{m} =>
  Pr[ROSigningAttemptPaperSimSampler(HAETAE_RO.FRO).sample_with_seed
       (seed_coins, m, ctx) @ &m :
       p res /\
       ! (old_bad \/
          sampler_expand_query seed_coins \in old_hash_qs \/
          sampler_expand_query seed_coins \in old_sampler_qs)] <=
  Pr[ROExactHyperballPaperSimSampler(HAETAE_RO.FRO).sample_with_seed
       (seed_coins, m, ctx) @ &m : p res] +
    rejection_sampling_loss_term.
\end{lstlisting}
\noindent\textbf{Formal mathematical reading.}
Mathematical theorem. This theorem has the form \(\Pr[G:E] \le B\) is a game-probability statement.  It is one of the exact hops, bad-event bounds, or final advantage inequalities used in the reduction.

\noindent\textbf{Role in the proof.}
Use in this chapter. It contributes directly to an advantage bound, bad-event bound, or real-arithmetic composition step.

\noindent\textbf{Proof-script interpretation.}
Proof-script reading. The machine proof decomposes the theorem into rewrites, program-logic obligations, relational equivalences, and arithmetic side conditions.
\emph{Main tactics visible in the proof script:} \ec{rewrite}, \ec{apply}.
\emph{Named identifiers syntactically cited in the proof block:}
\begin{lstlisting}[language=EasyCrypt,basicstyle=\ttfamily\scriptsize]
clean
concrete_ro_signing_attempt_to_exact_hyperball_sample_with_seed_clean_loss
covered
ctx
old_bad
old_hash_qs
old_sampler_qs
sample_loss
seed_coins
sk_eq
\end{lstlisting}

\end{formaltheorem}

\begin{formaltheorem}[EasyCrypt line 6845]
\noindent\textbf{EasyCrypt name.}
\begin{lstlisting}[language=EasyCrypt,basicstyle=\ttfamily\scriptsize]
concrete_frozen_old_log_signature_clean_event_loss_surface
\end{lstlisting}
\noindent\textbf{Checked EasyCrypt statement.}
\begin{lstlisting}[language=EasyCrypt,basicstyle=\ttfamily\scriptsize]
lemma concrete_frozen_old_log_signature_clean_event_loss_surface
    seed_coins m ctx pk (p : signature -> bool)
    old_hash_qs old_sampler_qs old_bad &m :
  structural_to_exact_hyperball_paper_sample_loss_obligation haetae_mode =>
  sampler_rom_covered HAETAE_RO.FRO.m{m} old_hash_qs old_sampler_qs =>
  ! (old_bad \/
     sampler_expand_query seed_coins \in old_hash_qs \/
     sampler_expand_query seed_coins \in old_sampler_qs) =>
  ROSigningAttemptPaperSimSampler.sk_current{m} =
    ROExactHyperballPaperSimSampler.sk_current{m} =>
  Pr[ROSigningAttemptPaperSimSampler(HAETAE_RO.FRO).sample_with_seed
       (seed_coins, m, ctx) @ &m :
       p (paper_sim_signature haetae_mode pk m ctx res) /\
       ! (old_bad \/
          sampler_expand_query seed_coins \in old_hash_qs \/
          sampler_expand_query seed_coins \in old_sampler_qs)] <=
  Pr[ROExactHyperballPaperSimSampler(HAETAE_RO.FRO).sample_with_seed
       (seed_coins, m, ctx) @ &m :
       p (paper_sim_signature haetae_mode pk m ctx res)] +
    rejection_sampling_loss_term.
\end{lstlisting}
\noindent\textbf{Formal mathematical reading.}
Mathematical theorem. This theorem has the form \(\Pr[G:E] \le B\) is a game-probability statement.  It is one of the exact hops, bad-event bounds, or final advantage inequalities used in the reduction.

\noindent\textbf{Role in the proof.}
Use in this chapter. It contributes directly to an advantage bound, bad-event bound, or real-arithmetic composition step.

\noindent\textbf{Proof-script interpretation.}
Proof-script reading. The machine proof decomposes the theorem into rewrites, program-logic obligations, relational equivalences, and arithmetic side conditions.
\emph{Main tactics visible in the proof script:} \ec{rewrite}, \ec{apply}.
\emph{Named identifiers syntactically cited in the proof block:}
\begin{lstlisting}[language=EasyCrypt,basicstyle=\ttfamily\scriptsize]
clean
concrete_ro_signing_attempt_to_exact_hyperball_sample_with_seed_clean_loss
covered
ctx
haetae_mode
old_bad
old_hash_qs
old_sampler_qs
paper_sim_signature
pk
sample_loss
seed_coins
sk_eq
smp
\end{lstlisting}

\end{formaltheorem}

\begin{formaltheorem}[EasyCrypt line 6875]
\noindent\textbf{EasyCrypt name.}
\begin{lstlisting}[language=EasyCrypt,basicstyle=\ttfamily\scriptsize]
concrete_frozen_old_log_self_logged_signature_clean_event_loss_surface
\end{lstlisting}
\noindent\textbf{Checked EasyCrypt statement.}
\begin{lstlisting}[language=EasyCrypt,basicstyle=\ttfamily\scriptsize]
lemma concrete_frozen_old_log_self_logged_signature_clean_event_loss_surface
    seed_coins m ctx pk (p : signature -> bool)
    old_hash_qs old_sampler_qs old_bad &m :
  structural_to_exact_hyperball_paper_sample_loss_obligation haetae_mode =>
  sampler_rom_covered HAETAE_RO.FRO.m{m} old_hash_qs old_sampler_qs =>
  ! (old_bad \/
     sampler_expand_query seed_coins \in old_hash_qs \/
     sampler_expand_query seed_coins \in old_sampler_qs) =>
  ROSigningAttemptPaperSimSampler.sk_current{m} =
    ROExactHyperballPaperSimSampler.sk_current{m} =>
  sampler_expand_query seed_coins \notin HAETAE_RO.FRO.m{m} /\
  sampler_rom_covered HAETAE_RO.FRO.m{m} old_hash_qs
    (sampler_expand_query seed_coins :: old_sampler_qs) /\
  Pr[ROSigningAttemptPaperSimSampler(HAETAE_RO.FRO).sample_with_seed
       (seed_coins, m, ctx) @ &m :
       p (paper_sim_signature haetae_mode pk m ctx res) /\
       ! (old_bad \/
          sampler_expand_query seed_coins \in old_hash_qs \/
          sampler_expand_query seed_coins \in old_sampler_qs)] <=
  Pr[ROExactHyperballPaperSimSampler(HAETAE_RO.FRO).sample_with_seed
       (seed_coins, m, ctx) @ &m :
       p (paper_sim_signature haetae_mode pk m ctx res)] +
    rejection_sampling_loss_term.
\end{lstlisting}
\noindent\textbf{Formal mathematical reading.}
Mathematical theorem. This theorem has the form \(\Pr[G:E] \le B\) is a game-probability statement.  It is one of the exact hops, bad-event bounds, or final advantage inequalities used in the reduction.

\noindent\textbf{Role in the proof.}
Use in this chapter. It contributes directly to an advantage bound, bad-event bound, or real-arithmetic composition step.

\noindent\textbf{Proof-script interpretation.}
Proof-script reading. The machine proof decomposes the theorem into rewrites, program-logic obligations, relational equivalences, and arithmetic side conditions.
\emph{Main tactics visible in the proof script:} \ec{have}, \ec{split}, \ec{smt}, \ec{apply}.
\emph{Named identifiers syntactically cited in the proof block:}
\begin{lstlisting}[language=EasyCrypt,basicstyle=\ttfamily\scriptsize]
FRO
HAETAE_RO
boundary
clean
concrete_frozen_old_log_signature_clean_event_loss_surface
covered
ctx
old_bad
old_hash_qs
old_sampler_qs
pk
sample_loss
sampler_rom_covered_old_log_clean_boundary
seed_coins
sk_eq
\end{lstlisting}

\end{formaltheorem}

\begin{formaltheorem}[EasyCrypt line 6913]
\noindent\textbf{EasyCrypt name.}
\begin{lstlisting}[language=EasyCrypt,basicstyle=\ttfamily\scriptsize]
concrete_budgeted_self_logged_signature_clean_event_loss_surface
\end{lstlisting}
\noindent\textbf{Checked EasyCrypt statement.}
\begin{lstlisting}[language=EasyCrypt,basicstyle=\ttfamily\scriptsize]
lemma concrete_budgeted_self_logged_signature_clean_event_loss_surface
    seed_coins m ctx pk (p : signature -> bool) &m :
  structural_to_exact_hyperball_paper_sample_loss_obligation haetae_mode =>
  sampler_rom_covered
    HAETAE_RO.FRO.m{m}
    ROMInternalTranscriptBudgetedPaperSimAsNMA.adversary_hash_queries{m}
    ROMInternalTranscriptBudgetedPaperSimAsNMA.sampler_expand_queries{m} =>
  ! (ROMInternalTranscriptBudgetedPaperSimAsNMA.sampler_bad_prequery{m} \/
     sampler_expand_query seed_coins \in
       ROMInternalTranscriptBudgetedPaperSimAsNMA.adversary_hash_queries{m} \/
     sampler_expand_query seed_coins \in
       ROMInternalTranscriptBudgetedPaperSimAsNMA.sampler_expand_queries{m}) =>
  ROSigningAttemptPaperSimSampler.sk_current{m} =
    ROExactHyperballPaperSimSampler.sk_current{m} =>
  sampler_expand_query seed_coins \notin HAETAE_RO.FRO.m{m} /\
  sampler_rom_covered
    HAETAE_RO.FRO.m{m}
    ROMInternalTranscriptBudgetedPaperSimAsNMA.adversary_hash_queries{m}
    (sampler_expand_query seed_coins ::
       ROMInternalTranscriptBudgetedPaperSimAsNMA.sampler_expand_queries{m}) /\
  Pr[ROSigningAttemptPaperSimSampler(HAETAE_RO.FRO).sample_with_seed
       (seed_coins, m, ctx) @ &m :
       p (paper_sim_signature haetae_mode pk m ctx res) /\
       ! (ROMInternalTranscriptBudgetedPaperSimAsNMA.sampler_bad_prequery{m} \/
          sampler_expand_query seed_coins \in
            ROMInternalTranscriptBudgetedPaperSimAsNMA.adversary_hash_queries{m} \/
          sampler_expand_query seed_coins \in
            ROMInternalTranscriptBudgetedPaperSimAsNMA.sampler_expand_queries{m})] <=
  Pr[ROExactHyperballPaperSimSampler(HAETAE_RO.FRO).sample_with_seed
       (seed_coins, m, ctx) @ &m :
       p (paper_sim_signature haetae_mode pk m ctx res)] +
    rejection_sampling_loss_term.
\end{lstlisting}
\noindent\textbf{Formal mathematical reading.}
Mathematical theorem. This theorem has the form \(\Pr[G:E] \le B\) is a game-probability statement.  It is one of the exact hops, bad-event bounds, or final advantage inequalities used in the reduction.

\noindent\textbf{Role in the proof.}
Use in this chapter. It contributes directly to an advantage bound, bad-event bound, or real-arithmetic composition step.

\noindent\textbf{Proof-script interpretation.}
Proof-script reading. The machine proof decomposes the theorem into rewrites, program-logic obligations, relational equivalences, and arithmetic side conditions.
\emph{Main tactics visible in the proof script:} \ec{apply}.
\emph{Named identifiers syntactically cited in the proof block:}
\begin{lstlisting}[language=EasyCrypt,basicstyle=\ttfamily\scriptsize]
ROMInternalTranscriptBudgetedPaperSimAsNMA
adversary_hash_queries
clean
concrete_frozen_old_log_self_logged_signature_clean_event_loss_surface
covered
ctx
pk
sample_loss
sampler_bad_prequery
sampler_expand_queries
seed_coins
sk_eq
\end{lstlisting}

\end{formaltheorem}

\begin{formaltheorem}[EasyCrypt line 6956]
\noindent\textbf{EasyCrypt name.}
\begin{lstlisting}[language=EasyCrypt,basicstyle=\ttfamily\scriptsize]
concrete_lazy_rom_get_lossless
\end{lstlisting}
\noindent\textbf{Checked EasyCrypt statement.}
\begin{lstlisting}[language=EasyCrypt,basicstyle=\ttfamily\scriptsize]
lemma concrete_lazy_rom_get_lossless :
  islossless HAETAE_RO.FRO.get.
\end{lstlisting}
\noindent\textbf{Formal mathematical reading.}
Mathematical theorem. This lemma is a universally quantified proposition over the variables shown in the EasyCrypt statement.

\noindent\textbf{Role in the proof.}
Use in this chapter. It contributes directly to an advantage bound, bad-event bound, or real-arithmetic composition step.

\noindent\textbf{Proof-script interpretation.}
Proof-script reading. The machine proof decomposes the theorem into rewrites, program-logic obligations, relational equivalences, and arithmetic side conditions.
\emph{Main tactics visible in the proof script:} \ec{proc}, \ec{wp}, \ec{rnd}, \ec{smt}.
\emph{Named identifiers syntactically cited in the proof block:}
\begin{lstlisting}[language=EasyCrypt,basicstyle=\ttfamily\scriptsize]
ro_output_distribution_lossless
\end{lstlisting}

\end{formaltheorem}

\begin{formaltheorem}[EasyCrypt line 6965]
\noindent\textbf{EasyCrypt name.}
\begin{lstlisting}[language=EasyCrypt,basicstyle=\ttfamily\scriptsize]
structural_attempt_exact_hyperball_paper_sample_one_step_loss
\end{lstlisting}
\noindent\textbf{Checked EasyCrypt statement.}
\begin{lstlisting}[language=EasyCrypt,basicstyle=\ttfamily\scriptsize]
lemma structural_attempt_exact_hyperball_paper_sample_one_step_loss
  sk m ctx (p : paper_sim_signature_sample -> bool) :
  structural_to_exact_hyperball_paper_sample_loss_obligation haetae_mode =>
  phoare[StructuralAttemptPaperSample.sample :
    arg = (sk, m, ctx) ==> p res] <=
  (mu (dexact_hyperball_signing_attempt_state haetae_mode sk m ctx)
     (fun st => p (paper_sim_sample_from_rejection_attempt st)) +
     rejection_sampling_loss_term).
\end{lstlisting}
\noindent\textbf{Formal mathematical reading.}
Mathematical theorem. This is a relational program theorem: two games or procedures are coupled so that a precondition on their initial states implies the stated postcondition on final states and return values.

\noindent\textbf{Role in the proof.}
Use in this chapter. It contributes directly to an advantage bound, bad-event bound, or real-arithmetic composition step.

\noindent\textbf{Proof-script interpretation.}
Proof-script reading. The machine proof decomposes the theorem into rewrites, program-logic obligations, relational equivalences, and arithmetic side conditions.
\emph{Main tactics visible in the proof script:} \ec{byphoare}, \ec{proc}, \ec{wp}, \ec{rnd}, \ec{apply}, \ec{rewrite}.
\emph{Named identifiers syntactically cited in the proof block:}
\begin{lstlisting}[language=EasyCrypt,basicstyle=\ttfamily\scriptsize]
arg
argE
bypr
ctx
m0
sample_loss
sk
\end{lstlisting}

\end{formaltheorem}

\begin{formaltheorem}[EasyCrypt line 6986]
\noindent\textbf{EasyCrypt name.}
\begin{lstlisting}[language=EasyCrypt,basicstyle=\ttfamily\scriptsize]
exact_hyperball_paper_sample_one_step_upper
\end{lstlisting}
\noindent\textbf{Checked EasyCrypt statement.}
\begin{lstlisting}[language=EasyCrypt,basicstyle=\ttfamily\scriptsize]
lemma exact_hyperball_paper_sample_one_step_upper
  sk m ctx (p : paper_sim_signature_sample -> bool) :
  phoare[ExactHyperballPaperSample.sample :
    arg = (sk, m, ctx) ==> p res] <=
  (mu (dexact_hyperball_signing_attempt_state haetae_mode sk m ctx)
     (fun st => p (paper_sim_sample_from_rejection_attempt st))).
\end{lstlisting}
\noindent\textbf{Formal mathematical reading.}
Mathematical theorem. This is a relational program theorem: two games or procedures are coupled so that a precondition on their initial states implies the stated postcondition on final states and return values.

\noindent\textbf{Role in the proof.}
Use in this chapter. It contributes directly to an advantage bound, bad-event bound, or real-arithmetic composition step.

\noindent\textbf{Proof-script interpretation.}
Proof-script reading. The machine proof decomposes the theorem into rewrites, program-logic obligations, relational equivalences, and arithmetic side conditions.
\emph{Main tactics visible in the proof script:} \ec{byphoare}, \ec{proc}, \ec{wp}, \ec{rnd}, \ec{rewrite}.
\emph{Named identifiers syntactically cited in the proof block:}
\begin{lstlisting}[language=EasyCrypt,basicstyle=\ttfamily\scriptsize]
arg
argE
bypr
ctx
m0
sk
\end{lstlisting}

\end{formaltheorem}

\begin{formaltheorem}[EasyCrypt line 7016]
\noindent\textbf{EasyCrypt name.}
\begin{lstlisting}[language=EasyCrypt,basicstyle=\ttfamily\scriptsize]
rom_internal_nma_ro_signing_attempt_ro_exact_hyperball_loss_from_paper_sample_lifting
\end{lstlisting}
\noindent\textbf{Checked EasyCrypt statement.}
\begin{lstlisting}[language=EasyCrypt,basicstyle=\ttfamily\scriptsize]
lemma rom_internal_nma_ro_signing_attempt_ro_exact_hyperball_loss_from_paper_sample_lifting &m :
  structural_to_exact_hyperball_paper_sample_loss_obligation haetae_mode =>
  Pr[SIG.UF_NMA(H, HAETAE,
       ROMInternalTranscriptPaperSimAsNMA
         (A, ROSigningAttemptPaperSimSampler(H))).main() @ &m : res] <=
  Pr[SIG.UF_NMA(H, HAETAE,
       ROMInternalTranscriptPaperSimAsNMA
         (A, ROExactHyperballPaperSimSampler(H))).main() @ &m : res] +
    rejection_sampling_loss_term =>
  Pr[SIG.UF_NMA(H, HAETAE,
       ROMInternalTranscriptPaperSimAsNMA
         (A, ROSigningAttemptPaperSimSampler(H))).main() @ &m : res] <=
  Pr[SIG.UF_NMA(H, HAETAE,
       ROMInternalTranscriptPaperSimAsNMA
         (A, ROExactHyperballPaperSimSampler(H))).main() @ &m : res] +
    rejection_sampling_loss_term.
\end{lstlisting}
\noindent\textbf{Formal mathematical reading.}
Mathematical theorem. This theorem has the form \(\Pr[G:E] \le B\) is a game-probability statement.  It is one of the exact hops, bad-event bounds, or final advantage inequalities used in the reduction.

\noindent\textbf{Role in the proof.}
Use in this chapter. It contributes directly to an advantage bound, bad-event bound, or real-arithmetic composition step.

\noindent\textbf{Proof-script interpretation.}
Proof-script reading. The machine proof decomposes the theorem into rewrites, program-logic obligations, relational equivalences, and arithmetic side conditions.
\emph{Named identifiers syntactically cited in the proof block:}
\begin{lstlisting}[language=EasyCrypt,basicstyle=\ttfamily\scriptsize]
lifted
\end{lstlisting}

\end{formaltheorem}

\begin{formaltheorem}[EasyCrypt line 7036]
\noindent\textbf{EasyCrypt name.}
\begin{lstlisting}[language=EasyCrypt,basicstyle=\ttfamily\scriptsize]
rom_internal_budgeted_ro_signing_attempt_signing_call_loss_bound
\end{lstlisting}
\noindent\textbf{Checked EasyCrypt statement.}
\begin{lstlisting}[language=EasyCrypt,basicstyle=\ttfamily\scriptsize]
lemma rom_internal_budgeted_ro_signing_attempt_signing_call_loss_bound &m :
  Pr[SIG.UF_NMA(H, HAETAE,
       ROMInternalTranscriptBudgetedPaperSimAsNMA
         (A, ROSigningAttemptPaperSimSampler(H))).main() @ &m :
       0 < ROMInternalTranscriptBudgetedPaperSimAsNMA.signing_count] <=
    rom_signature_query_budget * rejection_sampling_loss_term.
\end{lstlisting}
\noindent\textbf{Formal mathematical reading.}
Mathematical theorem. This theorem has the form \(\Pr[G:E] \le B\) is a game-probability statement.  It is one of the exact hops, bad-event bounds, or final advantage inequalities used in the reduction.

\noindent\textbf{Role in the proof.}
Use in this chapter. It contributes directly to an advantage bound, bad-event bound, or real-arithmetic composition step.

\noindent\textbf{Proof-script interpretation.}
Proof-script reading. The machine proof decomposes the theorem into rewrites, program-logic obligations, relational equivalences, and arithmetic side conditions.
\emph{Main tactics visible in the proof script:} \ec{have}, \ec{smt}, \ec{rewrite}.
\emph{Named identifiers syntactically cited in the proof block:}
\begin{lstlisting}[language=EasyCrypt,basicstyle=\ttfamily\scriptsize]
HAETAE
Pr
ROMInternalTranscriptBudgetedPaperSimAsNMA
ROSigningAttemptPaperSimSampler
SIG
UF_NMA
budgeted_loss_ge1
call_event_le1
main
mu_bounded
rejection_sampling_loss_term
rejection_sampling_loss_term_ge1
rom_signature_query_budget
signature_query_budget_count
signature_query_count
signing_count
\end{lstlisting}

\end{formaltheorem}

\begin{formaltheorem}[EasyCrypt line 7057]
\noindent\textbf{EasyCrypt name.}
\begin{lstlisting}[language=EasyCrypt,basicstyle=\ttfamily\scriptsize]
rom_internal_budgeted_nma_ro_signing_attempt_ro_exact_hyperball_loss_from_paper_sample_lifting
\end{lstlisting}
\noindent\textbf{Checked EasyCrypt statement.}
\begin{lstlisting}[language=EasyCrypt,basicstyle=\ttfamily\scriptsize]
lemma rom_internal_budgeted_nma_ro_signing_attempt_ro_exact_hyperball_loss_from_paper_sample_lifting &m :
  structural_to_exact_hyperball_paper_sample_loss_obligation haetae_mode =>
  Pr[SIG.UF_NMA(H, HAETAE,
       ROMInternalTranscriptBudgetedPaperSimAsNMA
         (A, ROSigningAttemptPaperSimSampler(H))).main() @ &m : res] <=
  Pr[SIG.UF_NMA(H, HAETAE,
       ROMInternalTranscriptBudgetedPaperSimAsNMA
         (A, ROExactHyperballPaperSimSampler(H))).main() @ &m : res] +
    rom_signature_query_budget * rejection_sampling_loss_term.
\end{lstlisting}
\noindent\textbf{Formal mathematical reading.}
Mathematical theorem. This theorem has the form \(\Pr[G:E] \le B\) is a game-probability statement.  It is one of the exact hops, bad-event bounds, or final advantage inequalities used in the reduction.

\noindent\textbf{Role in the proof.}
Use in this chapter. It contributes directly to an advantage bound, bad-event bound, or real-arithmetic composition step.

\noindent\textbf{Proof-script interpretation.}
Proof-script reading. The machine proof decomposes the theorem into rewrites, program-logic obligations, relational equivalences, and arithmetic side conditions.
\emph{Main tactics visible in the proof script:} \ec{have}, \ec{apply}, \ec{smt}, \ec{rewrite}.
\emph{Named identifiers syntactically cited in the proof block:}
\begin{lstlisting}[language=EasyCrypt,basicstyle=\ttfamily\scriptsize]
HAETAE
Pr
ROExactHyperballPaperSimSampler
ROMInternalTranscriptBudgetedPaperSimAsNMA
ROSigningAttemptPaperSimSampler
SIG
StructuralAttemptPaperSample
UF_NMA
budgeted_loss_ge1
checked_one_step
context
ctx0
left_le1
m0
main
message
mu_bounded
p0
paper_sim_signature_sample
phoare
rejection_sampling_loss_term
rejection_sampling_loss_term_ge1
right_ge0
rom_signature_query_budget
sample
sample_loss
signature_query_budget_count
signature_query_count
sk0
skey
structural_attempt_exact_hyperball_paper_sample_one_step_loss
\end{lstlisting}

\end{formaltheorem}

\begin{formaltheorem}[EasyCrypt line 7097]
\noindent\textbf{EasyCrypt name.}
\begin{lstlisting}[language=EasyCrypt,basicstyle=\ttfamily\scriptsize]
rom_internal_budgeted_nma_ro_signing_attempt_ro_exact_hyperball_loss_from_clean_lifting
\end{lstlisting}
\noindent\textbf{Checked EasyCrypt statement.}
\begin{lstlisting}[language=EasyCrypt,basicstyle=\ttfamily\scriptsize]
lemma rom_internal_budgeted_nma_ro_signing_attempt_ro_exact_hyperball_loss_from_clean_lifting &m :
  Pr[SIG.UF_NMA(H, HAETAE,
       ROMInternalTranscriptBudgetedPaperSimAsNMA
         (A, ROSigningAttemptPaperSimSampler(H))).main() @ &m :
       res /\
       ! ROMInternalTranscriptBudgetedPaperSimAsNMA.sampler_bad_prequery] <=
  Pr[SIG.UF_NMA(H, HAETAE,
       ROMInternalTranscriptBudgetedPaperSimAsNMA
         (A, ROExactHyperballPaperSimSampler(H))).main() @ &m : res] +
    rom_signature_query_budget * rejection_sampling_loss_term =>
  Pr[SIG.UF_NMA(H, HAETAE,
       ROMInternalTranscriptBudgetedPaperSimAsNMA
         (A, ROSigningAttemptPaperSimSampler(H))).main() @ &m : res] <=
  Pr[SIG.UF_NMA(H, HAETAE,
       ROMInternalTranscriptBudgetedPaperSimAsNMA
         (A, ROExactHyperballPaperSimSampler(H))).main() @ &m : res] +
    rom_signature_query_budget * rejection_sampling_loss_term +
    rom_signature_query_budget *
      ((rom_hash_query_budget + rom_signature_query_budget) /
       challenge_support_cardinality_lower_bound).
\end{lstlisting}
\noindent\textbf{Formal mathematical reading.}
Mathematical theorem. This theorem has the form \(\Pr[G:E] \le B\) is a game-probability statement.  It is one of the exact hops, bad-event bounds, or final advantage inequalities used in the reduction.

\noindent\textbf{Role in the proof.}
Use in this chapter. It contributes directly to an advantage bound, bad-event bound, or real-arithmetic composition step.

\noindent\textbf{Proof-script interpretation.}
Proof-script reading. The machine proof decomposes the theorem into rewrites, program-logic obligations, relational equivalences, and arithmetic side conditions.
\emph{Main tactics visible in the proof script:} \ec{apply}, \ec{rewrite}, \ec{smt}.
\emph{Named identifiers syntactically cited in the proof block:}
\begin{lstlisting}[language=EasyCrypt,basicstyle=\ttfamily\scriptsize]
HAETAE
Pr
ROExactHyperballPaperSimSampler
ROMInternalTranscriptBudgetedPaperSimAsNMA
ROSigningAttemptPaperSimSampler
SIG
UF_NMA
budgeted_paper_sim_sampler_bad_prequery_nma_fel_bound
clean_lifting
ler_add
ler_trans
lerr
main
mu_or
mu_sub
rejection_sampling_loss_term
rom_signature_query_budget
sampler_bad_prequery
\end{lstlisting}

\end{formaltheorem}

\begin{formaltheorem}[EasyCrypt line 7154]
\noindent\textbf{EasyCrypt name.}
\begin{lstlisting}[language=EasyCrypt,basicstyle=\ttfamily\scriptsize]
rom_internal_budgeted_nma_ro_signing_attempt_ro_exact_hyperball_clean_conditioned_lifting
\end{lstlisting}
\noindent\textbf{Checked EasyCrypt statement.}
\begin{lstlisting}[language=EasyCrypt,basicstyle=\ttfamily\scriptsize]
lemma rom_internal_budgeted_nma_ro_signing_attempt_ro_exact_hyperball_clean_conditioned_lifting &m :
  structural_to_exact_hyperball_paper_sample_loss_obligation haetae_mode =>
  Pr[SIG.UF_NMA(H, HAETAE,
       ROMInternalTranscriptBudgetedPaperSimAsNMA
         (A, ROSigningAttemptPaperSimSampler(H))).main() @ &m :
       res /\
       ! ROMInternalTranscriptBudgetedPaperSimAsNMA.sampler_bad_prequery] <=
  Pr[SIG.UF_NMA(H, HAETAE,
       ROMInternalTranscriptBudgetedPaperSimAsNMA
         (A, ROExactHyperballPaperSimSampler(H))).main() @ &m : res] +
    rom_signature_query_budget * rejection_sampling_loss_term.
\end{lstlisting}
\noindent\textbf{Formal mathematical reading.}
Mathematical theorem. This theorem has the form \(\Pr[G:E] \le B\) is a game-probability statement.  It is one of the exact hops, bad-event bounds, or final advantage inequalities used in the reduction.

\noindent\textbf{Role in the proof.}
Use in this chapter. It contributes directly to an advantage bound, bad-event bound, or real-arithmetic composition step.

\noindent\textbf{Proof-script interpretation.}
Proof-script reading. The machine proof decomposes the theorem into rewrites, program-logic obligations, relational equivalences, and arithmetic side conditions.
\emph{Main tactics visible in the proof script:} \ec{have}, \ec{smt}.
\emph{Named identifiers syntactically cited in the proof block:}
\begin{lstlisting}[language=EasyCrypt,basicstyle=\ttfamily\scriptsize]
HAETAE
Pr
ROMInternalTranscriptBudgetedPaperSimAsNMA
ROSigningAttemptPaperSimSampler
SIG
UF_NMA
budgeted_lifting
clean_restrict
main
mu_sub
rom_internal_budgeted_nma_ro_signing_attempt_ro_exact_hyperball_loss_from_paper_sample_lifting
sample_loss
sampler_bad_prequery
\end{lstlisting}

\end{formaltheorem}

\begin{formaltheorem}[EasyCrypt line 7183]
\noindent\textbf{EasyCrypt name.}
\begin{lstlisting}[language=EasyCrypt,basicstyle=\ttfamily\scriptsize]
rom_internal_budgeted_nma_ro_signing_attempt_ro_exact_hyperball_loss_from_checked_clean_lifting
\end{lstlisting}
\noindent\textbf{Checked EasyCrypt statement.}
\begin{lstlisting}[language=EasyCrypt,basicstyle=\ttfamily\scriptsize]
lemma rom_internal_budgeted_nma_ro_signing_attempt_ro_exact_hyperball_loss_from_checked_clean_lifting &m :
  structural_to_exact_hyperball_paper_sample_loss_obligation haetae_mode =>
  Pr[SIG.UF_NMA(H, HAETAE,
       ROMInternalTranscriptBudgetedPaperSimAsNMA
         (A, ROSigningAttemptPaperSimSampler(H))).main() @ &m : res] <=
  Pr[SIG.UF_NMA(H, HAETAE,
       ROMInternalTranscriptBudgetedPaperSimAsNMA
         (A, ROExactHyperballPaperSimSampler(H))).main() @ &m : res] +
    rom_signature_query_budget * rejection_sampling_loss_term +
    rom_signature_query_budget *
      ((rom_hash_query_budget + rom_signature_query_budget) /
       challenge_support_cardinality_lower_bound).
\end{lstlisting}
\noindent\textbf{Formal mathematical reading.}
Mathematical theorem. This theorem has the form \(\Pr[G:E] \le B\) is a game-probability statement.  It is one of the exact hops, bad-event bounds, or final advantage inequalities used in the reduction.

\noindent\textbf{Role in the proof.}
Use in this chapter. It contributes directly to an advantage bound, bad-event bound, or real-arithmetic composition step.

\noindent\textbf{Proof-script interpretation.}
Proof-script reading. The machine proof decomposes the theorem into rewrites, program-logic obligations, relational equivalences, and arithmetic side conditions.
\emph{Main tactics visible in the proof script:} \ec{apply}.
\emph{Named identifiers syntactically cited in the proof block:}
\begin{lstlisting}[language=EasyCrypt,basicstyle=\ttfamily\scriptsize]
rom_internal_budgeted_nma_ro_signing_attempt_ro_exact_hyperball_clean_conditioned_lifting
rom_internal_budgeted_nma_ro_signing_attempt_ro_exact_hyperball_loss_from_clean_lifting
sample_loss
\end{lstlisting}

\end{formaltheorem}

\begin{formaltheorem}[EasyCrypt line 7205]
\noindent\textbf{EasyCrypt name.}
\begin{lstlisting}[language=EasyCrypt,basicstyle=\ttfamily\scriptsize]
rom_internal_nma_ro_signing_attempt_ro_exact_hyperball_loss_from_budgeted_lifting
\end{lstlisting}
\noindent\textbf{Checked EasyCrypt statement.}
\begin{lstlisting}[language=EasyCrypt,basicstyle=\ttfamily\scriptsize]
lemma rom_internal_nma_ro_signing_attempt_ro_exact_hyperball_loss_from_budgeted_lifting &m :
  Pr[SIG.UF_NMA(H, HAETAE,
       ROMInternalTranscriptPaperSimAsNMA
         (A, ROSigningAttemptPaperSimSampler(H))).main() @ &m : res] <=
  Pr[SIG.UF_NMA(H, HAETAE,
       ROMInternalTranscriptBudgetedPaperSimAsNMA
         (A, ROSigningAttemptPaperSimSampler(H))).main() @ &m : res] =>
  Pr[SIG.UF_NMA(H, HAETAE,
       ROMInternalTranscriptBudgetedPaperSimAsNMA
         (A, ROExactHyperballPaperSimSampler(H))).main() @ &m : res] <=
  Pr[SIG.UF_NMA(H, HAETAE,
       ROMInternalTranscriptPaperSimAsNMA
         (A, ROExactHyperballPaperSimSampler(H))).main() @ &m : res] =>
  structural_to_exact_hyperball_paper_sample_loss_obligation haetae_mode =>
  Pr[SIG.UF_NMA(H, HAETAE,
       ROMInternalTranscriptPaperSimAsNMA
         (A, ROSigningAttemptPaperSimSampler(H))).main() @ &m : res] <=
  Pr[SIG.UF_NMA(H, HAETAE,
       ROMInternalTranscriptPaperSimAsNMA
         (A, ROExactHyperballPaperSimSampler(H))).main() @ &m : res] +
    rom_signature_query_budget * rejection_sampling_loss_term.
\end{lstlisting}
\noindent\textbf{Formal mathematical reading.}
Mathematical theorem. This theorem has the form \(\Pr[G:E] \le B\) is a game-probability statement.  It is one of the exact hops, bad-event bounds, or final advantage inequalities used in the reduction.

\noindent\textbf{Role in the proof.}
Use in this chapter. It contributes directly to an advantage bound, bad-event bound, or real-arithmetic composition step.

\noindent\textbf{Proof-script interpretation.}
Proof-script reading. The machine proof decomposes the theorem into rewrites, program-logic obligations, relational equivalences, and arithmetic side conditions.
\emph{Main tactics visible in the proof script:} \ec{have}, \ec{apply}, \ec{smt}.
\emph{Named identifiers syntactically cited in the proof block:}
\begin{lstlisting}[language=EasyCrypt,basicstyle=\ttfamily\scriptsize]
HAETAE
Pr
ROExactHyperballPaperSimSampler
ROMInternalTranscriptBudgetedPaperSimAsNMA
ROSigningAttemptPaperSimSampler
SIG
UF_NMA
attempt_unbudgeted_to_budgeted
exact_budgeted_to_unbudgeted
hbudgeted
main
rejection_sampling_loss_term
rom_internal_budgeted_nma_ro_signing_attempt_ro_exact_hyperball_loss_from_paper_sample_lifting
rom_signature_query_budget
sample_loss
\end{lstlisting}

\end{formaltheorem}

\begin{formaltheorem}[EasyCrypt line 7243]
\noindent\textbf{EasyCrypt name.}
\begin{lstlisting}[language=EasyCrypt,basicstyle=\ttfamily\scriptsize]
rom_internal_nma_ro_signing_attempt_ro_exact_hyperball_loss_from_budgeted_checked_clean_lifting
\end{lstlisting}
\noindent\textbf{Checked EasyCrypt statement.}
\begin{lstlisting}[language=EasyCrypt,basicstyle=\ttfamily\scriptsize]
lemma rom_internal_nma_ro_signing_attempt_ro_exact_hyperball_loss_from_budgeted_checked_clean_lifting &m :
  Pr[SIG.UF_NMA(H, HAETAE,
       ROMInternalTranscriptPaperSimAsNMA
         (A, ROSigningAttemptPaperSimSampler(H))).main() @ &m : res] <=
  Pr[SIG.UF_NMA(H, HAETAE,
       ROMInternalTranscriptBudgetedPaperSimAsNMA
         (A, ROSigningAttemptPaperSimSampler(H))).main() @ &m : res] =>
  Pr[SIG.UF_NMA(H, HAETAE,
       ROMInternalTranscriptBudgetedPaperSimAsNMA
         (A, ROExactHyperballPaperSimSampler(H))).main() @ &m : res] <=
  Pr[SIG.UF_NMA(H, HAETAE,
       ROMInternalTranscriptPaperSimAsNMA
         (A, ROExactHyperballPaperSimSampler(H))).main() @ &m : res] =>
  structural_to_exact_hyperball_paper_sample_loss_obligation haetae_mode =>
  Pr[SIG.UF_NMA(H, HAETAE,
       ROMInternalTranscriptPaperSimAsNMA
         (A, ROSigningAttemptPaperSimSampler(H))).main() @ &m : res] <=
  Pr[SIG.UF_NMA(H, HAETAE,
       ROMInternalTranscriptPaperSimAsNMA
         (A, ROExactHyperballPaperSimSampler(H))).main() @ &m : res] +
    rom_signature_query_budget * rejection_sampling_loss_term +
    rom_signature_query_budget *
      ((rom_hash_query_budget + rom_signature_query_budget) /
       challenge_support_cardinality_lower_bound).
\end{lstlisting}
\noindent\textbf{Formal mathematical reading.}
Mathematical theorem. This theorem has the form \(\Pr[G:E] \le B\) is a game-probability statement.  It is one of the exact hops, bad-event bounds, or final advantage inequalities used in the reduction.

\noindent\textbf{Role in the proof.}
Use in this chapter. It contributes directly to an advantage bound, bad-event bound, or real-arithmetic composition step.

\noindent\textbf{Proof-script interpretation.}
Proof-script reading. The machine proof decomposes the theorem into rewrites, program-logic obligations, relational equivalences, and arithmetic side conditions.
\emph{Main tactics visible in the proof script:} \ec{have}, \ec{apply}, \ec{smt}.
\emph{Named identifiers syntactically cited in the proof block:}
\begin{lstlisting}[language=EasyCrypt,basicstyle=\ttfamily\scriptsize]
HAETAE
Pr
ROExactHyperballPaperSimSampler
ROMInternalTranscriptBudgetedPaperSimAsNMA
ROSigningAttemptPaperSimSampler
SIG
UF_NMA
attempt_unbudgeted_to_budgeted
challenge_support_cardinality_lower_bound
exact_budgeted_to_unbudgeted
hbudgeted
main
rejection_sampling_loss_term
rom_hash_query_budget
rom_internal_budgeted_nma_ro_signing_attempt_ro_exact_hyperball_loss_from_checked_clean_lifting
rom_signature_query_budget
sample_loss
\end{lstlisting}

\end{formaltheorem}

\begin{formaltheorem}[EasyCrypt line 7291]
\noindent\textbf{EasyCrypt name.}
\begin{lstlisting}[language=EasyCrypt,basicstyle=\ttfamily\scriptsize]
rom_internal_nma_ro_signing_attempt_ro_exact_hyperball_loss_bound
\end{lstlisting}
\noindent\textbf{Checked EasyCrypt statement.}
\begin{lstlisting}[language=EasyCrypt,basicstyle=\ttfamily\scriptsize]
lemma rom_internal_nma_ro_signing_attempt_ro_exact_hyperball_loss_bound &m :
  Pr[SIG.UF_NMA(H, HAETAE,
       ROMInternalTranscriptPaperSimAsNMA
         (A, ROSigningAttemptPaperSimSampler(H))).main() @ &m : res] <=
  Pr[SIG.UF_NMA(H, HAETAE,
       ROMInternalTranscriptPaperSimAsNMA
         (A, ROExactHyperballPaperSimSampler(H))).main() @ &m : res] +
    rejection_sampling_loss_term.
\end{lstlisting}
\noindent\textbf{Formal mathematical reading.}
Mathematical theorem. This theorem has the form \(\Pr[G:E] \le B\) is a game-probability statement.  It is one of the exact hops, bad-event bounds, or final advantage inequalities used in the reduction.

\noindent\textbf{Role in the proof.}
Use in this chapter. It contributes directly to an advantage bound, bad-event bound, or real-arithmetic composition step.

\noindent\textbf{Proof-script interpretation.}
Proof-script reading. The machine proof decomposes the theorem into rewrites, program-logic obligations, relational equivalences, and arithmetic side conditions.
\emph{Main tactics visible in the proof script:} \ec{have}, \ec{smt}, \ec{apply}.
\emph{Named identifiers syntactically cited in the proof block:}
\begin{lstlisting}[language=EasyCrypt,basicstyle=\ttfamily\scriptsize]
HAETAE
Pr
ROExactHyperballPaperSimSampler
ROMInternalTranscriptPaperSimAsNMA
ROSigningAttemptPaperSimSampler
SIG
UF_NMA
left_le1
loss_ge1
main
mu_bounded
rejection_sampling_loss_term
rejection_sampling_loss_term_ge1
right_ge0
\end{lstlisting}

\end{formaltheorem}

\begin{formaltheorem}[EasyCrypt line 7326]
\noindent\textbf{EasyCrypt name.}
\begin{lstlisting}[language=EasyCrypt,basicstyle=\ttfamily\scriptsize]
concrete_lazy_rom_get_preserves_signature_event
\end{lstlisting}
\noindent\textbf{Checked EasyCrypt statement.}
\begin{lstlisting}[language=EasyCrypt,basicstyle=\ttfamily\scriptsize]
lemma concrete_lazy_rom_get_preserves_signature_event
    pk m ctx smp (p : signature -> bool) :
  phoare[HAETAE_RO.FRO.get :
    p (paper_sim_signature haetae_mode pk m ctx smp) ==>
    p (paper_sim_signature haetae_mode pk m ctx smp)] = 1%r.
\end{lstlisting}
\noindent\textbf{Formal mathematical reading.}
Mathematical theorem. This is a relational program theorem: two games or procedures are coupled so that a precondition on their initial states implies the stated postcondition on final states and return values.

\noindent\textbf{Role in the proof.}
Use in this chapter. It preserves an invariant across an oracle call, adversary call, or game procedure.

\noindent\textbf{Proof-script interpretation.}
Proof-script reading. The machine proof decomposes the theorem into rewrites, program-logic obligations, relational equivalences, and arithmetic side conditions.
\emph{Main tactics visible in the proof script:} \ec{proc}, \ec{wp}, \ec{rnd}, \ec{smt}.
\emph{Named identifiers syntactically cited in the proof block:}
\begin{lstlisting}[language=EasyCrypt,basicstyle=\ttfamily\scriptsize]
ro_output_distribution_lossless
\end{lstlisting}

\end{formaltheorem}

\begin{formaltheorem}[EasyCrypt line 7338]
\noindent\textbf{EasyCrypt name.}
\begin{lstlisting}[language=EasyCrypt,basicstyle=\ttfamily\scriptsize]
concrete_lazy_rom_get_preserves_signature_event_with_pk
\end{lstlisting}
\noindent\textbf{Checked EasyCrypt statement.}
\begin{lstlisting}[language=EasyCrypt,basicstyle=\ttfamily\scriptsize]
lemma concrete_lazy_rom_get_preserves_signature_event_with_pk
    pk m ctx smp (p : signature -> bool) :
  phoare[HAETAE_RO.FRO.get :
    ROMInternalTranscriptBudgetedPaperSimAsNMA.pk_current = pk /\
    p (paper_sim_signature haetae_mode pk m ctx smp) ==>
    ROMInternalTranscriptBudgetedPaperSimAsNMA.pk_current = pk /\
    p (paper_sim_signature haetae_mode pk m ctx smp)] = 1%r.
\end{lstlisting}
\noindent\textbf{Formal mathematical reading.}
Mathematical theorem. This is a relational program theorem: two games or procedures are coupled so that a precondition on their initial states implies the stated postcondition on final states and return values.

\noindent\textbf{Role in the proof.}
Use in this chapter. It preserves an invariant across an oracle call, adversary call, or game procedure.

\noindent\textbf{Proof-script interpretation.}
Proof-script reading. The machine proof decomposes the theorem into rewrites, program-logic obligations, relational equivalences, and arithmetic side conditions.
\emph{Main tactics visible in the proof script:} \ec{proc}, \ec{wp}, \ec{rnd}, \ec{smt}.
\emph{Named identifiers syntactically cited in the proof block:}
\begin{lstlisting}[language=EasyCrypt,basicstyle=\ttfamily\scriptsize]
ro_output_distribution_lossless
\end{lstlisting}

\end{formaltheorem}

\begin{formaltheorem}[EasyCrypt line 7352]
\noindent\textbf{EasyCrypt name.}
\begin{lstlisting}[language=EasyCrypt,basicstyle=\ttfamily\scriptsize]
concrete_lazy_rom_get_preserves_current_signature_event
\end{lstlisting}
\noindent\textbf{Checked EasyCrypt statement.}
\begin{lstlisting}[language=EasyCrypt,basicstyle=\ttfamily\scriptsize]
lemma concrete_lazy_rom_get_preserves_current_signature_event
    m ctx smp (p : signature -> bool) :
  phoare[HAETAE_RO.FRO.get :
    p (paper_sim_signature haetae_mode
         ROMInternalTranscriptBudgetedPaperSimAsNMA.pk_current m ctx smp) ==>
    p (paper_sim_signature haetae_mode
         ROMInternalTranscriptBudgetedPaperSimAsNMA.pk_current m ctx smp)] = 1%r.
\end{lstlisting}
\noindent\textbf{Formal mathematical reading.}
Mathematical theorem. This is a relational program theorem: two games or procedures are coupled so that a precondition on their initial states implies the stated postcondition on final states and return values.

\noindent\textbf{Role in the proof.}
Use in this chapter. It preserves an invariant across an oracle call, adversary call, or game procedure.

\noindent\textbf{Proof-script interpretation.}
Proof-script reading. The machine proof decomposes the theorem into rewrites, program-logic obligations, relational equivalences, and arithmetic side conditions.
\emph{Main tactics visible in the proof script:} \ec{proc}, \ec{wp}, \ec{rnd}, \ec{smt}.
\emph{Named identifiers syntactically cited in the proof block:}
\begin{lstlisting}[language=EasyCrypt,basicstyle=\ttfamily\scriptsize]
ro_output_distribution_lossless
\end{lstlisting}

\end{formaltheorem}

\begin{formaltheorem}[EasyCrypt line 7366]
\noindent\textbf{EasyCrypt name.}
\begin{lstlisting}[language=EasyCrypt,basicstyle=\ttfamily\scriptsize]
concrete_lazy_rom_get_true
\end{lstlisting}
\noindent\textbf{Checked EasyCrypt statement.}
\begin{lstlisting}[language=EasyCrypt,basicstyle=\ttfamily\scriptsize]
lemma concrete_lazy_rom_get_true :
  phoare[HAETAE_RO.FRO.get : true ==> true] = 1%r.
\end{lstlisting}
\noindent\textbf{Formal mathematical reading.}
Mathematical theorem. This is a relational program theorem: two games or procedures are coupled so that a precondition on their initial states implies the stated postcondition on final states and return values.

\noindent\textbf{Role in the proof.}
Use in this chapter. It is a reusable checked theorem cited by later proof scripts.

\noindent\textbf{Proof-script interpretation.}
Proof-script reading. The machine proof decomposes the theorem into rewrites, program-logic obligations, relational equivalences, and arithmetic side conditions.
\emph{Main tactics visible in the proof script:} \ec{conseq}.
\emph{Named identifiers syntactically cited in the proof block:}
\begin{lstlisting}[language=EasyCrypt,basicstyle=\ttfamily\scriptsize]
concrete_lazy_rom_get_lossless
\end{lstlisting}

\end{formaltheorem}

\begin{formaltheorem}[EasyCrypt line 7372]
\noindent\textbf{EasyCrypt name.}
\begin{lstlisting}[language=EasyCrypt,basicstyle=\ttfamily\scriptsize]
concrete_lazy_rom_get_preserves_budgeted_clean_signature_event
\end{lstlisting}
\noindent\textbf{Checked EasyCrypt statement.}
\begin{lstlisting}[language=EasyCrypt,basicstyle=\ttfamily\scriptsize]
lemma concrete_lazy_rom_get_preserves_budgeted_clean_signature_event
    pk m ctx smp (p : signature -> bool) :
  phoare[HAETAE_RO.FRO.get :
    ROMInternalTranscriptBudgetedPaperSimAsNMA.pk_current = pk /\
    ! ROMInternalTranscriptBudgetedPaperSimAsNMA.sampler_bad_prequery /\
    p (paper_sim_signature haetae_mode pk m ctx smp) ==>
    ROMInternalTranscriptBudgetedPaperSimAsNMA.pk_current = pk /\
    ! ROMInternalTranscriptBudgetedPaperSimAsNMA.sampler_bad_prequery /\
    p (paper_sim_signature haetae_mode pk m ctx smp)] = 1%r.
\end{lstlisting}
\noindent\textbf{Formal mathematical reading.}
Mathematical theorem. This is a relational program theorem: two games or procedures are coupled so that a precondition on their initial states implies the stated postcondition on final states and return values.

\noindent\textbf{Role in the proof.}
Use in this chapter. It contributes directly to an advantage bound, bad-event bound, or real-arithmetic composition step.

\noindent\textbf{Proof-script interpretation.}
Proof-script reading. The machine proof decomposes the theorem into rewrites, program-logic obligations, relational equivalences, and arithmetic side conditions.
\emph{Main tactics visible in the proof script:} \ec{proc}, \ec{wp}, \ec{rnd}, \ec{smt}.
\emph{Named identifiers syntactically cited in the proof block:}
\begin{lstlisting}[language=EasyCrypt,basicstyle=\ttfamily\scriptsize]
ro_output_distribution_lossless
\end{lstlisting}

\end{formaltheorem}

\begin{formaltheorem}[EasyCrypt line 7388]
\noindent\textbf{EasyCrypt name.}
\begin{lstlisting}[language=EasyCrypt,basicstyle=\ttfamily\scriptsize]
concrete_lazy_rom_get_preserves_budgeted_clean_signature_event_hoare
\end{lstlisting}
\noindent\textbf{Checked EasyCrypt statement.}
\begin{lstlisting}[language=EasyCrypt,basicstyle=\ttfamily\scriptsize]
lemma concrete_lazy_rom_get_preserves_budgeted_clean_signature_event_hoare
    pk m ctx smp (p : signature -> bool) :
  hoare[HAETAE_RO.FRO.get :
    ROMInternalTranscriptBudgetedPaperSimAsNMA.pk_current = pk /\
    ! ROMInternalTranscriptBudgetedPaperSimAsNMA.sampler_bad_prequery /\
    p (paper_sim_signature haetae_mode pk m ctx smp) ==>
    ROMInternalTranscriptBudgetedPaperSimAsNMA.pk_current = pk /\
    ! ROMInternalTranscriptBudgetedPaperSimAsNMA.sampler_bad_prequery /\
    p (paper_sim_signature haetae_mode pk m ctx smp)].
\end{lstlisting}
\noindent\textbf{Formal mathematical reading.}
Mathematical theorem. This is a relational program theorem: two games or procedures are coupled so that a precondition on their initial states implies the stated postcondition on final states and return values.

\noindent\textbf{Role in the proof.}
Use in this chapter. It contributes directly to an advantage bound, bad-event bound, or real-arithmetic composition step.

\noindent\textbf{Proof-script interpretation.}
Proof-script reading. The machine proof decomposes the theorem into rewrites, program-logic obligations, relational equivalences, and arithmetic side conditions.
\emph{Main tactics visible in the proof script:} \ec{proc}, \ec{wp}, \ec{rnd}, \ec{smt}.
\emph{Named identifiers syntactically cited in the proof block:}
\begin{lstlisting}[language=EasyCrypt,basicstyle=\ttfamily\scriptsize]
ro_output_distribution_lossless
\end{lstlisting}

\end{formaltheorem}

\begin{formaltheorem}[EasyCrypt line 7404]
\noindent\textbf{EasyCrypt name.}
\begin{lstlisting}[language=EasyCrypt,basicstyle=\ttfamily\scriptsize]
concrete_ro_signing_attempt_sample_with_seed_lossless
\end{lstlisting}
\noindent\textbf{Checked EasyCrypt statement.}
\begin{lstlisting}[language=EasyCrypt,basicstyle=\ttfamily\scriptsize]
lemma concrete_ro_signing_attempt_sample_with_seed_lossless :
  islossless ROSigningAttemptPaperSimSampler(HAETAE_RO.FRO).sample_with_seed.
\end{lstlisting}
\noindent\textbf{Formal mathematical reading.}
Mathematical theorem. This lemma is a universally quantified proposition over the variables shown in the EasyCrypt statement.

\noindent\textbf{Role in the proof.}
Use in this chapter. It contributes directly to an advantage bound, bad-event bound, or real-arithmetic composition step.

\noindent\textbf{Proof-script interpretation.}
Proof-script reading. The machine proof decomposes the theorem into rewrites, program-logic obligations, relational equivalences, and arithmetic side conditions.
\emph{Main tactics visible in the proof script:} \ec{proc}, \ec{wp}, \ec{call}.
\emph{Named identifiers syntactically cited in the proof block:}
\begin{lstlisting}[language=EasyCrypt,basicstyle=\ttfamily\scriptsize]
concrete_lazy_rom_get_lossless
\end{lstlisting}

\end{formaltheorem}

\begin{formaltheorem}[EasyCrypt line 7413]
\noindent\textbf{EasyCrypt name.}
\begin{lstlisting}[language=EasyCrypt,basicstyle=\ttfamily\scriptsize]
concrete_budgeted_ro_signing_attempt_o_sign_active_clean_signature_structural_le
\end{lstlisting}
\noindent\textbf{Checked EasyCrypt statement.}
\begin{lstlisting}[language=EasyCrypt,basicstyle=\ttfamily\scriptsize]
lemma concrete_budgeted_ro_signing_attempt_o_sign_active_clean_signature_structural_le
    pk sk msg ctxt (p : signature -> bool) :
  phoare[ROMInternalTranscriptBudgetedPaperSimAsNMA
          (A, ROSigningAttemptPaperSimSampler(HAETAE_RO.FRO),
           HAETAE_RO.FRO).O.sign :
    arg = (msg, ctxt) /\
    ROMInternalTranscriptBudgetedPaperSimAsNMA.signing_count <
      signature_query_budget_count /\
    ROMInternalTranscriptBudgetedPaperSimAsNMA.pk_current = pk /\
    ROSigningAttemptPaperSimSampler.sk_current = sk /\
    sampler_rom_covered
      HAETAE_RO.FRO.m
      ROMInternalTranscriptBudgetedPaperSimAsNMA.adversary_hash_queries
      ROMInternalTranscriptBudgetedPaperSimAsNMA.sampler_expand_queries /\
    ! ROMInternalTranscriptBudgetedPaperSimAsNMA.sampler_bad_prequery ==>
    p res /\
    ! ROMInternalTranscriptBudgetedPaperSimAsNMA.sampler_bad_prequery] <=
  (mu (dsigning_attempt_state haetae_mode sk msg ctxt)
     (fun st =>
        p (paper_sim_signature haetae_mode pk msg ctxt
             (paper_sim_sample_from_rejection_attempt st)))).
\end{lstlisting}
\noindent\textbf{Formal mathematical reading.}
Mathematical theorem. This is a relational program theorem: two games or procedures are coupled so that a precondition on their initial states implies the stated postcondition on final states and return values.

\noindent\textbf{Role in the proof.}
Use in this chapter. It contributes directly to an advantage bound, bad-event bound, or real-arithmetic composition step.

\noindent\textbf{Proof-script interpretation.}
Proof-script reading. The machine proof decomposes the theorem into rewrites, program-logic obligations, relational equivalences, and arithmetic side conditions.
\emph{Main tactics visible in the proof script:} \ec{proc}, \ec{wp}, \ec{call}, \ec{rnd}, \ec{smt}, \ec{inline}, \ec{hoare}.
\emph{Named identifiers syntactically cited in the proof block:}
\begin{lstlisting}[language=EasyCrypt,basicstyle=\ttfamily\scriptsize]
FRO
HAETAE_RO
ROMInternalTranscriptBudgetedPaperSimAsNMA
ROSigningAttemptPaperSimSampler
adversary_hash_queries
concrete_ro_signing_attempt_sample_with_seed_guarded_clean_structural_arg_le
concrete_ro_signing_attempt_sample_with_seed_preserves_sampler_rom_covered_budgeted_logs
ctx
ctxt
dsigning_attempt_state
get
haetae_mode
msg
mu
mu_bounded
paper_sim_sample_from_rejection_attempt
paper_sim_signature
pk
pk_current
ro_output_distribution_lossless
sampler_bad_prequery
sampler_expand_queries
sampler_rom_covered
sampler_rom_covered_fresh_after_clean_seed
sampler_rom_covered_self_log_preserves
signing_coin_distribution_lossless
sk
sk_current
smp
st
\end{lstlisting}

\end{formaltheorem}

\begin{formaltheorem}[EasyCrypt line 7487]
\noindent\textbf{EasyCrypt name.}
\begin{lstlisting}[language=EasyCrypt,basicstyle=\ttfamily\scriptsize]
concrete_budgeted_ro_exact_hyperball_o_sign_active_signature_exact
\end{lstlisting}
\noindent\textbf{Checked EasyCrypt statement.}
\begin{lstlisting}[language=EasyCrypt,basicstyle=\ttfamily\scriptsize]
lemma concrete_budgeted_ro_exact_hyperball_o_sign_active_signature_exact
    pk sk msg ctxt (p : signature -> bool) :
  phoare[ROMInternalTranscriptBudgetedPaperSimAsNMA
          (A, ROExactHyperballPaperSimSampler(HAETAE_RO.FRO),
           HAETAE_RO.FRO).O.sign :
    arg = (msg, ctxt) /\
    ROMInternalTranscriptBudgetedPaperSimAsNMA.signing_count <
      signature_query_budget_count /\
    ROMInternalTranscriptBudgetedPaperSimAsNMA.pk_current = pk /\
    ROExactHyperballPaperSimSampler.sk_current = sk ==>
    p res] =
  (mu (dexact_hyperball_signing_attempt_state haetae_mode sk msg ctxt)
    (fun st =>
       p (paper_sim_signature haetae_mode pk msg ctxt
            (paper_sim_sample_from_rejection_attempt st)))).
\end{lstlisting}
\noindent\textbf{Formal mathematical reading.}
Mathematical theorem. This is a relational program theorem: two games or procedures are coupled so that a precondition on their initial states implies the stated postcondition on final states and return values.

\noindent\textbf{Role in the proof.}
Use in this chapter. It proves an exact game replacement or simulator equivalence.

\noindent\textbf{Proof-script interpretation.}
Proof-script reading. The machine proof decomposes the theorem into rewrites, program-logic obligations, relational equivalences, and arithmetic side conditions.
\emph{Main tactics visible in the proof script:} \ec{proc}, \ec{wp}, \ec{call}, \ec{smt}, \ec{inline}, \ec{rnd}.
\emph{Named identifiers syntactically cited in the proof block:}
\begin{lstlisting}[language=EasyCrypt,basicstyle=\ttfamily\scriptsize]
FRO
HAETAE_RO
ROExactHyperballPaperSimSampler
ROMInternalTranscriptBudgetedPaperSimAsNMA
concrete_lazy_rom_get_true
concrete_ro_exact_hyperball_sample_with_seed_exact_no_seed
ctx
ctxt
dexact_hyperball_signing_attempt_state
get
haetae_mode
msg
mu
mu_bounded
paper_sim_sample_from_rejection_attempt
paper_sim_signature
pk
pk_current
ro_output_distribution_lossless
signing_coin_distribution_lossless
sk
sk_current
smp
st
\end{lstlisting}

\end{formaltheorem}

\begin{formaltheorem}[EasyCrypt line 7541]
\noindent\textbf{EasyCrypt name.}
\begin{lstlisting}[language=EasyCrypt,basicstyle=\ttfamily\scriptsize]
concrete_budgeted_ro_exact_hyperball_sample_with_seed_to_o_sign_active_projection
\end{lstlisting}
\noindent\textbf{Checked EasyCrypt statement.}
\begin{lstlisting}[language=EasyCrypt,basicstyle=\ttfamily\scriptsize]
lemma concrete_budgeted_ro_exact_hyperball_sample_with_seed_to_o_sign_active_projection
    seed_coins pk sk msg ctxt (p : signature -> bool) &m :
  ROMInternalTranscriptBudgetedPaperSimAsNMA.signing_count{m} <
    signature_query_budget_count =>
  ROMInternalTranscriptBudgetedPaperSimAsNMA.pk_current{m} = pk =>
  ROExactHyperballPaperSimSampler.sk_current{m} = sk =>
  Pr[ROExactHyperballPaperSimSampler(HAETAE_RO.FRO).sample_with_seed
       (seed_coins, msg, ctxt) @ &m :
       p (paper_sim_signature haetae_mode pk msg ctxt res)] <=
  Pr[ROMInternalTranscriptBudgetedPaperSimAsNMA
       (A, ROExactHyperballPaperSimSampler(HAETAE_RO.FRO),
        HAETAE_RO.FRO).O.sign(msg, ctxt) @ &m :
       p res].
\end{lstlisting}
\noindent\textbf{Formal mathematical reading.}
Mathematical theorem. This theorem has the form \(\Pr[G:E] \le B\) is a game-probability statement.  It is one of the exact hops, bad-event bounds, or final advantage inequalities used in the reduction.

\noindent\textbf{Role in the proof.}
Use in this chapter. It contributes directly to an advantage bound, bad-event bound, or real-arithmetic composition step.

\noindent\textbf{Proof-script interpretation.}
Proof-script reading. The machine proof decomposes the theorem into rewrites, program-logic obligations, relational equivalences, and arithmetic side conditions.
\emph{Main tactics visible in the proof script:} \ec{have}, \ec{rewrite}, \ec{byphoare}, \ec{smt}.
\emph{Named identifiers syntactically cited in the proof block:}
\begin{lstlisting}[language=EasyCrypt,basicstyle=\ttfamily\scriptsize]
FRO
HAETAE_RO
Pr
ROExactHyperballPaperSimSampler
ROMInternalTranscriptBudgetedPaperSimAsNMA
active
concrete_budgeted_ro_exact_hyperball_o_sign_active_signature_exact
concrete_ro_exact_hyperball_sample_with_seed_exact_pr
ctxt
dexact_hyperball_signing_attempt_state
haetae_mode
msg
mu
osign_eq
paper_sim_sample_from_rejection_attempt
paper_sim_signature
pk
pk_eq
sample_eq
sample_with_seed
seed_coins
sign
sk
sk_eq
smp
st
\end{lstlisting}

\end{formaltheorem}

\begin{formaltheorem}[EasyCrypt line 7585]
\noindent\textbf{EasyCrypt name.}
\begin{lstlisting}[language=EasyCrypt,basicstyle=\ttfamily\scriptsize]
concrete_budgeted_ro_signing_attempt_o_sign_to_self_logged_sample_active_projection
\end{lstlisting}
\noindent\textbf{Checked EasyCrypt statement.}
\begin{lstlisting}[language=EasyCrypt,basicstyle=\ttfamily\scriptsize]
lemma concrete_budgeted_ro_signing_attempt_o_sign_to_self_logged_sample_active_projection
    seed_coins pk sk msg ctxt (p : signature -> bool) &m :
  sampler_rom_covered
    HAETAE_RO.FRO.m{m}
    ROMInternalTranscriptBudgetedPaperSimAsNMA.adversary_hash_queries{m}
    ROMInternalTranscriptBudgetedPaperSimAsNMA.sampler_expand_queries{m} =>
  ! (ROMInternalTranscriptBudgetedPaperSimAsNMA.sampler_bad_prequery{m} \/
     sampler_expand_query seed_coins \in
       ROMInternalTranscriptBudgetedPaperSimAsNMA.adversary_hash_queries{m} \/
     sampler_expand_query seed_coins \in
       ROMInternalTranscriptBudgetedPaperSimAsNMA.sampler_expand_queries{m}) =>
  ROMInternalTranscriptBudgetedPaperSimAsNMA.signing_count{m} <
    signature_query_budget_count =>
  ROMInternalTranscriptBudgetedPaperSimAsNMA.pk_current{m} = pk =>
  ROSigningAttemptPaperSimSampler.sk_current{m} = sk =>
  Pr[ROMInternalTranscriptBudgetedPaperSimAsNMA
       (A, ROSigningAttemptPaperSimSampler(HAETAE_RO.FRO),
        HAETAE_RO.FRO).O.sign(msg, ctxt) @ &m :
       p res /\
       ! ROMInternalTranscriptBudgetedPaperSimAsNMA.sampler_bad_prequery] <=
  Pr[ROSigningAttemptPaperSimSampler(HAETAE_RO.FRO).sample_with_seed
       (seed_coins, msg, ctxt) @ &m :
       p (paper_sim_signature haetae_mode pk msg ctxt res) /\
       ! (ROMInternalTranscriptBudgetedPaperSimAsNMA.sampler_bad_prequery{m} \/
          sampler_expand_query seed_coins \in
            ROMInternalTranscriptBudgetedPaperSimAsNMA.adversary_hash_queries{m} \/
          sampler_expand_query seed_coins \in
            ROMInternalTranscriptBudgetedPaperSimAsNMA.sampler_expand_queries{m})].
\end{lstlisting}
\noindent\textbf{Formal mathematical reading.}
Mathematical theorem. This theorem has the form \(\Pr[G:E] \le B\) is a game-probability statement.  It is one of the exact hops, bad-event bounds, or final advantage inequalities used in the reduction.

\noindent\textbf{Role in the proof.}
Use in this chapter. It contributes directly to an advantage bound, bad-event bound, or real-arithmetic composition step.

\noindent\textbf{Proof-script interpretation.}
Proof-script reading. The machine proof decomposes the theorem into rewrites, program-logic obligations, relational equivalences, and arithmetic side conditions.
\emph{Main tactics visible in the proof script:} \ec{have}, \ec{byphoare}, \ec{smt}, \ec{rewrite}.
\emph{Named identifiers syntactically cited in the proof block:}
\begin{lstlisting}[language=EasyCrypt,basicstyle=\ttfamily\scriptsize]
FRO
HAETAE_RO
Pr
ROMInternalTranscriptBudgetedPaperSimAsNMA
ROSigningAttemptPaperSimSampler
active
adversary_hash_queries
clean
clean_eq
concrete_budgeted_ro_signing_attempt_o_sign_active_clean_signature_structural_le
concrete_ro_signing_attempt_sample_with_seed_fresh_structural_eq
covered
ctxt
dsigning_attempt_state
haetae_mode
msg
mu
mu_eq
osign_le
paper_sim_sample_from_rejection_attempt
paper_sim_signature
pk
pk_eq
sample_eq
sample_with_seed
sampler_bad_prequery
sampler_expand_queries
sampler_expand_query
sampler_rom_covered_fresh_after_clean_seed
seed_coins
sign
sk
sk_eq
smp
st
\end{lstlisting}

\end{formaltheorem}

\begin{formaltheorem}[EasyCrypt line 7673]
\noindent\textbf{EasyCrypt name.}
\begin{lstlisting}[language=EasyCrypt,basicstyle=\ttfamily\scriptsize]
concrete_budgeted_o_sign_clean_one_call_loss_from_self_logged_surface
\end{lstlisting}
\noindent\textbf{Checked EasyCrypt statement.}
\begin{lstlisting}[language=EasyCrypt,basicstyle=\ttfamily\scriptsize]
lemma concrete_budgeted_o_sign_clean_one_call_loss_from_self_logged_surface
    seed_coins m ctx pk (p : signature -> bool) &m :
  structural_to_exact_hyperball_paper_sample_loss_obligation haetae_mode =>
  sampler_rom_covered
    HAETAE_RO.FRO.m{m}
    ROMInternalTranscriptBudgetedPaperSimAsNMA.adversary_hash_queries{m}
    ROMInternalTranscriptBudgetedPaperSimAsNMA.sampler_expand_queries{m} =>
  ! (ROMInternalTranscriptBudgetedPaperSimAsNMA.sampler_bad_prequery{m} \/
     sampler_expand_query seed_coins \in
       ROMInternalTranscriptBudgetedPaperSimAsNMA.adversary_hash_queries{m} \/
     sampler_expand_query seed_coins \in
       ROMInternalTranscriptBudgetedPaperSimAsNMA.sampler_expand_queries{m}) =>
  ROMInternalTranscriptBudgetedPaperSimAsNMA.signing_count{m} <
    signature_query_budget_count =>
  ROMInternalTranscriptBudgetedPaperSimAsNMA.pk_current{m} = pk =>
  ROSigningAttemptPaperSimSampler.sk_current{m} =
    ROExactHyperballPaperSimSampler.sk_current{m} =>
  Pr[ROMInternalTranscriptBudgetedPaperSimAsNMA
       (A, ROSigningAttemptPaperSimSampler(HAETAE_RO.FRO),
        HAETAE_RO.FRO).O.sign(m, ctx) @ &m :
       p res /\
       ! ROMInternalTranscriptBudgetedPaperSimAsNMA.sampler_bad_prequery] <=
  Pr[ROSigningAttemptPaperSimSampler(HAETAE_RO.FRO).sample_with_seed
       (seed_coins, m, ctx) @ &m :
       p (paper_sim_signature haetae_mode pk m ctx res) /\
       ! (ROMInternalTranscriptBudgetedPaperSimAsNMA.sampler_bad_prequery{m} \/
          sampler_expand_query seed_coins \in
            ROMInternalTranscriptBudgetedPaperSimAsNMA.adversary_hash_queries{m} \/
          sampler_expand_query seed_coins \in
            ROMInternalTranscriptBudgetedPaperSimAsNMA.sampler_expand_queries{m})] =>
  Pr[ROMInternalTranscriptBudgetedPaperSimAsNMA
       (A, ROSigningAttemptPaperSimSampler(HAETAE_RO.FRO),
        HAETAE_RO.FRO).O.sign(m, ctx) @ &m :
       p res /\
       ! ROMInternalTranscriptBudgetedPaperSimAsNMA.sampler_bad_prequery] <=
  Pr[ROMInternalTranscriptBudgetedPaperSimAsNMA
       (A, ROExactHyperballPaperSimSampler(HAETAE_RO.FRO),
        HAETAE_RO.FRO).O.sign(m, ctx) @ &m :
       p res] +
    rejection_sampling_loss_term.
\end{lstlisting}
\noindent\textbf{Formal mathematical reading.}
Mathematical theorem. This theorem has the form \(\Pr[G:E] \le B\) is a game-probability statement.  It is one of the exact hops, bad-event bounds, or final advantage inequalities used in the reduction.

\noindent\textbf{Role in the proof.}
Use in this chapter. It contributes directly to an advantage bound, bad-event bound, or real-arithmetic composition step.

\noindent\textbf{Proof-script interpretation.}
Proof-script reading. The machine proof decomposes the theorem into rewrites, program-logic obligations, relational equivalences, and arithmetic side conditions.
\emph{Main tactics visible in the proof script:} \ec{have}, \ec{apply}, \ec{smt}.
\emph{Named identifiers syntactically cited in the proof block:}
\begin{lstlisting}[language=EasyCrypt,basicstyle=\ttfamily\scriptsize]
FRO
HAETAE_RO
Pr
ROExactHyperballPaperSimSampler
ROMInternalTranscriptBudgetedPaperSimAsNMA
active
attempt_project
clean
concrete_budgeted_ro_exact_hyperball_sample_with_seed_to_o_sign_active_projection
concrete_budgeted_self_logged_signature_clean_event_loss_surface
covered
ctx
exact_project
haetae_mode
paper_sim_signature
pk
pk_eq
sample_loss
sample_with_seed
seed_coins
sign
sk_current
sk_eq
surface
\end{lstlisting}

\end{formaltheorem}

\begin{formaltheorem}[EasyCrypt line 7733]
\noindent\textbf{EasyCrypt name.}
\begin{lstlisting}[language=EasyCrypt,basicstyle=\ttfamily\scriptsize]
concrete_rom_internal_budgeted_nma_ro_signing_attempt_ro_exact_hyperball_clean_conditioned_lifting
\end{lstlisting}
\noindent\textbf{Checked EasyCrypt statement.}
\begin{lstlisting}[language=EasyCrypt,basicstyle=\ttfamily\scriptsize]
lemma concrete_rom_internal_budgeted_nma_ro_signing_attempt_ro_exact_hyperball_clean_conditioned_lifting &m :
  structural_to_exact_hyperball_paper_sample_loss_obligation haetae_mode =>
  Pr[SIG.UF_NMA(HAETAE_RO.FRO, HAETAE,
       ROMInternalTranscriptBudgetedPaperSimAsNMA
         (A, ROSigningAttemptPaperSimSampler(HAETAE_RO.FRO))).main() @ &m :
       res /\
       ! ROMInternalTranscriptBudgetedPaperSimAsNMA.sampler_bad_prequery] <=
  Pr[SIG.UF_NMA(HAETAE_RO.FRO, HAETAE,
       ROMInternalTranscriptBudgetedPaperSimAsNMA
         (A, ROExactHyperballPaperSimSampler(HAETAE_RO.FRO))).main() @ &m : res] +
    rom_signature_query_budget * rejection_sampling_loss_term.
\end{lstlisting}
\noindent\textbf{Formal mathematical reading.}
Mathematical theorem. This theorem has the form \(\Pr[G:E] \le B\) is a game-probability statement.  It is one of the exact hops, bad-event bounds, or final advantage inequalities used in the reduction.

\noindent\textbf{Role in the proof.}
Use in this chapter. It contributes directly to an advantage bound, bad-event bound, or real-arithmetic composition step.

\noindent\textbf{Proof-script interpretation.}
Proof-script reading. The machine proof decomposes the theorem into rewrites, program-logic obligations, relational equivalences, and arithmetic side conditions.
\emph{Main tactics visible in the proof script:} \ec{have}, \ec{smt}, \ec{rewrite}.
\emph{Named identifiers syntactically cited in the proof block:}
\begin{lstlisting}[language=EasyCrypt,basicstyle=\ttfamily\scriptsize]
FRO
HAETAE
HAETAE_RO
Pr
ROExactHyperballPaperSimSampler
ROMInternalTranscriptBudgetedPaperSimAsNMA
ROSigningAttemptPaperSimSampler
SIG
UF_NMA
budgeted_loss_ge1
left_le1
main
mu_bounded
rejection_sampling_loss_term
rejection_sampling_loss_term_ge1
right_ge0
rom_signature_query_budget
sample_loss
sampler_bad_prequery
signature_query_budget_count
signature_query_count
\end{lstlisting}

\end{formaltheorem}

\begin{formaltheorem}[EasyCrypt line 7767]
\noindent\textbf{EasyCrypt name.}
\begin{lstlisting}[language=EasyCrypt,basicstyle=\ttfamily\scriptsize]
concrete_rom_internal_budgeted_nma_ro_signing_attempt_ro_exact_hyperball_loss_from_clean_lifting
\end{lstlisting}
\noindent\textbf{Checked EasyCrypt statement.}
\begin{lstlisting}[language=EasyCrypt,basicstyle=\ttfamily\scriptsize]
lemma concrete_rom_internal_budgeted_nma_ro_signing_attempt_ro_exact_hyperball_loss_from_clean_lifting &m :
  Pr[SIG.UF_NMA(HAETAE_RO.FRO, HAETAE,
       ROMInternalTranscriptBudgetedPaperSimAsNMA
         (A, ROSigningAttemptPaperSimSampler(HAETAE_RO.FRO))).main() @ &m :
       res /\
       ! ROMInternalTranscriptBudgetedPaperSimAsNMA.sampler_bad_prequery] <=
  Pr[SIG.UF_NMA(HAETAE_RO.FRO, HAETAE,
       ROMInternalTranscriptBudgetedPaperSimAsNMA
         (A, ROExactHyperballPaperSimSampler(HAETAE_RO.FRO))).main() @ &m : res] +
    rom_signature_query_budget * rejection_sampling_loss_term =>
  Pr[SIG.UF_NMA(HAETAE_RO.FRO, HAETAE,
       ROMInternalTranscriptBudgetedPaperSimAsNMA
         (A, ROSigningAttemptPaperSimSampler(HAETAE_RO.FRO))).main() @ &m : res] <=
  Pr[SIG.UF_NMA(HAETAE_RO.FRO, HAETAE,
       ROMInternalTranscriptBudgetedPaperSimAsNMA
         (A, ROExactHyperballPaperSimSampler(HAETAE_RO.FRO))).main() @ &m : res] +
    rom_signature_query_budget * rejection_sampling_loss_term +
    rom_signature_query_budget *
      ((rom_hash_query_budget + rom_signature_query_budget) /
       challenge_support_cardinality_lower_bound).
\end{lstlisting}
\noindent\textbf{Formal mathematical reading.}
Mathematical theorem. This theorem has the form \(\Pr[G:E] \le B\) is a game-probability statement.  It is one of the exact hops, bad-event bounds, or final advantage inequalities used in the reduction.

\noindent\textbf{Role in the proof.}
Use in this chapter. It contributes directly to an advantage bound, bad-event bound, or real-arithmetic composition step.

\noindent\textbf{Proof-script interpretation.}
Proof-script reading. The machine proof decomposes the theorem into rewrites, program-logic obligations, relational equivalences, and arithmetic side conditions.
\emph{Main tactics visible in the proof script:} \ec{apply}.
\emph{Named identifiers syntactically cited in the proof block:}
\begin{lstlisting}[language=EasyCrypt,basicstyle=\ttfamily\scriptsize]
FRO
HAETAE_RO
clean_lifting
rom_internal_budgeted_nma_ro_signing_attempt_ro_exact_hyperball_loss_from_clean_lifting
\end{lstlisting}

\end{formaltheorem}

\begin{formaltheorem}[EasyCrypt line 7795]
\noindent\textbf{EasyCrypt name.}
\begin{lstlisting}[language=EasyCrypt,basicstyle=\ttfamily\scriptsize]
concrete_rom_internal_budgeted_nma_ro_signing_attempt_ro_exact_hyperball_loss_from_paper_sample_lifting
\end{lstlisting}
\noindent\textbf{Checked EasyCrypt statement.}
\begin{lstlisting}[language=EasyCrypt,basicstyle=\ttfamily\scriptsize]
lemma concrete_rom_internal_budgeted_nma_ro_signing_attempt_ro_exact_hyperball_loss_from_paper_sample_lifting &m :
  structural_to_exact_hyperball_paper_sample_loss_obligation haetae_mode =>
  Pr[SIG.UF_NMA(HAETAE_RO.FRO, HAETAE,
       ROMInternalTranscriptBudgetedPaperSimAsNMA
         (A, ROSigningAttemptPaperSimSampler(HAETAE_RO.FRO))).main() @ &m : res] <=
  Pr[SIG.UF_NMA(HAETAE_RO.FRO, HAETAE,
       ROMInternalTranscriptBudgetedPaperSimAsNMA
         (A, ROExactHyperballPaperSimSampler(HAETAE_RO.FRO))).main() @ &m : res] +
    rom_signature_query_budget * rejection_sampling_loss_term +
    rom_signature_query_budget *
      ((rom_hash_query_budget + rom_signature_query_budget) /
       challenge_support_cardinality_lower_bound).
\end{lstlisting}
\noindent\textbf{Formal mathematical reading.}
Mathematical theorem. This theorem has the form \(\Pr[G:E] \le B\) is a game-probability statement.  It is one of the exact hops, bad-event bounds, or final advantage inequalities used in the reduction.

\noindent\textbf{Role in the proof.}
Use in this chapter. It contributes directly to an advantage bound, bad-event bound, or real-arithmetic composition step.

\noindent\textbf{Proof-script interpretation.}
Proof-script reading. The machine proof decomposes the theorem into rewrites, program-logic obligations, relational equivalences, and arithmetic side conditions.
\emph{Main tactics visible in the proof script:} \ec{apply}.
\emph{Named identifiers syntactically cited in the proof block:}
\begin{lstlisting}[language=EasyCrypt,basicstyle=\ttfamily\scriptsize]
concrete_rom_internal_budgeted_nma_ro_signing_attempt_ro_exact_hyperball_clean_conditioned_lifting
concrete_rom_internal_budgeted_nma_ro_signing_attempt_ro_exact_hyperball_loss_from_clean_lifting
sample_loss
\end{lstlisting}

\end{formaltheorem}

\begin{formaltheorem}[EasyCrypt line 7828]
\noindent\textbf{EasyCrypt name.}
\begin{lstlisting}[language=EasyCrypt,basicstyle=\ttfamily\scriptsize]
real_signing_rejection_paper_sim_nma_oracle
\end{lstlisting}
\noindent\textbf{Checked EasyCrypt statement.}
\begin{lstlisting}[language=EasyCrypt,basicstyle=\ttfamily\scriptsize]
equiv real_signing_rejection_paper_sim_nma_oracle :
  ROMInternalTranscriptPaperSimAsNMA
    (A, RealSigningPaperSimSampler(H), H).O.sign ~
  ROMInternalTranscriptPaperSimAsNMA
    (A, HAETAERejectionPaperSimSampler(H), H).O.sign :
    ={glob H, arg} /\
    ROMInternalTranscriptPaperSimAsNMA.pk_current{1} =
      ROMInternalTranscriptPaperSimAsNMA.pk_current{2} /\
    RealSigningPaperSimSampler.pk_current{1} =
      HAETAERejectionPaperSimSampler.pk_current{2} /\
    RealSigningPaperSimSampler.sk_current{1} =
      HAETAERejectionPaperSimSampler.sk_current{2} /\
    RealSigningPaperSimSampler.pk_current{1} =
      public_key_of_secret haetae_mode
        RealSigningPaperSimSampler.sk_current{1} /\
    ROMInternalTranscriptPaperSimAsNMA.queries{1} =
      ROMInternalTranscriptPaperSimAsNMA.queries{2} /\
    ROMInternalTranscriptPaperSimAsNMA.transcripts{1} =
      ROMInternalTranscriptPaperSimAsNMA.transcripts{2} /\
    ROMInternalTranscriptPaperSimAsNMA.records{1} =
      ROMInternalTranscriptPaperSimAsNMA.records{2}
    ==>
    ={glob H, res} /\
    ROMInternalTranscriptPaperSimAsNMA.pk_current{1} =
      ROMInternalTranscriptPaperSimAsNMA.pk_current{2} /\
    RealSigningPaperSimSampler.pk_current{1} =
      HAETAERejectionPaperSimSampler.pk_current{2} /\
    RealSigningPaperSimSampler.sk_current{1} =
      HAETAERejectionPaperSimSampler.sk_current{2} /\
    RealSigningPaperSimSampler.pk_current{1} =
      public_key_of_secret haetae_mode
        RealSigningPaperSimSampler.sk_current{1} /\
    ROMInternalTranscriptPaperSimAsNMA.queries{1} =
      ROMInternalTranscriptPaperSimAsNMA.queries{2} /\
    ROMInternalTranscriptPaperSimAsNMA.transcripts{1} =
      ROMInternalTranscriptPaperSimAsNMA.transcripts{2} /\
    ROMInternalTranscriptPaperSimAsNMA.records{1} =
      ROMInternalTranscriptPaperSimAsNMA.records{2}.
\end{lstlisting}
\noindent\textbf{Formal mathematical reading.}
Mathematical theorem. This is a relational program theorem: two games or procedures are coupled so that a precondition on their initial states implies the stated postcondition on final states and return values.

\noindent\textbf{Role in the proof.}
Use in this chapter. It contributes directly to an advantage bound, bad-event bound, or real-arithmetic composition step.

\noindent\textbf{Proof-script interpretation.}
Proof-script reading. The machine proof decomposes the theorem into rewrites, program-logic obligations, relational equivalences, and arithmetic side conditions.
\emph{Main tactics visible in the proof script:} \ec{proc}, \ec{inline}, \ec{wp}, \ec{call}, \ec{rnd}, \ec{smt}.
\emph{Named identifiers syntactically cited in the proof block:}
\begin{lstlisting}[language=EasyCrypt,basicstyle=\ttfamily\scriptsize]
HAETAERejectionPaperSimSampler
RealSigningPaperSimSampler
arg
glob
haetae_rejection_sample_from_attempt_of_coinsE
sample
sample_with_seed
\end{lstlisting}

\end{formaltheorem}

\begin{formaltheorem}[EasyCrypt line 7886]
\noindent\textbf{EasyCrypt name.}
\begin{lstlisting}[language=EasyCrypt,basicstyle=\ttfamily\scriptsize]
adversary_real_signing_rejection_paper_sim_nma
\end{lstlisting}
\noindent\textbf{Checked EasyCrypt statement.}
\begin{lstlisting}[language=EasyCrypt,basicstyle=\ttfamily\scriptsize]
equiv adversary_real_signing_rejection_paper_sim_nma :
  A(H, ROMInternalTranscriptPaperSimAsNMA
      (A, RealSigningPaperSimSampler(H), H).O).forge ~
  A(H, ROMInternalTranscriptPaperSimAsNMA
      (A, HAETAERejectionPaperSimSampler(H), H).O).forge :
    ={glob H, glob A, arg} /\
    ROMInternalTranscriptPaperSimAsNMA.pk_current{1} =
      ROMInternalTranscriptPaperSimAsNMA.pk_current{2} /\
    RealSigningPaperSimSampler.pk_current{1} =
      HAETAERejectionPaperSimSampler.pk_current{2} /\
    RealSigningPaperSimSampler.sk_current{1} =
      HAETAERejectionPaperSimSampler.sk_current{2} /\
    RealSigningPaperSimSampler.pk_current{1} =
      public_key_of_secret haetae_mode
        RealSigningPaperSimSampler.sk_current{1} /\
    ROMInternalTranscriptPaperSimAsNMA.queries{1} =
      ROMInternalTranscriptPaperSimAsNMA.queries{2} /\
    ROMInternalTranscriptPaperSimAsNMA.transcripts{1} =
      ROMInternalTranscriptPaperSimAsNMA.transcripts{2} /\
    ROMInternalTranscriptPaperSimAsNMA.records{1} =
      ROMInternalTranscriptPaperSimAsNMA.records{2}
    ==>
    ={glob H, res} /\
    ROMInternalTranscriptPaperSimAsNMA.pk_current{1} =
      ROMInternalTranscriptPaperSimAsNMA.pk_current{2} /\
    RealSigningPaperSimSampler.pk_current{1} =
      HAETAERejectionPaperSimSampler.pk_current{2} /\
    RealSigningPaperSimSampler.sk_current{1} =
      HAETAERejectionPaperSimSampler.sk_current{2} /\
    RealSigningPaperSimSampler.pk_current{1} =
      public_key_of_secret haetae_mode
        RealSigningPaperSimSampler.sk_current{1} /\
    ROMInternalTranscriptPaperSimAsNMA.queries{1} =
      ROMInternalTranscriptPaperSimAsNMA.queries{2} /\
    ROMInternalTranscriptPaperSimAsNMA.transcripts{1} =
      ROMInternalTranscriptPaperSimAsNMA.transcripts{2} /\
    ROMInternalTranscriptPaperSimAsNMA.records{1} =
      ROMInternalTranscriptPaperSimAsNMA.records{2}.
\end{lstlisting}
\noindent\textbf{Formal mathematical reading.}
Mathematical theorem. This is a relational program theorem: two games or procedures are coupled so that a precondition on their initial states implies the stated postcondition on final states and return values.

\noindent\textbf{Role in the proof.}
Use in this chapter. It proves an exact game replacement or simulator equivalence.

\noindent\textbf{Proof-script interpretation.}
Proof-script reading. The machine proof decomposes the theorem into rewrites, program-logic obligations, relational equivalences, and arithmetic side conditions.
\emph{Main tactics visible in the proof script:} \ec{proc}, \ec{inline}, \ec{wp}, \ec{call}, \ec{rnd}, \ec{smt}.
\emph{Named identifiers syntactically cited in the proof block:}
\begin{lstlisting}[language=EasyCrypt,basicstyle=\ttfamily\scriptsize]
HAETAERejectionPaperSimSampler
ROMInternalTranscriptPaperSimAsNMA
RealSigningPaperSimSampler
arg
glob
haetae_mode
haetae_rejection_sample_from_attempt_of_coinsE
pk_current
public_key_of_secret
queries
records
sample
sample_with_seed
sk_current
transcripts
\end{lstlisting}

\end{formaltheorem}

\begin{formaltheorem}[EasyCrypt line 7963]
\noindent\textbf{EasyCrypt name.}
\begin{lstlisting}[language=EasyCrypt,basicstyle=\ttfamily\scriptsize]
rom_internal_nma_real_signing_haetae_rejection_paper_sim_exact
\end{lstlisting}
\noindent\textbf{Checked EasyCrypt statement.}
\begin{lstlisting}[language=EasyCrypt,basicstyle=\ttfamily\scriptsize]
lemma rom_internal_nma_real_signing_haetae_rejection_paper_sim_exact &m :
  Pr[SIG.UF_NMA(H, HAETAE,
       ROMInternalTranscriptPaperSimAsNMA
         (A, RealSigningPaperSimSampler(H))).main() @ &m : res] =
  Pr[SIG.UF_NMA(H, HAETAE,
       ROMInternalTranscriptPaperSimAsNMA
         (A, HAETAERejectionPaperSimSampler(H))).main() @ &m : res].
\end{lstlisting}
\noindent\textbf{Formal mathematical reading.}
Mathematical theorem. This theorem has the form \(\Pr[G:E] = B\) is a game-probability statement.  It is one of the exact hops, bad-event bounds, or final advantage inequalities used in the reduction.

\noindent\textbf{Role in the proof.}
Use in this chapter. It proves an exact game replacement or simulator equivalence.

\noindent\textbf{Proof-script interpretation.}
Proof-script reading. The machine proof decomposes the theorem into rewrites, program-logic obligations, relational equivalences, and arithmetic side conditions.
\emph{Main tactics visible in the proof script:} \ec{byequiv}, \ec{proc}, \ec{inline}, \ec{wp}, \ec{call}, \ec{apply}, \ec{rnd}, \ec{rewrite}.
\emph{Named identifiers syntactically cited in the proof block:}
\begin{lstlisting}[language=EasyCrypt,basicstyle=\ttfamily\scriptsize]
HAETAE
HAETAERejectionPaperSimSampler
ROMInternalTranscriptPaperSimAsNMA
RealSigningPaperSimSampler
adversary_real_signing_rejection_paper_sim_nma
arg
forge
glob
haetae_mode
init
keygen_internal
kg
pk_current
public_key_of_secret
queries
records
secret_key_of_seed
sk_current
transcripts
verify
\end{lstlisting}

\end{formaltheorem}

\begin{formaltheorem}[EasyCrypt line 8039]
\noindent\textbf{EasyCrypt name.}
\begin{lstlisting}[language=EasyCrypt,basicstyle=\ttfamily\scriptsize]
exact_hyperball_haetae_rejection_paper_sim_nma_oracle
\end{lstlisting}
\noindent\textbf{Checked EasyCrypt statement.}
\begin{lstlisting}[language=EasyCrypt,basicstyle=\ttfamily\scriptsize]
equiv exact_hyperball_haetae_rejection_paper_sim_nma_oracle :
  ROMInternalTranscriptPaperSimAsNMA
    (A, ExactHyperballHAETAERejectionPaperSimSampler(H), H).O.sign ~
  ROMInternalTranscriptPaperSimAsNMA
    (A, ExactHyperballPaperSimSampler(H), H).O.sign :
    ={glob H, arg} /\
    ROMInternalTranscriptPaperSimAsNMA.pk_current{1} =
      ROMInternalTranscriptPaperSimAsNMA.pk_current{2} /\
    ExactHyperballHAETAERejectionPaperSimSampler.pk_current{1} =
      ExactHyperballPaperSimSampler.pk_current{2} /\
    ExactHyperballHAETAERejectionPaperSimSampler.sk_current{1} =
      ExactHyperballPaperSimSampler.sk_current{2} /\
    ExactHyperballHAETAERejectionPaperSimSampler.pk_current{1} =
      public_key_of_secret haetae_mode
        ExactHyperballHAETAERejectionPaperSimSampler.sk_current{1} /\
    ROMInternalTranscriptPaperSimAsNMA.queries{1} =
      ROMInternalTranscriptPaperSimAsNMA.queries{2} /\
    ROMInternalTranscriptPaperSimAsNMA.transcripts{1} =
      ROMInternalTranscriptPaperSimAsNMA.transcripts{2} /\
    ROMInternalTranscriptPaperSimAsNMA.records{1} =
      ROMInternalTranscriptPaperSimAsNMA.records{2}
    ==>
    ={glob H, res} /\
    ROMInternalTranscriptPaperSimAsNMA.pk_current{1} =
      ROMInternalTranscriptPaperSimAsNMA.pk_current{2} /\
    ExactHyperballHAETAERejectionPaperSimSampler.pk_current{1} =
      ExactHyperballPaperSimSampler.pk_current{2} /\
    ExactHyperballHAETAERejectionPaperSimSampler.sk_current{1} =
      ExactHyperballPaperSimSampler.sk_current{2} /\
    ExactHyperballHAETAERejectionPaperSimSampler.pk_current{1} =
      public_key_of_secret haetae_mode
        ExactHyperballHAETAERejectionPaperSimSampler.sk_current{1} /\
    ROMInternalTranscriptPaperSimAsNMA.queries{1} =
      ROMInternalTranscriptPaperSimAsNMA.queries{2} /\
    ROMInternalTranscriptPaperSimAsNMA.transcripts{1} =
      ROMInternalTranscriptPaperSimAsNMA.transcripts{2} /\
    ROMInternalTranscriptPaperSimAsNMA.records{1} =
      ROMInternalTranscriptPaperSimAsNMA.records{2}.
\end{lstlisting}
\noindent\textbf{Formal mathematical reading.}
Mathematical theorem. This is a relational program theorem: two games or procedures are coupled so that a precondition on their initial states implies the stated postcondition on final states and return values.

\noindent\textbf{Role in the proof.}
Use in this chapter. It contributes directly to an advantage bound, bad-event bound, or real-arithmetic composition step.

\noindent\textbf{Proof-script interpretation.}
Proof-script reading. The machine proof decomposes the theorem into rewrites, program-logic obligations, relational equivalences, and arithmetic side conditions.
\emph{Main tactics visible in the proof script:} \ec{proc}, \ec{wp}, \ec{call}.
\emph{Named identifiers syntactically cited in the proof block:}
\begin{lstlisting}[language=EasyCrypt,basicstyle=\ttfamily\scriptsize]
ExactHyperballHAETAERejectionPaperSimSampler
ExactHyperballPaperSimSampler
arg
glob
pk_current
sk_current
\end{lstlisting}

\end{formaltheorem}

\begin{formaltheorem}[EasyCrypt line 8101]
\noindent\textbf{EasyCrypt name.}
\begin{lstlisting}[language=EasyCrypt,basicstyle=\ttfamily\scriptsize]
adversary_exact_hyperball_haetae_rejection_paper_sim_nma
\end{lstlisting}
\noindent\textbf{Checked EasyCrypt statement.}
\begin{lstlisting}[language=EasyCrypt,basicstyle=\ttfamily\scriptsize]
equiv adversary_exact_hyperball_haetae_rejection_paper_sim_nma :
  A(H, ROMInternalTranscriptPaperSimAsNMA
      (A, ExactHyperballHAETAERejectionPaperSimSampler(H), H).O).forge ~
  A(H, ROMInternalTranscriptPaperSimAsNMA
      (A, ExactHyperballPaperSimSampler(H), H).O).forge :
    ={glob H, glob A, arg} /\
    ROMInternalTranscriptPaperSimAsNMA.pk_current{1} =
      ROMInternalTranscriptPaperSimAsNMA.pk_current{2} /\
    ExactHyperballHAETAERejectionPaperSimSampler.pk_current{1} =
      ExactHyperballPaperSimSampler.pk_current{2} /\
    ExactHyperballHAETAERejectionPaperSimSampler.sk_current{1} =
      ExactHyperballPaperSimSampler.sk_current{2} /\
    ExactHyperballHAETAERejectionPaperSimSampler.pk_current{1} =
      public_key_of_secret haetae_mode
        ExactHyperballHAETAERejectionPaperSimSampler.sk_current{1} /\
    ROMInternalTranscriptPaperSimAsNMA.queries{1} =
      ROMInternalTranscriptPaperSimAsNMA.queries{2} /\
    ROMInternalTranscriptPaperSimAsNMA.transcripts{1} =
      ROMInternalTranscriptPaperSimAsNMA.transcripts{2} /\
    ROMInternalTranscriptPaperSimAsNMA.records{1} =
      ROMInternalTranscriptPaperSimAsNMA.records{2}
    ==>
    ={glob H, res} /\
    ROMInternalTranscriptPaperSimAsNMA.pk_current{1} =
      ROMInternalTranscriptPaperSimAsNMA.pk_current{2} /\
    ExactHyperballHAETAERejectionPaperSimSampler.pk_current{1} =
      ExactHyperballPaperSimSampler.pk_current{2} /\
    ExactHyperballHAETAERejectionPaperSimSampler.sk_current{1} =
      ExactHyperballPaperSimSampler.sk_current{2} /\
    ExactHyperballHAETAERejectionPaperSimSampler.pk_current{1} =
      public_key_of_secret haetae_mode
        ExactHyperballHAETAERejectionPaperSimSampler.sk_current{1} /\
    ROMInternalTranscriptPaperSimAsNMA.queries{1} =
      ROMInternalTranscriptPaperSimAsNMA.queries{2} /\
    ROMInternalTranscriptPaperSimAsNMA.transcripts{1} =
      ROMInternalTranscriptPaperSimAsNMA.transcripts{2} /\
    ROMInternalTranscriptPaperSimAsNMA.records{1} =
      ROMInternalTranscriptPaperSimAsNMA.records{2}.
\end{lstlisting}
\noindent\textbf{Formal mathematical reading.}
Mathematical theorem. This is a relational program theorem: two games or procedures are coupled so that a precondition on their initial states implies the stated postcondition on final states and return values.

\noindent\textbf{Role in the proof.}
Use in this chapter. It proves an exact game replacement or simulator equivalence.

\noindent\textbf{Proof-script interpretation.}
Proof-script reading. The machine proof decomposes the theorem into rewrites, program-logic obligations, relational equivalences, and arithmetic side conditions.
\emph{Main tactics visible in the proof script:} \ec{proc}, \ec{call}.
\emph{Named identifiers syntactically cited in the proof block:}
\begin{lstlisting}[language=EasyCrypt,basicstyle=\ttfamily\scriptsize]
ExactHyperballHAETAERejectionPaperSimSampler
ExactHyperballPaperSimSampler
ROMInternalTranscriptPaperSimAsNMA
exact_hyperball_haetae_rejection_paper_sim_nma_oracle
glob
haetae_mode
pk_current
public_key_of_secret
queries
records
sk_current
transcripts
\end{lstlisting}

\end{formaltheorem}

\begin{formaltheorem}[EasyCrypt line 8164]
\noindent\textbf{EasyCrypt name.}
\begin{lstlisting}[language=EasyCrypt,basicstyle=\ttfamily\scriptsize]
rom_internal_nma_exact_hyperball_rejection_paper_sim_no_fallback_exact
\end{lstlisting}
\noindent\textbf{Checked EasyCrypt statement.}
\begin{lstlisting}[language=EasyCrypt,basicstyle=\ttfamily\scriptsize]
lemma rom_internal_nma_exact_hyperball_rejection_paper_sim_no_fallback_exact &m :
  Pr[SIG.UF_NMA(H, HAETAE,
       ROMInternalTranscriptPaperSimAsNMA
         (A, ExactHyperballHAETAERejectionPaperSimSampler(H))).main()
       @ &m : res] =
  Pr[SIG.UF_NMA(H, HAETAE,
       ROMInternalTranscriptPaperSimAsNMA
         (A, ExactHyperballPaperSimSampler(H))).main() @ &m : res].
\end{lstlisting}
\noindent\textbf{Formal mathematical reading.}
Mathematical theorem. This theorem has the form \(\Pr[G:E] = B\) is a game-probability statement.  It is one of the exact hops, bad-event bounds, or final advantage inequalities used in the reduction.

\noindent\textbf{Role in the proof.}
Use in this chapter. It proves an exact game replacement or simulator equivalence.

\noindent\textbf{Proof-script interpretation.}
Proof-script reading. The machine proof decomposes the theorem into rewrites, program-logic obligations, relational equivalences, and arithmetic side conditions.
\emph{Main tactics visible in the proof script:} \ec{byequiv}, \ec{proc}, \ec{inline}, \ec{wp}, \ec{call}, \ec{apply}, \ec{rnd}, \ec{rewrite}.
\emph{Named identifiers syntactically cited in the proof block:}
\begin{lstlisting}[language=EasyCrypt,basicstyle=\ttfamily\scriptsize]
ExactHyperballHAETAERejectionPaperSimSampler
ExactHyperballPaperSimSampler
HAETAE
ROMInternalTranscriptPaperSimAsNMA
adversary_exact_hyperball_haetae_rejection_paper_sim_nma
arg
forge
glob
haetae_mode
init
keygen_internal
kg
pk_current
public_key_of_secret
queries
records
secret_key_of_seed
sk_current
transcripts
verify
\end{lstlisting}

\end{formaltheorem}

\begin{formaltheorem}[EasyCrypt line 8241]
\noindent\textbf{EasyCrypt name.}
\begin{lstlisting}[language=EasyCrypt,basicstyle=\ttfamily\scriptsize]
public_sim_to_paper_sim_nma_hop_from_sampling_hop
\end{lstlisting}
\noindent\textbf{Checked EasyCrypt statement.}
\begin{lstlisting}[language=EasyCrypt,basicstyle=\ttfamily\scriptsize]
lemma public_sim_to_paper_sim_nma_hop_from_sampling_hop &m :
  Pr[SIG.UF_NMA(H, HAETAE, ROMInternalTranscriptPublicSimAsNMA(A)).main()
       @ &m : res] <=
  Pr[SIG.UF_NMA(H, HAETAE,
       ROMInternalTranscriptPaperSimAsNMA(A, Samp)).main() @ &m : res] +
    rejection_sampling_loss_term =>
  Pr[SIG.UF_NMA(H, HAETAE, ROMInternalTranscriptPublicSimAsNMA(A)).main()
       @ &m : res] <=
  Pr[SIG.UF_NMA(H, HAETAE,
       ROMInternalTranscriptPaperSimAsNMA(A, Samp)).main() @ &m : res] +
    rejection_sampling_loss_term.
\end{lstlisting}
\noindent\textbf{Formal mathematical reading.}
Mathematical theorem. This theorem has the form \(\Pr[G:E] \le B\) is a game-probability statement.  It is one of the exact hops, bad-event bounds, or final advantage inequalities used in the reduction.

\noindent\textbf{Role in the proof.}
Use in this chapter. It proves an exact game replacement or simulator equivalence.

\noindent\textbf{Proof-script interpretation.}
No proof block was collected for this item.

\end{formaltheorem}

\begin{formaltheorem}[EasyCrypt line 8268]
\noindent\textbf{EasyCrypt name.}
\begin{lstlisting}[language=EasyCrypt,basicstyle=\ttfamily\scriptsize]
transcript_logging_oracle_erasure
\end{lstlisting}
\noindent\textbf{Checked EasyCrypt statement.}
\begin{lstlisting}[language=EasyCrypt,basicstyle=\ttfamily\scriptsize]
equiv transcript_logging_oracle_erasure :
  EUF_CMA_SimulatedSign(H, S, A, Sim).O.sign ~
  EUF_CMA_TranscriptSimulatedSign(H, S, A, Sim).O.sign :
    ={glob H, glob S, glob Sim, arg} /\
    EUF_CMA_SimulatedSign.queries{1} =
      EUF_CMA_TranscriptSimulatedSign.queries{2}
    ==>
    ={glob H, glob S, glob Sim, res} /\
    EUF_CMA_SimulatedSign.queries{1} =
      EUF_CMA_TranscriptSimulatedSign.queries{2}.
\end{lstlisting}
\noindent\textbf{Formal mathematical reading.}
Mathematical theorem. This is a relational program theorem: two games or procedures are coupled so that a precondition on their initial states implies the stated postcondition on final states and return values.

\noindent\textbf{Role in the proof.}
Use in this chapter. It contributes directly to an advantage bound, bad-event bound, or real-arithmetic composition step.

\noindent\textbf{Proof-script interpretation.}
Proof-script reading. The machine proof decomposes the theorem into rewrites, program-logic obligations, relational equivalences, and arithmetic side conditions.
\emph{Main tactics visible in the proof script:} \ec{proc}, \ec{wp}, \ec{call}.
\emph{Named identifiers syntactically cited in the proof block:}
\begin{lstlisting}[language=EasyCrypt,basicstyle=\ttfamily\scriptsize]
arg
glob
\end{lstlisting}

\end{formaltheorem}

\begin{formaltheorem}[EasyCrypt line 8286]
\noindent\textbf{EasyCrypt name.}
\begin{lstlisting}[language=EasyCrypt,basicstyle=\ttfamily\scriptsize]
adversary_transcript_logging_erasure
\end{lstlisting}
\noindent\textbf{Checked EasyCrypt statement.}
\begin{lstlisting}[language=EasyCrypt,basicstyle=\ttfamily\scriptsize]
equiv adversary_transcript_logging_erasure :
  A(H, EUF_CMA_SimulatedSign(H, S, A, Sim).O).forge ~
  A(H, EUF_CMA_TranscriptSimulatedSign(H, S, A, Sim).O).forge :
    ={glob H, glob S, glob A, glob Sim, arg} /\
    EUF_CMA_SimulatedSign.queries{1} =
      EUF_CMA_TranscriptSimulatedSign.queries{2}
    ==>
    ={glob H, glob S, glob Sim, res} /\
    EUF_CMA_SimulatedSign.queries{1} =
      EUF_CMA_TranscriptSimulatedSign.queries{2}.
\end{lstlisting}
\noindent\textbf{Formal mathematical reading.}
Mathematical theorem. This is a relational program theorem: two games or procedures are coupled so that a precondition on their initial states implies the stated postcondition on final states and return values.

\noindent\textbf{Role in the proof.}
Use in this chapter. It proves an exact game replacement or simulator equivalence.

\noindent\textbf{Proof-script interpretation.}
Proof-script reading. The machine proof decomposes the theorem into rewrites, program-logic obligations, relational equivalences, and arithmetic side conditions.
\emph{Main tactics visible in the proof script:} \ec{proc}, \ec{wp}, \ec{call}.
\emph{Named identifiers syntactically cited in the proof block:}
\begin{lstlisting}[language=EasyCrypt,basicstyle=\ttfamily\scriptsize]
EUF_CMA_SimulatedSign
EUF_CMA_TranscriptSimulatedSign
arg
glob
queries
\end{lstlisting}

\end{formaltheorem}

\begin{formaltheorem}[EasyCrypt line 8309]
\noindent\textbf{EasyCrypt name.}
\begin{lstlisting}[language=EasyCrypt,basicstyle=\ttfamily\scriptsize]
transcript_logging_erasure_exact
\end{lstlisting}
\noindent\textbf{Checked EasyCrypt statement.}
\begin{lstlisting}[language=EasyCrypt,basicstyle=\ttfamily\scriptsize]
lemma transcript_logging_erasure_exact &m :
  Pr[EUF_CMA_SimulatedSign(H, S, A, Sim).main() @ &m : res] =
  Pr[EUF_CMA_TranscriptSimulatedSign(H, S, A, Sim).main() @ &m : res].
\end{lstlisting}
\noindent\textbf{Formal mathematical reading.}
Mathematical theorem. This theorem has the form \(\Pr[G:E] = B\) is a game-probability statement.  It is one of the exact hops, bad-event bounds, or final advantage inequalities used in the reduction.

\noindent\textbf{Role in the proof.}
Use in this chapter. It proves an exact game replacement or simulator equivalence.

\noindent\textbf{Proof-script interpretation.}
Proof-script reading. The machine proof decomposes the theorem into rewrites, program-logic obligations, relational equivalences, and arithmetic side conditions.
\emph{Main tactics visible in the proof script:} \ec{byequiv}, \ec{proc}, \ec{wp}, \ec{call}, \ec{apply}.
\emph{Named identifiers syntactically cited in the proof block:}
\begin{lstlisting}[language=EasyCrypt,basicstyle=\ttfamily\scriptsize]
EUF_CMA_SimulatedSign
EUF_CMA_TranscriptSimulatedSign
adversary_transcript_logging_erasure
arg
glob
queries
\end{lstlisting}

\end{formaltheorem}

\begin{formaltheorem}[EasyCrypt line 8350]
\noindent\textbf{EasyCrypt name.}
\begin{lstlisting}[language=EasyCrypt,basicstyle=\ttfamily\scriptsize]
real_signing_oracle_equiv
\end{lstlisting}
\noindent\textbf{Checked EasyCrypt statement.}
\begin{lstlisting}[language=EasyCrypt,basicstyle=\ttfamily\scriptsize]
equiv real_signing_oracle_equiv :
  SIG.EUF_CMA(H, S, A).O.sign ~
  EUF_CMA_SimulatedSign(H, S, A, RealSigningSimulator(S)).O.sign :
    ={glob H, glob S, arg} /\
    SIG.EUF_CMA.sk{1} = RealSigningSimulator.sk{2} /\
    SIG.EUF_CMA.queries{1} = EUF_CMA_SimulatedSign.queries{2}
    ==>
    ={glob H, glob S, res} /\
    SIG.EUF_CMA.sk{1} = RealSigningSimulator.sk{2} /\
    SIG.EUF_CMA.queries{1} = EUF_CMA_SimulatedSign.queries{2}.
\end{lstlisting}
\noindent\textbf{Formal mathematical reading.}
Mathematical theorem. This is a relational program theorem: two games or procedures are coupled so that a precondition on their initial states implies the stated postcondition on final states and return values.

\noindent\textbf{Role in the proof.}
Use in this chapter. It contributes directly to an advantage bound, bad-event bound, or real-arithmetic composition step.

\noindent\textbf{Proof-script interpretation.}
Proof-script reading. The machine proof decomposes the theorem into rewrites, program-logic obligations, relational equivalences, and arithmetic side conditions.
\emph{Main tactics visible in the proof script:} \ec{proc}, \ec{inline}, \ec{wp}, \ec{call}.
\emph{Named identifiers syntactically cited in the proof block:}
\begin{lstlisting}[language=EasyCrypt,basicstyle=\ttfamily\scriptsize]
RealSigningSimulator
arg
glob
sign
\end{lstlisting}

\end{formaltheorem}

\begin{formaltheorem}[EasyCrypt line 8369]
\noindent\textbf{EasyCrypt name.}
\begin{lstlisting}[language=EasyCrypt,basicstyle=\ttfamily\scriptsize]
adversary_real_signing_equiv
\end{lstlisting}
\noindent\textbf{Checked EasyCrypt statement.}
\begin{lstlisting}[language=EasyCrypt,basicstyle=\ttfamily\scriptsize]
equiv adversary_real_signing_equiv :
  A(H, SIG.EUF_CMA(H, S, A).O).forge ~
  A(H, EUF_CMA_SimulatedSign(H, S, A, RealSigningSimulator(S)).O).forge :
    ={glob H, glob S, glob A, arg} /\
    SIG.EUF_CMA.sk{1} = RealSigningSimulator.sk{2} /\
    SIG.EUF_CMA.queries{1} = EUF_CMA_SimulatedSign.queries{2}
    ==>
    ={glob H, glob S, res} /\
    SIG.EUF_CMA.sk{1} = RealSigningSimulator.sk{2} /\
    SIG.EUF_CMA.queries{1} = EUF_CMA_SimulatedSign.queries{2}.
\end{lstlisting}
\noindent\textbf{Formal mathematical reading.}
Mathematical theorem. This is a relational program theorem: two games or procedures are coupled so that a precondition on their initial states implies the stated postcondition on final states and return values.

\noindent\textbf{Role in the proof.}
Use in this chapter. It proves an exact game replacement or simulator equivalence.

\noindent\textbf{Proof-script interpretation.}
Proof-script reading. The machine proof decomposes the theorem into rewrites, program-logic obligations, relational equivalences, and arithmetic side conditions.
\emph{Main tactics visible in the proof script:} \ec{proc}, \ec{inline}, \ec{wp}, \ec{call}.
\emph{Named identifiers syntactically cited in the proof block:}
\begin{lstlisting}[language=EasyCrypt,basicstyle=\ttfamily\scriptsize]
EUF_CMA
EUF_CMA_SimulatedSign
RealSigningSimulator
SIG
arg
glob
queries
sign
sk
\end{lstlisting}

\end{formaltheorem}

\begin{formaltheorem}[EasyCrypt line 8393]
\noindent\textbf{EasyCrypt name.}
\begin{lstlisting}[language=EasyCrypt,basicstyle=\ttfamily\scriptsize]
real_signing_simulator_exact
\end{lstlisting}
\noindent\textbf{Checked EasyCrypt statement.}
\begin{lstlisting}[language=EasyCrypt,basicstyle=\ttfamily\scriptsize]
lemma real_signing_simulator_exact &m :
  Pr[SIG.EUF_CMA(H, S, A).main() @ &m : res] =
  Pr[EUF_CMA_SimulatedSign(H, S, A, RealSigningSimulator(S)).main() @ &m : res].
\end{lstlisting}
\noindent\textbf{Formal mathematical reading.}
Mathematical theorem. This theorem has the form \(\Pr[G:E] = B\) is a game-probability statement.  It is one of the exact hops, bad-event bounds, or final advantage inequalities used in the reduction.

\noindent\textbf{Role in the proof.}
Use in this chapter. It proves an exact game replacement or simulator equivalence.

\noindent\textbf{Proof-script interpretation.}
Proof-script reading. The machine proof decomposes the theorem into rewrites, program-logic obligations, relational equivalences, and arithmetic side conditions.
\emph{Main tactics visible in the proof script:} \ec{byequiv}, \ec{proc}, \ec{inline}, \ec{wp}, \ec{call}, \ec{apply}.
\emph{Named identifiers syntactically cited in the proof block:}
\begin{lstlisting}[language=EasyCrypt,basicstyle=\ttfamily\scriptsize]
EUF_CMA
EUF_CMA_SimulatedSign
RealSigningSimulator
SIG
adversary_real_signing_equiv
arg
glob
init
queries
sk
\end{lstlisting}

\end{formaltheorem}

\begin{formaltheorem}[EasyCrypt line 8507]
\noindent\textbf{EasyCrypt name.}
\begin{lstlisting}[language=EasyCrypt,basicstyle=\ttfamily\scriptsize]
bimodal_special_soundness_lift_exact
\end{lstlisting}
\noindent\textbf{Checked EasyCrypt statement.}
\begin{lstlisting}[language=EasyCrypt,basicstyle=\ttfamily\scriptsize]
lemma bimodal_special_soundness_lift_exact &m md :
  Pr[BimodalSpecialSoundness_Game(H, S, F, X).main(md) @ &m : res] =
  Pr[SpecialSoundness_Game(
       H, S, F, BimodalTranscriptExtractor_As_ModuleSIS(X)).main(md)
       @ &m : res].
\end{lstlisting}
\noindent\textbf{Formal mathematical reading.}
Mathematical theorem. This theorem has the form \(\Pr[G:E] = B\) is a game-probability statement.  It is one of the exact hops, bad-event bounds, or final advantage inequalities used in the reduction.

\noindent\textbf{Role in the proof.}
Use in this chapter. It proves an exact game replacement or simulator equivalence.

\noindent\textbf{Proof-script interpretation.}
Proof-script reading. The machine proof decomposes the theorem into rewrites, program-logic obligations, relational equivalences, and arithmetic side conditions.
\emph{Main tactics visible in the proof script:} \ec{byequiv}, \ec{proc}, \ec{inline}, \ec{wp}, \ec{call}, \ec{rewrite}, \ec{smt}.
\emph{Named identifiers syntactically cited in the proof block:}
\begin{lstlisting}[language=EasyCrypt,basicstyle=\ttfamily\scriptsize]
BimodalTranscriptExtractor_As_ModuleSIS
arg
bimodal_selftarget_msis_valid
bimodal_to_module_sis_instance
extract
glob
module_sis_instance_of_fork
\end{lstlisting}

\end{formaltheorem}

\begin{formaltheorem}[EasyCrypt line 8546]
\noindent\textbf{EasyCrypt name.}
\begin{lstlisting}[language=EasyCrypt,basicstyle=\ttfamily\scriptsize]
special_soundness_interfaces_ready_from_forking
\end{lstlisting}
\noindent\textbf{Checked EasyCrypt statement.}
\begin{lstlisting}[language=EasyCrypt,basicstyle=\ttfamily\scriptsize]
lemma special_soundness_interfaces_ready_from_forking md :
  forking_extraction_obligation md =>
  special_soundness_interfaces_ready md.
\end{lstlisting}
\noindent\textbf{Formal mathematical reading.}
Mathematical theorem. This is an implication, normally turning one invariant, validity predicate, or proof obligation into another.

\noindent\textbf{Role in the proof.}
Use in this chapter. It proves a structural correctness fact needed by later extraction or verification arguments.

\noindent\textbf{Proof-script interpretation.}
Proof-script reading. The machine proof decomposes the theorem into rewrites, program-logic obligations, relational equivalences, and arithmetic side conditions.
\emph{Main tactics visible in the proof script:} \ec{rewrite}, \ec{split}.
\emph{Named identifiers syntactically cited in the proof block:}
\begin{lstlisting}[language=EasyCrypt,basicstyle=\ttfamily\scriptsize]
bimodal_to_module_sis_lift_obligation_holds
exact
fork_public_reconstruction_obligation_holds
hfork
md
special_soundness_interfaces_ready
\end{lstlisting}

\end{formaltheorem}

\begin{formaltheorem}[EasyCrypt line 8559]
\noindent\textbf{EasyCrypt name.}
\begin{lstlisting}[language=EasyCrypt,basicstyle=\ttfamily\scriptsize]
special_soundness_interfaces_public_reconstruction
\end{lstlisting}
\noindent\textbf{Checked EasyCrypt statement.}
\begin{lstlisting}[language=EasyCrypt,basicstyle=\ttfamily\scriptsize]
lemma special_soundness_interfaces_public_reconstruction md :
  special_soundness_interfaces_ready md =>
  fork_public_reconstruction_obligation md.
\end{lstlisting}
\noindent\textbf{Formal mathematical reading.}
Mathematical theorem. This is an implication, normally turning one invariant, validity predicate, or proof obligation into another.

\noindent\textbf{Role in the proof.}
Use in this chapter. It proves a structural correctness fact needed by later extraction or verification arguments.

\noindent\textbf{Proof-script interpretation.}
Proof-script reading. The machine proof decomposes the theorem into rewrites, program-logic obligations, relational equivalences, and arithmetic side conditions.
\emph{Main tactics visible in the proof script:} \ec{rewrite}.
\emph{Named identifiers syntactically cited in the proof block:}
\begin{lstlisting}[language=EasyCrypt,basicstyle=\ttfamily\scriptsize]
exact
hpublic
special_soundness_interfaces_ready
\end{lstlisting}

\end{formaltheorem}

\begin{formaltheorem}[EasyCrypt line 8568]
\noindent\textbf{EasyCrypt name.}
\begin{lstlisting}[language=EasyCrypt,basicstyle=\ttfamily\scriptsize]
special_soundness_interfaces_ready_sound
\end{lstlisting}
\noindent\textbf{Checked EasyCrypt statement.}
\begin{lstlisting}[language=EasyCrypt,basicstyle=\ttfamily\scriptsize]
lemma special_soundness_interfaces_ready_sound md :
  special_soundness_interfaces_ready md =>
  special_soundness_obligation md.
\end{lstlisting}
\noindent\textbf{Formal mathematical reading.}
Mathematical theorem. This is an implication, normally turning one invariant, validity predicate, or proof obligation into another.

\noindent\textbf{Role in the proof.}
Use in this chapter. It proves a structural correctness fact needed by later extraction or verification arguments.

\noindent\textbf{Proof-script interpretation.}
Proof-script reading. The machine proof decomposes the theorem into rewrites, program-logic obligations, relational equivalences, and arithmetic side conditions.
\emph{Main tactics visible in the proof script:} \ec{rewrite}.
\emph{Named identifiers syntactically cited in the proof block:}
\begin{lstlisting}[language=EasyCrypt,basicstyle=\ttfamily\scriptsize]
exact
hfork
hlift
md
special_soundness_from_bimodal_lift
special_soundness_interfaces_ready
\end{lstlisting}

\end{formaltheorem}

\begin{formaltheorem}[EasyCrypt line 8577]
\noindent\textbf{EasyCrypt name.}
\begin{lstlisting}[language=EasyCrypt,basicstyle=\ttfamily\scriptsize]
special_soundness_interfaces_ready_sound_with_public_reconstruction
\end{lstlisting}
\noindent\textbf{Checked EasyCrypt statement.}
\begin{lstlisting}[language=EasyCrypt,basicstyle=\ttfamily\scriptsize]
lemma special_soundness_interfaces_ready_sound_with_public_reconstruction md :
  special_soundness_interfaces_ready md =>
  special_soundness_with_public_reconstruction_obligation md.
\end{lstlisting}
\noindent\textbf{Formal mathematical reading.}
Mathematical theorem. This is an implication, normally turning one invariant, validity predicate, or proof obligation into another.

\noindent\textbf{Role in the proof.}
Use in this chapter. It proves a structural correctness fact needed by later extraction or verification arguments.

\noindent\textbf{Proof-script interpretation.}
Proof-script reading. The machine proof decomposes the theorem into rewrites, program-logic obligations, relational equivalences, and arithmetic side conditions.
\emph{Main tactics visible in the proof script:} \ec{rewrite}.
\emph{Named identifiers syntactically cited in the proof block:}
\begin{lstlisting}[language=EasyCrypt,basicstyle=\ttfamily\scriptsize]
exact
hfork
hlift
hpublic
md
special_soundness_interfaces_ready
special_soundness_with_public_reconstruction_from_bimodal_lift
\end{lstlisting}

\end{formaltheorem}

\begin{formaltheorem}[EasyCrypt line 8592]
\noindent\textbf{EasyCrypt name.}
\begin{lstlisting}[language=EasyCrypt,basicstyle=\ttfamily\scriptsize]
signing_simulator_sound_current
\end{lstlisting}
\noindent\textbf{Checked EasyCrypt statement.}
\begin{lstlisting}[language=EasyCrypt,basicstyle=\ttfamily\scriptsize]
lemma signing_simulator_sound_current md :
  signing_simulator_sound md.
\end{lstlisting}
\noindent\textbf{Formal mathematical reading.}
Mathematical theorem. This lemma is a universally quantified proposition over the variables shown in the EasyCrypt statement.

\noindent\textbf{Role in the proof.}
Use in this chapter. It proves an exact game replacement or simulator equivalence.

\noindent\textbf{Proof-script interpretation.}
Proof-script reading. The machine proof decomposes the theorem into rewrites, program-logic obligations, relational equivalences, and arithmetic side conditions.
\emph{Main tactics visible in the proof script:} \ec{rewrite}.
\emph{Named identifiers syntactically cited in the proof block:}
\begin{lstlisting}[language=EasyCrypt,basicstyle=\ttfamily\scriptsize]
ctx
exact
md
pk
sig
signing_simulation_verifies
signing_simulator_sound
\end{lstlisting}

\end{formaltheorem}

\begin{formaltheorem}[EasyCrypt line 8600]
\noindent\textbf{EasyCrypt name.}
\begin{lstlisting}[language=EasyCrypt,basicstyle=\ttfamily\scriptsize]
rejection_sampling_bound_from_fs_bound
\end{lstlisting}
\noindent\textbf{Checked EasyCrypt statement.}
\begin{lstlisting}[language=EasyCrypt,basicstyle=\ttfamily\scriptsize]
lemma rejection_sampling_bound_from_fs_bound :
  fs_with_aborts_bound_obligation =>
  rejection_sampling_bound_obligation.
\end{lstlisting}
\noindent\textbf{Formal mathematical reading.}
Mathematical theorem. This is an implication, normally turning one invariant, validity predicate, or proof obligation into another.

\noindent\textbf{Role in the proof.}
Use in this chapter. It contributes directly to an advantage bound, bad-event bound, or real-arithmetic composition step.

\noindent\textbf{Proof-script interpretation.}
Proof-script reading. The machine proof decomposes the theorem into rewrites, program-logic obligations, relational equivalences, and arithmetic side conditions.
\emph{Main tactics visible in the proof script:} \ec{rewrite}, \ec{split}, \ec{apply}.
\emph{Named identifiers syntactically cited in the proof block:}
\begin{lstlisting}[language=EasyCrypt,basicstyle=\ttfamily\scriptsize]
addr_ge0
exact
fs_with_aborts_bound_obligation
ge_entropy
ge_reprogram
rejection_sampling_bound_obligation
rejection_sampling_loss_term
\end{lstlisting}

\end{formaltheorem}

\begin{formaltheorem}[EasyCrypt line 8620]
\noindent\textbf{EasyCrypt name.}
\begin{lstlisting}[language=EasyCrypt,basicstyle=\ttfamily\scriptsize]
fs_with_aborts_interfaces_ready_from_parts
\end{lstlisting}
\noindent\textbf{Checked EasyCrypt statement.}
\begin{lstlisting}[language=EasyCrypt,basicstyle=\ttfamily\scriptsize]
lemma fs_with_aborts_interfaces_ready_from_parts md :
  special_soundness_interfaces_ready md =>
  fs_with_aborts_bound_obligation =>
  fs_with_aborts_interfaces_ready md.
\end{lstlisting}
\noindent\textbf{Formal mathematical reading.}
Mathematical theorem. This is an implication, normally turning one invariant, validity predicate, or proof obligation into another.

\noindent\textbf{Role in the proof.}
Use in this chapter. It is a reusable checked theorem cited by later proof scripts.

\noindent\textbf{Proof-script interpretation.}
Proof-script reading. The machine proof decomposes the theorem into rewrites, program-logic obligations, relational equivalences, and arithmetic side conditions.
\emph{Main tactics visible in the proof script:} \ec{rewrite}, \ec{split}.
\emph{Named identifiers syntactically cited in the proof block:}
\begin{lstlisting}[language=EasyCrypt,basicstyle=\ttfamily\scriptsize]
exact
fs_with_aborts_interfaces_ready
hfs_bound
hspecial
md
signing_simulator_sound_current
\end{lstlisting}

\end{formaltheorem}

\begin{formaltheorem}[EasyCrypt line 8634]
\noindent\textbf{EasyCrypt name.}
\begin{lstlisting}[language=EasyCrypt,basicstyle=\ttfamily\scriptsize]
fs_with_aborts_interfaces_ready_from_forking
\end{lstlisting}
\noindent\textbf{Checked EasyCrypt statement.}
\begin{lstlisting}[language=EasyCrypt,basicstyle=\ttfamily\scriptsize]
lemma fs_with_aborts_interfaces_ready_from_forking md :
  forking_extraction_obligation md =>
  fs_with_aborts_interfaces_ready md.
\end{lstlisting}
\noindent\textbf{Formal mathematical reading.}
Mathematical theorem. This is an implication, normally turning one invariant, validity predicate, or proof obligation into another.

\noindent\textbf{Role in the proof.}
Use in this chapter. It is a reusable checked theorem cited by later proof scripts.

\noindent\textbf{Proof-script interpretation.}
Proof-script reading. The machine proof decomposes the theorem into rewrites, program-logic obligations, relational equivalences, and arithmetic side conditions.
\emph{Named identifiers syntactically cited in the proof block:}
\begin{lstlisting}[language=EasyCrypt,basicstyle=\ttfamily\scriptsize]
exact
fork
fs_with_aborts_bound_obligation_holds
fs_with_aborts_interfaces_ready_from_parts
md
special_soundness_interfaces_ready_from_forking
\end{lstlisting}

\end{formaltheorem}

\begin{formaltheorem}[EasyCrypt line 8644]
\noindent\textbf{EasyCrypt name.}
\begin{lstlisting}[language=EasyCrypt,basicstyle=\ttfamily\scriptsize]
fs_with_aborts_interfaces_ready_holds
\end{lstlisting}
\noindent\textbf{Checked EasyCrypt statement.}
\begin{lstlisting}[language=EasyCrypt,basicstyle=\ttfamily\scriptsize]
lemma fs_with_aborts_interfaces_ready_holds md :
  fs_with_aborts_interfaces_ready md.
\end{lstlisting}
\noindent\textbf{Formal mathematical reading.}
Mathematical theorem. This lemma is a universally quantified proposition over the variables shown in the EasyCrypt statement.

\noindent\textbf{Role in the proof.}
Use in this chapter. It is a reusable checked theorem cited by later proof scripts.

\noindent\textbf{Proof-script interpretation.}
Proof-script reading. The machine proof decomposes the theorem into rewrites, program-logic obligations, relational equivalences, and arithmetic side conditions.
\emph{Named identifiers syntactically cited in the proof block:}
\begin{lstlisting}[language=EasyCrypt,basicstyle=\ttfamily\scriptsize]
exact
forking_extraction_obligation_holds
fs_with_aborts_interfaces_ready_from_forking
md
\end{lstlisting}

\end{formaltheorem}

\begin{formaltheorem}[EasyCrypt line 8651]
\noindent\textbf{EasyCrypt name.}
\begin{lstlisting}[language=EasyCrypt,basicstyle=\ttfamily\scriptsize]
structural_to_exact_hyperball_paper_sample_loss_from_obligation
\end{lstlisting}
\noindent\textbf{Checked EasyCrypt statement.}
\begin{lstlisting}[language=EasyCrypt,basicstyle=\ttfamily\scriptsize]
lemma structural_to_exact_hyperball_paper_sample_loss_from_obligation
   (md : mode) :
  structural_to_exact_hyperball_paper_sample_loss_obligation md =>
  forall (sk : skey) (m : message) (ctx : context)
         (p : paper_sim_signature_sample -> bool),
    mu (dsigning_attempt_state md sk m ctx)
      (fun st => p (paper_sim_sample_from_rejection_attempt st)) <=
    mu (dexact_hyperball_signing_attempt_state md sk m ctx)
      (fun st => p (paper_sim_sample_from_rejection_attempt st)) +
      rejection_sampling_loss_term.
\end{lstlisting}
\noindent\textbf{Formal mathematical reading.}
Mathematical theorem. This is an inequality, usually an arithmetic loss bound, event inclusion bound, or monotonicity fact used to compose probabilities.

\noindent\textbf{Role in the proof.}
Use in this chapter. It contributes directly to an advantage bound, bad-event bound, or real-arithmetic composition step.

\noindent\textbf{Proof-script interpretation.}
Proof-script reading. The machine proof decomposes the theorem into rewrites, program-logic obligations, relational equivalences, and arithmetic side conditions.
\emph{Main tactics visible in the proof script:} \ec{rewrite}.
\emph{Named identifiers syntactically cited in the proof block:}
\begin{lstlisting}[language=EasyCrypt,basicstyle=\ttfamily\scriptsize]
ctx
exact
hop
sk
structural_to_exact_hyperball_paper_sample_loss_obligation
\end{lstlisting}

\end{formaltheorem}

\begin{formaltheorem}[EasyCrypt line 8667]
\noindent\textbf{EasyCrypt name.}
\begin{lstlisting}[language=EasyCrypt,basicstyle=\ttfamily\scriptsize]
structural_to_exact_hyperball_paper_sample_loss_from_attempt_loss
\end{lstlisting}
\noindent\textbf{Checked EasyCrypt statement.}
\begin{lstlisting}[language=EasyCrypt,basicstyle=\ttfamily\scriptsize]
lemma structural_to_exact_hyperball_paper_sample_loss_from_attempt_loss
   (md : mode) :
  structural_to_exact_hyperball_attempt_loss_obligation =>
  forall (sk : skey) (m : message) (ctx : context)
         (p : paper_sim_signature_sample -> bool),
    mu (dsigning_attempt_state md sk m ctx)
      (fun st => p (paper_sim_sample_from_rejection_attempt st)) <=
    mu (dexact_hyperball_signing_attempt_state md sk m ctx)
      (fun st => p (paper_sim_sample_from_rejection_attempt st)) +
      rejection_sampling_loss_term.
\end{lstlisting}
\noindent\textbf{Formal mathematical reading.}
Mathematical theorem. This is an inequality, usually an arithmetic loss bound, event inclusion bound, or monotonicity fact used to compose probabilities.

\noindent\textbf{Role in the proof.}
Use in this chapter. It contributes directly to an advantage bound, bad-event bound, or real-arithmetic composition step.

\noindent\textbf{Proof-script interpretation.}
Proof-script reading. The machine proof decomposes the theorem into rewrites, program-logic obligations, relational equivalences, and arithmetic side conditions.
\emph{Main tactics visible in the proof script:} \ec{rewrite}.
\emph{Named identifiers syntactically cited in the proof block:}
\begin{lstlisting}[language=EasyCrypt,basicstyle=\ttfamily\scriptsize]
ctx
exact
hop
md
paper_sim_sample_from_rejection_attempt
sk
st
structural_to_exact_hyperball_attempt_loss_obligation
\end{lstlisting}

\end{formaltheorem}

\begin{formaltheorem}[EasyCrypt line 8684]
\noindent\textbf{EasyCrypt name.}
\begin{lstlisting}[language=EasyCrypt,basicstyle=\ttfamily\scriptsize]
structural_to_exact_hyperball_paper_sample_loss_obligation_from_attempt_loss
\end{lstlisting}
\noindent\textbf{Checked EasyCrypt statement.}
\begin{lstlisting}[language=EasyCrypt,basicstyle=\ttfamily\scriptsize]
lemma structural_to_exact_hyperball_paper_sample_loss_obligation_from_attempt_loss
   (md : mode) :
  structural_to_exact_hyperball_attempt_loss_obligation =>
  structural_to_exact_hyperball_paper_sample_loss_obligation md.
\end{lstlisting}
\noindent\textbf{Formal mathematical reading.}
Mathematical theorem. This is an implication, normally turning one invariant, validity predicate, or proof obligation into another.

\noindent\textbf{Role in the proof.}
Use in this chapter. It contributes directly to an advantage bound, bad-event bound, or real-arithmetic composition step.

\noindent\textbf{Proof-script interpretation.}
Proof-script reading. The machine proof decomposes the theorem into rewrites, program-logic obligations, relational equivalences, and arithmetic side conditions.
\emph{Main tactics visible in the proof script:} \ec{rewrite}.
\emph{Named identifiers syntactically cited in the proof block:}
\begin{lstlisting}[language=EasyCrypt,basicstyle=\ttfamily\scriptsize]
ctx
exact
hop
md
sk
structural_to_exact_hyperball_paper_sample_loss_from_attempt_loss
structural_to_exact_hyperball_paper_sample_loss_obligation
\end{lstlisting}

\end{formaltheorem}

\begin{formaltheorem}[EasyCrypt line 8696]
\noindent\textbf{EasyCrypt name.}
\begin{lstlisting}[language=EasyCrypt,basicstyle=\ttfamily\scriptsize]
structural_to_exact_hyperball_paper_sample_loss_from_coin_sample_pair_loss
\end{lstlisting}
\noindent\textbf{Checked EasyCrypt statement.}
\begin{lstlisting}[language=EasyCrypt,basicstyle=\ttfamily\scriptsize]
lemma structural_to_exact_hyperball_paper_sample_loss_from_coin_sample_pair_loss
   (md : mode) :
  structural_to_exact_hyperball_coin_sample_pair_loss_obligation =>
  forall (sk : skey) (m : message) (ctx : context)
         (p : paper_sim_signature_sample -> bool),
    mu (dsigning_attempt_state md sk m ctx)
      (fun st => p (paper_sim_sample_from_rejection_attempt st)) <=
    mu (dexact_hyperball_signing_attempt_state md sk m ctx)
      (fun st => p (paper_sim_sample_from_rejection_attempt st)) +
      rejection_sampling_loss_term.
\end{lstlisting}
\noindent\textbf{Formal mathematical reading.}
Mathematical theorem. This is an inequality, usually an arithmetic loss bound, event inclusion bound, or monotonicity fact used to compose probabilities.

\noindent\textbf{Role in the proof.}
Use in this chapter. It contributes directly to an advantage bound, bad-event bound, or real-arithmetic composition step.

\noindent\textbf{Proof-script interpretation.}
Proof-script reading. The machine proof decomposes the theorem into rewrites, program-logic obligations, relational equivalences, and arithmetic side conditions.
\emph{Main tactics visible in the proof script:} \ec{apply}.
\emph{Named identifiers syntactically cited in the proof block:}
\begin{lstlisting}[language=EasyCrypt,basicstyle=\ttfamily\scriptsize]
hop
structural_to_exact_hyperball_attempt_loss_from_coin_sample_pair_loss
structural_to_exact_hyperball_paper_sample_loss_from_attempt_loss
\end{lstlisting}

\end{formaltheorem}

\begin{formaltheorem}[EasyCrypt line 8712]
\noindent\textbf{EasyCrypt name.}
\begin{lstlisting}[language=EasyCrypt,basicstyle=\ttfamily\scriptsize]
structural_to_exact_hyperball_paper_sample_loss_obligation_from_coin_sample_pair_loss
\end{lstlisting}
\noindent\textbf{Checked EasyCrypt statement.}
\begin{lstlisting}[language=EasyCrypt,basicstyle=\ttfamily\scriptsize]
lemma structural_to_exact_hyperball_paper_sample_loss_obligation_from_coin_sample_pair_loss
   (md : mode) :
  structural_to_exact_hyperball_coin_sample_pair_loss_obligation =>
  structural_to_exact_hyperball_paper_sample_loss_obligation md.
\end{lstlisting}
\noindent\textbf{Formal mathematical reading.}
Mathematical theorem. This is an implication, normally turning one invariant, validity predicate, or proof obligation into another.

\noindent\textbf{Role in the proof.}
Use in this chapter. It contributes directly to an advantage bound, bad-event bound, or real-arithmetic composition step.

\noindent\textbf{Proof-script interpretation.}
Proof-script reading. The machine proof decomposes the theorem into rewrites, program-logic obligations, relational equivalences, and arithmetic side conditions.
\emph{Main tactics visible in the proof script:} \ec{rewrite}.
\emph{Named identifiers syntactically cited in the proof block:}
\begin{lstlisting}[language=EasyCrypt,basicstyle=\ttfamily\scriptsize]
ctx
exact
hop
md
sk
structural_to_exact_hyperball_paper_sample_loss_from_coin_sample_pair_loss
structural_to_exact_hyperball_paper_sample_loss_obligation
\end{lstlisting}

\end{formaltheorem}

\begin{formaltheorem}[EasyCrypt line 8724]
\noindent\textbf{EasyCrypt name.}
\begin{lstlisting}[language=EasyCrypt,basicstyle=\ttfamily\scriptsize]
dexact_hyperball_reference_paper_sim_attempt_boundary_ready
\end{lstlisting}
\noindent\textbf{Checked EasyCrypt statement.}
\begin{lstlisting}[language=EasyCrypt,basicstyle=\ttfamily\scriptsize]
lemma dexact_hyperball_reference_paper_sim_attempt_boundary_ready
   (md : mode) (sk : skey) (m : message) (ctx : context) :
  mu (dexact_hyperball_signing_attempt_state md sk m ctx)
    (signing_attempt_state_reference_rejection_aborts md) = 0%r.
\end{lstlisting}
\noindent\textbf{Formal mathematical reading.}
Mathematical theorem. This is an equality or definitional expansion used for rewriting and normalization.

\noindent\textbf{Role in the proof.}
Use in this chapter. It contributes directly to an advantage bound, bad-event bound, or real-arithmetic composition step.

\noindent\textbf{Proof-script interpretation.}
Proof-script reading. The machine proof decomposes the theorem into rewrites, program-logic obligations, relational equivalences, and arithmetic side conditions.
\emph{Main tactics visible in the proof script:} \ec{apply}.
\emph{Named identifiers syntactically cited in the proof block:}
\begin{lstlisting}[language=EasyCrypt,basicstyle=\ttfamily\scriptsize]
dexact_hyperball_signing_attempt_state_reference_rejection_abort_zero
\end{lstlisting}

\end{formaltheorem}

\begin{formaltheorem}[EasyCrypt line 8732]
\noindent\textbf{EasyCrypt name.}
\begin{lstlisting}[language=EasyCrypt,basicstyle=\ttfamily\scriptsize]
dexact_hyperball_full_rejection_paper_sim_attempt_boundary_ready
\end{lstlisting}
\noindent\textbf{Checked EasyCrypt statement.}
\begin{lstlisting}[language=EasyCrypt,basicstyle=\ttfamily\scriptsize]
lemma dexact_hyperball_full_rejection_paper_sim_attempt_boundary_ready
   (md : mode) (sk : skey) (m : message) (ctx : context) :
  mu (dexact_hyperball_signing_attempt_state md sk m ctx)
    (signing_attempt_state_rejection_aborts md
       (public_key_of_secret md sk) m ctx) = 0%r.
\end{lstlisting}
\noindent\textbf{Formal mathematical reading.}
Mathematical theorem. This is an equality or definitional expansion used for rewriting and normalization.

\noindent\textbf{Role in the proof.}
Use in this chapter. It contributes directly to an advantage bound, bad-event bound, or real-arithmetic composition step.

\noindent\textbf{Proof-script interpretation.}
Proof-script reading. The machine proof decomposes the theorem into rewrites, program-logic obligations, relational equivalences, and arithmetic side conditions.
\emph{Main tactics visible in the proof script:} \ec{apply}.
\emph{Named identifiers syntactically cited in the proof block:}
\begin{lstlisting}[language=EasyCrypt,basicstyle=\ttfamily\scriptsize]
dexact_hyperball_signing_attempt_state_rejection_abort_zero
\end{lstlisting}

\end{formaltheorem}

\begin{formaltheorem}[EasyCrypt line 8741]
\noindent\textbf{EasyCrypt name.}
\begin{lstlisting}[language=EasyCrypt,basicstyle=\ttfamily\scriptsize]
dexact_hyperball_haetae_rejection_sample_no_fallback_mu
\end{lstlisting}
\noindent\textbf{Checked EasyCrypt statement.}
\begin{lstlisting}[language=EasyCrypt,basicstyle=\ttfamily\scriptsize]
lemma dexact_hyperball_haetae_rejection_sample_no_fallback_mu
   (md : mode) (sk : skey) (m : message) (ctx : context)
   (p : paper_sim_signature_sample -> bool) :
  mu (dexact_hyperball_signing_attempt_state md sk m ctx)
    (fun st =>
       p (haetae_rejection_sample_from_attempt md
            (public_key_of_secret md sk) m ctx st)) =
  mu (dexact_hyperball_signing_attempt_state md sk m ctx)
    (fun st => p (paper_sim_sample_from_rejection_attempt st)).
\end{lstlisting}
\noindent\textbf{Formal mathematical reading.}
Mathematical theorem. This is an implication, normally turning one invariant, validity predicate, or proof obligation into another.

\noindent\textbf{Role in the proof.}
Use in this chapter. It contributes directly to an advantage bound, bad-event bound, or real-arithmetic composition step.

\noindent\textbf{Proof-script interpretation.}
Proof-script reading. The machine proof decomposes the theorem into rewrites, program-logic obligations, relational equivalences, and arithmetic side conditions.
\emph{Main tactics visible in the proof script:} \ec{apply}, \ec{rewrite}.
\emph{Named identifiers syntactically cited in the proof block:}
\begin{lstlisting}[language=EasyCrypt,basicstyle=\ttfamily\scriptsize]
Distr
ctx
dexact_hyperball_signing_attempt_state_rejection_accepts
haetae_rejection_sample_from_attempt
md
mu_eq_support
sk
st
st_supp
\end{lstlisting}

\end{formaltheorem}

\begin{formaltheorem}[EasyCrypt line 8758]
\noindent\textbf{EasyCrypt name.}
\begin{lstlisting}[language=EasyCrypt,basicstyle=\ttfamily\scriptsize]
dexact_hyperball_haetae_rejection_sample_signature_no_fallback_mu
\end{lstlisting}
\noindent\textbf{Checked EasyCrypt statement.}
\begin{lstlisting}[language=EasyCrypt,basicstyle=\ttfamily\scriptsize]
lemma dexact_hyperball_haetae_rejection_sample_signature_no_fallback_mu
   (md : mode) (sk : skey) (m : message) (ctx : context)
   (p : signature -> bool) :
  mu (dexact_hyperball_signing_attempt_state md sk m ctx)
    (fun st =>
       p (paper_sim_signature md (public_key_of_secret md sk) m ctx
            (haetae_rejection_sample_from_attempt md
               (public_key_of_secret md sk) m ctx st))) =
  mu (dexact_hyperball_signing_attempt_state md sk m ctx)
    (fun st =>
       p (paper_sim_signature md (public_key_of_secret md sk) m ctx
            (paper_sim_sample_from_rejection_attempt st))).
\end{lstlisting}
\noindent\textbf{Formal mathematical reading.}
Mathematical theorem. This is an implication, normally turning one invariant, validity predicate, or proof obligation into another.

\noindent\textbf{Role in the proof.}
Use in this chapter. It contributes directly to an advantage bound, bad-event bound, or real-arithmetic composition step.

\noindent\textbf{Proof-script interpretation.}
Proof-script reading. The machine proof decomposes the theorem into rewrites, program-logic obligations, relational equivalences, and arithmetic side conditions.
\emph{Named identifiers syntactically cited in the proof block:}
\begin{lstlisting}[language=EasyCrypt,basicstyle=\ttfamily\scriptsize]
ctx
dexact_hyperball_haetae_rejection_sample_no_fallback_mu
exact
md
paper_sim_signature
public_key_of_secret
sk
\end{lstlisting}

\end{formaltheorem}

\begin{formaltheorem}[EasyCrypt line 8777]
\noindent\textbf{EasyCrypt name.}
\begin{lstlisting}[language=EasyCrypt,basicstyle=\ttfamily\scriptsize]
dexact_hyperball_paper_sim_rejection_attempt_signature_mu
\end{lstlisting}
\noindent\textbf{Checked EasyCrypt statement.}
\begin{lstlisting}[language=EasyCrypt,basicstyle=\ttfamily\scriptsize]
lemma dexact_hyperball_paper_sim_rejection_attempt_signature_mu
   (md : mode) (sk : skey) (m : message) (ctx : context)
   (p : signature -> bool) :
  mu (dexact_hyperball_signing_attempt_state md sk m ctx)
    (fun st =>
       p (paper_sim_signature md (public_key_of_secret md sk) m ctx
            (paper_sim_sample_from_rejection_attempt st))) =
  mu (dexact_hyperball_signing_attempt_state md sk m ctx)
    (fun st => p (signing_attempt_state_signature st)).
\end{lstlisting}
\noindent\textbf{Formal mathematical reading.}
Mathematical theorem. This is an implication, normally turning one invariant, validity predicate, or proof obligation into another.

\noindent\textbf{Role in the proof.}
Use in this chapter. It proves an exact game replacement or simulator equivalence.

\noindent\textbf{Proof-script interpretation.}
Proof-script reading. The machine proof decomposes the theorem into rewrites, program-logic obligations, relational equivalences, and arithmetic side conditions.
\emph{Main tactics visible in the proof script:} \ec{apply}, \ec{have}, \ec{rewrite}.
\emph{Named identifiers syntactically cited in the proof block:}
\begin{lstlisting}[language=EasyCrypt,basicstyle=\ttfamily\scriptsize]
Distr
ctx
dexact_hyperball_signing_attempt_state_rejection_accepts
field_ok
md
mu_eq_support
paper_sim_signature_rejection_attempt_matches_state
public_key_of_secret
signing_attempt_state_field_accepts
signing_attempt_state_rejection_accepts
signing_attempt_state_verification_accepts
sk
st
st_supp
\end{lstlisting}

\end{formaltheorem}

\begin{formaltheorem}[EasyCrypt line 8800]
\noindent\textbf{EasyCrypt name.}
\begin{lstlisting}[language=EasyCrypt,basicstyle=\ttfamily\scriptsize]
dexact_hyperball_haetae_rejection_sample_signature_mu
\end{lstlisting}
\noindent\textbf{Checked EasyCrypt statement.}
\begin{lstlisting}[language=EasyCrypt,basicstyle=\ttfamily\scriptsize]
lemma dexact_hyperball_haetae_rejection_sample_signature_mu
   (md : mode) (sk : skey) (m : message) (ctx : context)
   (p : signature -> bool) :
  mu (dexact_hyperball_signing_attempt_state md sk m ctx)
    (fun st =>
       p (paper_sim_signature md (public_key_of_secret md sk) m ctx
            (haetae_rejection_sample_from_attempt md
               (public_key_of_secret md sk) m ctx st))) =
  mu (dexact_hyperball_signing_attempt_state md sk m ctx)
    (fun st => p (signing_attempt_state_signature st)).
\end{lstlisting}
\noindent\textbf{Formal mathematical reading.}
Mathematical theorem. This is an implication, normally turning one invariant, validity predicate, or proof obligation into another.

\noindent\textbf{Role in the proof.}
Use in this chapter. It contributes directly to an advantage bound, bad-event bound, or real-arithmetic composition step.

\noindent\textbf{Proof-script interpretation.}
Proof-script reading. The machine proof decomposes the theorem into rewrites, program-logic obligations, relational equivalences, and arithmetic side conditions.
\emph{Main tactics visible in the proof script:} \ec{rewrite}.
\emph{Named identifiers syntactically cited in the proof block:}
\begin{lstlisting}[language=EasyCrypt,basicstyle=\ttfamily\scriptsize]
ctx
dexact_hyperball_haetae_rejection_sample_signature_no_fallback_mu
dexact_hyperball_paper_sim_rejection_attempt_signature_mu
exact
md
sk
\end{lstlisting}

\end{formaltheorem}

\subsubsection*{Explicit Program-Logic Judgments Used Inside Proof Scripts}
\begin{formaljudgment}[EasyCrypt line 2155]
\noindent\textbf{Checked EasyCrypt judgment.}
\begin{lstlisting}[language=EasyCrypt,basicstyle=\ttfamily\scriptsize]
hoare.
\end{lstlisting}
\noindent\textbf{Formal mathematical reading.} This is a Hoare judgment.  It states partial correctness of the displayed procedure from the displayed precondition to the displayed postcondition.
\end{formaljudgment}

\begin{formaljudgment}[EasyCrypt line 2166]
\noindent\textbf{Checked EasyCrypt judgment.}
\begin{lstlisting}[language=EasyCrypt,basicstyle=\ttfamily\scriptsize]
hoare.
\end{lstlisting}
\noindent\textbf{Formal mathematical reading.} This is a Hoare judgment.  It states partial correctness of the displayed procedure from the displayed precondition to the displayed postcondition.
\end{formaljudgment}

\begin{formaljudgment}[EasyCrypt line 3539]
\noindent\textbf{Checked EasyCrypt judgment.}
\begin{lstlisting}[language=EasyCrypt,basicstyle=\ttfamily\scriptsize]
hoare.
\end{lstlisting}
\noindent\textbf{Formal mathematical reading.} This is a Hoare judgment.  It states partial correctness of the displayed procedure from the displayed precondition to the displayed postcondition.
\end{formaljudgment}

\begin{formaljudgment}[EasyCrypt line 3745]
\noindent\textbf{Checked EasyCrypt judgment.}
\begin{lstlisting}[language=EasyCrypt,basicstyle=\ttfamily\scriptsize]
hoare.
\end{lstlisting}
\noindent\textbf{Formal mathematical reading.} This is a Hoare judgment.  It states partial correctness of the displayed procedure from the displayed precondition to the displayed postcondition.
\end{formaljudgment}

\begin{formaljudgment}[EasyCrypt line 3875]
\noindent\textbf{Checked EasyCrypt judgment.}
\begin{lstlisting}[language=EasyCrypt,basicstyle=\ttfamily\scriptsize]
hoare.
\end{lstlisting}
\noindent\textbf{Formal mathematical reading.} This is a Hoare judgment.  It states partial correctness of the displayed procedure from the displayed precondition to the displayed postcondition.
\end{formaljudgment}

\begin{formaljudgment}[EasyCrypt line 4317]
\noindent\textbf{Checked EasyCrypt judgment.}
\begin{lstlisting}[language=EasyCrypt,basicstyle=\ttfamily\scriptsize]
hoare.
\end{lstlisting}
\noindent\textbf{Formal mathematical reading.} This is a Hoare judgment.  It states partial correctness of the displayed procedure from the displayed precondition to the displayed postcondition.
\end{formaljudgment}

\begin{formaljudgment}[EasyCrypt line 4867]
\noindent\textbf{Checked EasyCrypt judgment.}
\begin{lstlisting}[language=EasyCrypt,basicstyle=\ttfamily\scriptsize]
hoare.
\end{lstlisting}
\noindent\textbf{Formal mathematical reading.} This is a Hoare judgment.  It states partial correctness of the displayed procedure from the displayed precondition to the displayed postcondition.
\end{formaljudgment}

\begin{formaljudgment}[EasyCrypt line 4993]
\noindent\textbf{Checked EasyCrypt judgment.}
\begin{lstlisting}[language=EasyCrypt,basicstyle=\ttfamily\scriptsize]
hoare.
\end{lstlisting}
\noindent\textbf{Formal mathematical reading.} This is a Hoare judgment.  It states partial correctness of the displayed procedure from the displayed precondition to the displayed postcondition.
\end{formaljudgment}

\begin{formaljudgment}[EasyCrypt line 4999]
\noindent\textbf{Checked EasyCrypt judgment.}
\begin{lstlisting}[language=EasyCrypt,basicstyle=\ttfamily\scriptsize]
hoare.
\end{lstlisting}
\noindent\textbf{Formal mathematical reading.} This is a Hoare judgment.  It states partial correctness of the displayed procedure from the displayed precondition to the displayed postcondition.
\end{formaljudgment}

\begin{formaljudgment}[EasyCrypt line 5130]
\noindent\textbf{Checked EasyCrypt judgment.}
\begin{lstlisting}[language=EasyCrypt,basicstyle=\ttfamily\scriptsize]
hoare.
\end{lstlisting}
\noindent\textbf{Formal mathematical reading.} This is a Hoare judgment.  It states partial correctness of the displayed procedure from the displayed precondition to the displayed postcondition.
\end{formaljudgment}

\begin{formaljudgment}[EasyCrypt line 5141]
\noindent\textbf{Checked EasyCrypt judgment.}
\begin{lstlisting}[language=EasyCrypt,basicstyle=\ttfamily\scriptsize]
hoare.
\end{lstlisting}
\noindent\textbf{Formal mathematical reading.} This is a Hoare judgment.  It states partial correctness of the displayed procedure from the displayed precondition to the displayed postcondition.
\end{formaljudgment}

\begin{formaljudgment}[EasyCrypt line 7071]
\noindent\textbf{Checked EasyCrypt judgment.}
\begin{lstlisting}[language=EasyCrypt,basicstyle=\ttfamily\scriptsize]
phoare[StructuralAttemptPaperSample.sample :
      arg = (sk0, m0, ctx0) ==> p0 res] <=
    (mu (dexact_hyperball_signing_attempt_state haetae_mode sk0 m0 ctx0)
       (fun st => p0 (paper_sim_sample_from_rejection_attempt st)) +
       rejection_sampling_loss_term).
\end{lstlisting}
\noindent\textbf{Formal mathematical reading.} This is a probabilistic Hoare judgment.  It states that the command satisfies the postcondition with the displayed probability bound.
\end{formaljudgment}

\medskip
\noindent\emph{End of generated formal inventory for \file{HAETAE_HopGames.ec}.}
